\documentclass[11pt]{article}
\usepackage[a4paper,margin=26mm]{geometry}
\usepackage{fontspec}
\setmainfont{Latin Modern Roman}
\usepackage{amsmath,amssymb,amsthm,mathtools}
\usepackage{longtable,booktabs,array,enumitem,xurl}
\usepackage[hidelinks]{hyperref}
\usepackage{fancyhdr}
\pagestyle{fancy}\fancyhf{}\fancyhead[L]{GCT\_KRON: checked research record}\fancyhead[R]{\thepage}
\setlength{\headheight}{15pt}
\setlength{\emergencystretch}{3em}
\newtheorem{theorem}{Theorem}
\newtheorem{lemma}[theorem]{Lemma}
\newtheorem{proposition}[theorem]{Proposition}
\newtheorem{corollary}[theorem]{Corollary}
\title{Kronecker coefficients through Grothendieck's\linebreak
\(\lambda\)-operations on finite group sets}
\author{Research run GCT\_KRON\_20261120}
\date{Working paper, 21 November 2026 (Europe/Berlin)}
\begin{document}
\maketitle
\begin{abstract}
This working paper audits the supplied archive transcriptions, derives
exact formulas and restricted nonvanishing statements, and distinguishes
support positivity from comparison of actual orbit-closure multiplicities.
It corrects the subset interpretation of the Burnside operations and proves
explicit limits on several proposed positivity-to-saturation inferences.
Projection, coefficient-support and jet-dominance arguments compare every
multiplicity copy in explicit ranges. At \(n\geq m^{25}\), a remaining
actual obstruction requires at least fifty-two rows and satisfies
\(\lambda_{2j}\geq28-j\) for \(1\leq j\leq26\), as well as a
joint later-row restriction. Prefix and order-ideal extensions couple
later rows with the second row; explicit positive-copy examples show
strict improvement beyond exact support-cover methods. A grouped differential refinement of the
known canonical Foulkes--Howe argument excludes every degree with
\(n\geq P_d(dm)\), where the explicit recurrence below has
\(P_d(dm)\leq\min\{(d-1)^2,\lfloor d(d+m-1)/2\rfloor\}\).
An explicit normal polynomial highest-weight algebra explains ordinary
saturation failure through its nonprimitive weight lattice. A separate
explicit family has unbounded determinant-orbit multiplicity while its
closure multiplicity is one. For the unpadded permanent, a Hessian
bracket gives an actual multiplicity-one obstruction at every size
\(n\geq3\), together with permanent-positive saturation holes.
A separate explicit coefficient projection gives the actual comparison
\(K=5<q=a=6<\operatorname{sk}=10<g=11\) at
\(n=m=3\), \(d=12\), and \(\lambda=(19,7,2^5)\).
Its determinant-closure value is a precisely cited published input.
The general large-padding
higher-multiplicity comparison remains open in this run. This is a
checkpoint working paper, not a claim to have settled that comparison.
\end{abstract}

\section*{Conventions, source status, and checks}
The attached Ikenmeyer--Panova and B\"urgisser--Ikenmeyer--Panova papers
use the internal padding \(X_{11}^{n-m}\operatorname{perm}_m\).
We use that convention for \(Z_{n,m}\) and \(q^m_\lambda(d[n])\),
so their literal constants and the bound \(\ell(\lambda)\leq m^2\)
apply. The explicit determinant constructions also work with an external
padding variable, but the stated cited GCT constants are not silently
transferred between conventions. Coordinate functions carry the polynomial
\(\mathrm{GL}(V)\)-representations; forms themselves lie in
\(\operatorname{Sym}^n V^*\).

Finding identifiers label arguments and refinements established in this
run, including short dictionary versions of known results. They do not
assert historical priority. Known inputs are identified with exact work
and statement locators. The archive inputs were transcriptions. The official cote~133 facsimile
was subsequently retrieved; A-029 records precisely which original
scans were visually read and recovers the named power-operation definition.
All thirteen attached archive pages and nine further cote~133 pages
were subsequently verified against official facsimiles. Conditional
readings, deletions and the unused uncertain name remain explicit. Full proof-reading events, source-use events
and finite computational checks are recorded separately in
\texttt{CHECK\_EVENTS.tsv}; finite experiments are not promoted to general
proofs. The model is GPT-6 Astra Pro; its reasoning-effort setting was not
exposed and is not invented. The run began at 2026-11-20 23:00:17 UTC.

The prompt stipulates Alp\"oge's RH and Dirichlet-GRH statements as
theorems. None of the arguments here uses them. The supplied four-page
\emph{Remark} was read; the named Lean files and the full scan of Lean
sources were not in the attached ZIP. No Lean verification is claimed.
No Masullo statement is used. The archive, formulas, schemes and
cohomological statements below do not depend on those analytic inputs.
\clearpage
\tableofcontents
\clearpage
\section*{Results and the remaining mathematical question}
\addcontentsline{toc}{section}{Results and the remaining mathematical question}

The supplied \(\lambda\)-operations give exact virtual identities,
positive formulas in specified families, and concrete restrictions on
possible multiplicity obstructions. They do not imply ordinary
Kronecker saturation or determine actual orbit-closure multiplicities
from support alone. The distinction is witnessed by explicit
counterexamples, rather than by an absence of a proposed proof.

\paragraph{The three multiplicities that must remain distinct.}
For \(\lambda\vdash nd\), write \(a_\lambda\) for its ambient
plethysm multiplicity, \(q^m_\lambda\) for its multiplicity on the
internally padded permanent closure, and \(K_n(\lambda)\) for
its multiplicity on the determinant closure. The relevant inequalities
are
\[
 q^m_\lambda\leq a_\lambda,\qquad
 K_n(\lambda)\leq\operatorname{sk}(\lambda,(d^n))
                         \leq g(\lambda,(d^n),(d^n)).
\]
The symmetric Kronecker term is the determinant \emph{orbit}
multiplicity. The closure can lose copies. Thus \(q^m_\lambda>K_n\)
is an actual multiplicity obstruction, whereas \(q^m_\lambda>g\)
is a stronger sufficient test. A theorem that \(g>0\), or even
\(K_n>0\), does not compare positive multiplicities. A-028 computes an unbounded strict gap explicitly: for even \(h\geq4\),
\(K_h(h,h,h)=a=1\), while the orbit and ordinary Kronecker
multiplicities both equal \(\#\{2a+3b+4c=h\}\geq2\).
The exact remaining data in general are ranks of evaluation maps,
as proved in F-010.

\paragraph{Explicit actual obstructions.}
The comparison can now be made with both sides positive. F-025 proves
\[
 \begin{gathered}
 K_3(19,7,2^5)=5
 <q^3_{(19,7,2^5)}(12[3])=a_{(19,7,2^5)}(12[3])=6,\\
 6<\operatorname{sk}=10<g=11.
 \end{gathered}
\]
The determinant-closure multiplicity five is the published calculation
of Landsberg--Manivel--Ressayre, Section~3.2, p.~476. An explicit
integer pencil with a nonzero \(49\)-by-\(49\) Jacobian minor proves
the permanent equality with all six ambient copies. Exact character
sums give the two larger orbit bounds. Thus neither \(q>g\) nor
\(q>\operatorname{sk}\) detects this actual closure obstruction.
More generally, the same projection retains every ambient permanent
copy at \(n=m=3\), \(\ell(\lambda)\leq7\), \(\lambda_3\leq2\).
A-031 extends this equality to \(\ell(\lambda)\leq2m+2\),
\(\lambda_3\leq2\), at each of the four sizes \(m=3,4,5,6\).
It gives further actual obstructions at sizes four, five and six,
with positive determinant multiplicities at least \(7,63,594\),
respectively. Their exact larger permanent counts are not asserted.
The finite witnesses are completely specified integer minors.

O-027 proves, for every determinant size \(n\), that
\(K_n=a_\lambda\) whenever \(\ell(\lambda)\leq2n\) and
\(\lambda_3\leq2\). The universal row threshold is sharp:
the known LMR equations lose an ambient copy at length \(2n+1\).
F-026 preserves six permanent copies in the internally padded shapes
\((12n-17,7,2^5)\), but O-027 shows that the determinant has
enough copies for every \(n\geq4\). Thus the displayed size-three
obstruction does not persist along those first-row shifts.

A-030 supplies a separate infinite family at \(n=m\geq3\): put
\(d=n^2+1\) and
\(\lambda=(d(n-2)+3,3,2^{d-3})\). Then \(q=a=1\) and
\(K=\operatorname{sk}=g=0\). The nonzero permanent section has
nonzero powers, whereas
\[
 g(s\lambda,((sd)^n),((sd)^n))
             =[z^s]\prod_{j=2}^n(1-z^j)^{-1}.
\]
This is zero at \(s=1\) and positive at every \(s\geq2\),
with genuine permanent positivity at every stretch. The determinant
closure multiplicities at those later stretches are not inferred
from this ordinary coefficient. Both examples have zero padding
exponent; their proofs do not establish a large-padding obstruction.

\paragraph{Archive corrections.}
The unsigned subset operation in cote~133, scan~019, is not the
exterior operation after linearization. The canonical opposite-symmetric
Burnside operations do linearize to exterior powers, but are not
special. Nontrivial irreducibles cannot be effective permutation
representations. The monomial Cauchy expansion in scan~028 is valid
in exact \(K_0\), while its proposed termwise effectivity is false.
The literal split quadratic group in scan~006 must be distinguished
from the modern metabolic Grothendieck--Witt group. A-001--A-017
prove the corrected statements or provide the explicit counterexamples.
The exhaustive assertion inventory is included as an appendix and TSV.

\paragraph{Positive families and saturation.}
F-004--F-005 give positive generating series for three two-row
and three hook shapes. F-002--F-003 give the rectangle/two-row
partition difference and rectangle/hook subset-sum product, with an
exact support description and a corrected reading of the cited
unimodality assertion. A-018 gives an explicit integer raising-map
kernel for every fixed-height family. O-001 classifies ordinary
saturation failures in the stated two-row-input range, while O-002
gives holes at every rectangle height. O-004 proves positivity of
every stretch at least two in the precise stable-tail range.
These statements do not convert the revised GCT~VI saturation-index
hypothesis into ordinary saturation. The normal, positive, pure-Tate
model in O-003 and the relative-lattice distinction in A-019
identify why positivity, normality and weights do not repair that
inference on their own. F-019 extracts the known complete rule for
the square of a two-row rectangle. A-024 realizes all its multiplicity
spaces as the monomial basis of an explicit polynomial highest-weight
algebra: its generated weight lattice has index four in the ordinary
even-size ambient lattice. Thus its positive basis and normality
coexist with precisely described ordinary saturation holes.
A-026 extracts a known polynomial invariant-ring theorem into the
positive count \(2a+3b+4c=h\) for three equal height-three
rectangles. Here the scalar weight lattice is primitive, while
degree one is missing and higher multiplicities are unbounded.
The two actual Kronecker algebras distinguish the two separate
mechanisms that invalidate a general saturation inference.
M-020 also records the known BIH theorem on saturation of actual
determinant support. When \(n\geq m^2\), every possible internally
padded permanent type belongs to that saturation. A short finite
reduction gives a uniform stretching integer for each fixed \(n\),
without an effective bound. Separately, known boundary localization
implies that all determinant-orbit copies extend after sufficiently
large full-rectangle translations. Those translations have \(n^2\)
rows and hence leave internally padded permanent support when \(m<n\).
Neither statement compares the original positive multiplicities.

\paragraph{Checked comparisons of actual multiplicities.}
Several arguments preserve all copies, rather than merely one
nonzero vector. Projection containment gives the ideal-space
comparison M-006. Explicit integral determinant constructions
then give row and shape bounds M-007--M-011. Cauchy flag minors
give the all-copy second-row stabilization F-014 and the nested
coefficient support F-016; restricting to those coefficient
families yields \(K_n=a_\lambda\) in their stated ranges.
The 2026 canonical Foulkes--Howe theorem is a known input, whose
complete written proof was checked here; M-010 transfers it to
the determinant closure. A-020--A-023 refine the differential
argument by tracking the partition's second row and treating powers
of a differential operator on their actual source spaces.
M-013 replaces exact containment by dominance of the coefficient
jets seen by a highest-weight function; its explicit integral
Jacobian preserves all copies even in a range where the exact
universal construction cannot work.

Here are useful uniform necessary conditions for a remaining
actual obstruction. All of them are restrictions, not an
assertion that a shape satisfying them is an obstruction.
\begin{center}
\renewcommand{\arraystretch}{1.2}
\begin{tabular}{@{}p{0.29\textwidth}p{0.29\textwidth}p{0.30\textwidth}@{}}
\toprule
Condition & \(n>3m^4\) & \(n\geq m^{25}\)\\
\midrule
Permanent size & \(m\geq19\) & \(m\geq190\)\\
Number of rows & \(\ell(\lambda)\geq10\) & \(\ell(\lambda)\geq52\)\\
Second row & \(\lambda_2\geq6\) & \(\lambda_2\geq27\)\\
Later row & \(\lambda_6\geq4\) & \(\lambda_{27}\geq14\)\\
Degree & \((d-1)^2>n\) & \((d-1)^2>n\)\\
Padding-refined degree & \(\lfloor d(d+m-1)/2\rfloor>n\) & \(\lfloor d(d+m-1)/2\rfloor>n\)\\
Permanent multiplicity & \(q^m_\lambda\geq1\) & \(q^m_\lambda\geq2\)\\
\bottomrule
\end{tabular}
\end{center}
The first four rows follow from F-009, M-008, F-014, F-017 and M-013.
The strongest uniform row restrictions are the staircase inequalities
\(\lambda_{12-2s}\geq s+1\) for \(1\leq s\leq5\) in
the IP range, and
\[
                  \lambda_{2j}\geq28-j\qquad(1\leq j\leq26)
\]
in the BIP range. M-014 gives the corresponding staircase
\(\lambda_{2j}\geq c+3-j\), \(1\leq j\leq c+1\),
for every integral \(c\geq6\) at \(n\geq m^c\).
The BIP case starts with
\(\lambda_2\geq27,\lambda_4\geq26,\lambda_6\geq25\)
and ends with \(\lambda_{52}\geq2\). The degree rows follow
from A-023 and the padding constraint
\(\lambda_2\leq|\bar\lambda|\leq md\). The final BIP
row also uses the cited no-occurrence theorem: positive
\(q^m_\lambda\) forces positive \(K_n\) in that range.
Every candidate must additionally fail the exact nonuniform
inequalities in M-009, A-023, F-015--F-018, M-011, M-013 and M-016--M-019,
and satisfy \(n<2^m-1\), \(\ell(\lambda)\leq m^2\),
\(\ell(\lambda)\leq d\), and \(|\bar\lambda|\leq md\).
The separate known inequality \(d>n/m\) is only asserted
for the stronger test \(q^m_\lambda>g\).

The strongest current shape-independent degree gate is the optimized
recurrence A-023:
\[
 n<P_d(dm),\qquad
 P_d(dm)=\binom d2+
 \max\!\left(0,\max_{2\leq j\leq d-1}
       \min\left\{\binom j2,dm-\binom j2\right\}\right).
\]
The inner maximum is checked at only the two integers bracketing
the real solution of \(j(j-1)=dm\), clipped to \(d-1\).
The empty maximum at \(d=2\) is zero. This implies the simpler
degree rows in the table and improves them in additional ranges.
Shape by shape, A-023 also gives the stronger necessary condition
\(n<P_d(\lambda_2)\), using its recurrence, for a type with a
nonzero second row.

The full nested support can improve on every individual row cut.
O-018 and A-025 prove the additional BIP necessary condition
\(\lambda_3\geq557\) or \(\lambda_{52}\geq3\).
F-018 expresses the general remaining support condition as
\(n<\kappa(\Gamma_{m,\lambda})\), where \(\kappa\) is an
explicit finite optimization over two monomial order ideals.
The prefix-extended jet construction M-017 gives a further
coupled restriction. Put \(C_*=478126755\). If
\(\lambda_2\leq26+t\), an actual BIP obstruction must have
\(\lambda_k\geq27\) whenever
\(1+\binom{k-2+t}{t}\leq C_*\). In particular,
\[
 \lambda_3\geq27\quad\text{or}\quad\lambda_2\geq478126780.
\]
At \(\lambda_2=27\), at least \(478126755\) rows are
necessary, and rows two through that index all equal twenty-seven.
This is a restriction on candidates, without asserting their existence.
The more general one-sided and two-sided order-ideal extensions
M-018--M-019 give further finite support criteria. Two convenient
consequences are
\[
 \begin{gathered}
 \lambda_2\geq159375611\ \text{or}\ \lambda_3\geq29
                         \ \text{or}\ \lambda_4\geq27,\\
 \lambda_2\geq239063403\ \text{or}\ \lambda_4\geq27.
 \end{gathered}
\]
O-025 supplies an actual \(q\geq4\) family excluded by M-019
although every M-018 parameter choice fails. This establishes a
strict additional reduction of the general gate. O-024 and O-026
give exact finite counts for these criteria; evaluating an existing
criterion is distinguished from proving a new reduction.
F-022 and M-016 also complete the second-row power rule:
for every integer \(c\geq1\), \(n\geq m^c\) requires
\(\lambda_2\geq c+2\). Equivalently, every actual obstruction
satisfies \(n<m^{\lambda_2-1}\).
The A-030 family attains \(\lambda_2=3\) at \(c=1\),
so that endpoint is sharp.

O-016, O-017 and O-019 prove separate limitations for exact
universal determinants, entire single-cut families and full jet
dominance. Those limitations do not imply an actual obstruction.

\paragraph{Examples with actual positive multiplicity.}
O-011 and O-013 determine a Hessian family exactly:
\(q=K=a=1\) within the stated permanent rank range.
M-012 gives \(q=K=a=4\) for \((nd-6,6)\) when
\(n>m\geq6\), \(d\geq6\). Multiplying its four independent
classes by nonzero Hessian sections gives higher-row families
with \(q\geq4\). These examples distinguish genuine
permanent multiplicity from a merely positive ambient coefficient.
O-020 proves \(K_n=a_\lambda\) for \(\lambda_3\leq2\)
at \(n\geq\ell(\lambda)-1\), with arbitrary second row.
F-021 proves full permanent multiplicity in its jet range and
exhibits a type with at least two actual copies beyond every
exact two-tree support cover. M-015's cyclic-star jet construction
also proves \(\lambda_2\geq6\) for actual obstructions already
at \(n\geq m^4\); this factor removal concerns that second-row
condition, not all the separate IP staircase conditions.
O-018 supplies such an example satisfying the whole uniform
staircase, with \(\lambda_2>n\), for which every enlarged
single-cut family fails containment by dimension but the full
nested support still proves \(4\leq q\leq K\).

\paragraph{What remains open in this run.}
Outside the proved conditions there is no general comparison
of the two evaluation ranks in the large-padding range, and no
proof of a universal \(q\leq K\) inequality there. The explicit
unpadded obstructions above do not resolve that remaining question. The present
results do not prove an algebraic-complexity separation or a
Boolean P-versus-NP conclusion. The dated ledger records the
strict reductions of this gate and the approaches still being
tested. Every attached archive page has now been verified against
its official facsimile; the remaining source qualifications and
mathematical questions are recorded separately.
\clearpage

\part{Archive audit and the exact Burnside-to-representation passage}

% BEGIN agents/archive/ARCHIVE_FINDINGS.tex
% Run GCT_KRON_20261120; author/checker GPT-6 Astra Pro, effort not exposed.
% Source scope updated after official cote133 facsimile verification; see A-029.
\section{Audit of the archive: conventions and the positivity gate}
The supplied sources are transcriptions, including explicit uncertain and
deleted readings. Every attached archive page has now been compared visually
with its official Montpellier facsimile: cote133 scans006,008,009,019,028;
cote2 scan017; cote12 scan041; cote43 scan037; cote119 scan217; and the
scope-exclusion pages cote4 scan007, cote16 scan038, cote79 scan051,
cote89 scan018. Nine additional cote133 pages020--027,029 were also read.
The source manifests and check events record the exact scope. An uncertain
surname on cote133 scan008 remains uncertified and is not used in any proof.
The statements below retain their explicit mathematical interpretations;
a deleted sentence or a question is not turned into an endorsed theorem.
The secondary study SA17 remains distinguished from the original pages,
and every relevant assertion, definition, question and scope exclusion is
recorded in the inventory.

For a finite group $G$, $B(G)$ is the group completion of finite $G$-sets under
disjoint union, with product induced by Cartesian product.  We distinguish
\[
 p_r(E)=[\{A\subset E:|A|=r\}],\qquad
 \sigma_r(E)=[\operatorname{Sym}^rE],\qquad
 \lambda_t^B(E)=\sigma_{-t}(E)^{-1}.
\]
Both $p_t$ and $\lambda_t^B$ are pre-$\lambda$ structures, but they are different.
The second is the conventional Burnside structure compatible with exterior
powers after linearization.  A ``special $\lambda$-ring'' means that the
universal product and iteration identities also hold (SGA 6, V, Definition
2.4, p.~320).  This is stronger than the addition identity for a pre-structure.

\subsection*{A-001. The subset reading of scan 019 contradicts its boxed identity}
\textbf{Source and status.} Cote 133, scan 019, lines 4--10 and 43--54;
SA17, lines 491--496.  Under the study's interpretation of $\mathcal P_r(E)$
as $r$-element subsets, the boxed identities are false in $B(G)$ and after
linearization.  The corrected identities are
\[
 \sigma_2(E)=p_2(E)+E,\qquad
 \lambda_2^B(E)=E^2-\sigma_2(E).
\]
\textbf{Proof.} A multiset of size two has either one repeated member or two
distinct members.  This gives the first equivariant bijection.  Taking the
coefficient of $t^2$ in $\sigma_t\lambda^B_{-t}=1$ gives the second equality.
Let $G=C_2$ and let $E$ be its free orbit.  Then $E^2=2E$ and $p_2(E)=1$.
Thus the proposed equality $2p_2(E)+E=E^2$ is $2+E=2E$.  Applying the number
of $G$-fixed points gives $2=0$, a contradiction.  Under linearization over
$\mathbb Q$, $E$ becomes $1+\varepsilon$, where $\varepsilon$ is the sign
character.  The two sides become $3+\varepsilon$ and $2+2\varepsilon$,
and are distinct.  In contrast, $\lambda_2^B(E)=E-1$ linearizes to
$\varepsilon=\bigwedge^2(1+\varepsilon)$.
The cardinality equality $2\binom d2+d=d^2$ is correct; forgetting the group
action conceals the error.  This also refutes SA17's assertion that the
Burnside identity was verified in degree two. \hfill$\square$

\paragraph{Dictionary.}
Finite $G$-sets give finite etale schemes with action of the constant group
scheme $G_{\mathbb Z}$ and permutation locally free sheaves on
$BG_{\mathbb Z}$.  The subset construction forgets the signs on an exterior
basis: a stabilizer can permute a chosen subset with odd sign.  The exterior
power is therefore not the permutation sheaf on subsets.  Constant group
schemes are constructed in SGA 3, I, 4.1, original p.~23; the representation
and comodule dictionary is I, Proposition 4.7.2, pp.~34--35.

\subsection*{A-002. Canonical Burnside operations are not special}
\textbf{Source and status.} The question $\lambda^i(xy)=?$ in cote 133,
scan 019, line 50 has a negative answer for the canonical inverse-symmetric
operations.  The general non-specialness and the signed subset formula are
known results; see Bouc--R\"okaeus, \emph{A note on the $\lambda$-structure on
the Burnside ring}, Theorem 1.1, formula (1.2).  We give a short explicit
counterexample to locate the failure relevant to the archive.

\textbf{Proof of the counterexample.}
Take $G=C_2\times C_2$, let $H_1,H_2,H_3$ be its three subgroups of order
two, put $X_i=G/H_i$, and put $E=G/1$.  Orbit counting gives
\[
 X_i^2=2X_i,\qquad X_iX_j=E\quad(i\ne j),\qquad
 \lambda_2^B(X_i)=X_i-1.
\]
If $H\leq G$, a multiset is $H$-fixed exactly when its multiplicity is
constant on each $H$-orbit.  For the mark $m_H$ this proves
\[
 m_H(\sigma_t(S))=\prod_{O\in H\backslash S}(1-t^{|O|})^{-1},\qquad
 m_H(\lambda_t^B(S))=\prod_{O\in H\backslash S}(1-(-t)^{|O|}).
\]
Consequently $\lambda_2^B(E)$ has marks $0,-2,-2,-2,6$ at
$G,H_1,H_2,H_3,1$, respectively, and
\[
 \lambda_2^B(E)=3E-X_1-X_2-X_3.
\]
The special degree-two identity would instead give
\[
 \lambda^2(X_1X_2)
 =X_1^2\lambda^2(X_2)+X_2^2\lambda^2(X_1)
       -2\lambda^2(X_1)\lambda^2(X_2)=2E-2.
\]
Their difference $E-X_1-X_2-X_3+2$ has $G$-mark $2$, hence is nonzero.
This difference linearizes to zero, illustrating why a character calculation
alone cannot verify specialness of the Burnside structure. \hfill$\square$

For completeness the subset pre-structure also fails the special
product identity. For the free two-point $C_2$-set $E$ one has
$p_2(E)=1$ and $E^2=2E$. The disjoint-union subset formula gives
$p_2(E^2)=p_2(2E)=2+E^2=2+2E$, whereas the special product
polynomial gives $E^2+E^2-2=4E-2$. Their difference has
$C_2$-fixed-point mark four. Both subset and canonical operations
do satisfy the pre-structure addition identity: subsets of a
disjoint union split by their two sizes, and the same argument
holds for multisets before inverting their total series.

\paragraph{The correct signed formula.}
Bouc--R\"okaeus, Theorem 1.1 (1.2), states
\[
 \lambda_n^B(S)=(-1)^n\sum_{\mu\vdash n}(-1)^{\ell(\mu)}
 \frac{\ell(\mu)!}{\prod_j\alpha_j(\mu)!}[\mathcal P_\mu(S)],
\]
where $\mathcal P_\mu(S)$ is the $G$-set of ordered pairwise disjoint subsets
with sizes given by the parts of $\mu$, and $\alpha_j(\mu)$ counts repeated
parts.  It is a formula in the group completion, not a positive formula.
Linearization commutes with these operations because multisets form a basis
of symmetric powers of a permutation module and the symmetric/exterior
generating-series identity holds in the representation ring.  The special
structure on the latter is SGA 6, VI, Theorem 3.3, pp.~392--393; it does not
transfer specialness backwards across a noninjective homomorphism.

\subsection*{A-003. Effective permutation lifts cannot supply nontrivial irreducibles}
\textbf{Claim.} For a field of characteristic zero and a finite group $G$,
no nontrivial irreducible representation is the linearization of an effective
finite $G$-set.  In particular, a nontrivial Specht representation has no
effective finite symmetric-group-set lift.

\textbf{Proof.} For a finite $G$-set $X$, the invariant subspace of $k[X]$
has basis the sums of the elements of each orbit.  Hence its dimension is
the number of orbits.  If $X$ is nonempty this dimension is positive.  A
nontrivial irreducible representation has no nonzero invariants; if it did,
its invariant subspace would be the whole representation and the
representation would be trivial.  The empty set linearizes to zero.
These exhaust the effective possibilities. \hfill$\square$

\paragraph{Consequence and dictionary.}
Virtual permutation resolutions, including Jacobi--Trudi constructions from
Young subgroups, can still yield exact signed Kronecker formulas.  This
argument rules out deriving a positive irreducible-basis formula merely by
treating each Specht class as an effective permutation class.  It does not
rule out an independent positive rule or a geometric realization by
cohomology.  On $BG$, the invariant global sections of a nonempty
permutation bundle are nonzero, whereas a nontrivial irreducible bundle has
none.  The group-completion step is the precise loss of effectiveness.

\subsection*{A-004. The monomial Cauchy coefficients on scan 028 need not be effective}
\textbf{Source and status.} Cote 133, scan 028, lines 16--24 asks whether
the summands in the displayed expansion are actual direct summands when
$\xi$ is effective.  Under its explicitly given monomial symmetric
coefficient convention, the answer is no, already over $\mathbb Q$.

\textbf{Proof.} In degree two the displayed expansion is
\[
 \tau^2(E)=\lambda^2(E)\rho_1^2+
             (E^2-2\lambda^2(E))\rho_2.
\]
Here $\rho_1^2$ is the regular representation of $S_2$ and $\rho_2$ is its
trivial representation; $E^2-2\lambda^2(E)=\psi^2(E)$.
Let $E$ be the two-dimensional standard representation of $S_3$, regarded
as a bundle on $BS_3$.  Its characters on the classes of the identity,
transpositions and three-cycles are $(2,0,-1)$.  Thus
\[
 \bigwedge^2E=\varepsilon,\quad
 \operatorname{Sym}^2E=1+E,\quad
 \psi^2(E)=1+E-\varepsilon.
\]
The coefficient of $\varepsilon$ in the last expression is negative, so it
is not the class of an effective representation.  In contrast the actual
$S_3\times S_2$ decomposition is
\[
 E^{\otimes2}=(1+E)\boxtimes1+\varepsilon\boxtimes\operatorname{sgn}.
\]
Expanding $\rho_1^2=1+\operatorname{sgn}$ cancels the negative coefficient
and recovers this genuine decomposition.  The monomial expansion is a
valid virtual identity in this example, but not a decomposition into
effective terms. \hfill$\square$

\paragraph{Dictionary.}
The example is a locally free sheaf on the ringed topos $BS_{3,\mathbb Q}$,
so it directly lies in the scope of scan 028.  Schur functors provide the
characteristic-zero direct-summand decomposition of a tensor power; changing
from Schur coefficients to monomial symmetric coefficients destroys
termwise effectiveness.  A monomial Cauchy identity therefore gives no
saturation statement.

\subsection*{A-005. The split quadratic Grothendieck group on scan 006 need not be special}
\textbf{Source and distinction.} Cote 133, scan 006, lines 6--10 defines
the group completion of the isometry monoid, without imposing relations
from metabolic spaces.  Lines 20--32 ask whether its exterior operations
are special when $2$ is invertible.  For this literal split definition the
answer is no.  This does not contradict the theorem for the modern
Grothendieck--Witt group: Zibrowius, \emph{Symmetric representation rings
are $\lambda$-rings}, introduction theorem, p.~1056, uses Definition 1.2,
p.~1060, which also imposes metabolic relations.

\textbf{Proof.} Work on the ringed classifying topos $BC_3$ over
$k=\mathbb F_3$.  Write $1,a,b$ for the classes of the indecomposable
modules $J_1,J_2,J_3$ on which a generator acts by one unipotent Jordan
block.  Since $k[C_3]=k[t]/(t^3)$, Jordan normal form implies that its
split representation ring is free as an abelian group on $1,a,b$.
The tensor and exterior identities in that ring are
\[
 a^2=b+1,\qquad ab=2b,\qquad b^2=3b,
 \qquad \lambda^2(a)=1,\quad\lambda^2(b)=b,\quad\lambda^3(b)=1.
\]
For completeness, on $J_2\otimes J_2$ the nilpotent operator $g-1$ has
ranks $2,1,0$ in powers one, two, three, so its blocks are $J_3+J_1$.
The module $J_3=k[C_3]$ is regular.  The map
$g\otimes v\mapsto g\otimes g^{-1}v$ identifies its tensor product with
$V$ under diagonal action with $\dim(V)$ copies of the regular module.
This proves the other tensor identities.  The determinants of unipotent
blocks are one; $J_3$ is self-dual and
$\bigwedge^2J_3\simeq J_3^*\otimes\det J_3\simeq J_3$, proving the
exterior identities.

In any pre-$\lambda$-ring the third Newton operation is
\[
 \psi^3(z)=z^3-3z\lambda^2(z)+3\lambda^3(z).
\]
It is additive, as follows by comparing coefficients of the formal
logarithmic derivative of $\lambda_t(z+w)=\lambda_t(z)\lambda_t(w)$.
The preceding identities give
\[
 \psi^3(a)=2b-2a,\qquad \psi^3(b)=3.
\]
Now take the symmetric hyperbolic modules $x=H(J_2)$ and $y=H(J_3)$.
Since $2$ is invertible, symmetric bilinear forms and quadratic forms
are equivalent.  Forgetting a form defines a ring homomorphism that
commutes with exterior powers from the split quadratic group to the
split representation ring.  The images of $x,y,xy$ are $2a,2b,8b$.
Consequently the images of $\psi^3(xy)$ and $\psi^3(x)\psi^3(y)$ are
respectively
\[
 24\cdot1\quad\hbox{and}\quad24b-24a.
\]
They are unequal in the free abelian group on $1,a,b$.  In a special
$\lambda$-ring every Newton operation is multiplicative: the universal
product identities imply this polynomial identity, which can be checked
on formal sums of line elements, where $\psi^3(\sum_i l_i)=\sum_i l_i^3$.
Thus the split quadratic group in this example is not special.
\hfill$\square$
\paragraph{Dictionary and significance.}
This is a counterexample on a classifying ringed topos, exactly the
framework of scan 006.  Quotienting the split group by metabolic relations
changes the mathematical object.  It is therefore essential to specify
the exact category and relations when importing a modern positivity or
$\lambda$-ring theorem into the manuscript's notation.  The counterexample
does not address a semisimple characteristic-zero restriction of the
split group; it refutes the stated generality with $2$ invertible.

\subsection*{A-006. Without invertibility of two, strict nonsingularity is not preserved by exterior powers}
\textbf{Source and status.} Cote 133, scan 006, lines 12--19 claims a
pre-$\lambda$ structure before assuming that $2$ is invertible and states
that ranks are even when $2$ is not invertible.  The rank assertion is
correct for a connected ringed topos with locally free modules of
constant rank; the claimed exterior-power pre-structure is not.

\textbf{Proof.} The polar form of a quadratic form has diagonal entries
$2q(e_i)$.  Modulo two its matrix is alternating.  An alternating matrix
of odd size has determinant zero over every ring, by the integral
alternating determinant identity.  Therefore the determinant of an odd
rank polar matrix belongs to the ideal $(2)$.  If it is a unit, then
$2$ is a unit.  Applied locally and then glued, this proves the even-rank
assertion whenever $2$ is not globally invertible and the rank is
constant.

Over a field of characteristic two the hyperbolic plane is a strictly
nonsingular quadratic module of rank two.  Its second exterior power
has rank one.  A rank-one quadratic module in characteristic two has
zero polar form, so cannot be strictly nonsingular.  Thus exterior
powers cannot define the claimed pre-structure on this collection of
objects.  When $2$ is invertible, exterior powers of the associated
nonsingular symmetric forms do give a pre-$\lambda$ structure, and
their ranks are binomial coefficients.  Specialness remains false for
the split group by A-005. \hfill$\square$

\subsection*{A-007. Exterior operations do not descend to the Witt quotient}
\textbf{Source and status.} Cote 133, scan 008, lines 32--38 asks whether
the pre-$\lambda$ structure descends modulo hyperbolic modules.  It does
not, even over a field of characteristic different from two.

\textbf{Proof.} The hyperbolic plane $H=\langle1,-1\rangle$ is zero in
the Witt group, while
\[
 \lambda^2(H)=\langle-1\rangle
\]
has odd rank and is nonzero by the dimension-modulo-two homomorphism.
Thus the hyperbolic ideal is not stable under exterior operations.
For $2$ invertible it is indeed an ideal: for a nonsingular symmetric
module $(M,b)$, the identification $M\simeq M^*$ induced by $b$ gives
\[
 H(E)\otimes(M,b)\simeq H(E\otimes M).
\]
Taking products with group-completed classes preserves the image of
the hyperbolic homomorphism.  Rank modulo two descends because
$\operatorname{rank}H(E)=2\operatorname{rank}E$; integral rank does
not descend, since $H(k)$ has rank two but is zero in the quotient.
\hfill$\square$

\subsection*{A-008. The general topos caveat on scan 008 is essential}
\textbf{Source and status.} Cote 133, scan 008, lines 2--8 asserts that
$q\perp(-q)$ is hyperbolic for modules over a ring, with a caveat for
general ringed topoi.  Lines 11--20 suggest a stable-hyperbolic
description of equality in the quotient.  The ring assertion holds for
finite projective modules with nonsingular polar form; the unrestricted
topos description fails in characteristic two, including its proposed
stable criterion.

\textbf{Proof over a ring.} Let $M$ be finite projective and let $b$ be
the nonsingular polar form of $q$.  The surjection
$M^*\otimes M^*\longrightarrow\operatorname{Sym}^2M^*$ splits as a
map of modules because its target is finite projective.  Hence choose
a bilinear form $c$ with $c(u,u)=q(u)$, and use $b$ to choose an
endomorphism $a$ satisfying $b(u,a(u))=q(u)$.  In the coordinates
$u=x-y,v=y$ one has
\[
 q(x)-q(y)=q(u)+b(u,v)=b(u,v+a(u)).
\]
The change of coordinates is invertible, and the last form is the
hyperbolic form after $b$ identifies the second copy of $M$ with $M^*$.
If $2$ is invertible, the choice $c=b/2$ is canonical and the argument
works on every ringed topos.  Without that hypothesis the global choice
of $c$ can fail. 

\textbf{Counterexample on a classifying topos.}
Let $k$ be algebraically closed of characteristic two, and let
$G=\mathbb G_m\rtimes C_2$, where the nontrivial element $s$ acts on
$\mathbb G_m$ by inversion.  Let $V=ke\oplus kf$ have torus weights
$1,-1$, with $s$ exchanging $e,f$, and set $q(xe+yf)=xy$.
This is a nonsingular invariant quadratic form, denoted $Q$.
In fact $G$ is its split orthogonal group: a matrix preserving $xy$
has $ac=bd=0$ and $ad+bc=1$, and hence lies in the diagonal or
anti-diagonal component, each isomorphic to $\mathbb G_m$.
Every $G$-endomorphism of $V$ is scalar, and $V$ is simple: a torus
stable line has one of the two weights and is not stable under $s$.
Thus a $G$-linear copy of $V$ in $V\oplus V$ has the form
$v\mapsto(av,bv)$, and the restriction of $q\perp q$ is
$(a+b)^2q$.  The only isotropic copy is the diagonal.  In contrast a
hyperbolic module has two complementary invariant isotropic submodules.
If $H(E)$ had underlying representation $V\oplus V$, semisimplicity
of this particular representation forces $E\simeq V$.  Hence
$Q\perp(-Q)=Q\perp Q$ is not hyperbolic.

We now prove failure even of the displayed stable criterion.  In any
$G$-representation the weight pair $\{1,-1\}$ is a canonical direct
summand, and an invariant nonsingular form restricts to a nonsingular
form there.  Write this summand as $V\otimes k^r$.  Torus invariance
forces its quadratic form to be $x^TAy$, where $A$ is an invertible
$r\times r$ matrix; $s$-invariance forces $A$ to be symmetric.  Changes
of equivariant coordinates act on $A$ by congruence.  The restriction
of a hyperbolic module $H(E)$ to this weight pair gives an alternating
matrix: if the multiplicity of $V$ in this pair of $E$ is $r$, this
matrix is congruent to
$\left(\begin{smallmatrix}0&I_r\\I_r&0\end{smallmatrix}\right)$.
The matrix for $2Q$ is $I_2$, which is nonalternating.  Adding an
alternating matrix preserves nonalternation.  Consequently no
\[
 2Q\perp H(F)\simeq H(E)
\]
can exist, for any representations $E,F$.

Nevertheless $3Q\simeq Q\perp H(V)$.  Indeed, for the bilinear form
with matrix $I_3$, the basis
\[
 (1,1,1),\quad(1,1,0),\quad(0,1,1)
\]
has Gram matrix $1\oplus
\left(\begin{smallmatrix}0&1\\1&0\end{smallmatrix}\right)$ in
characteristic two.  Thus in the group completion
$2[Q]=[H(V)]$, so $2Q$ is zero in the quotient by the hyperbolic image,
although it is not stably hyperbolic in the displayed sense.
This disproves the unrestricted criterion on scan 008.  Over a field
with no group action, the ordinary Witt cancellation statement quoted
there is compatible with the ring argument; the counterexample concerns
equivariant objects on a classifying topos. \hfill$\square$

\subsection*{A-009. Correct interpretation of the Stiefel--Whitney statements}
\textbf{Source.} Cote 133, scan 009, lines 2--39; SA17, lines 574--575.
The existence, naturality, normalization on lines, and orthogonal-sum
formula are known statements in the etale setting for schemes over
$\mathbb Z[1/2]$. Exact primary source: Esnault--Kahn--Viehweg,
\emph{Coverings with odd ramification and Stiefel--Whitney classes},
\S5.1--5.3, pp.~179--182. The hyperbolic formula is the Proposition in
\S5.5, pp.~183--184. Thus the formula called ``new in this reading''
by SA17 was already in that 1993 source.

The target of the total class on group-completed objects is
$1+\prod_{i>0}H^i_{\mathrm{et}}(X,\mathbb F_2)$, with multiplication of
series. A virtual object can have infinitely many nonzero classes,
because total classes of negatives are formal inverses. The degree-zero
component is one. For ordinary bundles the classes above the rank vanish.
The crossed-out assertion $w\circ H=1$ is false. The cited formula is
\[
 w(H(E))=\sum_{i=0}^{r}(1+\alpha)^{r-i}\bar c_i(E),
 \qquad \alpha=(-1)\in H^1_{\mathrm{et}}(X,\mathbb F_2),
\]
where $\bar c_i$ is the Chern class reduced modulo two in degree $2i$.
For a line bundle, $w_2(H(L))=\bar c_1(L)$. For instance on
$\mathbb P^1_{\mathbb C}$ with $L=\mathcal O(1)$ this is nonzero.

\textbf{Verification of the universal formulas used in the archive.}
Use the cohomologically injective orthogonal splitting principle
(Zibrowius, \emph{The gamma-filtration on the Witt ring of a scheme},
Proposition 3.3, pp.~19--20 of arXiv:1411.3198). If the first
Stiefel--Whitney roots of $M,N$ are $u_i,v_j$, respectively, line-bundle
normalization and the tensor product of the corresponding double-cover
torsors give
\[
 w(M\otimes N)=\prod_{i,j}(1+u_i+v_j),\qquad
 w(\bigwedge^sM)=\prod_{|I|=s}\left(1+\sum_{i\in I}u_i\right).
\]
These are symmetric polynomials in each root alphabet. The fundamental
theorem on symmetric polynomials expresses them universally in the
elementary symmetric polynomials, namely the $w_i$, and the numbers of
roots, namely the ranks. Injectivity of the splitting pullback proves
the original identities. The target of the asserted $\lambda$
homomorphism must be endowed with these Chern-root operations, not with
componentwise cup product as its ring product. The manuscript's
shorthand about an arbitrary strictly local ringed topos does not state
enough topology or splitting hypotheses to identify an equally broad
theorem; the verified statement here is the etale scheme formulation.

\subsection*{A-010. An integral proof of the monomial Cauchy identity}
\textbf{Source and conventions.} Cote 133, scan 028, lines 2--18 and
26--48. Let $K_0(X)$ mean the exact Grothendieck group of locally free
modules, with relations from short exact sequences; use bounded finite
ranks initially; A-017 below removes the boundedness restriction.
On $\bigoplus_{n\geq0}K_0(B_{S_nX})$ use induction along
$S_i\times S_j\subset S_{i+j}$ for the graded product. With these
conventions the displayed identity is integral; it does not require
division by $n!$. On an arbitrary ringed topos, ``bounded finite ranks''
includes all constant-rank modules and their finite virtual combinations.
This is enough for every representation considered in this work.

\textbf{Proof.} Tensor powers of a short exact sequence of locally
free modules have a filtration, preserved by $S_n$, whose associated
graded is the tensor power of the direct sum of the submodule and
quotient. Filter by the number of tensor factors lying in the submodule.
Its quotients are direct sums indexed by subsets of the tensor positions
and are locally free. Consequently
\[
 \tau_t([E]+[F])=\tau_t([E])\tau_t([F]),
 \quad \tau_n([E])=[E^{\otimes n}],
\]
where the right-hand product uses induction. This defines $\tau_t$
on $K_0(X)$ by group completion, with values in units of the completed
graded ring.

For an actual module of constant rank $d$, check the identity
universally on the locally ringed classifying topos of
$GL_{d,\mathbb Z}\times S_n$. The general bundle and its tensor
permutation action are obtained from this universal bundle by its frame
torsor and base change. Pull back to the full flag bundle of the
$GL_d$ bundle. Each successive projective-bundle pullback on $K_0$
is injective by SGA 6, VI, Theorem 1.1, and its locally ringed topos
extension, Remark 1.14, p.~380; the universal classifying topos has
quasi-compact final object. Apply this with the $S_n$ factor present,
so injectivity holds in the equivariant $K_0$ of the desired identity.

The pulled-back bundle has a filtration with line quotients
$L_1,\ldots,L_d$. Exterior-power and tensor-power filtrations permit
replacement by $\bigoplus_iL_i$. For this direct sum expand the
tensor product by its $n$ tensor positions. A multiplicity list
$(\nu_1,\ldots,\nu_d)$ with sum $n$ contributes
\[
 L_1^{\nu_1}\cdots L_d^{\nu_d}\otimes
 \operatorname{Ind}_{S_{\nu_1}\times\cdots\times S_{\nu_d}}^{S_n}1.
\]
The stabilizer acts trivially on tensor powers of a line. If
$\alpha_i$ counts indices with $\nu_j=i>0$, the induction class is
$\rho_1^{\alpha_1}\cdots\rho_n^{\alpha_n}$. Summing the distinct
multiplicity lists of this type gives the monomial symmetric function
$m_\mu(L_1,\ldots,L_d)$, where $\mu$ has $\alpha_i$ parts equal
to $i$. The fundamental theorem on symmetric polynomials writes this
as the integral polynomial $\Sigma_\alpha$ in the elementary classes
$\lambda^i(E)$. Injectivity of the flag pullback and then pullback
from the universal bundle prove the formula for $E$.

The right-hand side as a total series is also exponential under
addition. For finite alphabets of formal line elements its expansion is
\[
 \prod_j\left(\sum_{r\geq0}L_j^r\rho_r\right),\qquad\rho_0=1,
\]
so union of alphabets is multiplication. Degree by degree this is a
polynomial identity in elementary symmetric functions over $\mathbb Z$,
and hence holds for arbitrary pre-$\lambda$ inputs satisfying the
addition rule. Both sides extend identically to virtual classes. Thus
\[
 \tau_n(\xi)=\sum_{\sum i\alpha_i=n}
 \Sigma_\alpha(\lambda^1\xi,\ldots,\lambda^n\xi)
              \rho_1^{\alpha_1}\cdots\rho_n^{\alpha_n}.
\]
Bounded nonconstant rank is handled by its finite open-and-closed rank
decomposition. A-017 below justifies the infinite rank decomposition
as well, so the conclusion extends to all locally finite ranks.
\hfill$\square$

The universal polynomial $\Sigma_\alpha$ is homogeneous for weights
$\deg\lambda^i=i$. Its extreme terms are $m_{(1^n)}=e_n$ and
$m_{(n)}=p_n$; the coefficients of the formal monomials $\rho_1^n$
and $\rho_n$ are therefore $\lambda^n$ and the $n$th Newton operation.
If the target has relations among the $\rho_i$, these are coefficients
of the universal polynomial expression, not uniquely recoverable
coefficients of the target element. A-004 refutes the separate question
about effective direct summands.

\subsection*{A-011. Injection powers and the corrected fiber-type decomposition}
\textbf{Source status.} SA17, lines 498--515, describes cote 133 scans
20--23, which were not attached but have now been retrieved from the
official Montpellier facsimile and visually read (A-029). The
reported statements and the deletion on scan022 are verified against
those originals. For a finite $G$-set $E$, set
$\tau_n(E)=\operatorname{Inj}([n],E)$ and
$T_n(E)=\operatorname{Map}([n],E)$, with commuting $G$ and $S_n$ actions.
Then
\[
 \tau_n(E\amalg F)=\coprod_{i+j=n}
 \operatorname{Ind}_{G\times S_i\times S_j}^{G\times S_n}
       (\tau_i(E)\times\tau_j(F)).
\]
Indeed the inverse image of $E$ is an $i$-element subset of $[n]$;
choosing it is the coset space $S_n/(S_i\times S_j)$, and restrictions
of the injection give the remaining factors. This proves the
exponential property and its extension to $B(G)$ by group completion.

For a fiber type $\alpha=(\alpha_1,\ldots,\alpha_n)$ put
$m=\sum\alpha_i$, $n=\sum i\alpha_i$, and
$W_\alpha=\prod_i(S_i\wr S_{\alpha_i})$. The canonical surjection
$W_\alpha\to\prod_iS_{\alpha_i}\subset S_m$ permutes blocks of
each size. The summand of $T_n(E)$ of this type is
\[
 \operatorname{Ind}_{G\times W_\alpha}^{G\times S_n}
 \operatorname{Infl}_{G\times\prod_iS_{\alpha_i}}^{G\times W_\alpha}
 \operatorname{Res}_{G\times\prod_iS_{\alpha_i}}^{G\times S_m}
 \tau_m(E).
\]
\textbf{Proof.} A map determines its partition into nonempty fibers.
The stabilizer of a chosen partition of type $\alpha$ is $W_\alpha$.
Labeling its $m$ blocks by distinct elements of $E$ is an injection
from the block set, and the stabilizer acts through its permutation
quotient. Inducing over the orbit of the partition proves the formula.
The number of partitions of type $\alpha$ is
$n!/\prod_i\alpha_i!(i!)^{\alpha_i}$. Already for $n=2,m=1$,
induction from $S_1$ is wrong: diagonal maps have trivial $S_2$ action,
whereas induction from $S_1$ gives a free action. Thus the deletion
reported by SA17 is justified. \hfill$\square$

The operators $q_H^n$ were only named in the supplied study.
The original manuscript scan021 has now been retrieved and visually
read; A-029 gives its quotient definition, full extension proof and
precise polynomial-degree qualification. This resolves that earlier
source-availability gate.

\subsection*{A-012. Correcting the reported finite-injection-functor equivalence}
\textbf{Source status.} SA17, lines 517--532, describes cote133
scans24--27. These originally unattached pages have now been
visually read in the official facsimile (A-029); scan025 confirms
the single-map-functor assertion below, even for a sufficiently
large indexing set. Put
$h_I(E)=\operatorname{Inj}(I,E)$ and $H_J(E)=\operatorname{Map}(J,E)$
on finite sets and injections. The assertion that finite sums of
$h_I$ are equivalently direct summands of a single $H_J$ is false
as written.

\textbf{Proof and verification of the closure assertions.}
The functor $h_1\amalg h_1$ has cardinality two at a singleton,
whereas every $H_J$ has cardinality one there. It is neither a
coproduct summand nor a retract of a single $H_J$. The corrected
statement uses a finite sum of the $H_J$: each $h_I$ is a component
of $H_I$, consisting of its injective maps.

The category of elements of $h_I$ has initial object
$(I,\operatorname{id}_I)$, proving connectedness. The components
of $H_J$ are indexed by kernel equivalence relations on $J$; the
component with quotient $I$ is a copy of $h_I$. Products of
representables split into finitely many representables, according
to overlap equivalence relations of the two injection images.
By Yoneda, a natural map from $h_I$ to a finite sum of representables
is determined by a target component and an injection from its
representing set to $I$. The equalizer of two such maps is all
of $h_I$ if the choices agree, and empty otherwise. Thus finite
sums of representables are closed under finite limits, including
the terminal $h_\varnothing$.

An element of $h_I\circ h_J$ is an injective list indexed by $I$
of injections $J\to E$. Its equality pattern on $I\times J$ has
finitely many possibilities. Admissible patterns have no equalities
within a row and no two identical rows. Each pattern gives one
representable, by injecting its quotient set into $E$. Finite sums
in either argument are handled by the finitely many choices of
target component for each row. This proves closure under composition.

Finite sums of $h_I/H$ correspond to expressions
$\coprod_n h_{[n]}\times_{S_n}X_n$ with each $X_n$ finite and with
\emph{finite degree support}: decompose $X_n$ into $S_n$-orbits.
Dropping finite degree support would no longer describe a finite
sum. Eventual bijectivity satisfies two-out-of-three by taking the
maximum of the two thresholds; pullback preserves it pointwise.
Finally functors with injective transition maps are closed under
sums, products and composition, since each constituent transition
in the composite is an injection. These verify the substantive
closure assertions SA17 had grouped as ``definitions''.
\hfill$\square$

% The following block audits the separately reported scan 029 question.

\subsection*{A-013. Polynomial induction rings: fields, local rational algebras, lattices, and a singular counterexample}
\textbf{Source status.} SA17, lines 534--556, reports a question on
cote133 scan029, which was not attached but has now been visually
read in the official facsimile (A-029). Its local-ring question
is confirmed. In this paragraph $R_A(S_n)$
always means exact $K_0$ of representations finite projective over $A$.
For a local ring these underlying modules are finite free. The product
across $n$ is induction; $\rho_n$ is the trivial representation in degree
$n$. This is different from the split group used in A-005.

Over a characteristic-zero field the induction ring is the symmetric
function ring $\mathbb Z[h_1,h_2,\ldots]$, with $\rho_i$ corresponding to
$h_i$. This is the classical Frobenius characteristic theorem (Macdonald, second edition, I, (7.3), p.~113, and I, (2.8), p.~22). The same
conclusion holds for every local $\mathbb Q$-algebra $A$: the split
algebra $\mathbb Q[S_n]=\prod_{\lambda\vdash n}\operatorname{Mat}_{f^\lambda}(\mathbb Q)$
base-changes to the analogous matrix-algebra product over $A$.
Morita equivalence identifies the representation exact category with a
product of finite projective $A$-module categories. Such modules are
free because $A$ is local. The rational Specht bases, and the induction
matrices between them, therefore base-change unchanged.

For $A=\mathbb F_2$, the regular $S_2$ module has a trivial submodule
and trivial quotient, giving $\rho_1^2=2\rho_2$. Thus the polynomial
assertion fails. More generally $R_{\mathbb F_p}(S_p)$ has rank
$p(p)-1$, since simples are indexed by $p$-regular partitions and
$(1^p)$ is the unique excluded partition of $p$ (James, \emph{The
Representation Theory of the Symmetric Groups}, Theorem 11.5, p.~40 (with Definition 11.2, p.~39)). A polynomial ring with one generator in each
positive degree would have rank $p(p)$ in this degree.

Characteristic zero alone is also insufficient for local rings.
Let
\[
 A=\mathbb Z_{(2)}[\epsilon]/(2\epsilon,\epsilon^2).
\]
This is a local ring of characteristic zero, with nonzero $\epsilon$.
The unit $1+\epsilon$ has square one, so defines a rank-one free
$A[S_2]$ module $L$ on which the transposition acts by $1+\epsilon$.
Trace of this transposition is additive in exact sequences of
$A$-free representations, since these sequences split as $A$-modules.
It therefore defines a homomorphism $R_A(S_2)\to A$. Its values on
$\rho_1^2$ and $\rho_2$ are $0$ and $1$, while its value on $L$ is
$1+\epsilon$. The latter is not in the integer subgroup of $A$.
Thus $[L]$ is not in the degree-two image of
$\mathbb Z[\rho_1,\rho_2,\ldots]$, proving failure of surjectivity.

\paragraph{The specific $\mathbb Z_p$-lattice question has an affirmative answer.}
Let $O=\mathbb Z_p$, $K=\mathbb Q_p$, and $R=O[S_n]$. The generic-fiber
map
\[
 K_0(\text{$O$-free finite $R$-modules})\longrightarrow R_K(S_n)
\]
is an isomorphism. Here is a proof spelling out the torsion issue.
Let $G_0(R)$ use all finitely generated $R$-modules. Any such module
$M$ has a finite free $R$-module cover $P\to M$ whose kernel $L$ is
$O$-torsion-free, hence $O$-free of finite rank. Set
$\chi(M)=[P]-[L]$ in the lattice $K_0$. Schanuel's lemma shows
independence of the cover. For a short exact sequence, choose free
covers of the outside terms and lift the quotient cover through the
middle term. The resulting direct-sum cover has kernels in a short
exact sequence; this proves additivity of $\chi$. On a lattice it
is its original class. Hence the lattice $K_0$ equals $G_0(R)$.

Every $O$-torsion $R$-module has a filtration by powers of $p$, with
quotients that are finite $\mathbb F_p[S_n]$ modules. The classes of
these modular modules are integer combinations of reductions of
Specht lattices: the submatrix indexed by $p$-regular partitions is
unitriangular (James, Corollary 12.3, p.~43; compatibility with
integral reduction is Corollary 8.11, p.~31). For each Specht lattice
$S_O^\lambda$, the exact sequence
\[
 0\longrightarrow S_O^\lambda\xrightarrow{p}S_O^\lambda
 \longrightarrow S_O^\lambda/pS_O^\lambda\longrightarrow0
\]
shows that its reduction is zero in $G_0(R)$. Therefore all torsion
module classes are zero there.

Any $K[S_n]$ module has an $S_n$-stable $O$-lattice, obtained by
summing translates of an arbitrary full lattice, proving surjectivity.
If two lattice classes have the same generic class, semisimplicity
over $K$ makes their generic representations isomorphic. Multiplying
such an isomorphism by a power of $p$ embeds one lattice into the
other with torsion cokernel. That cokernel is zero in $G_0(R)$, so
the lattice classes agree. This proves injectivity. Generic fiber
commutes with induction, hence
\[
 \bigoplus_n R_{\mathbb Z_p}(S_n)\simeq
 \bigoplus_n R_{\mathbb Q_p}(S_n)=\mathbb Z[\rho_1,\rho_2,\ldots].
\]
The argument is specific to these exact lattice categories and the
surjective decomposition map for symmetric groups. It does not extend
to arbitrary singular local rings, as the explicit $A$ above shows.
A classification of all local coefficient rings is not asserted.


\subsection*{A-014. The big-Witt calculations of cote 43}
\textbf{Source.} Cote 43, scan 037, lines 3--50. These are statements
about the big Witt ring in exterior-series convention. They do not
describe subset operations in a Burnside ring. With this convention
the calculations are valid, with continuity and naturality understood
in the final uniqueness assertion.

\textbf{Verification.} Put $W^+(A)=1+tA[[t]]$, with addition given
by ordinary multiplication of series. On formal roots its other
product is
\[
 \prod_i(1+x_it)*\prod_j(1+y_jt)=\prod_{i,j}(1+x_iy_jt).
\]
Each coefficient is symmetric in both alphabets and therefore is
an integral polynomial in their elementary symmetric functions.
These polynomials define the product for every ring and give an
affine ring scheme over $\mathbb Z$ with underlying scheme
$\operatorname{Spec}\mathbb Z[a_1,a_2,\ldots]$. Associativity and
distributivity are identities on formal roots, hence polynomial
identities over $\mathbb Z$. The standard special $\lambda$
structure is that of SGA 6, V, 2.3--2.5, especially Example 2.5(b), pp.~321--322; the line identity is
\[
 (1+at)*f(t)=f(at).
\]
If $A$ has characteristic $p$, coefficientwise Frobenius
$F(\sum a_it^i)=\sum a_i^pt^i$ is a ring endomorphism because
the operations are integral polynomials. The map $V(f)(t)=f(t^p)$
is injective and additive for Witt addition. The ordinary series
identity $f(t)^p=F(f)(t^p)$ proves
\[
 FV=VF=p
\]
as endomorphisms of the additive Witt group. If $x=a^{p^i}$, then
$1+xt^{p^i}=(1+at)^{p^i}$ and distributivity plus the line identity
give
\[
 (1+xt^{p^i})*f(t)=F^i(f)(xt^{p^i}).
\]
For arbitrary $x$, adjoin a root of $T^{p^i}-x$. This is a free,
faithfully flat extension, so the identity descends. This verifies
the characteristic-$p$ formulas without a perfectness assumption.

For the last map to $p$-typical Witt vectors, convert the exterior
series to the usual convention by $f(t)\mapsto f(-t)^{-1}$ and
take the usual $p$-typical projection, retaining the ghost coordinates
indexed by $1,p,p^2,\ldots$. Equivalently, on a finite formal root
alphabet set
\[
 \psi\left(\prod_i(1+x_it)\right)=\sum_i[x_i]\in W_p(A).
\]
The coordinates of the right side are symmetric integral
polynomials in the $x_i$, using the defining integral addition
polynomials of the Witt functor; its $j$th coordinate has degree
$p^j$, so depends on only finitely many coefficients of the input
series. This constructs a continuous natural map on all series.
The ghost identities show that it respects both ring operations
over the universal torsion-free polynomial ring, and hence, as
integral polynomial identities, over every ring. It sends
$1+at$ to the Teichmuller vector $(a,0,\ldots)$. A continuous
natural homomorphism with these line values is unique: split
universal finite polynomials into formal roots, use additivity,
and pass coefficientwise to formal series. Without this natural
continuous interpretation, the assertion is not a claim that the
line elements generate $W^+(A)$ as an abstract abelian group.
\hfill$\square$

\subsection*{A-015. Tensor functor identities and the finite-operator viewpoint}
\textbf{Source.} Cote 2, scan 017, lines 2--38; cote 12, scan 041,
lines 22--24; cote 119, scan 217, lines 25--31.
For group representations on finite locally free modules,
restriction preserves tensor products and exterior powers. External
tensor product is the tensor of the two pullbacks; pullback by the
diagonal gives the internal tensor product. These identities follow
directly from the actions on elementary tensors. Restrictions and
inductions are transitive, by the associativity isomorphism for
tensoring bimodules.

The projection formula for induction along $u:G\to H$ is the explicit
isomorphism
\[
 kH\otimes_{kG}(M\otimes\operatorname{Res}N)
 \longrightarrow(kH\otimes_{kG}M)\otimes N,
 \qquad h\otimes(m\otimes n)\longmapsto(h\otimes m)\otimes hn.
\]
Its inverse sends $(h\otimes m)\otimes n$ to
$h\otimes(m\otimes h^{-1}n)$; both maps respect the balancing
relations and the $H$ action. For finite groups these constructions
preserve finite locally free underlying modules in the usual
subgroup-induction setting. If $u_*$ means right-adjoint
coinduction instead, the projection formula holds when $N$ is
dualizable, by tensor-Hom adjunction. It is not asserted here for
arbitrary sheaves without this hypothesis.

The statement on cote 12 is a programme to determine $\lambda$
operations using finite operator groups, not a positivity theorem.
The historical remark in cote 119 supplies no further mathematical
identity. The precise familiar dictionary is that, if $r!$ is
invertible, $\bigwedge^r E$ is the image of the antisymmetrizer
$r!^{-1}\sum_{\pi\in S_r}\operatorname{sgn}(\pi)\pi$ on
$E^{\otimes r}$. Over $\mathbb Z$ the exterior quotient and the
exterior filtrations still exist, but this averaging idempotent
need not. SGA 6, V, Example 2.2(b) and Lemma 2.2.1, pp.~315--316,
give the exact-category construction. None of these formal
identities alone provides positive irreducible multiplicities.
\hfill$\square$

For the classifying-topos reduction on cote 133 scan 006, lines
34--45, add the local-standard-form hypothesis explicitly. If
$q_m,q_n$ are the chosen standard forms and the given forms are
locally isometric to them, their sheaves of isometries are
$O(q_m)$ and $O(q_n)$ torsors. The associated-bundle functor from
representations of their product preserves orthogonal sums,
tensor products, and exterior powers. It therefore pulls any
universal identity for the standard forms back to the given pair;
one factor suffices for an iteration identity. In the usual etale
setting with two invertible, nonsingular symmetric forms are
locally standard. An arbitrary ringed topos need not already have
that topology, so this step is conditional on the stated local
triviality, or must be formulated on the universal stack of
quadratic forms. A-005 still refutes specialness for the split
group in the archive's generality.

\subsection*{A-016. What the supplied Borger source does and does not transfer}
In the supplied \texttt{lrfoe13.tex}, \S1.1 (lines 774--807)
describes $\Lambda$-structures and, for flat schemes over
$\mathbb Z$, commuting Frobenius lifts. These are special
$\lambda$ structures. A-002 therefore blocks an automatic transfer
from the canonical Burnside pre-structure to that geometry.
\S2.5 (lines 956--977) gives the Chebychev example
$\psi_2(x)=x^2-2$ on the representation ring of $SL_2$;
the existence of a $\Lambda$-structure does not assert preservation
of an arbitrary coefficientwise positive cone. Its cohomological
theorem in \S6, labelled
\texttt{thm:lambda-varieties-are-abelian} (lines 1863--1910), and
the following corollary
\texttt{cor:lambda-varieties-are-Artin-Tate} (1913--1944), require
an actual finite-type $\Lambda$-scheme. They supply no such scheme
with fibers or cohomology equal to the required Kronecker
multiplicity spaces. That missing realization, including its
grading and a proof of support saturation, cannot be replaced
by identifying the coefficient ring as a $\lambda$-ring.

The section titled ``Relation to finite sets'', starting at line
2682, occurs after the first \verb|\end{document}| at line 2680.
It is inactive source material in this supplied draft; it must
not be cited as a numbered theorem of the compiled paper. The
source also describes itself as preliminary at lines 441--443.
These observations restrict the claimed dictionary; they do not
refute Borger's stated theorems or assert a general impossibility
of a future geometric Kronecker rule.

\subsection*{A-017. Removing the global rank bound in A-010}
\textbf{Claim.} If a locally ringed topos decomposes into open-and-closed
rank strata, the exact $K_0$ of its category of locally free modules
of finite local rank is the product of the corresponding exact
$K_0$ groups of the strata. Consequently A-010 holds for virtual
locally free classes without a global rank bound.

\textbf{Proof.} We prove the elementary exact-category statement
\[
 K_0\left(\prod_i\mathcal C_i\right)\simeq\prod_iK_0(\mathcal C_i).
\]
Every element of an exact $K_0$ is a difference of two object classes;
choosing these coordinatewise proves surjectivity. For injectivity
write an element as $[A]-[B]$ and suppose every coordinate is zero.
In the split group completion of $\mathcal C_i$, the difference
$[A_i]-[B_i]$ is a finite signed sum of short-exact-sequence
relations. Combine all positive relations by direct sum into
one sequence $U_i\to V_i\to W_i$, and all negative relations into
one sequence $U_i'\to V_i'\to W_i'$. Thus
\[
 [A_i]-[B_i]=[V_i]-[U_i]-[W_i]-[V_i']+[U_i']+[W_i']
\]
in the split group completion. By the defining equivalence of a
commutative-monoid group completion there is an object $D_i$ and
an isomorphism
\[
 A_i\oplus U_i\oplus W_i\oplus V_i'\oplus D_i
 \simeq B_i\oplus V_i\oplus U_i'\oplus W_i'\oplus D_i.
\]
The coordinate tuples form objects and exact sequences in the
product category. Their isomorphism, together with its two exact
relations, proves $[A]=[B]$ there. No uniform bound on the number
of original relations is needed after folding them into two
sequences. This proves the product assertion.

Now write a virtual class as $[E]-[F]$ and decompose the topos by
the pair of locally constant ranks of $E,F$. Modules on this
disjoint decomposition glue coordinatewise, so the exact category
of locally free modules of finite local rank is the product of
the categories on the strata. The same holds after adjoining the
finite $S_n$ action. The two sides of A-010 agree on each stratum
by its constant-rank proof, hence agree globally by the product
assertion. This establishes the claimed extension. \hfill$\square$

\subsection*{A-018. A positive kernel model for every fixed-height Kronecker family}
\textbf{Statement.} Let $\lambda,\mu,\nu\vdash N$ have at most
$p,q,r$ parts, respectively, and pad them by zeros to these lengths.
Let $\mathcal T(\lambda,\mu,\nu)$ be the set of nonnegative integer
$p\times q\times r$ arrays $a=(a_{ijk})$ with one-dimensional margins
$\lambda,\mu,\nu$. For $1\leq i<p$ define the integer matrix
\[
 D_i^{(1)} e_a=\sum_{j=1}^q\sum_{k=1}^r a_{i+1,j,k}
 e_{a-e_{i+1,j,k}+e_{i,j,k}},
\]
from the vector space with basis $\mathcal T(\lambda,\mu,\nu)$ to
that with margins $(\lambda+\epsilon_i-\epsilon_{i+1},\mu,\nu)$.
Terms with zero coefficient are omitted. Define $D_j^{(2)}$ and
$D_k^{(3)}$ by the identical rule on the other axes. The target
margins need not be partitions; an impossible margin has zero
target space. If $D$ is the vertical stack of all these matrices,
then
\[
 \boxed{\quad g(\lambda,\mu,\nu)=\dim_{\mathbb Q}\ker D
 =|\mathcal T(\lambda,\mu,\nu)|-\operatorname{rank}_{\mathbb Q}D.\quad}
\]
This gives explicit multiplicity spaces and an exact algorithm
polynomial in $N$ for fixed $p,q,r$. It does not assert polynomial
complexity in the binary encoding length of $N$, or when heights
vary. The binary-height instance is F-001; this statement includes
all fixed-height triples and fixed-height rectangular triples.

\textbf{Proof.} Take $U=\mathbb Q^p$, $V=\mathbb Q^q$,
$W=\mathbb Q^r$. In $\operatorname{Sym}^N(U\otimes V\otimes W)$
the monomials $\prod x_{ijk}^{a_{ijk}}$ form a basis, and the three
diagonal tori have weight equal to the three margins. Thus the
weight $(\lambda,\mu,\nu)$ space has exactly the indicated basis.
The elementary upper-triangular Lie algebra operator
$E_{i,i+1}$ on $U$ acts as
$\sum_{j,k}x_{ijk}\partial/\partial x_{i+1,j,k}$, giving precisely
$D_i^{(1)}$. The same computation gives the other two sets of
matrices. Simple positive root operators generate all positive
root operators by commutators. In characteristic zero their
common kernel is the subspace fixed by the upper unipotent group:
on a finite-dimensional polynomial representation each root
subgroup acts as the finite exponential of its nilpotent
derivative.

For clarity the multiplicity calculation is a finite-group
calculation. Reordering tensor factors and taking diagonal
$S_N$ invariants gives
\[
 \operatorname{Sym}^N(U\otimes V\otimes W)
 \simeq (U^{\otimes N}\otimes V^{\otimes N}
                  \otimes W^{\otimes N})^{S_N}.
\]
Characteristic-zero averaging identifies symmetric coinvariants
with these invariants. Apply Schur--Weyl decomposition separately
to the three tensor factors. This yields
\[
 \bigoplus_{\alpha,\beta,\gamma\vdash N}
 S_\alpha U\otimes S_\beta V\otimes S_\gamma W\otimes
 ([\alpha]\otimes[\beta]\otimes[\gamma])^{S_N}.
\]
Terms whose partition has too many rows vanish. Symmetric-group
irreducibles are self-dual, so the last space has dimension
$g(\alpha,\beta,\gamma)$. In the irreducible polynomial
$GL(U)$ module $S_\alpha U$, the simultaneous positive-root
kernel is its one-dimensional highest-weight line of weight
$\alpha$. One may check uniqueness directly: order products of
matrix Lie algebra generators using their commutator relations,
with lowering operators first, diagonal operators next, and
raising operators last. A weight vector killed by all raising
operators generates only weights below its own. Irreducibility
then forces its weight to be the highest weight $\alpha$, whose
weight space has dimension one. The same applies to $V,W$.
Consequently the joint kernel in weight $(\lambda,\mu,\nu)$ is
exactly the corresponding invariant multiplicity space. The
integer matrices have the same rank over $\mathbb Q$ and
$\mathbb C$, proving the formula. Schur--Weyl here is the standard
characteristic-zero equivalence in Macdonald, Chapter I,
Appendix A, (5.3)--(5.6), pp.~155--156.

There are
$D_0=pqr-p-q-r+2$ free table entries: after choosing all cells
except $(1,1,1)$, $(i,1,1)$ for $i>1$, $(1,j,1)$ for $j>1$,
and $(1,1,k)$ for $k>1$, the margins determine those remaining
cells uniquely. Each chosen entry lies between zero and $N$;
nonnegative completion can be checked directly. Thus every
source and target table set has at most $(N+1)^{D_0}$ elements.
For fixed $p,q,r$ the number of targets and the number of nonzero
entries per column are bounded independently of $N$, and matrix
entries have magnitude at most $N$. Fraction-free integer
elimination has polynomial bit complexity in these dimensions
and in $\log(N+1)$: intermediate determinants have bit lengths
bounded by the Hadamard bound, polynomial in the matrix dimension
and entry bit length. This proves the stated polynomial-in-$N$
algorithm, including zero and degenerate margins. \hfill$\square$

\textbf{Geometric and positivity scope.} Equivalently this is the
weight $(\lambda,\mu,\nu)$ part of
$H^0(\mathfrak u_p\oplus\mathfrak u_q\oplus\mathfrak u_r,
\operatorname{Sym}^N(U\otimes V\otimes W))$.
All spaces and matrices are explicitly defined over $\mathbb Z$,
using the coordinate affine scheme and its weight grading; the
diagonalizable-group/graded-module dictionary is SGA 3, I,
Proposition 4.7.3, pp.~35--36. The equality with ordinary
Kronecker coefficients is asserted only over characteristic zero.
This is a constructive positive vector-space realization and
finite rank algorithm, not a cancellation-free integer-point
count, a saturation theorem, or a uniformly efficient rule when
the determinant size and rectangle heights grow.

\paragraph{A direct finite-$G$-set construction of the same kernel.}
This model can be placed before linearization. Let $G=S_N$ and
let $X_\lambda$ be the finite $G$-set of ordered decompositions
$[N]=B_1\sqcup\cdots\sqcup B_p$ with $|B_i|=\lambda_i$.
For each $i<p$, the correspondence which moves one element
from $B_{i+1}$ to $B_i$ gives the equivariant linear map
\[
 \partial_i:\mathbb Q[X_\lambda]\longrightarrow
                 \mathbb Q[X_{\lambda+\epsilon_i-\epsilon_{i+1}}],
 \qquad [B]\longmapsto
                 \sum_{u\in B_{i+1}}[B\text{ with }u\text{ moved}].
\]
Indeed, identify $\mathbb Q[X_\lambda]$ with the weight-$\lambda$
subspace of $(\mathbb Q^p)^{\otimes N}$ by recording the block
of every tensor position. Then $\partial_i$ is exactly the
simple positive-root operator. Schur--Weyl and the highest-weight
argument above identify its joint kernel with the Specht module
$[\lambda]$, including its $S_N$ action.

Take the three corresponding systems for $\lambda,\mu,\nu$.
Tensoring over a field preserves kernels, and invariants preserve
kernels; hence the joint kernel inside
\[
 \bigl(\mathbb Q[X_\lambda]\otimes\mathbb Q[X_\mu]
                     \otimes\mathbb Q[X_\nu]\bigr)^{S_N}
\]
is $([\lambda]\otimes[\mu]\otimes[\nu])^{S_N}$.
The orbits of triples of ordered decompositions are indexed by
the intersection tables
$a_{ijk}=|B_i\cap C_j\cap E_k|$, giving precisely
$\mathcal T(\lambda,\mu,\nu)$. Use the normalized orbit sums
as a rational invariant basis. From any point of orbit $a$,
there are $a_{i+1,j,k}$ ways to move an element in the indicated
intersection to the new block. Equivariance distributes those
edges uniformly over the target orbit. Therefore the coefficient
on its normalized orbit sum is exactly $a_{i+1,j,k}$, which
recovers the matrix $D_i^{(1)}$, and likewise on the other axes.
This proves the formula directly through finite $G$-sets and
their incidence correspondences. It uses actual kernels after
linearization, without claiming that a nontrivial irreducible
itself has an effective permutation lift. The constant-group
scheme dictionary is SGA 3, I, 4.1, p.~23; the representation
comodule dictionary is I, Proposition 4.7.2, pp.~34--35.

\subsection*{A-019. A sufficient positive-basis hypothesis and its exact saturation scope}
\textbf{Statement.} Let $A=\bigoplus_{\alpha\in L_0}A_\alpha$ be
a commutative domain graded by a finite-rank free abelian group $L_0$, over
an algebraically closed field $k$. Suppose $A$ is integrally
closed in its fraction field and each $A_\alpha$ has dimension
at most one over $k$. Put
$S=\{\alpha:A_\alpha\ne0\}$ and $L=\mathbb ZS\subseteq L_0$.
Then $S$ is saturated in $L$: if $\alpha\in L$ and
$h\alpha\in S$ for an integer $h\geq1$, then $\alpha\in S$.
Moreover there is a homogeneous basis $(u_\alpha)_{\alpha\in S}$
with $u_\alpha u_\beta=u_{\alpha+\beta}$; thus
$A\simeq k[S]$. Finite generation is not needed for these
assertions. If $A$ is finitely generated, this supplies an affine
semigroup algebra with an effective monomial basis. If in
addition $L$ is primitive in $L_0$, the same saturation implication
holds for every $\alpha\in L_0$.

\textbf{Proof of saturation.} The support is a monoid because
$A$ is a domain. Write $\alpha=\beta-\gamma$ with
$\beta,\gamma\in S$; every element of $\mathbb ZS$ has such an
expression by collecting the positive and negative summands.
Choose nonzero $f_\beta\in A_\beta$,
$f_\gamma\in A_\gamma$, and $f_{h\alpha}\in A_{h\alpha}$.
Both $f_\beta^h$ and $f_\gamma^h f_{h\alpha}$ are nonzero
elements of the same at-most-one-dimensional space $A_{h\beta}$.
Hence $z=f_\beta/f_\gamma$ satisfies
$z^h=c f_{h\alpha}$ for some $c\in k^\times$.
It is integral over $A$ and therefore lies in $A$ by normality.
Its degree is $\alpha$: in the homogeneous decomposition of
$z$, multiplying by the nonzero homogeneous $f_\gamma$ can
produce the single degree $\beta$ only from degree $\alpha$.
Thus $A_\alpha\ne0$. If $L$ is primitive and
$h\alpha\in S\subseteq L$ for $\alpha\in L_0$, then
$\alpha\in L$ by primitivity, proving the last assertion.

\textbf{Proof of the positive basis.} Localize $A$ by all its
nonzero homogeneous elements, obtaining a graded domain $B$.
Its degree-zero part is $k$: the numerator and denominator of a
homogeneous fraction of degree zero belong to a common
one-dimensional $A$-piece and hence are proportional. Every
degree in $L$ occurs, and every nonzero homogeneous fraction is
invertible. Consequently each $B_\alpha$ is one-dimensional.
Choose a $\mathbb Z$-basis $\ell_1,\ldots,\ell_s$ of $L$ and
a nonzero $z_i\in B_{\ell_i}$. Distinct Laurent monomials in
the $z_i$ have distinct degrees, so they are linearly independent.
They span $B$ because every homogeneous piece is a line. Hence
$B=k[z_1^{\pm1},\ldots,z_s^{\pm1}]$ as a graded algebra.
For $\alpha\in S$, the corresponding Laurent monomial
$u_\alpha$ spans $B_\alpha$ and belongs to $A_\alpha$, since
that is a nonzero line in $B_\alpha$. These monomials give the
claimed multiplicative basis and algebra isomorphism.
\hfill$\square$

The algebraic-closure assumption is harmless for the intended
complex representation application; the displayed proof in
fact needs only that $A_0=k$ and the stated one-dimensional
grading. It is explicit multiplication plus normality and a
multiplicity-free grading that produce the positive saturated
model. An abstract special $\lambda$-ring does not provide these
hypotheses. In a highest-weight algebra whose graded dimensions
are Kronecker coefficients, the one-dimensional hypothesis fails
as soon as some coefficient exceeds one; no global application
to that algebra is asserted here.

The lattice qualification cannot be dropped. The polynomial
algebra $k[t]$ with $\deg t=2$ is normal and multiplicity-free,
but its support $2\mathbb N$ is saturated only in $2\mathbb Z$,
not in the chosen ambient lattice $\mathbb Z$. Nor does a
saturated finer support guarantee a saturated projected grading:
$\mathbb N^2$ is saturated in $\mathbb Z^2$, while its image under
$(a,b)\mapsto2a+3b$ is $\langle2,3\rangle$, which misses $1$
although it generates all of $\mathbb Z$ and contains $2$.
A proposed toric degeneration must therefore also prove
saturation of the particular weight projection that records
the desired multiplicities.

% END agents/archive/ARCHIVE_FINDINGS.tex

% BEGIN agents/archive/ORIGINAL_POWER_OPERATIONS.tex
\section{Original manuscript verification and the named power operations}

\subsection*{A-029. The operators \(q_H^n\) recovered from the original scan}
The definition absent from the supplied study is explicit on the
original cote 133, scan 021, lower half (PDF page 21, printed
footer 20). Let \(G\) be a finite group, \(H\leq S_n\),
and write \(A(G)\) for its Burnside ring and
\(A_n(G)=A(G\times S_n)\). The manuscript first defines
\[
 q_H:A_n(G)\longrightarrow A(G),\qquad
                   q_H([E])=[H\backslash E],                 \tag{A-029.1}
\]
and then defines
\[
                       q_H^n(x)=q_H(\tau^n(x)).                \tag{A-029.2}
\]
The injection operation \(\tau^n\) and its completed exponential
extension are on scans 020--021 and were proved in A-011. Consequently,
for an actual finite \(G\)-set \(E\),
\[
             q_H^n([E])=[\operatorname{Inj}([n],E)/H].          \tag{A-029.3}
\]
Left versus right orbit notation is immaterial here: the \(S_n\)
action is precomposition, converted to a right action by inversion,
and commutes with the \(G\) action. These are effective operations
on actual sets and integral polynomial operations of degree \(n\)
on the additive Burnside group. They are generally neither additive
nor homogeneous under ordinary integer scaling.

\begin{proof}
Taking \(H\)-orbits commutes with disjoint union. It therefore
induces the additive homomorphism (A-029.1) on group completions.
The action of \(G\) descends because it commutes with that of \(H\).
A-011 constructs the unique homomorphism of groups
\[
 \tau^*:A(G)\longrightarrow
        1+\prod_{j\geq1}A_j(G)
\]
where multiplication is the induction product. Its degree-\(n\)
coefficient agrees with the class of injections for actual sets.
Composing that coefficient with (A-029.1) defines (A-029.2) on
all virtual classes and proves (A-029.3).

For completeness, polynomial degree is meant in ordinary additive
coordinates, with integer-valued polynomial coefficients. Choose
the finite basis \(u_1,\ldots,u_s\) of transitive \(G\)-sets
for the free abelian group \(A(G)\). For integers \(a_i\),
\[
             \tau^*\left(\sum_i a_i u_i\right)
                         =\prod_i\tau^*(u_i)^{a_i}.
\]
Write \(\tau^*(u_i)=1+v_i\), where \(v_i\) has positive
induction degree. The degree-\(n\) truncation of
\((1+v_i)^{a_i}\) is
\(\sum_{j=0}^n\binom{a_i}{j}v_i^j\), also for negative
\(a_i\), by the formal inverse identity in a complete graded ring.
A term contributing to degree \(n\) has \(\sum_i j_i\leq n\).
Its coefficient is a product of integer-valued binomial polynomials
of total degree at most \(n\). The additive map \(q_H\)
preserves this conclusion. No denominator is inverted in the
Burnside group: every binomial coefficient takes an integer value
at every integer argument.

To see that the degree is exactly \(n\), restrict to multiples
\(a\cdot1\) of the trivial point. The action of \(H\) on
injections into an \(a\)-element set is free, giving for
\(a\geq0\)
\[
 q_H^n(a\cdot1)=\frac{n!}{|H|}\binom an\,1.                    \tag{A-029.4}
\]
Both sides are polynomial functions of \(a\), so the equality
holds at every integer. The leading coefficient is \(1/|H|\),
which is nonzero over \(\mathbb Q\); thus the degree is \(n\).
For \(n=0\) the operation is the constant one, of degree zero.
For every \(n\geq2\), (A-029.4) gives \(q_H^n(1)=0\) but
\(q_H^n(n\cdot1)=n!/|H|>0\). It is therefore not homogeneous
of degree \(n\) under ordinary scaling, and is not additive.

The parenthetical wording on scan 021 appears to call the
operations polynomial and homogeneous of degree \(n\). Its
unambiguous graded meaning is that \(\tau^n\) lies in induction
degree \(n\). If ``homogeneous'' is instead interpreted as the
identity \(q_H^n(ax)=a^nq_H^n(x)\), the last calculation refutes
that interpretation. No additional homogeneous-scaling assertion
is needed for the construction.
\end{proof}

For \(H=S_n\), (A-029.3) is exactly the \(n\)-element-subset
operation; for \(H=1\), it is the ordered-injection operation.
Thus the original definition connects the supplied notation to the
subset operations already distinguished from canonical opposite
symmetric powers in A-001--A-002. Linearization gives actual
permutation representations of these orbit sets, not an automatic
exterior-power identity or a counting formula for every irreducible
Kronecker multiplicity. The previous saturation counterexamples
are unaffected.

\paragraph{Original-source verification and exact scope.}
The official source is
\url{https://grothendieck.umontpellier.fr/133.pdf}. It contains
54 PDF pages, including its cover; the catalogue's 53 is the number
of manuscript/offprint pages. In this file, the supplied scan number
is the PDF page number, while the printed footer is one smaller.
The complete PDF hash and byte count are recorded in
\path{agents/archive/GROTHENDIECK133_SOURCE_MANIFEST.json}.
In this run the following originals were visually read in full:
scans 006, 008, 009, and 019--029. Their material agrees with the
supplied transcriptions or secondary study at the specific points
used in A-001--A-015, subject to the mathematical corrections already
proved. In particular the deletions on scans 009 and 022 are visible,
the subset identities on scan 019 are not deleted, the direct-summand
claim on scan 025 really is stated for a single sufficiently large
\(H_I\), and scan 029 explicitly asks the local-ring polynomiality
question. The uncertain name on scan 008 is not authenticated.
The remaining attached originals were subsequently verified as well:
cote2 scan017, cote12 scan041, cote43 scan037, cote119 scan217,
and the scope-exclusion pages cote4 scan007, cote16 scan038,
cote79 scan051 and cote89 scan018. Their relevant assertions and
symbol interpretations agree with the scoped audit. In particular,
the typescript in cote119 explicitly contains the historical statement
about exterior operations and symmetric-group operations. This is a
verification of its wording, not an independent priority claim.

The large volumes were accessed through exact HTTP byte ranges and
only their requested pages were extracted. Their individual manifests
record the official URL, PDF page, full-volume size and server ETag,
every retrieved range hash, and the extracted-page hash; a full-volume
SHA256 is not claimed for those eight sources. Thus all thirteen
attached archive pages, and nine additional cote133 pages, have an
actual recorded facsimile check. Remaining paleographic uncertainty
is stated where relevant and is not converted into a theorem.

\paragraph{Grothendieck dictionary.}
For a finite \(G\)-set, its associated constant finite scheme carries
an action of the constant finite group scheme \(G_{\mathbb Z}\).
The injection subset and the finite orbit quotient in (A-029.3)
are actual finite schemes with the induced action. Disjoint union
supplies the additive Burnside group, and induction supplies the
graded multiplication. Linearization is the associated finite free
permutation comodule, in the representation formalism of SGA 3, I,
Proposition 4.7.2, pp.~34--35. The homogeneous induction degree
is a torus grading as in SGA 3, I, Proposition 4.7.3, pp.~35--36.
Exterior operations on locally free sheaves instead use SGA 6, V,
Definition 2.1, p.~314, and VI, Theorem 3.3, pp.~392--393.
That dictionary preserves the distinction between orbit-quotient
set operations and exterior Schur operations; it supplies no
missing specialness, positivity of virtual differences, or
Kronecker saturation theorem.

% END agents/archive/ORIGINAL_POWER_OPERATIONS.tex

% BEGIN proofs/root_findings.tex
\section{The exact passage from finite sets to Kronecker multiplicities}

\subsection*{Finding M-001: a universal signed orbit formula}
Fix \(N\geq 0\).  For a weak composition
\(\alpha=(\alpha_1,\ldots,\alpha_l)\) of \(N\), let
\(X_\alpha\) be the \(S_N\)-set of ordered decompositions
\([N]=A_1\sqcup\cdots\sqcup A_l\), with \(|A_i|=\alpha_i\).
If a coordinate of \(\alpha\) is negative, put \([X_\alpha]=0\).
For three weak compositions of \(N\), write
\[
 C(\alpha,\beta,\gamma)=
 \#\left\{(u_{ijk})\in\mathbb Z_{\geq0}^{l\times m\times r}:
 \sum_{j,k}u_{ijk}=\alpha_i,\ 
 \sum_{i,k}u_{ijk}=\beta_j,\ 
 \sum_{i,j}u_{ijk}=\gamma_k\right\}.
\]
The same convention makes \(C=0\) if any margin is negative.
For partitions \(\lambda,\mu,\nu\) of \(N\), of lengths \(l,m,r\),
respectively, define
\(\alpha(\pi)_i=\lambda_i-i+\pi(i)\), and similarly
\(\beta(\sigma)\), \(\gamma(\tau)\).  Then
\begin{equation}\label{eq:M001}
 g(\lambda,\mu,\nu)=
 \sum_{\pi\in S_l,\,\sigma\in S_m,\,\tau\in S_r}
 \operatorname{sgn}(\pi)\operatorname{sgn}(\sigma)
 \operatorname{sgn}(\tau)
 C\bigl(\alpha(\pi),\beta(\sigma),\gamma(\tau)\bigr).
\end{equation}
This is a derived form of the Jacobi--Trudi identity, not a claim
of a new formula in the literature.

\begin{proof}
Let \(\ell:B(S_N)\longrightarrow R_{\mathbb Q}(S_N)\) be linearization.
The Young permutation module \(\mathbb Q[X_\alpha]\) has Frobenius
characteristic \(h_{\alpha_1}\cdots h_{\alpha_l}\).
Jacobi--Trudi therefore gives
\[
 [S^\lambda]=
 \ell\left(\sum_{\pi\in S_l}
                  \operatorname{sgn}(\pi)[X_{\alpha(\pi)}]\right).
\]
The precise classical inputs are Macdonald, \emph{Symmetric Functions
and Hall Polynomials}, second edition, I, (3.4) and (3.4'), and I, \S7, equations (7.2)--(7.6);
the Frobenius map and the tensor-product inner product are also stated
in Rosas, \S2, equations (3)--(5).
All irreducible characters of \(S_N\) are real and self-dual.
Consequently \(g(\lambda,\mu,\nu)\) is the dimension of the
\(S_N\)-invariants in \(S^\lambda\otimes S^\mu\otimes S^\nu\).
The dimension of the invariants of a permutation module
\(\mathbb Q[Y]\) is the number of orbits of \(Y\).
On substituting the three displayed virtual expressions, it remains
to determine the orbits of
\(X_\alpha\times X_\beta\times X_\gamma\).
The integers \(u_{ijk}=|A_i\cap B_j\cap C_k|\) are orbit invariants.
Conversely, given two triples with the same integers, choose a
bijection separately on each corresponding triple intersection;
their disjoint union is a permutation carrying one triple to the
other.  Every nonnegative table with the required margins is realized
by partitioning \([N]\) into cells of the prescribed sizes.
Thus the orbit number is exactly \(C(\alpha,\beta,\gamma)\), proving
\eqref{eq:M001}.
\end{proof}

\paragraph{What the formula does and does not provide.}
It recovers every Kronecker coefficient from actual finite-set orbit
counts and integral changes of basis.  The determinant signs are
essential to this argument.  It is not a cancellation-free count,
and it supplies neither saturation nor a polynomial-time algorithm
for binary-encoded arbitrary partitions.  The effective-lift
obstruction in A-003 explains why simply replacing the virtual
Young expression by an effective lift cannot work.

\subsection*{Finding M-003: the rational gamma graded ring loses this information}
Let \(G\) be any finite group, let \(R=R_{\mathbb C}(G)\) have its
exterior-power special lambda structure and rank augmentation, and
let \(I=\ker(\dim:R\to\mathbb Z)\).
For the gamma filtration of SGA~6, Exp.~V, 3.10,
\begin{equation}\label{eq:M003}
 F_\gamma^rR\otimes\mathbb Q=I\otimes\mathbb Q
 \quad(r\geq1),\qquad
 \operatorname{gr}_\gamma^rR\otimes\mathbb Q=0
 \quad(r\geq1).
\end{equation}
This statement concerns rationalization of this filtration, not its
integral associated graded ring.

\begin{proof}
In \(R_{\mathbb Q}:=R\otimes_{\mathbb Z}\mathbb Q\) put
\(e=[\mathbb C[G]]/|G|\).
The character of \(e\) is one at the identity and zero elsewhere.
Characters determine representation-ring classes, so \(e^2=e\) and
\(xe=\dim(x)e\) for every \(x\in R_{\mathbb Q}\).
It follows that \(I_{\mathbb Q}=(1-e)R_{\mathbb Q}\), and hence
\(I_{\mathbb Q}^r=I_{\mathbb Q}\) for \(r\geq1\).
By definition \(F_\gamma^rR\) is generated by products
\(\prod_j\gamma^{i_j}(x_j)\), where
\(x_j\in I\), \(i_j>0\), and \(\sum_j i_j\geq r\).
Since \(\gamma^1(x)=x\), it contains \(I^r\).
The augmentation is a lambda homomorphism; for \(x\in I\),
\(\dim(\gamma_t(x))=1\).  Thus every positive gamma operation on
\(x\) has rank zero and \(F_\gamma^rR\subset I\).
Tensoring \(I^r\subset F_\gamma^rR\subset I\) with the flat
\(\mathbb Z\)-module \(\mathbb Q\) proves the first equality.
Exactness of rationalization proves the assertion about the
associated graded.
\end{proof}

\paragraph{Consequence.}
The degree-zero rational graded piece retains the augmentation,
namely dimension; every positive graded piece is zero. Passing
to this associated graded cannot retain the additional representation
data needed for shape-dependent Kronecker multiplicities. In particular, a
saturation argument that uses only that associated graded has
discarded the data it would need.  This is a precise loss-of-data
statement; it does not rule out arguments using the full lambda
ring or integral gamma information.

\subsection*{Finding M-004: an integral invariant scheme, with its good-prime qualification}
Assume \(N\geq1\). Choose stable full integral lattices in the rational Specht modules, denoted \(S^\lambda_{\mathbb Z}\),
\(S^\mu_{\mathbb Z}\), \(S^\nu_{\mathbb Z}\), and put
\[
 T=S^\lambda_{\mathbb Z}\otimes S^\mu_{\mathbb Z}
                  \otimes S^\nu_{\mathbb Z}.
\]
Let \(s_1,\ldots,s_{N-1}\) be the adjacent transpositions of \(S_N\).
The integral map
\[
 A:T\longrightarrow T^{\,N-1},\qquad
 v\longmapsto((s_i-1)v)_i
\]
has
\begin{equation}\label{eq:M004}
 \operatorname{rank}_{\mathbb Z}\ker A
       =g(\lambda,\mu,\nu).
\end{equation}
Moreover the affine group scheme of solutions \(Av=0\) over
\(\mathbb Z[1/N!]\) is a vector bundle of this rank.  Formation of
its module of sections commutes with arbitrary base change within
that open subscheme of \(\operatorname{Spec}\mathbb Z\).

\begin{proof}
The equations \(s_iv=v\) for generators are equivalent to invariance
under \(S_N\).  Such lattices exist: take the integer span of all group translates of a rational basis. They are finite free, and their rationalizations are the corresponding irreducible rational representations.  Since localization is exact,
\((\ker A)\otimes\mathbb Q=(T\otimes\mathbb Q)^{S_N}\).
The latter has dimension \(g\), by self-duality and the definition
of a tensor-product multiplicity.  A submodule of a finite free
\(\mathbb Z\)-module is free, proving \eqref{eq:M004}.
Over \(B=\mathbb Z[1/N!]\), the averaging operator
\[
 P=\frac{1}{N!}\sum_{\sigma\in S_N}\sigma
\]
is an idempotent with image precisely \(T_B^{S_N}\).
This image is a direct summand of a finite free module.  It is
therefore finite projective, and its rank is constant on the
connected scheme \(\operatorname{Spec}B\); the generic rank is \(g\).
The same idempotent exists after any \(B\)-algebra base change, so
its image commutes with that base change.  Finally the equations
\(Av=0\) represent the functor of fixed vectors, and over \(B\)
this functor is the vector bundle associated to the indicated
projective module.
\end{proof}

\paragraph{Limits of the model.}
At primes dividing \(N!\), invariant formation need not commute
with base change, so the same dimension assertion has not been
made there.  The constant group scheme \(S_{N,\mathbb Z}\) is finite
etale even at these primes; it is the failure of linear
reductivity, not a failure of etaleness, that matters.
Computing a kernel dimension is an exact nonnegative linear-algebra
description.  Its ambient lattice can be very large, and this
description alone is not a positive combinatorial formula of the
kind required in the GCT flip.

\subsection*{Finding M-005: pure Artin realizations still retain a negative trace}
Let \(k=\mathbb F_q\), \(X=\operatorname{Spec}\mathbb F_{q^2}\), and
\(\ell\ne\operatorname{char}k\).
On the two geometric points of \(X\), geometric Frobenius is a
transposition.  The permutation representation
\(V=H^0(X_{\overline k},\mathbb Q_\ell)\) has eigenvalues \(1,-1\);
therefore \(\Lambda^2V\) is one-dimensional, pure of weight zero,
and has Frobenius trace \(-1\).

\begin{proof}
Frobenius exchanges the two basis vectors of \(V\).
In the basis consisting of their sum and difference, its matrix
is \(\operatorname{diag}(1,-1)\).
On \(\Lambda^2V\) it consequently acts by its determinant, \(-1\).
Every complex absolute value of this eigenvalue is one, which is
exactly purity of weight zero.  The assertions follow directly.
\end{proof}

In the category of rational Artin motives the same argument is the
decomposition \(h^0(X)=\mathbf1\oplus\varepsilon\), with
\(\Lambda^2h^0(X)=\varepsilon\); the rational idempotents are
\((1\pm \mathrm{Frob})/2\) in the finite quotient \(C_2\).
This is the motivic form of the corrected two-element Burnside
example, not a conjectural statement about general motives.
Purity controls absolute values and does not turn exterior powers
into effective finite sets or force nonnegative traces.  The RH
and Dirichlet-GRH input stipulated in the prompt supplies no
additional step in any proof in this paper.

\section{Dictionary into Grothendieck's frameworks}

\paragraph{Finite sets, group schemes and classifying topoi.}
For a finite group \(G\), SGA~3, Exp.~I, 4.1 (original p.~23) constructs its constant
group scheme \(G_{\mathbb Z}\).  Exp.~I, Proposition 4.7.2 (original pp.~34--35) identifies
representations of affine group schemes with comodules.
Thus a finite \(G\)-set \(E\) gives the locally free permutation
sheaf \(\mathcal O[E]\) on the classifying topos \(BG\).
Disjoint union becomes direct sum and Cartesian product becomes
tensor product.  Linearization is the resulting map into the
Grothendieck group of representations.  With short-exact-sequence
relations this is \(K_0(\operatorname{Rep}(G))\);
it must not be confused with the split group of indecomposable
representations in modular characteristic.

\paragraph{Exterior powers, symmetric powers and their signs.}
SGA~6, Exp.~V, Definition 2.1 (p.~314), Example 2.2(b), and Lemma 2.2.1 (pp.~315--316)
construct the pre-lambda structure using exterior powers and
their exact-sequence filtration.  Exp.~VI, Theorem 3.3
(printed pp.~392--393) proves the special lambda identities for
the \(K\)-group of locally free modules of bounded rank on a
commutatively ringed topos.  The exterior basis associated with
an \(r\)-element subset has the orientation sign of the induced
permutation of that subset.  Forgetting that sign gives the
incorrect subset formula audited in cote~133, scan~019.
The correct opposite-symmetric Burnside operations linearize to
these exterior powers, while remaining non-special upstairs.
This dictionary carries the proofs of A-001--A-002 and A-003 and the
signed formula M-001 to sheaves on \(BS_N\).

\paragraph{Schur functors and tensor powers.}
In characteristic zero, or after inverting \(N!\) for degree \(N\),
the rational symmetric-group projectors on a tensor power give
the usual Schur isotypic summands, compatibly with base change.
Over \(\mathbb Z\) one may instead use integral Schur modules and
filtrations; a rational direct-sum decomposition is not silently
asserted integrally.  Cote~133, scan~028 is the archive locator for
the tensor-power class and its monomial-symmetric expansion.
It concerns the induction product across the tower of symmetric
groups.  The Kronecker product is the tensor product at a fixed
degree, equivalently restriction along the diagonal
\(S_N\to S_N\times S_N\); cote~2, scan~017 records that
external-product/diagonal mechanism.

\paragraph{Gamma filtration.}
SGA~6, Exp.~V, 3.10 (p.~331) defines the gamma filtration and
3.10.1--3.10.2 give its generators.  Formula M-003 shows exactly
what happens to this filtration on \(K_0(\operatorname{Rep}_{\mathbb C}G)\)
after rationalization.  This calculation has no implication that
the integral filtration, or the gamma filtration of an arbitrary
scheme, vanishes.


\paragraph{Borger's lambda geometry.}
The representation ring, with its special lambda structure, is torsion-free
as an abelian group. Its Adams operations give commuting Frobenius lifts
on its spectrum: on characters, \(\psi^p(\chi)(g)=\chi(g^p)\).
This is the flat affine case of Borger, \emph{Lambda-rings and the field
with one element}, \S1.1 (the supplied source, lines 774--797).
The non-special Burnside structure audited in A-001--A-002 does not qualify for
that conclusion by the same route. Nor does the Frobenius-lift
interpretation restore effectivity: on the standard representation
\(V\) of \(S_3\), one has
\(\psi^2(V)=1+V-\operatorname{sgn}\), as computed in A-004.
Thus a scheme endomorphism arising from an Adams operation can have
negative coordinates in the irreducible-representation basis.
This is a concrete compatibility statement and a concrete limitation,
not a claim that all methods of lambda geometry fail.

\paragraph{Orbit-closure schemes and parameter spaces.}
The determinant and padded permanent are homogeneous polynomials
with integral coefficients.  Their rational orbit closures admit
integral models by contracting their defining ideals from
\(\mathbb Q\) to the ambient polynomial algebra over \(\mathbb Z\).
Those quotients are torsion-free as \(\mathbb Z\)-modules and hence
flat over \(\mathbb Z\); their generic fibers are the prescribed
closures.  It is not asserted that every special fiber is the
orbit closure of the reduced polynomial.
Projectivizing these homogeneous models places flat families with
fixed Hilbert polynomial in a Hilbert scheme.
The precise FGA source is Grothendieck, Bourbaki exposé~221,
Theorems 3.1 and 3.2 (printed pp.~259--260), for the Quot functor
and its fixed-Hilbert-polynomial projective pieces; the Hilbert
case is \(F=\mathcal O_X\), explained immediately after the
construction.  These representability results organize possible
families.  They do not prove the required comparison of the
determinant and permanent isotypic multiplicities.

\paragraph{Sheaves, weights and motives.}
The permutation sheaves and their orientation summands have
finite monodromy and admit Artin realizations.  Finding M-005
exhibits the sign already at this level.  For Grassmannian
realizations of the rectangular/two-row formula, Schubert cells
give even Tate cohomology and hard Lefschetz identifies the
partition-number difference with a primitive multiplicity.
The exact cohomological locators used here are SGA~4,
Exp.~XVIII, Theorem 3.2.5 (duality), SGA~5, Exp.~III, 4.7
(Lefschetz--Verdier trace formula), and Deligne,
\emph{La conjecture de Weil II}, 1.2.2 (purity) and
Theorem 4.1.1 (hard Lefschetz).
Their use does not identify an arbitrary Kronecker coefficient
with cohomology concentrated in one degree.
Only Artin motives and Tate realizations are used here. The
Grassmannian primitive quotient is a statement about cohomological
dimensions; no canonical motivic identification with a Kronecker
multiplicity space is asserted.

% END proofs/root_findings.tex

\part{Positive formulas and restricted saturation}

% BEGIN agents/formulas/formulas.tex
% GCT_KRON_20261120, positive_formulas worker, GPT-6 Astra Pro.
% This fragment uses standard amsmath/amssymb/amsthm commands only.
% Finding means established in this run; no claim of literature originality.

\section{Positive formulas obtained from finite sets and exterior powers}
Throughout this section the ground field has characteristic zero. Write
$H_a^N=[N-a,1^a]$, $T_a^N=[N-a,a]$, and $R=(b^a)$, with $N=ab$ in the
rectangular case. Zero parts are omitted. The superscript on $H_a^N$ and
$T_a^N$ denotes the size, not a tensor power. Statements about the effective
cone of the Burnside ring are not implicit in any statement below.

\subsection{F-001: a positive kernel formula on binary contingency tables}
\textbf{Statement.}
For $0\leq r,s,u\leq N/2$, let $\mathcal T_N(r,s,u)$ be the set of eight-entry
nonnegative integral arrays $x=(x_{ijk})_{i,j,k\in\{0,1\}}$ satisfying
\[
 \sum x_{ijk}=N,\qquad \sum_{jk}x_{1jk}=r,\qquad
 \sum_{ik}x_{i1k}=s,\qquad \sum_{ij}x_{ij1}=u.
\]
Define three maps on the rational vector space with basis $e_x$ indexed by
$\mathcal T_N(r,s,u)$. The first map has target basis
$\mathcal T_N(r-1,s,u)$ and is
\[
 D_1e_x=\sum_{j,k}x_{1jk}
 e_{x-\mathbf1_{1jk}+\mathbf1_{0jk}}.                 \tag{F-001.1}
\]
A term of coefficient zero is omitted. Define $D_2,D_3$ by changing the
second and third indices in the same way. If a marginal is zero its target
space and map are zero. For the stacked matrix $D=(D_1,D_2,D_3)$,
\[
 g((N-r,r),(N-s,s),(N-u,u))
 =\dim_{\mathbb Q}\ker D
 =|\mathcal T_N(r,s,u)|-\operatorname{rank}_{\mathbb Q}D. \tag{F-001.2}
\]
In particular the coefficient is nonzero exactly when the specified integral
matrix has a nonzero rational null vector. This is a positive vector-space
formula and an exact rank criterion, not a cancellation-free enumeration of
set orbits.

\textbf{Proof.}
Let $M_v=\mathbb Q[\binom{[N]}v]$. On subset basis vectors define
$D_v(A)=\sum_{a\in A}(A\setminus\{a\})$ and
$U_v(A)=\sum_{b\notin A}(A\cup\{b\})$.
Counting a deletion followed by an insertion gives
\[
 D_{v+1}U_v-U_{v-1}D_v=(N-2v)\,\mathrm{id}_{M_v}.        \tag{F-001.3}
\]
Give subset basis vectors the standard orthonormal inner product. Then
$U_v^*=D_{v+1}$. If $v<N/2$, (F-001.3) implies
\[
 \|U_vw\|^2=\|D_vw\|^2+(N-2v)\|w\|^2,
\]
so $U_v$ is injective and $D_{v+1}$ is surjective. For
$1\leq v\leq N/2$, the characteristic of $M_v$ is $h_{N-v}h_v$ and Pieri's
rule (Macdonald, \emph{Symmetric Functions and Hall Polynomials}, second
edition, I~(5.16), p.~73, together with I~(7.3), p.~113) gives
$M_v\cong\bigoplus_{j=0}^v[N-j,j]$.
Consequently $K_v:=\ker D_v\cong[N-v,v]$; put $K_0=M_0$.
This identifies an actual irreducible kernel in a permutation module, instead
of merely subtracting two classes in its Grothendieck group.

In $M_r\otimes M_s\otimes M_u$, the intersection of the three partial
kernels is $K_r\otimes K_s\otimes K_u$. To see this directly, choose vector
space complements $M_v=K_v\oplus C_v$. The tensor product splits into eight
summands, and the first partial kernel consists of those having first factor
$K_r$, then the second and third partial kernels impose their corresponding
conditions. This argument also applies when one of the target spaces is
zero. Taking $S_N$-invariants therefore gives
\[
 \bigcap_{i=1}^3\ker(D_i|_{(M_r\otimes M_s\otimes M_u)^{S_N}})
  =(K_r\otimes K_s\otimes K_u)^{S_N}.
\]
The right side has dimension the required Kronecker coefficient.

The orbits of triples $(A,B,C)$ of subsets of prescribed cardinalities are
classified by their eight membership counts $x_{ijk}$. In the invariant
space use the \emph{average} of the orbit basis vectors,
$e_x=|O_x|^{-1}\sum_{(A,B,C)\in O_x}(A,B,C)$.
Deleting a point of type $(1,j,k)$ has $x_{1jk}$ possible choices.
Every target triple in the new orbit has the same number of preimages;
double counting gives
$|O_x|x_{1jk}=|O_{x'}|(x_{0jk}+1)$.
Thus the coefficient on orbit averages is $x_{1jk}$, exactly (F-001.1).
For unnormalized orbit sums the coefficient would instead be $x_{0jk}+1$.
This proves the formula. \hfill$\square$

\textbf{Complexity and dictionary.}
Four free entries determine the remaining entries of a binary table with
these margins, so there are at most $(N+1)^4$ tables. Rational matrix rank
yields a polynomial algorithm in $N$, with standard exact arithmetic; this
is not a claim of polynomial complexity in $\log N$.
The set of $v$-subsets is a finite constant scheme over $\mathbb Z$ carrying
an action of the constant finite group scheme $S_N$. Its free module is the
linearization of that scheme. The incidence maps are integral equivariant
maps. Over $\mathbb Q$ their kernels are representations, hence vector bundles
on the classifying stack $BS_N$, and (F-001.2) computes an invariant Hom rank
in this category. Exterior powers and Schur functors act on its $K_0$.
No claim of integral semisimplicity, integral direct-summand projectors, or a
finite-set basis for $K_v$ is made.

\subsection{F-004: a positive generating function for three two-row shapes}
\textbf{Statement.}
The following equality of formal power series holds:
\[
\begin{gathered}
 \sum_{N\geq0}\ \sum_{0\leq r,s,u\leq N/2}
 g((N-r,r),(N-s,s),(N-u,u))t^Nx^ry^sz^u\\
 =\frac{1+t^3xyz}{(1-t)(1-t^2xy)(1-t^2xz)(1-t^2yz)
 (1-t^4x^2y^2z^2)}.
\end{gathered}                                      \tag{F-004.1}
\]
The $N=0$ term is $1$. In particular, after sorting $0\leq u\leq s\leq r\leq
N/2$, put
\[
 L=\max\left(0,\left\lceil\frac{r+s+u-N}{2}\right\rceil\right),\qquad
 U=\left\lfloor\frac{s+u-r}{2}\right\rfloor.
\]
Then the coefficient is the positive integer-interval count
\[
 g((N-r,r),(N-s,s),(N-u,u))=\#([L,U]\cap\mathbb Z).       \tag{F-004.2}
\]
The interval formula is the classical result of Rosas, Corollary~5,
\emph{The Kronecker product of Schur functions indexed by two-row shapes or
hook shapes}, arXiv:math/0001084, PDF p.~11. The statement established here
is the explicit positive-series variant and its dictionary from subset
orbits; it is not a claim of originality for the coefficient formula.

\textbf{Proof of the generating variant.}
The orbit table count in F-001 has generating series
\[
 B(t,x,y,z)=\prod_{i,j,k\in\{0,1\}}(1-tx^iy^jz^k)^{-1}.
\]
For $r,s,u\leq N/2$, the coefficient of $t^Nx^ry^sz^u$ in
$(1-x)(1-y)(1-z)B$ is the Kronecker coefficient: in each factor use
$[N-r,r]=[M_r]-[M_{r-1}]$ and then take invariants. The finite differences
have additional coefficients outside this dominant chamber, so they are
not equal to (F-004.1) as an unrestricted series.

Write $G$ for the rational function on the right side of (F-004.1). For
$\epsilon=(\epsilon_1,\epsilon_2,\epsilon_3)\in\{0,1\}^3$ put
$P_\epsilon=x^{\epsilon_1}y^{\epsilon_2}z^{\epsilon_3}$, and define
\[
 \mathcal W(G)=\sum_\epsilon(-1)^{|\epsilon|}P_\epsilon
 G(tP_\epsilon,x^{1-2\epsilon_1},y^{1-2\epsilon_2},z^{1-2\epsilon_3}).
\]
There is an exact rational identity
\[
 \mathcal W(G)=(1-x)(1-y)(1-z)B.                       \tag{F-004.3}
\]
For completeness, here is a finite polynomial certificate of that identity,
which also specifies an elementary way to check every denominator and sign.
Put $K=t^2xyz$ and let $v_1=x,v_2=y,v_3=z$. Multiply (F-004.3) by
\[
 (1-K^2)\prod_\eta(1-tP_\eta)
 \prod_{i=1}^3(1-K/v_i)(1-Kv_i).
\]
All these factors are polynomials. The difference of the two sides becomes
\begin{align*}
 &\sum_\epsilon(-1)^{|\epsilon|}P_\epsilon(1+KtP_\epsilon)
  \prod_{\eta\ne\epsilon}(1-tP_\eta)
  \prod_{i=1}^3(1-Kv_i^{1-2\epsilon_i})\\
 &\hspace{10mm}-(1-x)(1-y)(1-z)(1-K^2)
   \prod_{i=1}^3(1-K/v_i)(1-Kv_i)=0.                 \tag{F-004.4}
\end{align*}
This polynomial identity follows by expanding the eight displayed summands
and collecting like monomials after $K=t^2xyz$; the supplied
\texttt{check\_formulas.py}, function \texttt{rational\_identity\_certificate},
performs exactly this multiplication in $\mathbb Z[t,x,y,z]$ and returns the
zero polynomial. Thus no convergence or numerical evaluation is involved.

Every monomial $t^Nx^ry^sz^u$ in $G$ satisfies $2r,2s,2u\leq N$: this holds
separately for every generator in its denominator and for its two numerator
monomials. A nonidentity summand of $\mathcal W$ replaces at least one such
exponent, say $r$, by $N-r+1>N/2$. Hence it contributes nothing in the
simultaneous dominant chamber. Extracting coefficients in that chamber from
(F-004.3) proves (F-004.1).

To extract its positive coefficient, choose $\varepsilon\in\{0,1\}$ from
the numerator and let $j\geq0$ be the exponent of
$t^4x^2y^2z^2$. The three exponents on pair generators are forced to be
\[
 \frac{r+s-u-2j-\varepsilon}{2},\quad
 \frac{r+u-s-2j-\varepsilon}{2},\quad
 \frac{s+u-r-2j-\varepsilon}{2}.
\]
The exponent on $t$ alone is $N-r-s-u+2j$. Integrality uniquely determines
$\varepsilon\equiv r+s+u\pmod2$. Nonnegativity of all four exponents says
exactly $L\leq j\leq U$. This proves (F-004.2). \hfill$\square$

\textbf{Geometric dictionary.}
Let $A,B,C$ be rank-two vector spaces. Schur--Weyl duality (Macdonald, I, Appendix~A, (5.5)--(5.6), p.~156) identifies the
coefficient above with the dimension of the highest-weight subspace of
$\operatorname{Sym}^N(A\otimes B\otimes C)$ of three weights
$(N-r,r),(N-s,s),(N-u,u)$. Choose rank-two free lattices with rational fibers $A,B,C$. The eight
variables are the weights of the coordinate functions on the affine dual
tensor scheme $(A\otimes B\otimes C)^*$, whose integral model is
$\mathbb A^8_{\mathbb Z}$ and whose coordinate ring is the symmetric
algebra on the tensor product of those lattices. The operator $\mathcal W$ is
rank-one Weyl antisymmetrization in each $\mathrm{GL}_2$ factor. Thus the
positive series is a multigraded highest-weight Hilbert series, while the
unrestricted table series is its torus-weight Hilbert series. This is the
precise distinction lost if one treats finite differences as positive orbit
counts.

\subsection{F-005: a symmetric positive formula for three hooks}
\textbf{Statement.}
For $0\leq r,s,u\leq N-1$,
\[
 g(H_r^N,H_s^N,H_u^N)=
 \begin{cases}
 1,&r\leq s+u,\ s\leq r+u,\ u\leq r+s,\quad r+s+u\leq2N-2,\\
 0,&\text{otherwise}.
 \end{cases}                                         \tag{F-005.1}
\]
The positive generating form is
\[
 \sum_{N\geq1}\sum_{r,s,u=0}^{N-1}g(H_r^N,H_s^N,H_u^N)t^Nx^ry^sz^u
 =\frac{t(1+t^2xyz)}{(1-t)(1-txy)(1-txz)(1-tyz)}.      \tag{F-005.2}
\]
This is a symmetric all-boundaries form of the classical hook formulas;
compare Rosas, Theorem~13(4), arXiv:math/0001084, PDF pp.~14--17. We give the
short exterior-power derivation because it exhibits exactly which
$\lambda$-operations are involved.

\textbf{Proof.}
Let $V_N=\mathbb Q^N$ be the permutation representation and
$W_N=V_N-\mathbf1$ its actual augmentation representation for $N\geq1$.
The hook identity is $\bigwedge^rW_N\cong H_r^N$: indeed,
$\bigwedge^rV_N$ is induced from $\operatorname{sgn}_{S_r}\boxtimes\mathbf1$,
so the vertical Pieri formula gives $H_r^N\oplus H_{r-1}^N$ in the
interior degrees; induction using $V_N=\mathbf1\oplus W_N$ gives the
claim, with the endpoints checked separately. The exact Pieri locator is
Macdonald, I~(5.17), p.~73.
A permutation with cycle lengths $m_1,\ldots,m_q$ satisfies
\[
 \lambda_x(V_N)(\sigma)=\det(1+x\sigma|V_N)
 =\prod_{i=1}^q(1-(-x)^{m_i}).
\]
For each cycle type, averaging the product of the three such determinants
computes the invariant dimension. The symmetric-group cycle-index identity
therefore gives
\begin{align*}
 R(t,x,y,z)
 &=\sum_{N\geq0}t^N\frac1{N!}\sum_{\sigma\in S_N}
   \det(1+x\sigma)\det(1+y\sigma)\det(1+z\sigma)\\
 &=\exp\left(\sum_{m\geq1}\frac{t^m}{m}
 (1-(-x)^m)(1-(-y)^m)(1-(-z)^m)\right)\\
 &=\frac{(1+tx)(1+ty)(1+tz)(1+txyz)}
 {(1-t)(1-txy)(1-txz)(1-tyz)}.
\end{align*}
The exponential identity follows directly by choosing the number of cycles
of each length; a cycle type with $c_m$ cycles of length $m$ has normalized
weight $\prod_m(m^{c_m}c_m!)^{-1}$. Thus it also proves the only cycle-index
fact used here.

Removing the trivial summand in each exterior-power series divides by
$(1+x)(1+y)(1+z)$ for $N\geq1$. The $N=0$ term must first be removed.
The elementary factorization
\begin{align*}
 &(1+tx)(1+ty)(1+tz)(1+txyz)
 -(1-t)(1-txy)(1-txz)(1-tyz)\\
 &\hspace{15mm}=t(1+x)(1+y)(1+z)(1+t^2xyz)
\end{align*}
proves (F-005.2).

For coefficient extraction, choose $\varepsilon=0$ or $1$ in the numerator.
The pair exponents are forced to be
$(r+s-u-\varepsilon)/2$, $(r+u-s-\varepsilon)/2$,
$(s+u-r-\varepsilon)/2$; their integrality uniquely determines
$\varepsilon\equiv r+s+u\pmod2$.
They are nonnegative exactly when the triangle inequalities hold. The
remaining exponent of $t$ is
$N-1-(r+s+u+\varepsilon)/2$; it is nonnegative exactly when
$r+s+u\leq2N-2$. Each admissible coefficient is therefore counted once.
The formula also shows $r,s,u\leq N-1$ automatically, and covers $N=1$.
\hfill$\square$

\textbf{Dictionary.}
Over $\mathbb Z$, $W_N$ is the locally free kernel of the augmentation map
$\mathbb Z^N\to\mathbb Z$. It is a representation of the constant group
scheme $S_N$, and $\lambda^r[W_N]=[\bigwedge^rW_N]$ in the representation
category's $K_0$. Over $\mathbb Q$ the invariant projection is the Reynolds
idempotent $N!^{-1}\sum\sigma$. Formula (F-005.2) computes ranks of these
invariant vector bundles. These are genuine exterior powers, with the
orientation signs of permuted wedges. They must not be replaced by the
unsigned set of $r$-subsets.

\subsection{F-002: the positive stable rectangle--two-row formula}
\textbf{Statement and source.}
Let $p_r(a,b)$ count partitions of $r$ in an $a$ by $b$ rectangle, and put
$p_{-1}=0$. The known rectangle identity is
\[
 g((b^a),(b^a),(ab-r,r))=p_r(a,b)-p_{r-1}(a,b),
 \qquad 0\leq r\leq ab/2.                            \tag{F-002.1}
\]
Its exact source is Pak--Panova, \emph{Strict unimodality of $q$-binomial
coefficients}, Lemma~3, p.~2. Suppose $a\leq b$. In the larger stable window
$r\leq b$ (still with $r\leq ab/2$), the positive refinement is
\[
 g((b^a),(b^a),(ab-r,r))
 =[q^r]\prod_{i=2}^a(1-q^i)^{-1}.                    \tag{F-002.2}
\]
Thus it counts partitions of $r$ into parts from $\{2,3,\ldots,a\}$.
For $a=1$ it is $1$ only at $r=0$; for $a=2$ it is $1$ precisely for even
$r$; for $a\geq3$ it is positive precisely for $r=0$ or $r\geq2$ in this
window.

\textbf{Proof of the refinement and the Burnside dictionary.}
Since $r\leq b$, the width condition in the rectangle imposes no restriction.
Partitions counted by $p_r(a,b)$ have at most $a$ parts. Conjugation
identifies them with partitions whose parts are at most $a$.
Multiplication of their generating series by $(1-q)$ removes the arbitrary
number of parts equal to $1$, proving (F-002.2). Its support statements
follow from the empty partition, parity for $a=2$, and the fact that every
integer $r\geq2$ is $2v+3w$ with $v,w\geq0$.
For the link with F-001, the two-row identity
$[N-r,r]=[M_r]-[M_{r-1}]$ and Frobenius reciprocity give
\[
 g(\alpha,\beta,(N-r,r))=A_r(\alpha,\beta)-A_{r-1}(\alpha,\beta),\qquad
 A_r=\sum_{\eta\vdash r,\theta\vdash N-r}
 c^\alpha_{\eta\theta}c^\beta_{\eta\theta}.
\]
For a rectangle $R=(b^a)$, the coefficient $c^R_{\eta\theta}$ is $1$ exactly
when $\theta$ is the rotated complement of $\eta$ in $R$, and is $0$
otherwise. One way to check this particular Littlewood--Richardson fact is
$S_R(\mathbb Q^a)=\det^b$: its multiplicity in
$S_\eta\otimes S_\theta$ is the multiplicity of $S_\theta$ in
$\det^b\otimes S_\eta^*$, an irreducible whose weight is that rotated
complement if it is polynomial. Thus $A_r(R,R)=p_r(a,b)$.
This recovers the known identity in precisely the form needed for the
finite-set linearization. \hfill$\square$

\subsection{F-003: rectangular hooks, with all endpoints}
\textbf{Statement and source.}
Put $m=\min(a,b)$. For every $0\leq k\leq ab-1$,
\[
 g((b^a),(b^a),(ab-k,1^k))
 =[q^k]\prod_{i=2}^m(1+q^{2i-1}).                    \tag{F-003.1}
\]
This is Ikenmeyer--Panova, \emph{Rectangular Kronecker coefficients and
plethysms in geometric complexity theory}, Proposition~6.4, PDF p.~15;
attached source lines 1380--1427. We use the empty-product convention:
if $m=1$, only $k=0$ has value $1$. The polynomial has degree $m^2-1$, so
all $k>m^2-1$ vanish, even when the rectangle is much longer in one direction.
The nonvanishing criterion is exact: $k$ must be a sum of distinct members
of $\{3,5,\ldots,2m-1\}$.

\textbf{Short exterior-power verification of the boundaries.}
The permutation representation satisfies
$\bigwedge^kV_N=H_k^N\oplus H_{k-1}^N$ for $1\leq k\leq N-1$,
with the obvious single summands at $k=0,N$ and $H_{-1}=H_N=0$.
The module $\bigwedge^kV_N$ is induced from the sign representation on
$S_k$ and the trivial representation on $S_{N-k}$, not the unsigned subset
module. Restricting the two rectangular factors and using the complement
rule from F-002 gives
\[
 g_k+g_{k-1}=\#\{\eta\vdash k:\eta=\eta',\ \eta\subseteq(b^a)\}.
\]
Self-conjugate diagrams inside the rectangle are exactly those inside the
$m$ by $m$ square. Their diagonal hook lengths are distinct odd integers in
$\{1,3,\ldots,2m-1\}$, and these lengths sum to the diagram's size.
This bijection gives $(1+q)\sum g_kq^k=\prod_{i=1}^m(1+q^{2i-1})$.
Cancelling $1+q$ proves the endpoint extension and records explicitly the
orientation needed in the $\lambda$-operation. \hfill$\square$

\subsection{F-006 and F-007: a corrected hook-positivity statement}
\textbf{F-006 (proved refinement).}
For $m\geq7$, the support of the product in (F-003.1) is
\[
 \{0,1,\ldots,m^2-1\}\setminus
 \{1,2,4,6,m^2-7,m^2-5,m^2-3,m^2-2\}.               \tag{F-006.1}
\]
Here is an elementary full proof of the support assertion, independent of
any strict-unimodality assertion. Let $A_m=\{3,5,\ldots,2m-1\}$ and
$D_m=\sum A_m=m^2-1$. Subset complementation makes its subset-sum support
symmetric about $D_m/2$. No subset sums to $1,2,4,6$; their complements are
also impossible. For $m=7$, all values $7$ through $24$ are represented by
\[
\begin{array}{c|rrrrrrrrr}
k&7&8&9&10&11&12&13&14&15\\ \hline
\text{sum}&7&3+5&9&3+7&11&3+9&13&3+11&3+5+7
\end{array}
\]
\[
\setlength{\arraycolsep}{3pt}
\begin{array}{c|rrrrrrrrr}
k&16&17&18&19&20&21&22&23&24\\ \hline
\text{sum}&3+13&3+5+9&5+13&3+5+11&7+13&3+5+13&9+13&3+7+13&11+13.
\end{array}
\]
Their complements cover $24$ through $41=D_7-7$. The additional values
$0,3,5,D_7-5,D_7-3,D_7$ are immediate. This proves the base case.
If the assertion holds for $m\geq7$, the new part $2m+1$ shows that the next
support contains both intervals
$[7,D_m-7]$ and $[2m+8,D_{m+1}-7]$. They overlap or are adjacent because
$2m+8\leq m^2-7$ for $m\geq7$. The boundary values $0,3,5$ and their
new complements are present. The eight impossible values already identified
are the only remaining values. This proves (F-006.1). \hfill$\square$

\textbf{F-007 (refutation of literal source wording).}
In the attached version of Ikenmeyer--Panova Theorem~6.1
(\texttt{thm:hook\_positive}, source lines 1340--1356), the last exceptional
value is printed as $d^1-2$. The correct last value is $d^2-2$, as follows
both from Proposition~6.4 and from the same source's later table.
More substantively, its unqualified sentence
$g_k(d,n)>g_{k-2}(d,n)$ for $k\leq d^2/2$ is false. For every $d\geq2$ and
$n\geq d$, (F-003.1) gives $g_0=1$ and $g_2=0$; for $d\geq2$, $k=2$ is
in the stated range. In particular the failure persists for $d\geq7$.
This refutes that sentence, not Proposition~6.4 or the corrected support
classification. We use F-006, not the false strictness claim, anywhere that
hook support is needed.

\subsection{F-008: the Grassmannian interpretation and its limit}
Let $X=\operatorname{Gr}(a,a+b)$ over $\mathbb Z$, of dimension $ab$.
Its Schubert cells give
$\dim H^{2r}(X_{\mathbb C},\mathbb Q)=p_r(a,b)$ and vanishing odd cohomology.
Cup product by the Pl\"ucker hyperplane class $L$ is injective from
$H^{2r-2}$ to $H^{2r}$ for $r\leq ab/2$, by hard Lefschetz. Hence
\[
 g((b^a),(b^a),(ab-r,r))
 =\dim\frac{H^{2r}(X_{\mathbb C},\mathbb Q)}{L H^{2r-2}(X_{\mathbb C},\mathbb Q)}.
                                                            \tag{F-008.1}
\]
The quotient is isomorphic to primitive cohomology after choosing the
Lefschetz decomposition. Formula (F-008.1) is a consequence of the known
rectangle identity and hard Lefschetz; it asserts equality of dimensions,
not a constructed canonical identification with a Kronecker multiplicity
space. Over a finite field the Schubert cell filtration gives the same
ranks in $\ell$-adic cohomology, and the even groups have Tate Frobenius
weights $2r$; the normalized quotient after a Tate twist has weight zero.
These facts provide the requested schemes, cohomology, weights, and Tate
motive dictionary in this family. They do not prove strict positivity of
every primitive degree or a saturation theorem: the parity example with
$a=2$ is an explicit obstruction. None of this argument uses RH or GRH.

% END agents/formulas/formulas.tex

% BEGIN agents/formulas/hooks_two_row.tex
\section{Two hooks and one two-row shape, with all overlap cases}

\subsection*{Finding F-020: a positive two-hook formula}
Let \(N\geq1\), \(H_e=(N-e,1^e)\),
\(0\leq e,f\leq N-1\), and \(0\leq k\leq N/2\).
Write \(\mathbf1_P\) for the indicator of a condition \(P\).
For the one-row output,
\[
                    g((N),H_e,H_f)=\mathbf1_{e=f}.
\]
For \(k=1\), which requires \(N\geq2\),
\begin{equation}\label{eq:F020hook}
 g((N-1,1),H_e,H_f)=
    \mathbf1_{|e-f|=1}
       +\mathbf1_{e=f}\mathbf1_{1\leq e\leq N-2}.
\end{equation}
For \(2\leq k\leq N/2\), put \(a=N-k\). Then
\begin{equation}\label{eq:F020}
 \begin{split}
 g((N-k,k),H_e,H_f)
   ={}&\mathbf1_{e=f}
       \bigl(\mathbf1_{k-1\leq e\leq a-1}
                    +\mathbf1_{k\leq e\leq a}\bigr)\\
     &+\mathbf1_{|e-f|=1}
                   \mathbf1_{k-1\leq\min(e,f)\leq a-1}.
 \end{split}
\end{equation}
These are nonnegative formulas with values in \(\{0,1,2\}\).
The genuine two-row case \(k\geq2\) is equivalent to
Rosas, Corollary~14, p.~18 of arXiv:math/0001084v1;
the restriction proof below also supplies all boundary cases.
The formula is recorded as a checked variant of known hook
rules, not as a claim of a new Kronecker formula.

\begin{proof}
Let \(U_t\) be the augmentation representation of \(S_t\).
The exterior power \(\Lambda^i U_t\), \(0\leq i\leq t-1\),
is the irreducible hook representation \([(t-i,1^i)]\).
One quick verification uses the signed subset module
\[
 \Lambda^i\mathbb C^t=
 \operatorname{Ind}_{S_{t-i}\times S_i}^{S_t}
                          (\mathbf1\boxtimes\operatorname{sgn}).
\]
Its characteristic is \(h_{t-i}e_i\). Vertical Pieri,
Macdonald I(5.17), and
\(\Lambda^i(\mathbf1\oplus U_t)
=\Lambda^i U_t\oplus\Lambda^{i-1}U_t\), prove the
assertion successively from \(i=0\). The endpoint \(t=1\)
has only \(\Lambda^0 U_1=\mathbf1\).

For \(1\leq j\leq N-1\), restrict to
\(S_{N-j}\times S_j\). The two block-constant vectors
in the permutation representation leave one trivial vector
in its augmentation, so
\[
 U_N\big|_{S_{N-j}\times S_j}
                  =U_{N-j}\oplus U_j\oplus\mathbf1.
\]
Consequently
\[
 \Lambda^e U_N\big|_{S_{N-j}\times S_j}
   =\bigoplus_{i+t+\delta=e\atop\delta\in\{0,1\}}
                        \Lambda^i U_{N-j}\boxtimes\Lambda^t U_j,
\]
with \(0\leq i\leq N-j-1\) and \(0\leq t\leq j-1\).
All displayed summands are pairwise nonisomorphic for fixed
\(e\): the pair \((i,t)\) determines \(\delta\).

Define
\[
 D_l(A,B)=\#\{(i,t):0\leq i\leq A-1,\ 0\leq t\leq B-1,
                                                 \ i+t=l\},
\]
with \(D_l=0\) outside the allowed range. Since the hook
modules are self-dual, the dimension of the invariant space
of \([H_e]\otimes[H_f]\) under this Young subgroup is
\[
 I_j(e,f)=
 \begin{cases}
 D_e(N-j,j)+D_{e-1}(N-j,j),&e=f,\\
 D_{\min(e,f)}(N-j,j),&|e-f|=1,\\
 0,&|e-f|>1.
 \end{cases}
\]
Indeed the two restrictions share a summand precisely when
the same pair \((i,t)\) can use \(\delta,\delta'\in\{0,1\}\)
to reach \(e,f\), respectively.

Let \(M_j\) be the permutation module on \(j\)-subsets
of \([N]\). Horizontal Pieri, Macdonald I(5.16), gives
\([M_j]-[M_{j-1}]=[(N-j,j)]\) for
\(1\leq j\leq N/2\). Frobenius reciprocity therefore yields
\[
                    g((N-k,k),H_e,H_f)=I_k(e,f)-I_{k-1}(e,f),
\]
where \(I_0(e,f)=\mathbf1_{e=f}\) by irreducibility.

For \(k\geq2\), both Young subgroups have two nonempty
blocks. The generating polynomial of \(D_l(A,B)\) is
\((1-x^A)(1-x^B)/(1-x)^2\), so, with \(a=N-k\geq k\),
\[
 \begin{split}
 \sum_l\bigl(D_l(a,k)-D_l(a+1,k-1)\bigr)x^l
   &=\frac{x^{k-1}-x^a}{1-x}\\
   &=x^{k-1}+x^k+\cdots+x^{a-1}.
 \end{split}
\]
Substitution into \(I_k-I_{k-1}\) gives
\eqref{eq:F020}, with no remaining subtraction.
For \(k=1\), one has
\(D_l(N-1,1)=\mathbf1_{0\leq l\leq N-2}\).
Subtracting \(I_0\) gives \eqref{eq:F020hook}, including
the endpoints \(e=0,N-1\). The case \(k=0\) was
Schur orthogonality. This proves every stated case.
\end{proof}

\paragraph{A source-overlap check.}
The displayed expression in Rosas's Corollary~14 is
\[
 \mathbf1_{k-1\leq e\leq N-k}\mathbf1_{e=f}
  +\mathbf1_{k\leq(e+f+1)/2\leq N-k}
                         \mathbf1_{|e-f|\leq1}.
\]
It agrees with \eqref{eq:F020} when \(k\geq2\).
Its proof invokes the non-hook double-hook formula,
Theorem~13(3), p.~14. Without separating the overlap where
the two-row output is itself a hook, that printed expression
cannot be applied literally: at
\(N=4,k=1,e=f=1\) it gives two, whereas
\eqref{eq:F020hook} gives
\(g((3,1),(3,1),(3,1))=1\).
Likewise the one-row case is Schur orthogonality.
This check concerns the stated arXiv version and does not
claim an error in an unexamined later version. The explicit
endpoint formulas above remove the ambiguity.

\paragraph{Grothendieck dictionary.}
The subset modules are linearizations of actual finite sets;
the signed subset modules are exterior powers of their
permutation vector bundles. Restriction to the constant
Young subgroup scheme and finite-group invariants recover
the displayed tensor multiplicities. In SGA~3 I4.7.2 these
are comodules; their symmetric/exterior operations and
exact-sequence identities are the representation \(K_0\)
operations of SGA~6 V2.2.1 and VI3.3. The cancellation
\(I_k-I_{k-1}\) becomes an explicit nonnegative interval
count by the proved generating-polynomial identity. No
effective set lift of the hook representation is presumed.
The result supplies a positive Kronecker rule in another
special family; it does not compare the evaluation kernels
for the determinant and permanent orbit closures.

% END agents/formulas/hooks_two_row.tex

% BEGIN agents/jet_review/hessian_bracket_multiplicity.tex
\section{A multiplicity-one family detected by a Hessian bracket}

\subsection*{Finding F-024: an exact near-rectangular plethysm coefficient}
For every \(d\geq3\) and \(n\geq3\), put
\[
                     \lambda=(d(n-2)+3,3,2^{d-3}).
                                                               \tag{F-024.1}
\]
Then, in characteristic zero,
\begin{equation}\label{eq:F024ambient}
             a_\lambda(d[n])=1.
\end{equation}
The ambient vector space must have dimension at least
\(d-1=\ell(\lambda)\); otherwise this Schur module is zero.
This exact extraction is derived below from standard Cauchy
and Pieri identities, without a claim of historical priority.
It concerns the complete ambient multiplicity, rather than
only the existence of one highest-weight vector.

\paragraph{The column representation and its balanced weight.}
Write \(a=n-2\geq1\), \(D=da\), and let
\(V=\mathbb C^d\) be the column space. The Cauchy
polynomial-matrix model used in F-014 and F-016 identifies
\(a_\lambda(d[n])\) with
\[
                 \dim\left((S_\lambda V)_{(n^d)}\right)^{S_d}.
                                                               \tag{F-024.2}
\]
Indeed the degree \((n,\ldots,n)\) column multigrading
is \((\operatorname{Sym}^n U)^{\otimes d}\), and
column-permutation invariants are
\(\operatorname{Sym}^d\operatorname{Sym}^n U\).
Cauchy identifies the complete row highest-weight space
with the column module \(S_\lambda V\), preserving that
multigrading and the column permutations.
The precise standard inputs are Macdonald, I(4.3), p.~63,
and Appendix~A(5.3)--(5.6), pp.~155--156.

Pad \(\lambda\) with a final zero and twist this column
module by \(\det(V)^{-2}\). Its highest weight becomes
\[
                   \theta=(D+1,1,0^{d-3},-2),
                                                               \tag{F-024.3}
\]
and the balanced weight becomes \((a^d)\).
Every column permutation has determinant \(\pm1\),
so its action on the determinant-square twist is trivial.
Therefore (F-024.2) is exactly
\(\dim(L_\theta{}_{(a^d)})^{S_d}\), where \(L_\theta\)
is the irreducible rational \(\mathrm{GL}(V)\)-module
of highest weight (F-024.3).

\paragraph{A four-term rational character identity.}
Write \(h_j\) for \(\operatorname{Sym}^j V\),
write \(V^*\) for the dual standard module, and interpret
the following identities in the characteristic-zero
representation Grothendieck group. Dual Pieri gives
\[
\begin{split}
 [S_{(D+1,1)}V\otimes\operatorname{Sym}^2V^*]
  ={}&[L_\theta]
       +[L_{(D,1,0^{d-3},-1)}]
       +[L_{(D+1,0^{d-2},-1)}]\\
     &+[S_{(D-1,1)}V]+[h_D],\\
 [S_{(D,1)}V\otimes V^*]
  ={}&[L_{(D,1,0^{d-3},-1)}]
       +[S_{(D-1,1)}V]+[h_D],\\
 [h_{D+1}\otimes V^*]
  ={}&[L_{(D+1,0^{d-2},-1)}]+[h_D].
\end{split}\tag{F-024.4}
\]
For clarity, these are the complete dual Pieri lists, not
selected summands. In the first line, a horizontal removal
of two boxes from \((D+1,1,0,\ldots,0)\), allowing a
negative last coordinate, has decrements in the first,
second and last coordinates whose sum is two. The second
coordinate can decrease by at most one. This gives exactly
\((0,0,2),(1,0,1),(0,1,1),(2,0,0),(1,1,0)\),
the five terms displayed. All intermediate zero coordinates
stay zero by the interlacing condition. The one-box removals
give the next two lines. The assumptions \(d\geq3\) and
\(D\geq3\) ensure all of these are distinct dominant
weights, including at \(d=3\).
Dual Pieri follows from ordinary Pieri by duality and a
determinant twist; the ordinary rule is Macdonald, I(5.16)--(5.17),
p.~73. Thus no positivity conjecture is used in these identities.

Subtracting the last two lines from the first and adding
\([h_D]\) isolates \([L_\theta]\). Ordinary one-box
Pieri also gives
\[
 S_{(D+1,1)}V=h_{D+1}\otimes V-h_{D+2},\qquad
 S_{(D,1)}V=h_D\otimes V-h_{D+1}
\]
in the same group. Substitution and cancellation yield
\begin{equation}\label{eq:F024character}
 \begin{split}
 [L_\theta]
  ={}&[h_{D+1}\otimes V\otimes\operatorname{Sym}^2V^*]
   -[h_{D+2}\otimes\operatorname{Sym}^2V^*]\\
   &-[h_D\otimes V\otimes V^*]+[h_D].
 \end{split}
\end{equation}
This is an exact character identity; its use below is
justified by exactness of balanced-weight extraction and
finite-group invariants in characteristic zero.

\begin{proof}[The positive coefficient equals one]
Let \(e_1,\ldots,e_d\) be the monomial basis of \(V\),
and let \(\mathbf a=(a,\ldots,a)\). The balanced-weight
subspaces of the four actual tensor modules in
\eqref{eq:F024character} have permutation bases under \(S_d\).
The invariant dimension of such a basis is the number of
its orbits.

For the first term, a basis vector is specified by a marked
index \(j\) in the factor \(V\) and a degree-two dual
multi-index \(\beta\). Its symmetric-power exponent is
forced to be
\[
                           \alpha=\mathbf a+\beta-e_j.
\]
All entries are nonnegative because \(a\geq1\), and
\(|\alpha|=D+1\). Thus every pair \((j,\beta)\)
is allowed. There are exactly four permutation orbits:
\(\beta=2e_j\); \(\beta=2e_k\) with \(k\ne j\);
\(\beta=e_j+e_k\) with \(k\ne j\); and
\(\beta=e_k+e_l\) with \(j,k,l\) all distinct.
The last orbit exists because \(d\geq3\).
Consequently this term contributes four invariant vectors.

For the second term, only \(\beta\) is chosen, and
\(\alpha=\mathbf a+\beta\) is forced. There are two
orbits, according as its two indices are equal or distinct.
For the third term, choose the standard index \(j\) and
the dual index \(k\); the forced exponent is
\(\alpha=\mathbf a+e_k-e_j\), again nonnegative.
There are two orbits, \(j=k\) and \(j\ne k\).
The last term has the single balanced monomial
\(e_1^a\cdots e_d^a\).

The balanced weight is invariant under the permutation
normalizer of the diagonal torus. Taking its subspace is
exact on rational torus modules, and then taking \(S_d\)
invariants is exact over \(\mathbb C\). We may therefore
apply these dimensions to the character identity, obtaining
\[
                        a_\lambda(d[n])=4-2-2+1=1.
\]
This proves \eqref{eq:F024ambient} for all stated parameters.
\end{proof}

\paragraph{The unique copy has an explicit Hessian bracket.}
Put \(s=d-1\geq2\), and let \(F\) be a homogeneous
form of degree \(n\) in \(s\) variables. Define
\[
 H_F(x)=\det\operatorname{Hess}_x F,
           \qquad q=s(n-2),\qquad L=n+q-1,
\]
and the polynomial one-form, written with an auxiliary vector
\(y\),
\[
 \Theta_F(x;y)=qH_F(x)\,dF_x(y)-nF(x)\,dH_{F,x}(y).
                                                               \tag{F-024.5}
\]
It has degree \(s+1=d\) in the coefficients of \(F\),
degree \(L\) in \(x\), and degree one in \(y\).
Euler's formula gives \(\Theta_F(x;x)=0\).
Its coefficient span is consequently the image of
\[
 \det(U)^2\otimes
 \ker\bigl(\operatorname{Sym}^L U\otimes U
                       \longrightarrow\operatorname{Sym}^{L+1}U\bigr)
       =\det(U)^2\otimes S_{(L,1)}U,
                                                               \tag{F-024.6}
\]
where \(U\) is the coordinate-function representation,
and the arrow symmetrizes the two factors. The equality is
ordinary one-box Pieri. More explicitly, the Euler identity
means that the coefficient map kills the fully symmetric
summand; its quotient is identified with the displayed
kernel by that Pieri decomposition. The determinant-square covariance
follows from
\[
 H_{F\circ g}(x)=\det(g)^2H_F(gx),\qquad
 \Theta_{F\circ g}(x;y)=\det(g)^2\Theta_F(gx;gy).
\]
Thus (F-024.6) is an equivariant coefficient map into
\(\operatorname{Sym}^d\operatorname{Sym}^n U\).
The source is irreducible with highest weight
\[
  (L+2,3,2^{s-2})=(d(n-2)+3,3,2^{d-3})=\lambda.
\]
A nonzero such coefficient map has image a complete copy
of this irreducible module.

To check nonvanishing, take \(F=\sum_{i=1}^s x_i^n\),
\(x=(2,1,\ldots,1)\), and \(y=e_1\). Then
\[
 H_F(x)=[n(n-1)]^s2^{n-2}\ne0,
\]
and direct substitution gives
\[
 \Theta_F(x;e_1)
   =\frac{n(n-2)(s-1)(2^n-1)}2 H_F(x)\ne0.
                                                               \tag{F-024.7}
\]
Hence the coefficient span of \(\Theta\) is a nonzero
copy of \(S_\lambda U\). By the ambient multiplicity
calculation it is the unique copy. In particular, to prove
that the corresponding ideal multiplicity is one, or that
the corresponding coordinate-ring multiplicity is one, it
suffices to prove respectively the vanishing or nonvanishing
of this explicit bracket on the \(\mathrm{GL}(U)\)-stable
variety in question.
That final geometric test is separate from the ambient
coefficient computation.

\paragraph{Exact computation and geometric dictionary.}
As a separate finite check, the supplied script
\path{agents/jet_review/check_hessian_bracket_plethysm.py}
expands \(h_d[h_3]\) by its wreath-product cycle index
and evaluates the Schur coefficients by exact
Murnaghan--Nakayama characters. It gives one for
\(s=2,\ldots,12\), including the nine-variable case
\((13,3,2^7)\) at outer degree ten. Its beta-bead
character engine was compared with the independent connected
border-strip engine on all 919 shape/cycle pairs of sizes
zero through eight. These computations supplement the full
arbitrary-parameter proof above.

The torus grading and finite-group action in (F-024.2)
are comodules in the sense of SGA~3, I,
Propositions~4.7.2--4.7.3, pp.~34--36. The Schur and
exterior constructions have the exact-functor dictionary
of SGA~6, VI, Theorem~3.3, pp.~392--393.
The intermediate rational twist is an invertible character
of \(\mathrm{GL}_d\); it is not an effective polynomial
representation on its own. Exactness is used only in
characteristic zero, where the permutation-group invariants
are exact. The Hessian bracket is an actual polynomial
coefficient morphism, rather than a virtual class, and its
vanishing is an ideal-comodule statement. These distinctions
make the passage from the calculated ambient multiplicity
to an actual orbit-closure comparison precise.

% END agents/jet_review/hessian_bracket_multiplicity.tex

% BEGIN agents/archive/rectangle_square_known.tex
\section{An exact multiplicity-free rectangular family}

\subsection*{F-019. The full square of a two-row rectangle}
This is an extracted form of a known positive formula, not a
claim of a new Kronecker rule. Let $h\geq1$ and
$\nu\vdash2h$. If $\ell(\nu)\leq4$, pad $\nu$ by zeros
to four parts. Then
\[
 g(\nu,(h,h),(h,h))=
 \begin{cases}
 1,&\nu_1,\nu_2,\nu_3,\nu_4\text{ have the same parity},\\
 0,&\text{otherwise}.
 \end{cases}                                          \tag{F-019.1}
\]
The coefficient is zero if $\ell(\nu)>4$. Thus the positive
cases are exactly all-even row lengths (with zero parts
allowed), and four positive odd row lengths. This includes
every output shape, rather than only two-row outputs.

\paragraph{Known source and short extraction.}
The exact supplied source is Blasiak--Mulmuley--Sohoni,
GCT IV, Proposition 18.5, label
\texttt{p Kronecker consequence l1 = 0}, in
\texttt{Apr13KroneckerGCT4.tex}, lines 7806--7813. The
equivalent positive sum immediately before it, label
\texttt{e BWZ main result}, is at lines 7778--7784; the
paper attributes that sum to its reference BWZ. The column
notation is defined at lines 1535--1536 and used at lines
7453--7466. Numbering is with the
shared theorem counter in the supplied version; a current
attempt to reopen the author-hosted PDF was unavailable,
so no new PDF page verification is claimed here.

Write $c_i=\nu_i-\nu_{i+1}$, $1\leq i\leq3$.
In that source $(h,h)$ is $[h,0]$, and the coefficient of
$x^k$ in its generating function after replacing $x$ by
$x^{1/2}$ corresponds to the second two-row factor having
row difference $2k$. For a second rectangle, $k=0$.
The positive sum at that coefficient has just the term
$j=0$ in its first sum; the second sum is empty. That term
is one exactly when $c_1,c_2,c_3$ are all even, and zero
otherwise. The three differences are even precisely when
all four padded row lengths have the same parity. This
proves (F-019.1) directly from the stated known formula.

For the length assertion, realize the coefficient by
Schur--Weyl as the multiplicity of
$S_{(h,h)}\mathbb C^2\otimes S_{(h,h)}\mathbb C^2$
in $S_\nu(\mathbb C^2\otimes\mathbb C^2)$.
The latter is zero when $\ell(\nu)>4$.
This uses the characteristic-zero Schur--Weyl equivalence,
Macdonald, I, Appendix A, (5.3)--(5.6), pp.~155--156.

\paragraph{Complete stretching in this family.}
For every positive integer $t$,
\[
 g(t\nu,(th,th),(th,th))=
 \begin{cases}
 1,&t\text{ even, or the four parts of }\nu
                              \text{ already have one parity},\\
 0,&t\text{ odd and those parities are mixed},
 \end{cases}                                          \tag{F-019.2}
\]
provided $\ell(\nu)\leq4$; for greater length every stretch
is zero. This follows by multiplying all three adjacent
row differences by $t$ in (F-019.1). Thus every mixed-parity
shape of length at most four is an ordinary saturation hole,
and stretch two has coefficient exactly one. The conclusion
is consistent with, and a special exact-value case of, O-001.

\paragraph{Grothendieck dictionary and GCT scope.}
Let $G$ be the constant group scheme $S_{2h}$ over
$\mathbb Z$, as in SGA 3, I, 4.1, p.~23. Over
$\mathbb Q$, the finite-group invariant space
\[
   \bigl([\nu]\otimes[(h,h)]\otimes[(h,h)]\bigr)^G
\]
has rank exactly zero or one as specified above.
The Specht modules may be constructed by the finite-set
incidence kernels in A-018. Their tensor and Schur functors
are operations in the exact representation $K_0$, with
SGA 3, I, Proposition 4.7.2, pp.~34--35, and SGA 6, VI,
Theorem 3.3, pp.~392--393. No effective permutation lift
of an irreducible is asserted. This is an explicit positive
multiplicity rule in a rectangular subfamily; it does not
compare the evaluation kernels defining the two GCT orbit
closures and therefore supplies no additional actual
multiplicity obstruction on its own.

% END agents/archive/rectangle_square_known.tex

% BEGIN agents/archive/POSITIVE_HW_ALGEBRA.tex
% Run GCT_KRON_20261120; GPT-6 Astra Pro; effort not exposed.
\section{A concrete positive algebra and its nonprimitive weight lattice}

\subsection*{A-024. Gram minors realize the rectangle-square rule}
This is a standard invariant-theoretic specialization, reconstructed
here from F-019 with an explicit proof. No priority claim is made.
It supplies an actual positive basis satisfying the hypotheses of
A-019 and identifies precisely why ordinary saturation still fails.

Let $E,F$ be two-dimensional complex vector spaces with chosen
nonzero alternating forms, let $W=E\otimes F$, and let
$H=\mathrm{SL}(E)\times\mathrm{SL}(F)$. For $r\geq4$, put
$U=\mathbb C^r$ and
\[
 B=\operatorname{Sym}(W\otimes U),\qquad A=B^H.
\]
The group $\mathrm{GL}(U)$ acts on $A$ and commutes with $H$.
Let $U_r$ denote its upper unitriangular subgroup. There are
explicit homogeneous highest-weight invariants
\[
 \Delta_1,\quad\Delta_2,\quad\Delta_3,\quad D_4\in A^{U_r}
\]
such that
\[
 \boxed{\quad A^{U_r}=\mathbb C[\Delta_1,\Delta_2,\Delta_3,D_4].\quad}
                                                        \tag{A-024.1}
\]
This is a polynomial algebra. Their weights, padded with zeros
to $r$ coordinates, are respectively
\[
 (2),\qquad(2,2),\qquad(2,2,2),\qquad(1,1,1,1).          \tag{A-024.2}
\]
The monomials in these four generators give one basis vector
for every nonzero rectangle-square Kronecker multiplicity space.

\paragraph{Construction of the generators.}
The product of the alternating forms is a nondegenerate symmetric
form on $W$. Its inverse gives an invariant tensor in
$\operatorname{Sym}^2W$. Choose bases and denote its matrix
by $J^{-1}$. For the standard basis $u_1,\ldots,u_r$ of $U$,
write $x_{ai}=w_a\otimes u_i$ for the generators of $B$,
where $1\leq a\leq4$. Set
\[
 Q_{ij}=\sum_{a,b=1}^4(J^{-1})_{ab}x_{ai}x_{bj},\qquad
 \Delta_j=\det(Q_{uv})_{1\leq u,v\leq j}\quad(1\leq j\leq3),
\]
and let $D_4$ be the determinant of the $4$-by-$4$ matrix
$(x_{ai})_{1\leq a,i\leq4}$. The first three generators are
$H$-invariant because the inverse form is an invariant tensor.
The fourth is invariant because
\[
        \det(g\otimes h)=(\det g)^2(\det h)^2=1
                            \quad((g,h)\in H).
\]
Each generator is a nonzero polynomial: after a complex change
of basis the nondegenerate symmetric form has an orthonormal
basis, and substitution of its first four vectors gives
nonzero Gram principal minors and a nonzero determinant.

For an upper triangular matrix acting on $U$, the first $j$
columns transform among themselves. Their Gram determinant
therefore transforms by the square of the determinant of that
$j$-by-$j$ upper block. The first four-column determinant
transforms by the determinant of its upper block. They are
thus fixed by $U_r$ and have exactly the weights (A-024.2).
This also fixes the Borel convention used in the formula.

\begin{proof}[Proof of the complete positive basis]
The characteristic-zero Cauchy decomposition, followed by
Schur--Weyl restriction to the tensor product, gives
\[
 A_N=\bigoplus_{\nu\vdash N}
       (S_\nu(E\otimes F))^H\otimes S_\nu U,
\]
with only types fitting both dimensions. An irreducible
$\mathrm{GL}_2$ module $S_\alpha E$ restricts trivially to
$\mathrm{SL}_2$ exactly when $\alpha=(h,h)$; its size is
then $2h$. Consequently the displayed multiplicity is zero
for odd $N$, and for $N=2h>0$ it is
\[
       \dim(S_\nu(E\otimes F))^H
                         =g(\nu,(h,h),(h,h)).           \tag{A-024.3}
\]
At degree zero it is one. By the known rule F-019, $A$ is
therefore multiplicity-free under $\mathrm{GL}(U)$, with
support the partitions of length at most four whose four
padded parts all have the same parity.

For such a partition write
\[
 a_1=\frac{\nu_1-\nu_2}{2},\quad
 a_2=\frac{\nu_2-\nu_3}{2},\quad
 a_3=\frac{\nu_3-\nu_4}{2},\quad a_4=\nu_4.
\]
These are nonnegative integers. The polynomial
\[
                    \Delta_1^{a_1}\Delta_2^{a_2}
                                  \Delta_3^{a_3}D_4^{a_4}             \tag{A-024.4}
\]
has weight $\nu$ and is nonzero because $B$ is a domain.
An irreducible characteristic-zero $\mathrm{GL}(U)$ module
has a one-dimensional $U_r$-fixed highest-weight line.
Thus (A-024.4) spans the complete $\nu$-weight space of
$A^{U_r}$. Conversely, every $U_r$-fixed vector is a sum
of these highest-weight vectors, since the graded pieces
have the direct-sum decomposition above. This proves
generation in (A-024.1).

The four weights in (A-024.2) are linearly independent over
$\mathbb Q$. Different monomials in the four generators
therefore have different torus weights. A polynomial relation
would be a linear relation among nonzero vectors of distinct
weights and is impossible. This proves algebraic independence
and the asserted monomial basis. Its structure constants are
one, and the exponent vector is recovered from $\nu$ by the
four displayed difference formulas. The proof identifies the
whole one-dimensional multiplicity space, not just occurrence.
\end{proof}

\paragraph{Relative saturation and the precise parity defect.}
Let $S$ be this weight semigroup, written in its first four
coordinates. Its generated lattice and its cone are
\[
 \begin{split}
 L_S&=\{v\in\mathbb Z^4:v_1\equiv v_2\equiv v_3\equiv v_4
                                                       \pmod2\},\\
 C&=\{v\in\mathbb R^4:v_1\geq v_2\geq v_3\geq v_4\geq0\}.
 \end{split}
\]
Indeed the same difference formulas, now allowing arbitrary
integer exponents, characterize the lattice generated by
(A-024.2), and allowing nonnegative real exponents gives $C$.
It follows directly that
\[
                         S=C\cap L_S.                 \tag{A-024.5}
\]
In particular $S$ is saturated in $L_S$: if $v\in L_S$
and $kv\in S$ for an integer $k\geq1$, the cone inequalities
give $v\in C$, so $v\in S$.

For this Kronecker family the ordinary ambient lattice is
\[
                L_{\mathrm{even}}=
                    \{v\in\mathbb Z^4:\textstyle\sum_i v_i
                                                        \text{ is even}\},
\]
because $h=\frac12\sum_i v_i$ must be integral. Reduction
modulo two gives eight possible vectors for $L_{\mathrm{even}}$
and only $0000,1111$ for $L_S$. Hence
\[
                         [L_{\mathrm{even}}:L_S]=4.    \tag{A-024.6}
\]
The lattice is not primitive. Every mixed-parity partition in
$C\cap L_{\mathrm{even}}$ lies outside $S$, while its double
lies in $S$. For example $\nu=(2,1,1,0)$ has coefficient
zero at $h=2$, whereas $2\nu=(4,2,2,0)$ has coefficient
one at $h=4$. These are failures of saturation in the larger
ambient lattice, not holes inside the generated lattice.

For completeness the original invariant algebra $A$ is also
normal. If $z\in\operatorname{Frac}(B^H)$ is integral over
$B^H$, it is integral over the polynomial domain $B$ and
belongs to $\operatorname{Frac}(B)$. Normality of $B$ implies
$z\in B$. Being a fraction of invariant elements, it is
$H$-invariant, so $z\in B^H$. The same argument for $U_r$
gives normality of $A^{U_r}$, also apparent from (A-024.1).
Thus even an explicit normal polynomial highest-weight algebra
does not force saturation in a nonprimitive ambient lattice.

\paragraph{Grothendieck dictionary and scope.}
The groups $\mathrm{SL}_2\times\mathrm{SL}_2$, $\mathrm{GL}_r$
and $U_r$, their tensor representation, the alternating-form
tensors, and the determinant polynomial are defined over
$\mathbb Z$. In the usual alternating bases the inverse-form
matrix is integral as well. They give polynomial comodule maps
in the sense of SGA 3, I, Proposition 4.7.2, pp.~34--35;
the displayed highest weights give the torus gradings of
Proposition 4.7.3, pp.~35--36. The Schur and exterior
functors are the operations of SGA 6, VI, Theorem 3.3,
pp.~392--393. The precise characteristic-zero decomposition
input is Macdonald, I(4.3), p.~63, and Appendix A(5.3)--(5.6),
pp.~155--156, together with the supplied GCT IV formula
located and extracted in F-019.

The equality (A-024.1) and its irreducible interpretation are
asserted over $\mathbb C$. No commutation of invariants with
arbitrary integral base change or integral irreducible splitting
is claimed. This gives an explicit positive model for the
additional multiplicative basis required in A-019. It does not
produce a positive rule for general Kronecker coefficients or
compare the determinant and permanent evaluation kernels.

% END agents/archive/POSITIVE_HW_ALGEBRA.tex

% BEGIN agents/archive/HEIGHT_THREE_RECTANGLES.tex
% Run GCT_KRON_20261120; GPT-6 Astra Pro; effort not exposed.
\section{A positive rule for three equal height-three rectangles}

\subsection*{A-026. An actual Kronecker polynomial ring with a saturation hole}
For every integer $h\geq0$, with the empty-partition convention
at $h=0$, one has the positive formula
\[
 \boxed{\quad
 g((h,h,h),(h,h,h),(h,h,h))
    =\#\{(a,b,c)\in\mathbb Z_{\geq0}^3:2a+3b+4c=h\}.
 \quad}                                                 \tag{A-026.1}
\]
In particular the coefficient can be evaluated by the closed form
\[
 \begin{cases}
 \left\lfloor\dfrac{h^2+12h+48}{48}\right\rfloor,&h\text{ even},\\[3pt]
 \left\lfloor\dfrac{h^2+6h+21}{48}\right\rfloor,&h\text{ odd}.
 \end{cases}                                             \tag{A-026.2}
\]
This is an extraction from a known invariant-ring theorem,
not a claim of a new general Kronecker rule. It supplies
positive monomial bases in an unbounded family, with
unbounded multiplicities, and gives a precise limit to the
multiplicity-free construction in A-024.

\paragraph{Exact primary input and hypotheses.}
Let $E,F,G$ be complex vector spaces of dimension three,
and let
\[
       \mathcal A=\mathbb C[E\otimes F\otimes G]
                  ^{\mathrm{SL}(E)\times\mathrm{SL}(F)\times\mathrm{SL}(G)}.
\]
The established invariant-theory result is
\[
                  \mathcal A=\mathbb C[I_6,I_9,I_{12}],           \tag{A-026.3}
\]
with algebraically independent generators of the displayed
ordinary polynomial degrees. A precise primary locator is
Domokos--Drensky, \emph{Defining relation for semi-invariants
of three by three matrix triples}, arXiv:1101.3178v1
(17 January 2011), Section 2, Theorem 1 and Remark 2(ii),
printed p.~6. The characteristic-zero hypothesis is stated
at the start of Section 2, p.~3. Its generators are
$I_6=H$ from (3), p.~4, $I_9=Q$ from (4), p.~5,
and $I_{12}=\widetilde S$, defined by the ternary-cubic
substitution on p.~5. All these pages were read in the
primary PDF; pp.~5--6 were also visually checked.

The algebraic step in that proof is transparent. Its invariant
algebra is free with basis $1,Q$ over
$\mathbb C[\widetilde S,\widetilde T,H]$, and Theorem 1 gives
\[
                 Q^2=H^3+27H\widetilde S-\frac{27}{4}\widetilde T.
\]
The free-module statement makes this the only relation.
Since the coefficient of $\widetilde T$ is nonzero in
characteristic zero, elimination of $\widetilde T$ gives
the polynomial algebra in $H,Q,\widetilde S$.
This checks the deduction used in (A-026.3). The paper's
matrix-semi-invariant presentation and classical ternary-cubic
invariant theorem are known inputs; an independent re-proof
of those earlier classification results is not claimed.

\begin{proof}[Proof of the positive Kronecker formula]
In characteristic zero, reorder the tensor factors in the
$N$-fold tensor power of $E^*\otimes F^*\otimes G^*$
and take diagonal $S_N$ invariants. Schur--Weyl decomposition
in the three factors gives
\[
 \operatorname{Sym}^N(E^*\otimes F^*\otimes G^*)
  =\bigoplus_{\alpha,\beta,\gamma\vdash N}
       (S_\alpha E^*\otimes S_\beta F^*\otimes S_\gamma G^*)
                      ^{\oplus g(\alpha,\beta,\gamma)},          \tag{A-026.4}
\]
where only partitions with at most three rows occur. Here
the coefficient is the invariant dimension of the three
Specht factors; their characteristic-zero self-duality gives
the usual Kronecker coefficient.

An irreducible polynomial $\mathrm{GL}_3$ module, or its
dual, restricts trivially to $\mathrm{SL}_3$ precisely when
its highest weight has all three parts equal. Such a module
is a one-dimensional determinant power and has degree $3h$.
Taking the three special-linear invariant spaces in
(A-026.4) therefore gives
\[
           \dim\mathcal A_{3h}
                      =g((h,h,h),(h,h,h),(h,h,h)),                \tag{A-026.5}
\]
and invariant spaces in all other degrees vanish.
The fixed choices of volume forms identify the determinant
lines with the scalar field; no dimension factors are lost.

By (A-026.3), the monomials
\[
                        I_6^a I_9^b I_{12}^c
\]
are a basis of $\mathcal A$. Their ordinary degree is
$6a+9b+12c=3(2a+3b+4c)$. Taking the degree-$3h$ basis
proves (A-026.1), including every multiplicity copy.
The positive indexing set is thus explicit, and multiplication
of its basis elements adds the three exponent vectors with
structure constant one. In particular this is more data
than an effective class in a representation $K_0$.
\end{proof}

\begin{proof}[Proof of the closed form]
Let $\epsilon\in\{0,1\}$ be the parity of $h$.
In (A-026.1), $b$ has that same parity. Write
$b=\epsilon+2j$ and put $M=(h-3\epsilon)/2$.
For $M\geq0$, the equation becomes
\[
                    a+2c+3j=M,
\]
so its number of solutions is
\[
           p(M)=\sum_{j=0}^{\lfloor M/3\rfloor}
                         \left(\left\lfloor\frac{M-3j}{2}\right\rfloor+1\right).
\]
Write $M=6q+s$, $0\leq s\leq5$. Pairing the terms
with indices $2i$ and $2i+1$ gives contribution $M-6i$.
For $s=0,1,2$ there are $q$ complete pairs and one
unpaired final term $\lfloor s/2\rfloor+1$; for
$s=3,4,5$ there are $q+1$ complete pairs. Consequently
\[
 p(M)=3q^2+(s+3)q+c_s,
           \qquad(c_0,c_1,c_2,c_3,c_4,c_5)=(1,1,2,3,4,5).
\]
This equals $\lfloor(M^2+6M+12)/12\rfloor$ in all
six residues. Substitution of $M=h/2$ for even $h$ and
$M=(h-3)/2$ for odd $h\geq3$ gives (A-026.2).
At $h=1$ the solution set is empty and the displayed odd
formula gives zero; at $h=0$ both formulas and the empty
Kronecker convention give one. This checks all endpoints.
The closed form requires only integer arithmetic on numbers
whose bit length is polynomial in the bit length of $h$.
\end{proof}

\paragraph{A primitive-lattice saturation failure inside an actual Kronecker ring.}
Grade $\mathcal A$ by $h=N/3$. Its three generator weights
are $2,3,4$. The support is exactly
\[
              S=\{0\}\cup\{2,3,4,\ldots\}.                     \tag{A-026.6}
\]
Indeed even $h\geq2$ uses $a=h/2,b=c=0$, and odd
$h\geq3$ uses $b=1,a=(h-3)/2,c=0$. No nonnegative
combination of $2,3,4$ is one. Since $3-2=1$, the
generated lattice is the full primitive lattice $\mathbb Z$.
Nevertheless $1\notin S$ while $2\cdot1\in S$.
Thus ordinary saturation fails even in this primitive lattice.
Equivalently
\[
                     g((1^3),(1^3),(1^3))=0,
             \qquad g((2^3),(2^3),(2^3))=1.
\]
The polynomial invariant algebra is a normal domain, and its
affine spectrum is smooth three-dimensional affine space.
These geometric properties do not repair the missing graded
piece. There is no conflict with A-019: its one-dimensional
graded-piece hypothesis fails, already at $h=4$, where
$I_6^2$ and $I_{12}$ are independent and the coefficient
is two. In fact at $h=4t$ the monomials with
$b=0,c=0,\ldots,t,a=2t-2c$ give at least $t+1$
independent vectors, so the multiplicities are unbounded.

In the finer exponent grading the support is the saturated
monoid $\mathbb Z_{\geq0}^3$. Its projection
$(a,b,c)\mapsto2a+3b+4c$ has the support (A-026.6).
This is an actual Kronecker realization of the failure of
saturation under a grading projection, not an unrelated
weighted-ring analogy.

\paragraph{What does and does not extend from A-024.}
For $r\geq3$, let
\[
       \mathcal R_r=\operatorname{Sym}
             (\mathbb C^3\otimes\mathbb C^3\otimes\mathbb C^r)
                                ^{\mathrm{SL}_3\times\mathrm{SL}_3}.
\]
Its $S_\nu\mathbb C^r$ multiplicity in degree $3h$ is
$g(\nu,(h^3),(h^3))$, by the same Cauchy argument.
Taking $h=4,\nu=(4,4,4)$ and using (A-026.1) gives
multiplicity two. Thus its highest-weight algebra has a
two-dimensional weight space and cannot satisfy the
multiplicity-free hypothesis of A-024. No classification
or nonpolynomiality claim about the full highest-weight
algebra is inferred. The equal-rectangle subalgebra in
(A-026.3) is precisely the positive specialization proved here.

\paragraph{Grothendieck dictionary and source limits.}
The three special-linear group schemes and their tensor
representation are defined over $\mathbb Z$; their invariant
spaces are kernels of comodule maps as in SGA 3, I,
Proposition 4.7.2, pp.~34--35. The determinant-line
condition is a Schur-functor calculation, with SGA 6, VI,
Theorem 3.3, pp.~392--393, and the exact Schur--Weyl
background in Macdonald, I, Appendix A(5.3)--(5.6),
pp.~155--156. The scalar and exponent gradings are the
torus gradings of SGA 3, I, Proposition 4.7.3, pp.~35--36.
The extra invariant-ring theorem (A-026.3) is essential;
it does not follow from these $K_0$ operations alone.

The primary source is available at
\url{https://arxiv.org/pdf/1101.3178}. Its characteristic-zero
statement is all that is used here. We make no assertion
that these invariant generators or their decomposition commute
with arbitrary integral base change. This positive family
and its saturation defect do not compare determinant and
permanent evaluation kernels and assert no additional
actual GCT multiplicity obstruction.

% END agents/archive/HEIGHT_THREE_RECTANGLES.tex

% BEGIN agents/archive/CASTLING_RECTANGLE_FAMILY.tex
% Run GCT_KRON_20261120; GPT-6 Astra Pro; effort not exposed.
\section{Transporting the positive rule to rectangles of growing heights}

\subsection*{A-027. A Markov-mutation family of rectangular coefficients}
Let $\mathcal M$ be the set of positive integer triples
obtained from $(1,1,1)$ by permutations and the mutations
\[
                         (x,y,z)\longmapsto(x,y,3xy-z).
\]
Every such mutation remains positive, as proved below.
For $(x,y,z)\in\mathcal M$ and $h\geq0$, one has
\[
 \boxed{\quad
 g\bigl(((yzh)^{3x}),((xzh)^{3y}),((xyh)^{3z})\bigr)
   =\#\{(u,v,w)\in\mathbb Z_{\geq0}^3:2u+3v+4w=h\}.
 \quad}                                                   \tag{A-027.1}
\]
The three partitions have the common size $3xyz h$.
Zero widths at $h=0$ mean the empty partition. The formula
is a consequence of elementary rectangular Schur duality
and the known invariant-theory specialization A-026; no
historical-priority claim is made.

There are triples in $\mathcal M$ for which all three
rectangle heights tend to infinity. Thus this positive
family does not impose a fixed height bound. For each
fixed triple it has an ordinary saturation failure at
$h=1$, since its coefficient is zero there and one at
$h=2$. The full value is given by A-026.2 at every $h$.

\paragraph{The rectangular complement identity.}
Let $a,b,c,N$ be positive integers such that $a,b,c$ divide
$N$ and $c<ab$. Set $t=N/c$, $c'=ab-c$, and $N'=tc'$.
Then $a,b,c'$ divide $N'$ and
\[
 g\bigl(((N/a)^a),((N/b)^b),(t^c)\bigr)
  =g\bigl(((N'/a)^a),((N'/b)^b),(t^{c'})\bigr).             \tag{A-027.2}
\]
The equality retains the entire multiplicity space by a
perfect duality, after choosing determinant-line trivializations.

\begin{proof}
Let $A,B$ be complex vector spaces of dimensions $a,b$,
let $W=A\otimes B$, and let $H=\mathrm{SL}(A)\times\mathrm{SL}(B)$.
Tensor Schur--Weyl decomposition identifies the left side of
(A-027.2) with
\[
                         \dim(S_{(t^c)}W)^H.             \tag{A-027.3}
\]
Indeed an $H$-fixed irreducible factor must be a determinant
power in each of $A,B$, so its partition labels are exactly
the two rectangles displayed on the left.

Write $D=ab$. Highest-weight duality gives
\[
            S_{(t^c)}W\simeq(\det W)^t
                                      \otimes(S_{(t^{D-c})}W)^*.
                                                               \tag{A-027.4}
\]
To verify this identity directly, the highest weight of
the dual on the right before the determinant twist is
$(0^c,(-t)^{D-c})$. Adding $(t^D)$ gives
$(t^c,0^{D-c})$, the highest weight on the left.
Both rational $\mathrm{GL}(W)$ representations are
irreducible, so the highest-weight classification gives
the claimed isomorphism, equivalently a perfect pairing
with values in $(\det W)^t$.

For $(g_1,g_2)\in H$ one has
$\det(g_1\otimes g_2)=(\det g_1)^b(\det g_2)^a=1$.
Thus the determinant line in (A-027.4) is trivial under $H$.
Complete reducibility in characteristic zero makes the
induced pairing of the $H$-invariant spaces perfect: a
nontrivial irreducible summand pairs trivially with a
trivial summand, and nondegeneracy on the full modules
therefore gives nondegeneracy on their trivial isotypic parts.
Consequently (A-027.3) equals
$\dim(S_{(t^{c'})}W)^H$.

Finally $N'=tab-N$, so $N'/a=tb-N/a$ and
$N'/b=ta-N/b$ are integers. They are positive since
$N'=tc'>0$. Applying the same tensor Schur--Weyl
calculation at size $N'$ gives the right side of
(A-027.2). This proves the coefficient identity and the
assertion about all multiplicity copies.
\end{proof}

\begin{proof}[Proof of the mutation formula]
Every triple in $\mathcal M$ satisfies the Markov equation
\[
                         x^2+y^2+z^2=3xyz.              \tag{A-027.5}
\]
It holds initially, and substituting $z'=3xy-z$ leaves
the difference between the two sides unchanged. Moreover
\[
                          z(3xy-z)=x^2+y^2>0.
\]
If the old triple is positive, $z'=3xy-z$ is therefore
a positive integer. This proves positivity at every step,
without using a classification of all Markov solutions.

For $h>0$, apply (A-027.2) with
\[
 a=3x,\quad b=3y,\quad c=3z,\quad
       N=3xyz h,\quad t=xyh.
\]
The preceding positivity also gives $c<ab$, and
\[
             c'=9xy-3z=3z',\qquad N'=3xyz'h.
\]
Thus the three new rectangle widths are $yz'h,xz'h,xyh$,
exactly those prescribed by the mutated triple. Permutations
of the triple preserve $g$ by its symmetric definition as
a triple Specht invariant dimension. Induction along the
mutation sequence reduces the coefficient to
$g((h^3),(h^3),(h^3))$, which is A-026.1.
The case $h=0$ is the constant empty-partition case.
This proves (A-027.1).

For an explicit unbounded sequence, start with a sorted
triple $1\leq x\leq y\leq z$ and replace its smallest
entry, giving the new sorted triple
\[
                            (y,z,3yz-x).
\]
The new last entry is larger than $z$: using $x\leq z$
and $y\geq1$, it is at least $2z>z$.
Thus the maximum increases without bound. Two steps later
the new minimum is the old maximum, so the minimum also
tends to infinity. All three entries, hence all three
heights $3x,3y,3z$, are unbounded. The initial triples are
\[
 (1,1,1),\ (1,1,2),\ (1,2,5),\ (2,5,29),\ (5,29,433).
\]
For example the third one gives
\[
       g(((10h)^3),((5h)^6),((2h)^{15}))
                         =\#\{2u+3v+4w=h\},
\]
and the fourth gives
$g(((145h)^6),((58h)^{15}),((10h)^{87}))$ with the
same value. These examples have the required common sizes
$30h$ and $870h$, respectively.
\end{proof}

\paragraph{Positive bases, with no multiplication claim.}
Fix a mutation path and determinant-line trivializations.
The basis $I_6^uI_9^vI_{12}^w$ in A-026 transfers through
the perfect pairings (A-027.4) by taking the uniquely dual
basis at each step, together with the Schur--Weyl
identifications. It gives bases of the required coefficient
spaces indexed by the same nonnegative triples as in
(A-027.1). This is a statement about graded vector spaces
and all their multiplicity copies. Compatibility with
multiplication has not been established here. In particular
neither a polynomial invariant-ring isomorphism for these
larger formats nor such an isomorphism from equality of
Hilbert series is asserted.

\paragraph{A second growing-height family with coefficient one.}
The same complement identity also gives a positive family
without the invariant-theory input A-026. Start with the
format $(a,b,c)=(1,2,2)$ and allow permutations and
complement mutations $c\mapsto ab-c$ whenever the new
coordinate is positive. Every resulting format has at least
two even coordinates: this holds initially, and a mutation
of an odd coordinate keeps it odd, while a mutation of an
even coordinate keeps it even. Therefore the three pairwise
products $ab,ac,bc$ are even, and for every $h\geq0$
\[
       g\left(((bch/2)^a),((ach/2)^b),((abh/2)^c)\right)=1.
                                                               \tag{A-027.6}
\]
At the starting format this is
$g((2h),(h,h),(h,h))=1$, since tensoring with the trivial
$S_{2h}$ representation preserves the irreducible $(h,h)$.
For a complement mutation the common size is $N=abc h/2$
and the width in the changed factor is $t=ab h/2$.
Thus (A-027.2) proves (A-027.6) by induction.

All three heights can again grow. Replacing the smallest
coordinate in sorted order gives
\[
              (1,2,2),\ (2,2,3),\ (2,3,4),\ (3,4,10),\ (4,10,37),\ldots.
\]
The first two changes increase the maximum. Thereafter the
middle entry $b$ is at least three and
$bc-a\geq(b-1)c\geq2c>c$. Positivity is therefore
automatic on this explicit sequence, and the same two-step
minimum argument proves that all heights tend to infinity.
For instance $(2,3,4)$ gives
$g(((6h)^2),((4h)^3),((3h)^4))=1$ for every $h$.
The scalar support in this family is all of
$\mathbb Z_{\geq0}$, in contrast with the missing first
degree in (A-027.1). Only allowed positive mutations are
claimed here; some other mutations of $(1,2,2)$ produce zero
and are excluded by the definition.

\paragraph{Grothendieck dictionary and GCT scope.}
The tensor representation and determinant-line pairing are
polynomial comodule constructions, as in SGA 3, I,
Propositions 4.7.2--4.7.3, pp.~34--36. The underlying
duality pairs complementary exterior powers into the
determinant line; (A-027.4) is its rectangular Schur-functor
form. The exact $K_0$/Schur background is SGA 6, VI,
Theorem 3.3, pp.~392--393, and the tensor Schur--Weyl
decomposition is Macdonald, I, Appendix A(5.3)--(5.6),
pp.~155--156. The proof above establishes the complement
identity directly, so no unverified historical castling
theorem number is used. Its semisimple invariant-space
argument is explicitly in characteristic zero.

The family answers a positive Kronecker question with
growing rectangle heights and transfers explicit saturation
holes. The two rectangular factors need not be equal,
as they are in the standard determinant-stabilizer
Kronecker coefficient. No determinant/permanent evaluation
kernel inclusion, actual multiplicity obstruction, or
complexity separation follows from these identities alone.

% END agents/archive/CASTLING_RECTANGLE_FAMILY.tex

% BEGIN agents/archive/ORBIT_CLOSURE_GAP.tex
% Run GCT_KRON_20261120; GPT-6 Astra Pro; effort not exposed.
\section{An explicit unbounded gap between the determinant orbit and its closure}

\subsection*{A-028. Rectangular examples where symmetric Kronecker multiplicities lose copies}
For an integer \(h\geq2\), put
\[
 \rho_h=(h,h,h),\qquad \mu_h=(\underbrace{3,\ldots,3}_{h}),\qquad
 p_h=\#\{(a,b,c)\in\mathbb Z_{\geq0}^3:2a+3b+4c=h\}.
\]
Here the determinant has size \(n=h\), and the outer polynomial
 degree is \(d=3\). Then
\begin{align}
 g(\rho_h,\mu_h,\mu_h)&=p_h,                         \tag{A-028.1}\\
 \operatorname{sk}(\rho_h,\mu_h)&=
 \begin{cases}p_h&h\text{ even},\\0&h\text{ odd},\end{cases}
                                                               \tag{A-028.2}\\
 K_h(\rho_h)=a_{\rho_h}(3[h])&=
 \begin{cases}1&h\text{ even},\\0&h\text{ odd}.\end{cases}
                                                               \tag{A-028.3}
\end{align}
The symmetric Kronecker term means the multiplicity of the Specht
module \([\rho_h]\) in \(\operatorname{Sym}^2[\mu_h]\).
It is the determinant orbit multiplicity in the established
orbit-stabilizer formula recalled in B\"urgisser--Ikenmeyer--Panova,
\emph{No occurrence obstructions in geometric complexity theory},
Section 1, subsection ``Related work: Kronecker coefficients,
paragraph defining the GCT-coefficients (supplied source
\texttt{main.tex}, lines 415--428). Consequently, for every
even \(h\geq4\),
\[
 K_h(\rho_h)=1<\operatorname{sk}(\rho_h,\mu_h)
                       =g(\rho_h,\mu_h,\mu_h)=p_h.              \tag{A-028.4}
\]
These gaps are unbounded: \(p_{4t}\geq t+1\).
This is a concrete extraction from known invariant and plethysm
inputs, with the transpose action checked below; historical priority
is not asserted. It exhibits why an orbit multiplicity cannot be
substituted for the actual closure multiplicity.

\paragraph{The transpose action on the known invariant generators.}
Use the algebra \(\mathcal A\) and generators \(H,Q,\widetilde S\)
from A-026. Identify all three three-dimensional vector spaces with
\(\mathbb C^3\), with fixed ordered bases. Regard a tensor as a
triple of \(3\)-by-\(3\) matrices \((A_1,A_2,A_3)\), and let
\(\tau\) simultaneously transpose the three matrices. This exchanges
the first two tensor factors and preserves \(\mathcal A\).
We claim
\[
              \tau H=H,\qquad\tau Q=-Q,\qquad
                     \tau\widetilde S=\widetilde S.             \tag{A-028.5}
\]
The degree-six and degree-nine components of the polynomial algebra
\(\mathcal A=\mathbb C[H,Q,\widetilde S]\) are respectively the
one-dimensional spaces spanned by \(H\) and by \(Q\). Thus \(\tau\)
acts on each by a scalar, necessarily \(\pm1\).

The primary source explicitly evaluates these two generators on
skew-symmetric matrix triples: Domokos--Drensky,
\emph{Defining relation for semi-invariants of three by three matrix
triples}, arXiv:1101.3178v1, Section 2, printed p.~5, the displayed
skew-symmetric triple immediately before the page break to p.~6.
Writing its three columns of skew coordinates as a matrix \(M\),
the displayed values are \(H=(\det M)^2\) and \(Q=(\det M)^3\).
Take \(M=I_3\). Both values are nonzero, while transposition sends
the entire triple to its negative. Homogeneity therefore multiplies
\(H\) by \((-1)^6=1\) and \(Q\) by \((-1)^9=-1\).
This determines the first two scalars in (A-028.5) on their entire
one-dimensional invariant spaces, not merely at this point.

For the last identity, each coefficient \(f_{i,j,k}\) of
\(\det(t_1A_1+t_2A_2+t_3A_3)\) is unchanged by simultaneous
transposition, since determinant is unchanged by transpose.
The definition of \(\widetilde S\) in the same source, Section 2,
printed p.~5, is a polynomial in these coefficients. It too is
unchanged. This completes the verification of (A-028.5).

\begin{proof}[Proof of all three multiplicity formulas]
By A-026, the degree-\(3h\) invariant space has basis
\[
                    H^a Q^b\widetilde S^c,
                    \qquad 2a+3b+4c=h.
\]
Every such \(b\) has the parity of \(h\), so (A-028.5) shows
that \(\tau\) acts on the entire space by \((-1)^h\).

In the three-factor Schur--Weyl decomposition (A-026.4), this
invariant space is the multiplicity space
\[
                  ([\rho_h]\otimes[\rho_h]\otimes[\rho_h])^{S_{3h}}.
                                                               \tag{A-028.6}
\]
The determinant lines in the three Schur factors are identified with
\(\mathbb C\) by the fixed ordered volume forms. Exchanging the
first two such lines acts trivially on their tensor product. These
are ordinary tensor spaces, with no graded or exterior sign in the
interchange. Thus the induced action on (A-028.6) is precisely the
flip of its first two Specht factors. Its fixed space has dimension
\(p_h\) for even \(h\), and zero for odd \(h\).

Conjugating a partition twists its Specht module by the sign
representation. Since \(\mu_h=\rho_h'\), one has
\([\mu_h]\cong[\rho_h]\otimes\operatorname{sgn}\).
The natural cancellation of the two sign factors gives an
\(S_{3h}\)-equivariant isomorphism
\[
 [\mu_h]\otimes[\mu_h]\ \cong
 [\rho_h]\otimes[\rho_h]
\]
which commutes with the flip. In particular their symmetric squares
are isomorphic as \(S_{3h}\)-modules. Taking the \([\rho_h]\)
multiplicity, using characteristic-zero self-duality of Specht modules,
proves (A-028.2). Taking the full tensor multiplicity proves
(A-028.1), also directly the usual double-conjugation invariance
of Kronecker coefficients. The precise classical sign-twist and
Schur--Weyl inputs are the character involution in Macdonald,
Chapter I, Section 7, equations (7.3)--(7.6), pp.~113--114, and
Appendix A(5.3)--(5.6), pp.~155--156.

For the ambient coefficient, apply F-011 with \(d=3,n=k=h\).
Its partition \((d(n-k)+k,k,k)\) is exactly \(\rho_h\),
and its value is one for even \(h\), zero for odd \(h\).
Since \(\rho_h\) has three rows and \(h\geq2\), O-007 gives
\(K_h(\rho_h)=a_{\rho_h}(3[h])\). This proves (A-028.3).
O-007 uses Beauville, \emph{Determinantal hypersurfaces},
Proposition 3.1(a), printed p.~46; all of that input's hypotheses
and the degree-two endpoint were checked there.

For even \(h\geq4\), the two solutions
\((a,b,c)=(h/2,0,0)\) and \(((h-4)/2,0,1)\) are distinct,
so \(p_h\geq2\). For \(h=4t\), the \(t+1\) choices
\(b=0,c=0,\ldots,t,a=2t-2c\) prove the unbounded gap.
All uses of the invariant-ring theorem, ordinary/symmetric tensor
operations, and ambient-to-closure comparison are thereby separate.
\end{proof}

\paragraph{Consequences and their scope.}
This family is stronger than an example with \(K<g\) caused only
by the antisymmetric summand: at every even index all the \(p_h\)
copies already belong to the symmetric Kronecker term, while exactly
one extends to the determinant closure. The difference
\(p_h-1\) is a loss of orbit functions at the boundary.
It supplies no actual permanent obstruction: \(q\leq a=1=K\)
for even \(h\), and \(q=0\) for odd \(h\). Thus ordinary and
symmetric Kronecker coefficients, their positive monomial bases,
and their saturation behavior still do not compute the rank needed
for a general actual multiplicity obstruction.

\paragraph{Grothendieck dictionary.}
The tensor representations and the transposition automorphism are
maps of affine group-scheme representations over \(\mathbb Z\).
Their invariant spaces are kernels of coactions, as in SGA 3, I,
Proposition 4.7.2, pp.~34--35. The Schur decompositions and
symmetric-square summand are the characteristic-zero realization
of Schur operations on locally free modules from SGA 6, VI,
Theorem 3.3, pp.~392--393; the displayed calculation is at the
level of actual modules and their explicit involution, not only
of virtual classes. The degree gradings are SGA 3, I,
Proposition 4.7.3, pp.~35--36.

Let \(\Omega_h\) be the scheme-theoretic determinant orbit closure
and \(\mathcal O_h\) its dense orbit. Restriction is the injective
map of algebras
\(\mathbb C[\Omega_h]\hookrightarrow\mathbb C[\mathcal O_h]\).
The multiplicity subspaces of type \(\rho_h\) have respective
dimensions one and \(p_h\) at even \(h\). Therefore \(p_h-1\)
independent orbit copies do not extend as regular functions across
the closure. This boundary-extension condition is additional
geometric data, invisible in a positive class of \(K_0\).
The orbit-closure ideal and its graded quotients can be placed in
the Hilbert/Quot framework of FGA, Bourbaki expos\'e 221,
Theorem 3.1, p.~259, and Theorem 3.2, p.~260; representability
alone does not force these restriction maps to be surjective.
All multiplicity assertions here are in characteristic zero;
no arbitrary integral base-change assertion is made for invariants.

% END agents/archive/ORBIT_CLOSURE_GAP.tex

% BEGIN proofs/saturation_and_boundary_translation.tex
\section{Saturation of determinant support and extension across the boundary}

\subsection*{M-020. Two precise statements that do not compare the original ranks}
There is a known saturation theorem for the actual determinant closure,
as well as a known localization theorem relating its orbit and closure.
Their short consequences below specify what they give in the report's
notation. They must be distinguished from ordinary Kronecker saturation,
and from a comparison of permanent and determinant multiplicities.
No historical priority is asserted for these elementary consequences.

\paragraph{Known support theorem and a finite uniform-stretch reduction.}
Let
\[
 S_n=\{\lambda:K_n(\lambda)>0\},\qquad
 T_n=\{\lambda:\ell(\lambda)\leq n,\ n\mid |\lambda|\}.
\]
Zero is included in both monoids. B\"urgisser--Ikenmeyer--H\"uttenhain,
\emph{Permanent versus determinant: not via saturations},
arXiv:1501.05528v2, Theorem~1, printed p.~3, proves
\[
                         T_n\subseteq\operatorname{Sat}(S_n)
                         \qquad(n>2).                         \tag{M-020.1}
\]
Here saturation is taken in the group generated by the monoid,
as defined immediately before that theorem. In particular each
\(\lambda\in T_n\) has some positive integral stretch in \(S_n\).
The primary source is \url{https://arxiv.org/pdf/1501.05528v2}.
This is a cited theorem, not a result derived here from Burnside
effectivity or K-theoretic positivity.

For clarity, there is even a uniform stretching integer for each
fixed \(n\):
\[
 \text{there exists }c_n\geq1\text{ such that }
            K_n(c_n\lambda)>0\quad\text{for every }\lambda\in T_n.
                                                               \tag{M-020.2}
\]
One finite reduction for choosing it uses
\[
 H_n=\{\theta\ne0:\ell(\theta)\leq n,\ \theta_1\leq n,
                                      \ n\mid |\theta|\}.
                                                               \tag{M-020.3}
\]
For each \(\theta\in H_n\), choose any \(c_\theta\geq1\)
with \(K_n(c_\theta\theta)>0\), which exists by (M-020.1).
Then \(c_n=\operatorname{lcm}_{\theta\in H_n}c_\theta\)
works. This specifies a finite gate; it supplies no numerical
bound on the choices \(c_\theta\).

\begin{proof}[The finite reduction]
Regard a partition in \(T_n\) as a multiset of columns with
heights in \(\{1,\ldots,n\}\). Any sequence of at least \(n\)
columns contains a nonempty consecutive subsequence of length
at most \(n\) whose total height is divisible by \(n\): among
the first \(n+1\) partial sums, including zero, two agree
modulo \(n\). Remove that subsequence. Its column heights,
sorted, define an element of \(H_n\); removal preserves
divisibility of the remaining total height. Continue until
fewer than \(n\) columns remain, which themselves give an
element of \(H_n\), unless the remainder is empty.
Thus every element of \(T_n\) is a sum of elements of \(H_n\).
The determinant closure is irreducible. Its coordinate ring
is therefore a domain, and products of nonzero highest-weight
vectors show that \(S_n\) is closed under addition. Multiples
of \(c_\theta\theta\) and then sums of those multiples prove
(M-020.2) for the stated least common multiple.
\end{proof}

For internal padding, \(q^m_\lambda(d[n])>0\) implies
\(\ell(\lambda)\leq m^2\) and \(|\lambda|=nd\).
Consequently, when \(n\geq m^2\) and \(n>2\), every possible
permanent type already belongs to \(\operatorname{Sat}(S_n)\).
Replacing determinant support by its saturation therefore erases
any occurrence distinction in this range. This statement gives
neither \(K_n(\lambda)>0\) at the original weight nor an inequality
between two positive multiplicities. It also does not say that
every sufficiently large integer stretch is positive: a uniform
multiple and all sufficiently large integers are different assertions.
In the stronger BIP range the separately cited no-occurrence theorem
does give original-weight positivity; the positive-rank comparison
still remains.

\paragraph{Known boundary localization and its all-copy consequence.}
Put \(R=n^2\), \(G=\operatorname{GL}_R\),
\(X=\overline{G\cdot\det_n}\), and \(O=G\cdot\det_n\).
Write \(A=\mathbb C[X]\) and \(B=\mathbb C[O]\).
B\"urgisser--Ikenmeyer, \emph{Fundamental invariants of orbit closures},
arXiv:1511.02927v2, Definition~3.8 and Proposition~3.9,
printed pp.~9--10, supply a homogeneous nonzero fundamental
\(\operatorname{SL}_R\)-invariant \(\Phi\in A\) such that
\[
             O=X\setminus V(\Phi),\qquad B=A[\Phi^{-1}].
                                                               \tag{M-020.4}
\]
The source's determinant polystability is Corollary~2.9,
printed pp.~6--7. The primary PDF is
\url{https://arxiv.org/pdf/1511.02927v2}.
Write \(e=\deg\Phi\). In the polynomial coordinate convention
used here, \(\Phi\) spans a one-dimensional representation
of highest weight \(\delta\mathbf1_R\), where
\[
                 \delta\in\mathbb Z_{>0},\qquad e=n\delta.
                                                               \tag{M-020.5}
\]
Indeed an \(\operatorname{SL}_R\)-invariant polynomial
\(\operatorname{GL}_R\)-character is a nonnegative determinant
power, and comparison with the scalar action gives
\(R\delta=ne\). No explicit value of \(e\) is assumed.

Fix \(\lambda\vdash nd\), and let
\(b_\lambda=[\mathbb S_\lambda\mathbb C^R:B]\), which is
the symmetric Kronecker orbit multiplicity in this report.
Then
\[
 K_n(\lambda+k\delta\mathbf1_R)
 \ \text{is nondecreasing in }k\geq0,\quad
 K_n(\lambda+k\delta\mathbf1_R)\leq b_\lambda,
                                                               \tag{M-020.6}
\]
and, for some \(k_0=k_0(n,\lambda,\Phi)\),
\[
             K_n(\lambda+k\delta\mathbf1_R)=b_\lambda
                         \qquad\text{for every }k\geq k_0.
                                                               \tag{M-020.7}
\]
The outer degree on the left is \(d+ke\). This is a
rectangular translation of the weight, not its dilation.

\begin{proof}[The complete multiplicity-space statement]
Let \(A^U_\mu\) and \(B^U_\mu\) be the highest-weight spaces
for a fixed maximal unipotent subgroup \(U\). Multiplication
by \(\Phi\) is injective because \(A\) is a domain, and shifts
the weight by \(\delta\mathbf1_R\). Division by \(\Phi^k\)
is defined in \(B\) and identifies the image of
\(A^U_{\lambda+k\delta\mathbf1_R}\) with a subspace of
\(B^U_\lambda\). These subspaces are nested as \(k\) grows.
They have dimensions equal to the left side of (M-020.6).

The orbit highest-weight space is finite-dimensional; in this
case its dimension is the known symmetric Kronecker coefficient.
Choose a basis \(f_1,\ldots,f_{b_\lambda}\). By (M-020.4),
for each basis vector some \(\Phi^{k_i}f_i\) belongs to \(A\).
It remains a highest-weight vector, of the shifted weight,
because \(\Phi\) is a semi-invariant. Taking
\(k_0=\max_i k_i\) makes every basis vector occur in the
divided image for every \(k\geq k_0\). The upper bound in
(M-020.6) then gives equality. When \(b_\lambda=0\), take
\(k_0=0\); both sides are zero for all \(k\).
\end{proof}

For example, F-025's type has orbit multiplicity ten and
closure multiplicity five. It follows that its translated
closure multiplicities eventually equal ten, although no
value of the first successful translation is supplied here.
This argument does not determine the dilates in A-030.

There is a direct padding limitation. If \(m<n\) and \(k>0\),
then \(\lambda+k\delta\mathbf1_{n^2}\) has exactly \(n^2\)
nonzero parts. Thus, at its corresponding degree,
\[
 q^m_{\lambda+k\delta\mathbf1_{n^2}}((d+ke)[n])=0,
                                                               \tag{M-020.8}
\]
by the permanent's \(m^2\) essential-variable bound. Hence
clearing determinant-boundary denominators in this manner
leaves the possible internally padded permanent support.
The localization theorem does not compare the original two
positive ranks and is not a way to transfer an unpadded
obstruction to the BIP range.

\paragraph{Grothendieck dictionary.}
The support statement concerns the character grading on
\(A^U\); its saturation is a lattice-and-cone operation,
not the class of a virtual representation in \(K_0\).
The coordinate-comodule and weight-space dictionary is
SGA~3, I, 4.7.2--4.7.3, pp.~34--36.
The boundary statement is the distinguished affine open
\(D(\Phi)\subseteq\operatorname{Spec}A\), with coordinate
ring \(A_\Phi\); it uses ordinary localization and the
finite-dimensionality of each Schur multiplicity space.
The Schur operations themselves are those of SGA~6, VI,
Theorem~3.3, pp.~392--393. The proof explicitly clears
denominators of a basis of actual sections. It is therefore
a statement about extension of sections across the boundary,
with an unspecified extension threshold, rather than an
effectivity consequence of a virtual \(\lambda\)-identity.

% END proofs/saturation_and_boundary_translation.tex

% BEGIN agents/obstructions/obstructions_findings.tex
% GCT_KRON_20261120, saturation/obstructions worker.
% All original arguments in this file are findings of this run, not claims of
% historical priority. Known results are identified as such.

\section{Saturation: precise negative and positive statements}

Throughout this part, coefficients and decomposition multiplicities are taken
over $\mathbb C$.  Write $R_{h,d}=(d^h)$ and
$g_{h,d}(\lambda)=g(\lambda,R_{h,d},R_{h,d})$.  Put
$\rho(N)=(N-|\rho|,\rho)$ only when this is a partition.
``Ordinary saturation'' means the implication
\[
 g(t\lambda,t\mu,t\nu)>0\text{ for some }t\geq1
 \quad\Longrightarrow\quad g(\lambda,\mu,\nu)>0.
\]
This must be distinguished from the residue-class saturation index in
Mulmuley's revised GCT VI, Definition~1.2.4 (arXiv:0704.0229v4,
printed p.~14; PDF p.~17), and Hypothesis~3.3.2 (printed p.~51;
PDF p.~54). Identically zero constituents of a quasipolynomial do not
violate that definition.  The revised hypothesis asks for a polynomial
bound on the index, and a generic index-zero statement; it does not assert
ordinary saturation.

\subsection{O-001: exact holes and a factor two for two two-row factors}

The input here is a known theorem, not a new positivity rule: Blasiak,
Mulmuley and Sohoni, GCT IV, Proposition~18.6, equations~(175)--(180),
supplied \texttt{Apr13KroneckerGCT4.tex} version, pp.~128--129 in
the matching author-hosted PDF. The proposition itself has no source
label; the six equation labels are \texttt{e g = 0 1} through
\texttt{e g = 0 6}. The proposition applies to
$\ell(\lambda),\ell(\mu)\leq2$ and $\ell(\nu)\leq4$.
Use column counts
\[
 A=\lambda_1-\lambda_2,\quad B=\mu_1-\mu_2,
 \quad c_i=\nu_i-\nu_{i+1}\ (1\leq i\leq4),\quad\nu_5=0.
\]
All three partitions have the same size.  The known classification says
that the coefficient vanishes if and only if at least one of the following
conditions holds:
\begin{align}
 &(A+B)/2>c_1+c_2+c_3; \tag{O-001.a}\\
 &|A-B|/2>\min(c_1,c_3)+c_2; \tag{O-001.b}\\
 &c_1=c_3=0,\qquad (A+B)/2\not\equiv c_2\pmod2; \tag{O-001.c}\\
 &\min(A,B)=0,\quad \min(c_1,c_2,c_3)=0,
       \quad\max(A,B)/2\not\equiv c_1+c_2\pmod2; \tag{O-001.d}\\
 &\{A,B\}=\{0,2\},\qquad c_1,c_2,c_3\text{ all even};\tag{O-001.e}\\
 &A=B=0,\qquad c_1\text{ or }c_2\text{ odd}. \tag{O-001.f}
\end{align}
The halves in the congruences are integers whenever the other conditions
in that line hold; for instance $A$ and $B$ have the parity of the common
size, and $\min(A,B)=0$ forces both to be even.

\paragraph{Finding O-001.}
In this entire $2\times2\times4$ family, a positive stretch exists if and
only if neither (O-001.a) nor (O-001.b) holds, and this is equivalent to
\[
                         g(2\lambda,2\mu,2\nu)>0.
\]
Consequently, the ordinary-saturation holes are exactly the triples
satisfying the two weak cone inequalities and at least one of
(O-001.c)--(O-001.f).

\begin{proof}
The quantities $A,B,c_i$ all multiply by the stretching factor.  A failure
of either homogeneous cone inequality therefore persists under every
stretch.  At stretch two, (O-001.c), (O-001.d), and (O-001.f) are false
because all relevant integers in their congruences are even.  Condition
(O-001.e) is false: it would require $\{A,B\}=\{0,1\}$ before stretching,
which is impossible because $A$ and $B$ have the same parity.  The known
classification now gives the asserted equivalence and the exact list of
holes.  This proves a factor-two statement, not saturation at factor one.
\end{proof}

For three two-row partitions an even simpler complete description is
available.  Write their second rows as $0\leq c\leq b\leq a\leq N/2$.
Rosas, \emph{The Kronecker product of Schur functions indexed by two-row
shapes or hook shapes}, Corollary~5 (arXiv:math/0001084, p.~11), gives
\[
 g=\#\left\{j\in\mathbb Z:
 \max(0,\lceil(a+b+c-N)/2\rceil)\leq j
                  \leq\lfloor(b+c-a)/2\rfloor\right\}.
\]
Thus a positive stretch exists exactly when $a\leq b+c$; within this cone
the holes are exactly $2a=N$ and $b+c-a$ odd.  To check the last assertion,
put $D=b+c-a\geq0$ and $E=N-2a\geq0$: the interval is
$[\max(0,\lceil(D-E)/2\rceil),\lfloor D/2\rfloor]$.
It is nonempty for $E\geq1$, and for $E=0$ precisely when $D$ is even.
The rectangular parity ray
\[
            g((k,k),(k,k),(k,k))=\begin{cases}1&k\text{ even},\\0&k\text{ odd}
                                      \end{cases}
\]
is already Rosas's Example~6, p.~12. Its revised GCT VI saturation index
is zero.  By contrast, Briand--Orellana--Rosas, Theorem~2.4,
equation~(2.5), gives the strong-saturation counterexamples
\[
\begin{gathered}
 \alpha=(k,k),\quad\beta=(k+1,k-1),\quad
 \gamma=(2k-2i-2j,2i,2j),\quad i>j>0,\ k>2i+j,\\
 g(t\gamma,t\alpha,t\beta)=
 \begin{cases}t/2+1&t\text{ even},\\(t-1)/2&t\text{ odd}.
 \end{cases}
\end{gathered}
\]
Its smallest example is $(\alpha,\beta,\gamma)=((6,6),(7,5),(6,4,2))$.
Here the constituent containing $t=1$ is not identically zero.

\subsection{O-002: exact stretching on long-row rectangular holes}

\paragraph{Finding O-002.}
For $h\geq2$, $d\geq1$, put $\eta=(hd-1,1)$.  For every integer
$t\geq1$,
\[
 g(t\eta,(td)^h,(td)^h)
       =[z^t]\prod_{j=2}^{h}\frac1{1-z^j}.                 \tag{O-002.1}
\]
In particular $g(\eta,R_{h,d},R_{h,d})=0$ and
$g(2\eta,R_{h,2d},R_{h,2d})=1$.  For $h=2$ the support of the ray is
the positive even integers; for $h\geq3$ it is every integer $t\geq2$.
The formal $t=0$ value is one.

\begin{proof}
Use the known rectangular/two-row identity (Pak--Panova, \emph{Strict
unimodality of $q$-binomial coefficients}, Lemma~3, p.~2), in the form
proved in F-002:
\[
 g((htd-t,t),(td)^h,(td)^h)=p_t(h,td)-p_{t-1}(h,td).
\]
The two-row shape is valid because $hd\geq2$. Here $p_s(h,td)$ counts
partitions of $s$ with at most $h$ rows, each of size at most $td$.
Since $td\geq t$, the width bound is inactive for $s=t,t-1$.
Conjugating partitions shows that the generating series is
$\prod_{j=1}^{h}(1-z^j)^{-1}$; multiplication by $1-z$ proves
(O-002.1).  The claimed supports follow from the additive semigroup
generated by $2,\ldots,h$. If $h\geq3$, an even integer $t\geq2$ is a
sum of twos, while an odd integer $t\geq3$ is three plus a sum of twos.
The coefficient at $t=2$ is one.
\end{proof}

This yields holes at arbitrarily large rectangle height, not only at
height two.  If $m\geq2$, the partition $\eta=(hd-1,1)$ satisfies both
coarse necessary padded-permanent conditions
\[
                 \ell(\eta)\leq m^2,
       \qquad \eta_1\geq (h-m)d.
\]
Indeed $\ell(\eta)=2$ and the second inequality is $md\geq1$.
These conditions therefore do not imply ordinary saturation even when
$h$ exceeds any prescribed polynomial in $m$.  They do not make $\eta$
an occurrence candidate: the ambient plethysm coefficient
$a_{\eta}(d[h])$ is zero, as O-005 proves for every nontrivial hook.

\subsection{O-003: a normal, positive, pure geometric model with a hole}

\paragraph{Finding O-003.}
For $h\geq2$, conjugation invariants of a traceless matrix satisfy
\[
 \mathbb C[\mathfrak{sl}_h]^{\mathrm{GL}_h}
       =\mathbb C[c_2,\ldots,c_h],\qquad \deg c_j=j,       \tag{O-003.1}
\]
where $c_j$ are the coefficients of its characteristic polynomial.
This gives a geometric realization of (O-002.1). For $h\geq3$ the
support monoid of the displayed grading is
$\{0,2,3,4,\ldots\}$; its generated group is $\mathbb Z$, and it is
not saturated in that group.  Nevertheless the invariant quotient is a
smooth affine space, its coordinate ring is normal, and multiplication
in the monomial basis has nonnegative structure constants.

\begin{proof}
Let $D$ be the diagonal traceless matrices.  Restriction of invariants to
$D$ is injective: the matrices with pairwise distinct eigenvalues form a
dense open subset of $\mathfrak{sl}_h$, and each is conjugate to a point
of $D$.  Its image lies in $\mathbb C[D]^{S_h}$.
The fundamental theorem of symmetric polynomials, followed by imposing
$e_1=0$, gives
\[
 \mathbb C[D]^{S_h}=\mathbb C[e_2,\ldots,e_h].
\]
For completeness, invariants commute with this quotient: averaging a
polynomial representative over $S_h$ gives a symmetric representative
of each invariant class, and a symmetric polynomial divisible by $e_1$
has symmetric quotient. The functions $e_2,\ldots,e_h$ are algebraically
independent after $e_1=0$ and are restrictions of $c_2,\ldots,c_h$.
This proves (O-003.1) and its graded Hilbert series.

Manivel's rectangular stability theorem (\emph{On rectangular Kronecker
coefficients}, Theorem~1; reproduced as Ikenmeyer--Panova Theorem~2.1)
identifies the coefficient on the left of (O-002.1) with
$\dim(\operatorname{Sym}^t\mathfrak{sl}_h)^{\mathrm{SL}_h}$: transpose
both rectangles if necessary, and use the inequality $td\geq t$.
The trace pairing identifies the adjoint representation with its dual,
and $\mathrm{GL}_h$ and $\mathrm{SL}_h$ have the same conjugation
invariants.  Hence (O-003.1) independently gives the Hilbert series in
(O-002.1).

For $h\geq3$, degrees two and three occur and degree one does not.
Their generated group contains $3-2=1$, so the support is not saturated
even in its own group. A polynomial ring is normal, and its affine
spectrum is smooth. The product of two monomials is a monomial with
coefficient one. These statements prove all the asserted properties.
\end{proof}

\paragraph{Integral and cohomological dictionary.}
The weighted affine scheme
\[
       Q_h=\operatorname{Spec}\mathbb Z[c_2,\ldots,c_h]
\]
and the $\mathbb G_m$-action $c_j\mapsto u^jc_j$ are defined over
$\mathbb Z$. SGA~3, Expos\'e~I, Proposition~4.7.3 (original pp.~35--36)
identifies the grading with this diagonalizable-group action; the
following remark makes the correspondence for algebras and affine
schemes explicit. Its complex generic quotient interpretation is
(O-003.1); no assertion about formation of matrix invariants in all bad
characteristics is needed. For a finite field $\mathbb F_q$ and
$\ell\nmid q$, its underlying affine space has
\[
 H_c^i((Q_h)_{\overline{\mathbb F}_q},\mathbb Q_\ell)=
 \begin{cases}\mathbb Q_\ell(-(h-1))&i=2(h-1),\\0&\text{otherwise}.
 \end{cases}
\]
This follows from SGA~4, Expos\'e~XV, Corollary~2.2, p.~192
(affine-space acyclicity for invertible torsion), and
Expos\'e~XVIII, Theorem~3.2.5, p.~585 (Poincar\'e duality),
followed by the compatible $\ell$-adic limit.
Thus even this pure Tate cohomology does not remove the degree-one hole.
The $\mathbb G_m$-weight of a regular function and a Frobenius weight
on cohomology are different gradings. The same example can be phrased
after base change to a field as $M_c(Q_h)=\mathbb L^{h-1}$ in a
theory of motives with compact supports and affine-space homotopy
invariance.

This is a counterexample to the inference ``normality, positive
structure constants, or purity alone force saturation of degree
support.'' It is not a counterexample to the correct normal-toric
statement about a semigroup algebra: the grading here identifies
distinct monomials, so $\mathbb C[c_2,\ldots,c_h]$ is not the semigroup
algebra of its one-dimensional degree support.

\subsection{O-004: all stretches at least two in a stable rectangular range}

\paragraph{Finding O-004.}
Let $h\geq7$, let $\rho$ be a partition with $\ell(\rho)<h^2$, and
let $d\geq\max(1,|\rho|)$.  Then $\lambda=\rho(hd)$ is a partition and
\[
          g(t\lambda,(td)^h,(td)^h)>0\qquad(t\geq2).       \tag{O-004.1}
\]
For $\rho=\varnothing$ the coefficient is one.  This is a short
strengthening of the even-tail conclusion stated in
Ikenmeyer--Panova Corollary~5.3, obtained from their
Proposition~5.2.  It is not claimed to be a new theorem of the literature.

\begin{proof}
Use $a_\sigma(h)$ for the stable rectangular invariant multiplicity
in Ikenmeyer--Panova Theorem~2.1.  Proposition~5.2 states that for
$h\geq7$, $0\leq k<h^2$, and every $t>1$,
\[
      g\bigl(t(h^2-k,1^k),(th)^h,(th)^h\bigr)>0.
\]
The tail of its first argument is $(t^k)$. By monotonicity as the
rectangle width increases, followed by stability, it follows that
$a_{(t^k)}(h)>0$ for every $k<h^2$.

Write $\rho$ as the rowwise sum of its columns:
$\rho=\sum_{j=1}^{\rho_1}(1^{\rho'_j})$.
The stable multiplicities have the semigroup property. One can see
this without an additional theorem: choose a width large enough that
the two positive coefficients are stable, apply the ordinary
Kronecker semigroup property to their first-row-extended shapes, and
then pass to the stable limit of the sum. Thus the already proved
positivity of every $a_{(t^{\rho'_j})}(h)$ implies
$a_{t\rho}(h)>0$.

Finally, $td\geq t|\rho|=|t\rho|$. Manivel stability, with the two
rectangle sides exchanged, therefore identifies
$g((t\rho)(htd),(td)^h,(td)^h)$ with $a_{t\rho}(h)$.
Since $(t\rho)(htd)=t\lambda$, this is (O-004.1).
\end{proof}

This concerns the stable-tail range $|\bar\lambda|\leq d$. The
padded-permanent condition only gives $|\bar\lambda|\leq md$, so this
is not a universal factor-two theorem for every GCT candidate.
The source's separate Corollary~5.4 asserts that for rectangle height
$h\geq7$, the group generated by the positive rectangular semigroup
contains every partition of size divisible by $h$ and length at most
$h^2$, the ambient representation-theoretic length bound. Together with its
full rational cone, this says its saturation is as large as possible;
it does not say the original semigroup is saturated. O-002 supplies
explicit holes in that maximal saturation.

\section{From formulas to genuine multiplicity statements}

Let $V=\mathbb C^{n^2}$, let
$\Omega_n=\overline{\mathrm{GL}(V)\cdot\det_n}$, and let
$Z_{n,m}=\overline{\mathrm{GL}(V)\cdot X_{11}^{n-m}\operatorname{per}_m}$.
The padding variable is internal: $X_{11}$ is one of the $m^2$
permanent variables, exactly as in both supplied IP and BIP sources.
An external padding variable would change the elementary length
bound to $m^2+1$ and is not silently substituted in their stated
thresholds. Put
\[
\begin{split}
 a_\lambda(d[n])&=\operatorname{mult}_{S_\lambda V}
                         \operatorname{Sym}^d(\operatorname{Sym}^nV),\\
 K_n(\lambda)&=\operatorname{mult}_{S_\lambda V}
                         \mathbb C[\Omega_n]_d,\\
 q^m_\lambda(d[n])&=\operatorname{mult}_{S_\lambda V}
                         \mathbb C[Z_{n,m}]_d.
\end{split}
\]
Here $|\lambda|=nd$.  The relevant chain is
\[
 K_n(\lambda)\leq\operatorname{sk}(\lambda,(d^n))
                  \leq g_{n,d}(\lambda),\qquad
 q^m_\lambda(d[n])\leq a_\lambda(d[n]).                    \tag{O-005.0}
\]
The middle term is the symmetric rectangular Kronecker coefficient,
the multiplicity in the coordinate ring of the determinant orbit;
it is not automatically the closure multiplicity. The chain and
interpretation are stated in the supplied B\"urgisser--Ikenmeyer--Panova
source, Section~1(b), and refer there to BLMW, equation~(5.2.7).

\subsection{O-005: exterior operations exclude every nontrivial hook}

The vanishing statement below is known: Giannelli,
\emph{On the decomposition of the Foulkes module}, Theorem~1.1;
the exterior-power identities are his Lemmas~3.2--3.3 and
Proposition~3.4. We record a short $\lambda$-ring form of its proof,
because this identifies exactly what the Burnside-to-representation
construction supplies here.

\paragraph{Finding O-005 (the $\lambda$-ring dictionary for hook vanishing).}
Let $P_N$ be the permutation representation of $S_N$ on $N=nd$
letters and $U_N=P_N-1$ its augmentation representation. If
$H=S_n\wr S_d$, then
\[
       \dim(\lambda^k U_N)^H=0\quad(1\leq k<N),
       \qquad \dim(\lambda^0 U_N)^H=1.                  \tag{O-005.1}
\]
By Schur--Weyl duality this says
\[
               a_{(nd-k,1^k)}(d[n])=0\quad(1\leq k<nd).
\]
Thus the nontrivial hook formulas for rectangular Kronecker coefficients
provide neither occurrence nor multiplicity obstructions for these
orbit closures: the permanent-side multiplicity already vanishes.

\begin{proof}[Proof of the $\lambda$-ring form]
Assume first $n,d\geq2$ and let $B=(S_n)^d\triangleleft H$.
As an $H$-module, $U_N$ is the direct sum of the block-constant
augmentation space $U_d$ and the direct sum $T$ of the $d$ within-block
augmentation spaces $U_n^{(i)}$.  The base group $B$ acts trivially on
$U_d$ and independently on the summands of $T$.

For $S_s$, the only invariant in the exterior algebra of $U_s$ has
degree zero. Indeed, $\Lambda^r P_s$ has no invariant for $r\geq2$:
a transposition of two members of an $r$-subset fixes that subset but
negates its exterior basis vector. For $r=1$ the invariant dimension
is one. The identity
$\lambda_t(P_s)=(1+t)\lambda_t(U_s)$ now proves the claim degree by
degree, using semisimplicity in characteristic zero.
The exterior-power direct-sum decomposition consequently gives
\[
       (\Lambda^\bullet U_N)^B=\Lambda^\bullet U_d.
\]
Taking invariants under $H/B=S_d$ gives (O-005.1).  If $n=1$ or
$d=1$, the wreath product is $S_N$ and the same claim applies directly.

Finally, $\lambda^k U_N=[nd-k,1^k]$, and the natural decomposition
\[
 \operatorname{Sym}^d\operatorname{Sym}^n V
       =(V^{\otimes nd})^{S_n\wr S_d}
       =\bigoplus_{\lambda\vdash nd}S_\lambda V
                         \otimes[\lambda]^H
\]
identifies its multiplicity with the left side of (O-005.1).
The quotient defining a coordinate ring cannot create new
characteristic-zero constituents, proving the obstruction conclusion.
\end{proof}

Two further known ambient restrictions are useful when testing
candidates. Giannelli's Proposition~2.7 gives
$a_\lambda(d[n])=0$ when $\ell(\lambda)>d$. His Theorem~1.2
also excludes many partitions near hooks: write
\[
 \lambda=(nd-k-|\alpha|,\alpha_1+1,\ldots,\alpha_t+1,1^{k-t}),
 \qquad t\leq k,\qquad b=\sum_{j=2}^{t}\alpha_j.
\]
For a nonempty partition $\alpha$, if $k>b$ and
$\alpha_1<(k-b)(k-b+1)/2$, then $a_\lambda(d[n])=0$.
These are known vanishing criteria, not new claims about saturation.

\subsection{O-006: full determinant multiplicities in two-row shapes}

\paragraph{Finding O-006.a (subspace criterion).}
Let $X\subseteq\operatorname{Sym}^n V^*$ be a closed
$\mathrm{GL}(V)$-stable subscheme over $\mathbb C$. Suppose it contains
the reduced linear space $\operatorname{Sym}^n U^*$ for some
$r$-dimensional quotient $V\twoheadrightarrow U$.
Then the coordinate-ring quotient loses no multiplicity with
$\ell(\lambda)\leq r$:
\[
 \operatorname{mult}_{S_\lambda V}\mathbb C[X]_d
                         =a_\lambda(d[n])
              \quad\text{if }\ell(\lambda)\leq r.       \tag{O-006.1}
\]

\begin{proof}
There is a decomposition natural in linear maps of vector spaces,
\[
 F(V):=\operatorname{Sym}^d\operatorname{Sym}^n V
       =\bigoplus_{\lambda\vdash nd}S_\lambda V\otimes M_\lambda,
 \qquad M_\lambda=[\lambda]^{S_n\wr S_d}.
\]
This follows by applying wreath-product invariants to Schur--Weyl
decomposition of $V^{\otimes nd}$; the factor $M_\lambda$ is independent
of $V$. The polynomial-functor decomposition and its naturality are
Macdonald, \emph{Symmetric Functions and Hall Polynomials}, second
edition, Chapter~I, Appendix~A, equations~(5.3)--(5.4), p.~155, and
(5.5)--(5.6), p.~156. Let $I(X)_d\subseteq F(V)$ be the degree-$d$ ideal.
Semisimplicity gives
$I(X)_d=\bigoplus_\lambda S_\lambda V\otimes J_\lambda$
for subspaces $J_\lambda\subseteq M_\lambda$.
Restriction to the displayed linear subspace is $F(\pi)$, where
$\pi:V\twoheadrightarrow U$. It acts on the $\lambda$ summand as
$S_\lambda(\pi)\otimes\operatorname{id}_{M_\lambda}$.
The map $S_\lambda(\pi)$ is surjective, since a linear section of
$\pi$ gives a right inverse after applying the Schur functor.
Its target is nonzero precisely when $\ell(\lambda)\leq r$.
Since every element of $I(X)_d$ restricts to zero, these facts force
$J_\lambda=0$ for such $\lambda$. This proves (O-006.1), including
its assertion about all copies of an irreducible, not just occurrence.
\end{proof}

\paragraph{Finding O-006.b.}
For $n\geq2$, $d\geq1$, and $\lambda\vdash nd$ with at most two
rows,
\[
 K_n(\lambda)=a_\lambda(d[n])
             =\operatorname{sk}(\lambda,(d^n))
             =g_{n,d}(\lambda).                         \tag{O-006.2}
\]
In particular, no such $\lambda$ can be a multiplicity obstruction
between the padded permanent and the determinant, for any $m<n$.

\begin{proof}
The polynomial map that substitutes an arbitrary linear endomorphism
of $V$ into $\det_n$ takes every endomorphism into $\Omega_n$:
$\mathrm{GL}(V)$ is dense in $\operatorname{End}(V)$, and a polynomial
that vanishes on its image also vanishes on the image of every
endomorphism. A binary homogeneous form of degree $n$ over
$\mathbb C$ factors into $n$ linear forms. It is the determinant of
the diagonal matrix with those forms on its diagonal, hence belongs
to $\Omega_n$. Consequently $\Omega_n$ contains every linear space
of binary forms, and (O-006.1) with $r=2$ proves $K_n=a_\lambda$.

For $\lambda=(nd-k,k)$ with $0\leq k\leq nd/2$, the weight
$(nd-k,k)$ in $\operatorname{Sym}^d\operatorname{Sym}^n\mathbb C^2$
has dimension $p_k(d,n)$: a monomial is a multiset of $d$ numbers
from $0,\ldots,n$ with sum $k$. The irreducible
$S_{(nd-j,j)}\mathbb C^2$ has each weight with second coordinate
$j,j+1,\ldots,nd-j$ once. Subtracting successive weight dimensions
therefore gives
\[
                  a_{(nd-k,k)}(d[n])=p_k(d,n)-p_{k-1}(d,n).
\]
The same difference is the rectangular Kronecker coefficient by
F-002 (the identity of Pak--Panova, Lemma~3). Thus $a_\lambda=g$.
The inequalities (O-005.0) imply the equality with the symmetric
coefficient as well. Finally $q^m_\lambda\leq a_\lambda=K_n$,
which excludes a multiplicity obstruction.
\end{proof}

\subsection{O-007: the same exclusion extends to three rows}

\paragraph{Finding O-007.}
For $n\geq2$ and $\ell(\lambda)\leq3$,
\[
                         K_n(\lambda)=a_\lambda(d[n]).   \tag{O-007.1}
\]
Thus a multiplicity obstruction for these two orbit closures must have
at least four rows. This strengthens the two-row exclusion; equality
with the ordinary rectangular Kronecker coefficient is not asserted
for three rows.

\begin{proof}
Beauville, \emph{Determinantal hypersurfaces}, Proposition~3.1(a),
printed p.~46, states that a smooth plane curve of degree $n$ with
a non-effective line bundle of degree $g-1$ has an $n\times n$
linear determinantal equation. Such a bundle exists over $\mathbb C$:
for $g\geq1$ the effective classes form the proper theta locus in
$\operatorname{Pic}^{g-1}(C)$; for $g=0$, take
$\mathcal O_{\mathbb P^1}(-1)$. The properness assertion also follows
from $\dim C^{(g-1)}=g-1<g=\dim\operatorname{Pic}^{g-1}(C)$.
For the relative Jacobian and theta locus see Beauville's following
paragraph~(3.2), on the same page.

Hence every smooth ternary form of degree $n$ is a determinant of
linear forms, up to a nonzero scalar that can be absorbed into a row
of the matrix. The substitution argument in O-006.b puts these
forms in $\Omega_n$. Smooth ternary forms are a dense open subset
of the space of all ternary forms (the discriminant is nonzero, as
the Fermat form verifies in characteristic zero). Since $\Omega_n$
is closed, it contains that entire space. Apply O-006.a with $r=3$.
\end{proof}

\subsection{O-008: full ambient multiplicity first fails in four rows}

\paragraph{Finding O-008.}
For every $n\geq4$, there exist $d\geq1$ and a partition
$\lambda\vdash nd$ with exactly four rows such that
\[
                         K_n(\lambda)<a_\lambda(d[n]).   \tag{O-008.1}
\]
Thus the universal equality in O-007 cannot be extended to every
four-row partition. This is an existence assertion about the
determinant ideal; it does not assert a permanent-versus-determinant
multiplicity obstruction.

\begin{proof}
Fix a four-dimensional quotient $V\twoheadrightarrow U$ and a linear
section. Let $L=\operatorname{Sym}^n U^*$, viewed as a linear subspace
of $\operatorname{Sym}^n V^*$, and let $r$ be the corresponding
linear retraction to $L$. Define
\[
 \delta_U:\operatorname{Mat}_n(\mathbb C)^4\longrightarrow L,
 \qquad (A_0,A_1,A_2,A_3)\longmapsto
                   \det(x_0A_0+x_1A_1+x_2A_2+x_3A_3),
\]
and put $D_U=\overline{\operatorname{im}\delta_U}$.
Every element of the image is an endomorphism substitution into
$\det_n$, so $D_U\subseteq\Omega_n\cap L$.
Conversely, restricting any endomorphism substitution to $L$ gives
a point of $\operatorname{im}\delta_U$. Since $r$ is continuous and
$\Omega_n$ is the closure of the endomorphism-substitution image,
$r(\Omega_n)\subseteq D_U$. A point of $\Omega_n\cap L$ is fixed by
$r$, proving
\[
                         \Omega_n\cap L=D_U.            \tag{O-008.2}
\]

We give an explicit parameter count for $D_U$, avoiding any assumption
that a dimension bound for an orbit quotient also controls its closure.
On the dense open set where $A_0$ is invertible, write
$B_i=A_0^{-1}A_i$ and $c=\det A_0$. On a further dense open set,
$B_1$ has distinct eigenvalues, so simultaneous conjugation makes it
diagonal. We can also assume that all entries $(B_2)_{1j}$ for
$j>1$ in an eigenbasis are nonzero. This condition is independent of
scaling eigenvectors, and holds on a nonempty open set after the
finite choice of ordering eigenvalues; for example take $B_1$
diagonal with distinct entries and $B_2$ with all entries nonzero.
Conjugation by a diagonal matrix then makes
$(B_2)_{1j}=1$ for $j>1$, while preserving the diagonal form of $B_1$.

Consequently a dense set of the image is parametrized by
\[
 c\det\bigl(x_0 I+x_1\operatorname{diag}(b_1,\ldots,b_n)
                                      +x_2 C+x_3 E\bigr),
 \qquad C_{1j}=1\ (j>1).
\]
The number of parameters is
$1+n+(n^2-n+1)+n^2=2n^2+2$.
The coefficients of these forms are polynomials in that many
parameters. Their generated field therefore has transcendence degree
at most $2n^2+2$, proving
$\dim D_U\leq2n^2+2$.
Equivalently, the original left-right action of
$\mathrm{SL}_n\times\mathrm{SL}_n$ has generic orbit dimension
$2n^2-2$: a tuple beginning with $I$, a diagonal matrix with distinct
entries, and a matrix with all entries nonzero has only a finite
simultaneous centralizer. The explicit parameter count already
suffices for the asserted bound.

For $n\geq4$,
\[
 \dim L-(2n^2+2)
 =\binom{n+3}{3}-2n^2-2
 =\frac{(n-1)(n-2)(n-3)}6>0.
\]
Thus $D_U$ is a proper closed cone. There is a nonzero homogeneous
polynomial $f\in\mathbb C[L]_d$ vanishing on it, with $d\geq1$.
The polynomial $f\circ r$ is nonzero and lies in $I(\Omega_n)_d$,
because $r(\Omega_n)\subseteq D_U$.
The injection
$\operatorname{Sym}^d\operatorname{Sym}^n U\hookrightarrow
\operatorname{Sym}^d\operatorname{Sym}^n V$ induced by the chosen
section has image only in isotypic components $S_\lambda V$ with
$\ell(\lambda)\leq4$, by the same functorial Schur decomposition as
O-006.a. At least one such component of $f\circ r$ is nonzero in the
ideal. O-007 excludes every component with at most three rows.
Therefore a component with exactly four rows occurs in the ideal,
which is precisely (O-008.1).
\end{proof}

\paragraph{Where the remaining multiplicity gate lies.}
The general sufficient inequality
$q^m_\lambda(d[n])>g_{n,d}(\lambda)$ implies a determinantal lower
bound (Ikenmeyer--Panova Theorem~1.4). It is stronger than the actual
multiplicity inequality $q^m_\lambda(d[n])>K_n(\lambda)$.
Ikenmeyer--Panova Proposition~2.8 already shows
\[
       q^m_\lambda(d[n])\leq a_\lambda(d[n])\leq g_{n,d}(\lambda)
                 \quad\text{if }n\geq|\bar\lambda|.
\]
Its Corollary~1.9 consequently requires $d>n/m$ for a
Kronecker-upper-bound multiplicity obstruction satisfying
$|\bar\lambda|\leq md$. This does not exclude an actual obstruction
caused by loss of multiplicity from the determinant orbit to its
closure. O-006 and O-007 do exclude that possibility in the specified
short-row families, because the closure has full ambient multiplicity.

\subsection{O-009: ambient-positive padded shapes already have positive $g$}

\paragraph{Finding O-009.}
Let $m\geq3$, $n>3m^4$, and $\lambda\vdash nd$. If
\[
 \ell(\lambda)\leq m^2,\qquad |\bar\lambda|\leq md,
                     \qquad a_\lambda(d[n])>0,
\]
then $g_{n,d}(\lambda)>0$.
This is a useful explicit variant of the known Ikenmeyer--Panova
no-occurrence theorem: membership in the actual permanent coordinate
ring is not required.

\begin{proof}
If $g_{n,d}(\lambda)=0$, then $a_\lambda(d[n])>g_{n,d}(\lambda)$.
Ikenmeyer--Panova Corollary~1.9 gives
$d>n/m>3m^3$. Their Theorem~1.10(a) excludes the exceptional tail set
\[
 X=\{(1),(1,1),(1,1,1,1),(1,1,1,1,1,1),(2,1),(3,1)\},
\]
because the ambient multiplicity is positive. All the hypotheses
of their Theorem~1.10(b) now hold, giving $g_{n,d}(\lambda)>0$,
a contradiction.
\end{proof}

Thus in this range adding ambient occurrence to the coarse padding
conditions eliminates the holes themselves. O-002 shows why that
additional hypothesis cannot be dropped. This statement does not
compare positive multiplicities on the two orbit closures, and it
does not prove ordinary saturation outside its stated region.

\subsection{O-010: a dimension bound on universal projection arguments}

\paragraph{Finding O-010.}
Let $r\geq3$, $n\geq m\geq1$, and fix a nonzero linear form $L$
on $\mathbb C^r$. If every form
$L^{n-m}f$, with $f$ an arbitrary homogeneous degree-$m$ form in
$r$ variables, is in the closure of the set of determinants of
$n\times n$ matrices of linear forms in those variables, then
\[
       \binom{m+r-1}{r-1}\leq(r-2)n^2+2.                \tag{O-010.1}
\]
Consequently, for a fixed positive integer $c$, a universal assertion
of this kind with $n=m^c$ and a fixed number $r$ of variables for
all sufficiently large $m$ requires $r\leq2c+1$. In particular
$r\leq51$ is necessary when $c=25$.

\begin{proof}
For $n\geq2$, apply the generic gauge from O-008 to an $r$-tuple
of matrices. After making its first matrix invertible and its
second diagonal, and normalizing $n-1$ entries of its third, the
remaining $r-3$ matrices are unrestricted. The number of parameters is
\[
       1+n+(n^2-n+1)+(r-3)n^2=(r-2)n^2+2.
\]
Thus the closed determinantal locus has dimension at most the
right side of (O-010.1). For $n=1$ the locus is the space of all
linear forms and has dimension $r=(r-2)n^2+2$, so the same bound
holds. Multiplication by the nonzero polynomial $L^{n-m}$ is an
injective linear map on the polynomial ring, so the space of forms
in the hypothesis is a linear space of dimension
$\binom{m+r-1}{r-1}$. Its containment proves (O-010.1).
For $n=m^c$, the left side is a polynomial in $m$ of degree $r-1$
with positive leading coefficient $1/(r-1)!$, while the right side
has degree $2c$. The inequality for all large $m$ therefore forces
$r-1\leq2c$.
\end{proof}

This bounds arguments that replace projected permanents by every
form of the same degree and number of variables. It does not bound
arguments using properties specific to the permanent, whose
projections occupy a much smaller parameter space. Combined with
the constructive bounds of this run, for $c\geq6$ it identifies
the optimal asymptotic fixed-variable threshold $r=2c+1$ for that
universal containment method. It is not a limitation theorem for
geometric complexity theory as a whole.

\subsection{O-011: an exact multiplicity-one family detected by Hessians}

\paragraph{Finding O-011.}
Let $n>m\geq2$, and, for $1\leq d\leq n^2$, put
\[
              \lambda_d=(d(n-2)+2,2^{d-1})\vdash nd.
\]
Then the actual closure and ambient multiplicities are
\[
 K_n(\lambda_d)=a_{\lambda_d}(d[n])=1,\qquad
 q^m_{\lambda_d}(d[n])=
 \begin{cases}1&d\leq m^2,\\0&m^2<d\leq n^2.\end{cases}       \tag{O-011.1}
\]
In particular this entire family, including outputs with as many
as $m^2$ rows and positive permanent multiplicity, has no actual
multiplicity obstruction. No bound such as $n\geq m^{25}$ or
the sufficient weight-box inequality is needed for this family.

\begin{proof}
For a form $F$ of degree $n$ on $V$, define the degree-$d$
polynomial function
\[
 W_d(F)=\det\left(\frac{\partial^2F}{\partial x_i\partial x_j}
                                  (e_1)\right)_{1\leq i,j\leq d}.
\]
For an upper triangular $g\in\mathrm{GL}(V)$, write $g_d$ for
its leading $d\times d$ block. Since $ge_1=g_{11}e_1$, the chain
rule and homogeneity give
\[
 W_d(g^{-1}\cdot F)
       =g_{11}^{d(n-2)}\det(g_d)^2 W_d(F).
\]
Here $g^{-1}\cdot F=F\circ g$, consistently with the action on
the coordinate ring. Hence $W_d$ is a highest-weight vector of
weight $\lambda_d$, provided it is nonzero. The ambient
multiplicity is one by F-011 with $k=2$: in the Cauchy
decomposition on $d$ columns, the relevant multiplicity module is
\[
 S_{\lambda_d}\mathbb C^d
       =\det^2\otimes\operatorname{Sym}^{d(n-2)}\mathbb C^d.
\]
Its column weight $(n,\ldots,n)$ is the one-dimensional space
spanned, after removing the determinant factor, by
$\prod_i z_i^{n-2}$, and column permutations fix it. The case
$d=1$ is the same calculation directly.

At the identity matrix, the Hessian of $\det_n$ is the bilinear form
\[
               (A,B)\longmapsto\operatorname{tr}(A)\operatorname{tr}(B)
                                      -\operatorname{tr}(AB).
\]
It is nondegenerate: the trace-zero subspace is nondegenerate for
$-\operatorname{tr}(AB)$, its orthogonal complement is the span
of $I_n$, and the latter has square $n(n-1)\ne0$.
A nondegenerate symmetric bilinear form and a vector with nonzero
square admit a nondegenerate $d$-dimensional subspace containing
that vector for every $1\leq d\leq n^2$; diagonalize its
orthogonal complement to see this. Choose coordinates taking
$e_1$ to $I_n$ and the first $d$ coordinate directions to such
a subspace. Then $W_d$ is nonzero on a translate of $\det_n$.
Thus $K_n(\lambda_d)\geq1$, and the ambient multiplicity-one
bound gives $K_n(\lambda_d)=1$.

It remains to prove nonvanishing on the internally padded permanent.
Set $s=n-m>0$, $P=\operatorname{per}_m$, and $F=X_{11}^sP$.
Work first at the all-ones matrix. Put $h=(m-2)!$, let $u$ be
the all-ones vector of length $m^2$, let $e$ be the coordinate
vector for the $(1,1)$ entry, and put
\[
 B=(J_m-I_m)\otimes(J_m-I_m).
\]
Direct differentiation of the permanent gives
$P=m!$, $\nabla P=(m-1)!u$, and $H_P=hB$ at that point.
Therefore
\[
 \frac{H_F}{h}
   =B+s(m-1)(eu^T+ue^T)+s(s-1)m(m-1)ee^T.              \tag{O-011.2}
\]
The eigenvalues of $B$ are $(m-1)^2$, $-(m-1)$ and $1$, with
multiplicities $1$, $2(m-1)$ and $(m-1)^2$, respectively, so it
is invertible. Its inverse gives the following $2\times2$ Gram matrix:
\[
 G=\begin{pmatrix}e^TB^{-1}e&e^TB^{-1}u\\
                  u^TB^{-1}e&u^TB^{-1}u\end{pmatrix}
   =\frac1{(m-1)^2}
       \begin{pmatrix}(m-2)^2&1\\1&m^2\end{pmatrix}.
\]
Apply the matrix determinant lemma to (O-011.2), with
\[
 C=\begin{pmatrix}s(s-1)m(m-1)&s(m-1)\\s(m-1)&0\end{pmatrix}.
\]
Expanding $\det(I_2+CG)$ and factoring yields the exact identity
\[
 \frac{\det H_F}{\det H_P}
  =\frac{m+s-1}{m-1}\bigl(1-(m^2-3m+1)s\bigr)
  =\frac{n-1}{m-1}\bigl(1-(m^2-3m+1)(n-m)\bigr).       \tag{O-011.3}
\]
For positive integral $s$ this is nonzero except when $(m,n)=(3,4)$:
the integer $m^2-3m+1$ is $-1$ for $m=2$, $1$ for $m=3$,
and at least $5$ for $m\geq4$.

The exceptional point does not cause an exceptional coefficient.
For $m=3,n=4$, evaluate instead at the matrix with $X_{11}=2$
and every other entry equal to one. Then $F=16$. In the
orthonormal basis of sums and differences under interchange of the
last two rows and last two columns, the Hessian has blocks
\[
 A=\begin{pmatrix}
 4&2\sqrt2&2\sqrt2&10\\
 2\sqrt2&0&4&2\sqrt2\\
 2\sqrt2&4&0&2\sqrt2\\
 10&2\sqrt2&2\sqrt2&4
 \end{pmatrix},\qquad
 C_0=\begin{pmatrix}0&-2\sqrt2\\-2\sqrt2&-4\end{pmatrix}
 \quad\text{twice},\qquad (4).
\]
For clarity the four invariant vectors for the first block are
$e_{11}$, $(e_{12}+e_{13})/\sqrt2$,
$(e_{21}+e_{31})/\sqrt2$, and
$(e_{22}+e_{23}+e_{32}+e_{33})/2$; replacing sums by the
row or column differences gives the two sign blocks and the
double-sign line. Direct differentiation gives the displayed
blocks. Their determinants are $576,-8,-8,4$, so
$\det H_F=147456\ne0$.

In every case we have found a point $v$ in the $m^2$-dimensional
permanent variable space with $F(v)\ne0$ and nondegenerate Hessian.
Euler's identity gives
\[
                          H_F(v,v)=n(n-1)F(v)\ne0.
\]
As for the determinant, choose any $d\leq m^2$ and a
nondegenerate $d$-subspace containing $v$, and then an invertible
change of ambient coordinates. This gives a translate of $F$
on which $W_d$ is nonzero. Hence $q^m_{\lambda_d}(d[n])\geq1$;
the ambient bound gives equality. For $d>m^2$, the length
$\ell(\lambda_d)=d$ violates the necessary internal-padding
length condition (IP Theorem~1.8), so $q^m_{\lambda_d}(d[n])=0$.
This proves every case of (O-011.1).
\end{proof}

The construction is the determinant of a restriction of the second
jet of the universal degree-$n$ section at a chosen point. It is
equivariant under the Borel preserving that point and flag; its
weight is the displayed determinant-line character. Its derivatives,
matrix entries and identities before the convenient orthonormal
change of basis are all defined over $\mathbb Z$. The multiplicity
interpretation uses the characteristic-zero Schur decomposition.
This places the explicit calculation in the equivariant locally
free-sheaf and Schur-functor dictionary of SGA~3, I, 4.7.2--4.7.3,
SGA~6, VI, Theorem~3.3, and cote~133, scan~028, without asserting
the same semisimple multiplicities in residue characteristics.

\section{Padding from a bounded degree preserves full multiplicity}

\subsection{O-012.a: a determinant bound obtained by splitting degree}

For $b\geq1$ and $r\geq2$, set
\[
 u=\left\lfloor\frac{b-1}{2}\right\rfloor,\qquad
 v=\left\lceil\frac{b-1}{2}\right\rceil,\qquad
 B(b,r)=\binom{r-1+u}{u}+\binom{r-1+v}{v}-1.
\]
Every degree-$b$ homogeneous form $f$ in $r$ variables, with a
chosen nonzero linear form $L$, admits a homogeneous linear
determinant of size $B(b,r)$ equal to
\[
                           L^{B(b,r)-b}f.               \tag{O-012.1}
\]
Consequently $L^{n-b}f\in\Omega_n$ whenever $n\geq B(b,r)$
and $r\leq n^2$. In one variable a diagonal determinant of size
$b$ suffices. These are integral identities in universal
coefficient variables.

\begin{proof}
Use $L$ and $z_1,\ldots,z_{r-1}$ as coordinates, and write
$f=\sum_{|\gamma|\leq b}c_\gamma z^\gamma L^{b-|\gamma|}$.
Make two copies of the monomial tree on the \emph{same} $r-1$
variables. The first has vertices of total degree at most $u$
and edges pointing away from its root; the second has vertices
of total degree at most $v$ and edges pointing toward its root.
The parent of a nonzero multi-index is obtained by subtracting
the last nonzero coordinate, so paths are unique. Label the
variable edges by $-z_j$, and identify the two roots into one
central vertex. Keep all other vertices distinct, even when
their multi-index labels are the same. There are $B=B(b,r)$ vertices.

Fix once and for all a rule that splits every multi-index of
degree at most $u+v=b-1$ as $\alpha+\beta$, with
$|\alpha|\leq u$, $|\beta|\leq v$; taking an initial segment
of its sorted variable word gives such a rule. For each term
with $|\gamma|\leq b-1$, insert the coefficient edge from
the prefix vertex $\alpha$ to the suffix vertex $\beta$ with
label $c_\gamma L$. For $|\gamma|=b$, remove one variable
$z_j$ by a fixed rule, split $\gamma-e_j=\alpha+\beta$,
and put label $c_\gamma z_j$ on that coefficient edge.
If several labels have the same endpoints, add them in the
matrix entry. All entries remain homogeneous linear forms.
Put an $L$-loop at every noncentral vertex; the central entry
is exactly the sum of its coefficient-loop labels.

As in M-007, every non-loop directed cycle uses the central
vertex and exactly one coefficient edge. Thus a determinant
term chooses one coefficient-edge summand and the unique
variable paths to and from it, with loops at all unused vertices.
For $|\gamma|=k<b$, the cycle has $k+1$ edges. Its sign
$(-1)^k$ cancels the sign on the $k$ variable edges, and its
contribution is $c_\gamma z^\gamma L^{B-k}$.
For $|\gamma|=b$, the cycle has $b$ edges, of which $b-1$
are tree edges and the last carries $c_\gamma z_j$.
The signs again cancel and the contribution is
$c_\gamma z^\gamma L^{B-b}$.
Summation gives (O-012.1). Since $r\geq2$, the two trees have
at least $u+1$ and $v+1$ vertices, so $B\geq u+v+1=b$;
all exponents are nonnegative. Adjoin the block $L I_{n-B}$
when a larger size $n$ is required.
\end{proof}

This complements the variable-splitting bound M-007. Splitting
degree is more efficient when the base degree is fixed and the
number of variables grows. One may use the smaller of the two
proved bounds. No assertion that either is an optimal general
determinantal representation is made.

\subsection{O-012.b: an arbitrary-rank padding inequality}

Let $U$ have dimension $r$, let $b\leq n$, and let
$\mu\vdash db$ have length at most $r$. Put
\[
                \lambda=\mu+d(n-b)e_1.
\]
If a closed $\mathrm{GL}(V)$-stable cone
$X\subseteq\operatorname{Sym}^n V^*$ contains every form
$x_1^{n-b}f$ with $f\in\operatorname{Sym}^b U^*$, then
\[
       \operatorname{mult}_{S_\lambda V}\mathbb C[X]_d
                                     \geq a_\mu(d[b]).  \tag{O-012.2}
\]
Here $U$ is the span of the first $r$ coordinate directions
in $V$. This is an arbitrary-rank form of the padding inequality
stated in Ikenmeyer--Panova Proposition~2.6(b), p.~6, citing
Kadish--Landsberg Theorem~1.3. The proof below makes the
normalization issue explicit.

\begin{proof}
Put $p=n-b$, and let
$T:\operatorname{Sym}^b U\to\operatorname{Sym}^n U$
be multiplication by $e_1^p$. For the upper Borel,
$gT=g_{11}^pTg$, so $\operatorname{Sym}^d T$ takes a
highest-weight vector of weight $\mu$ to one of weight
$\lambda$. It is injective.

Use the differential pairings
$\langle a,F\rangle=a(\partial)F|_0$ in each symmetric power.
The transpose $R$ of the padding map
$P:f\mapsto x_1^pf$ is $\partial_{e_1}^p$. On a monomial
$e^\alpha$ of degree $b$ one has
\[
                    RT(e^\alpha)
                    =\frac{(\alpha_1+p)!}{\alpha_1!}e^\alpha.
\]
Thus $RT$ is an invertible diagonal operator in characteristic
zero. Pullback to the padded linear space, after applying
$\operatorname{Sym}^dT$, is
$\operatorname{Sym}^d(RT)$, also invertible.
In particular, a basis of the degree-$d$ highest-weight space
of weight $\mu$ lifts to highest-weight vectors of weight
$\lambda$ whose restrictions to that padded space remain
linearly independent. Since that space lies in $X$, the lifted
vectors remain independent in $\mathbb C[X]_d$, proving
(O-012.2). A different conventional normalization of the
symmetric-power pairing changes the displayed diagonal factors
by nonzero constants and does not change this argument.
\end{proof}

Together O-012.a and O-012.b give
\[
             K_n(\mu+d(n-b)e_1)\geq a_\mu(d[b])
                 \quad\text{if }n\geq B(b,r).           \tag{O-012.3}
\]
The geometric statement concerns a whole padded linear space;
it does not assert identical evaluation of the lifted polynomial
on a single unchanged base form. Invertibility of the diagonal
pullback, rather than a false coefficient-normalization identity,
is what preserves all copies.

\subsection{O-012.c: stabilization after removing the last-row rectangle}

Let $d\geq2$ and $\lambda\vdash nd$ have at most $d$ rows;
pad it with zeros to length $d$. Define
\[
 k=\lambda_d,\qquad
 \rho=(\lambda_2-k,\ldots,\lambda_d-k),\qquad
 t=|\rho|,\qquad b=k+t.
\]
Zero parts of $\rho$ may be deleted. If $b\geq1$ and $n\geq b$,
put
\[
       \mu=(k+(d-1)t,\lambda_2,\ldots,\lambda_d)\vdash db.
\]
Then
\[
                         a_\lambda(d[n])=a_\mu(d[b]).  \tag{O-012.4}
\]
When $b=0$, the partition is the one-row shape and its ambient
multiplicity is one. For $\ell(\lambda)<d$ the definition gives
$k=0$, $b=|\bar\lambda|$; for $\ell(\lambda)=d$ it can
give a much smaller base degree than the total tail size.

\begin{proof}
Compute stable plethysm multiplicities using a $d$-dimensional
row space and a $d$-dimensional column space, with matrix
variables $Z=(z_{ij})$. By the Cauchy decomposition, the row
highest-weight space of weight $\lambda$ is the column module
$S_\lambda\mathbb C^d$. Its subspace of column multidegree
$(n,\ldots,n)$ fixed by column permutations has dimension
$a_\lambda(d[n])$.

Every polynomial in this row highest-weight space is divisible
by $\det(Z)^k$. To justify the assertion about every vector,
multiplication by that determinant power embeds the row
highest-weight space of weight $\lambda-k^d$ into the one
of weight $\lambda$. As column modules the map is the nonzero
map between the identical irreducibles
$\det^k\otimes S_{\lambda-k^d}\mathbb C^d$ and
$S_\lambda\mathbb C^d$, so it is an isomorphism.

After division, write $a=n-k$. The row weight is
$(da-t,\rho)$ and every column has degree $a$.
Each monomial has total degree $t$ outside the first row.
Consequently its exponent of $z_{1j}$ is at least $a-t$
in every column $j$. Since $a\geq t$, every such polynomial
is divisible by
\[
                         Q=\prod_{j=1}^d z_{1j}^{a-t}.
\]
This factor is fixed by the upper row unipotent group and by
column permutations. Division by it preserves the highest-weight
condition, subtracts $d(a-t)e_1$ from the row weight, and
reduces every column degree to $t$. Conversely multiplication
by $Q$ reverses the operation. Thus it is a bijection of the
two relevant column-permutation modules, including all their
copies and their signs.

Multiplying back by $\det(Z)^k$ gives row weight
$(k+(d-1)t,\lambda_2,\ldots,\lambda_d)=\mu$ and
column degree $b=k+t$. The determinant factor has the same
column sign character before and after, so the invariant
dimensions agree. This proves (O-012.4). The tuple $\mu$ is
a partition: $\rho_1\leq t\leq(d-1)t$, so its first
row is at least its second. Cauchy naturality then carries
the stable multiplicity to every actual vector-space dimension
in which $\lambda$ and $\mu$ fit.
\end{proof}

\subsection{O-012.d: full actual multiplicities and the four-copy family}

Under the hypotheses of O-012.c, set $r=\ell(\lambda)$.
If $r\leq n^2$ and
\[
                          n\geq\max\{b,B(b,r)\},       \tag{O-012.5}
\]
then
\[
                         K_n(\lambda)=a_\lambda(d[n]). \tag{O-012.6}
\]
For $r=1$ use the pure-power diagonal determinant instead of
$B(b,r)$. This excludes an actual multiplicity obstruction for
every padded-permanent closure, since $q\leq a=K_n$.

Indeed, $\lambda=\mu+d(n-b)e_1$, so O-012.3 gives
$K_n(\lambda)\geq a_\mu(d[b])=a_\lambda(d[n])$.
The reverse inequality is the ambient coordinate-ring quotient.

In particular, suppose $m\geq2$, $n\geq m^{25}$, and
$r\leq m^2$. Every shape satisfying $b\leq25$ in O-012.c
has full ambient determinant multiplicity. To verify the uniform
bound, use monotonicity in $b,r$:
\[
 B(b,r)\leq B(25,m^2)=2\binom{m^2+11}{12}-1.
\]
For $m=2$ the latter is $909<2^{25}$. For $m\geq3$,
\[
 2\binom{m^2+11}{12}
  \leq\frac{2(3m^2)^{12}}{12!}
  <m^{24}\leq m^{25},
\]
because $m^2+11\leq3m^2$ and $2\cdot3^{12}<12!$.
Also $n\geq m^{25}\geq25\geq b$.

For the explicit first higher-multiplicity example, take $d\geq6$,
$n\geq8$ and
\[
               \lambda=(d(n-2)-4,8,2^{d-2}).
\]
Here $k=2$, $t=6$, and $b=8$. Removing the column
$\det^2$ factor and using the binary weight count of F-013
(equivalently O-006.b) gives
\[
 a_\lambda(d[n])=p_6(d,n-2)-p_5(d,n-2)=11-7=4.
\]
Both rectangle bounds are inactive since $d,n-2\geq6$.
Therefore
\[
       n\geq\max\{8,B(8,d)\}\quad\Longrightarrow\quad
                         K_n(\lambda)=a_\lambda(d[n])=4. \tag{O-012.7}
\]
In the BIP range this holds for every $6\leq d\leq m^2$.
The statement settles all four determinant-side copies and
excludes $q>K_n$ there; it does not assert a value for the
permanent-side multiplicity between zero and four.

\paragraph{Geometric dictionary and scope.}
The determinant matrices define morphisms of affine coefficient
schemes over $\mathbb Z$. Multiplication by a chosen section
and contraction by its dual are maps of locally free symmetric
powers; the factorial diagonal calculation specifies their
characteristic-zero effect on all highest-weight copies.
The stabilization is a calculation in the coordinate algebra
of the matrix scheme: a determinant-line factor, followed by
division by a common first-row monomial, preserves the indicated
torus-weight and finite-group-invariant modules. The precise
background is the Cauchy identity, Macdonald I(4.3), p.~63,
and the functorial Schur decomposition I Appendix~A(5.3)--(5.6),
pp.~155--156, together with SGA~3 I4.7.2--4.7.3 and SGA~6
VI3.3. The related known padding statement is IP Proposition~2.6(b).
No claim is made that the factorial factors are invertible over
all of $\mathbb Z$. Integral polynomial identities and
characteristic-zero irreducible multiplicities remain separate.

The new gate is exact: a shape in this stabilized class can
only be an actual multiplicity obstruction if it fails
(O-012.5). In particular, at $n\geq m^{25}$ a surviving
maximal-row shape $\ell(\lambda)=d$ must have
$\lambda_d+\sum_{i=2}^d(\lambda_i-\lambda_d)>25$.
This is a full-multiplicity exclusion, including non-scalar
ambient multiplicities, rather than another occurrence test.
The later quadratic Chow bound M-010 excludes all maximal-row
shapes with $d\leq m^2$ throughout the BIP range as well.
Thus this particular BIP specialization is superseded by that
stronger degree result. O-012.a and O-012.b remain the inputs
to F-014 and F-015, which apply also at arbitrarily large outer
degree; the parameter-dependent statement (O-012.5) remains
valid in every stated range.

\subsection{O-013: the Hessian family at every outer degree}

Let $n\geq2$, $1\leq r\leq\min(d,n^2)$, and put
\[
          \lambda_{d,r}=(nd-2(r-1),2^{r-1})\vdash nd.
\]
Then
\[
           K_n(\lambda_{d,r})=a_{\lambda_{d,r}}(d[n])=1.
                                                        \tag{O-013.1}
\]
For internal padding and $n>m\geq2$, one moreover has
\[
 q^m_{\lambda_{d,r}}(d[n])=
 \begin{cases}1&r\leq m^2,\\0&r>m^2.\end{cases}          \tag{O-013.2}
\]
There is no upper bound on the outer degree $d$. Thus this
extends O-011 beyond the maximal-row case $d=r$ and gives
genuine positive-permanent examples even at degrees beyond
the independent quadratic Chow bound M-010.

\begin{proof}
The one-row case is the power of the first coefficient, so
assume $r\geq2$. Apply F-014 with second row $2$:
\[
 a_{\lambda_{d,r}}(d[n])=
 a_{\mu}(d[2]),\qquad
       \mu=(2(d-r+1),2^{r-1}).
                                                        \tag{O-013.3}
\]
We give the short multiplicity-free argument needed on the
right, rather than presuming a formula for all quadratic
plethysms. In the space of quadratic forms on $\mathbb C^r$,
the upper triangular group has a dense open orbit. Indeed,
on the open set where all leading principal minors of the
symmetric coefficient matrix are nonzero, symmetric Gaussian
elimination gives an upper-triangular congruence to a diagonal
matrix with nonzero entries. Over $\mathbb C$ a further
diagonal congruence changes that matrix to the identity.
Consequently the ratio of two nonzero polynomial highest-weight
functions of the same weight is a rational invariant of the
upper triangular group, and is constant on its open orbit.
The two functions are proportional. Every irreducible
multiplicity in the coordinate ring of quadratic forms is
therefore at most one.

For a degree-$n$ form $F$, let $W_r(F)$ be the Hessian minor
used in O-011, and define
\[
                  H_{d,r}(F)=W_r(F)F(e_1)^{d-r}.
                                                        \tag{O-013.4}
\]
This is a degree-$d$ polynomial function. The chain-rule
calculation in O-011 and the identity
$F(ge_1)=g_{11}^{n}F(e_1)$ for upper triangular $g$ show
that it is a highest-weight vector of weight
\[
  (r(n-2)+2+n(d-r),2^{r-1})=\lambda_{d,r}.
\]
For $n=2$, evaluation on a nonsingular diagonal quadratic
form with nonzero first coefficient makes (O-013.4) nonzero.
Thus the multiplicity at $\mu$ in (O-013.3) is exactly one,
and so is the ambient multiplicity at $\lambda_{d,r}$.

At the determinant point, O-011 supplies for every $r\leq n^2$
a change of coordinates under which $W_r$ is nonzero and
$F(e_1)$ is nonzero: the first direction corresponds to the
identity matrix, and the other directions span a nondegenerate
Hessian subspace containing it. Hence (O-013.4) is nonzero
on the determinant orbit, proving $K_n\geq1$. The ambient
upper bound gives (O-013.1).

For the internally padded permanent, O-011 supplies the same
two nonvanishing conditions for every $r\leq m^2$, including
the separately checked exceptional pair $(m,n)=(3,4)$.
Thus (O-013.4) is nonzero on that orbit as well, and
$q\geq1$. Again the ambient upper bound gives $q=1$.
If $r>m^2$, the permanent depends on only $m^2$ variables,
so the length restriction gives $q=0$.
\end{proof}

This is an exact scalar family at all outer degrees, with
an explicit surviving highest-weight function. It does not
claim that occurrence determines multiplicity when the ambient
coefficient exceeds one. The determinant-side proof uses
the actual orbit, so no replacement of closure multiplicities
by symmetric or ordinary rectangular coefficients occurs.

\paragraph{Geometric dictionary.}
Quadratic forms are points of the affine scheme of symmetric
bilinear matrices. Triangular congruence gives the displayed
open orbit after base change to $\mathbb C$; the resulting
character semi-invariants are torus-weight sections in the
grading/comodule correspondence of SGA~3, I, Proposition~4.7.3,
pp.~35--36. The Hessian and evaluation polynomials define
sections of the indicated Schur-type summands in the coordinate
algebra of forms. Their functorial decompositions are Macdonald,
I Appendix~A, (5.3)--(5.6), pp.~155--156, and the locally free
$K_0$ framework is SGA~6, VI, Theorem~3.3, pp.~392--393.
The scheme maps and derivative polynomials have integral
coefficients, but the Borel orbit argument, complete
reducibility, and Hessian nondegeneracy used here are
characteristic-zero assertions. Nonzero restriction to the
actual orbit means that the corresponding section survives
the coherent ideal quotient of its closure.

\section{A bounded-exponent determinant construction and a strict test case}

\subsection{O-014.a: degree splitting with exponent bounds}
Let $b\geq1$, $r\geq2$, and choose positive integer caps
\[
             c=(c_2,\ldots,c_r),\qquad c_2=b,
                    \quad c_i\leq b\quad(i\geq3).
\]
Write $u=\lfloor(b-1)/2\rfloor$, $v=\lceil(b-1)/2\rceil$,
and define
\[
 N_j(c)=\#\{\alpha\in\mathbb Z_{\geq0}^{r-1}:
                         \alpha\leq c,\ |\alpha|\leq j\},
 \qquad D(b,c)=N_u(c)+N_v(c)-1.                         \tag{O-014.1}
\]
Equivalently,
\[
 N_j(c)=[z^j]\frac{1}{1-z}\prod_{i=2}^r
                                      (1+z+\cdots+z^{c_i}).
\]
For every degree-$b$ form
\[
 f(L,z_2,\ldots,z_r)=
       \sum_{|\gamma|\leq b,\,\gamma\leq c}
                          a_\gamma z^\gamma L^{b-|\gamma|},
\]
there is a homogeneous linear matrix of size $D=D(b,c)$
whose determinant equals $L^{D-b}f$. Its entries are linear
forms in the form variables, with integral polynomial dependence
on the coefficient parameters.
In particular, every $L^{n-b}f$ in this bounded coefficient
space lies in $\Omega_n$ if $n\geq D$ and $r\leq n^2$.

\begin{proof}
Use the two trees of O-012.a, retaining only multi-indices
$\alpha\leq c$ in the first tree of height $u$ and the second
tree of height $v$. The retained sets are closed under taking
parents, so every root path remains present. They have $N_u(c)$
and $N_v(c)$ vertices; merge only their roots.
For a coefficient with $|\gamma|<b$, split its sorted word
into $\alpha+\beta=\gamma$ with $|\alpha|\leq u$ and
$|\beta|\leq v$, and insert the coefficient edge labeled
$a_\gamma L$. Both endpoints satisfy the caps because they
are coordinatewise bounded by $\gamma$. For $|\gamma|=b$,
remove one variable $z_i$, split $\gamma-e_i$ in the same
way, and label the coefficient edge by $a_\gamma z_i$.
Again both endpoints and every path are retained.

All tree edges carry the negative of their variable, all
noncentral unused vertices carry $L$-loops, and coincident
coefficient-edge labels are added. Every directed non-loop
cycle passes through the central vertex and exactly one
coefficient edge. The determinant signs cancel the tree-edge
signs just as in O-012.a, giving one term
$a_\gamma z^\gamma L^{D-|\gamma|}$ for each coefficient.
Their sum is $L^{D-b}f$. Since $c_2=b$, the first variable
alone supplies paths of full heights $u$ and $v$, so
$D\geq u+v+1=b$. No negative power of $L$ occurs.
Adding an $L I_{n-D}$ block gives the larger determinant.
\end{proof}

For a partition $\lambda$, use $b=\lambda_2$ and
$c=(\lambda_2,\ldots,\lambda_r)$. F-015 proves that this
whole restricted coefficient space is sufficient to preserve
every highest-weight copy, so that $n\geq D(b,c)$ implies
$K_n(\lambda)=a_\lambda(d[n])$. The argument uses the
second-row stabilization F-014 and the invertible diagonal
padding pullback of O-012.b. The size in (O-014.1) is at
most the unbounded size $B(b,r)$, and can be strictly smaller.

\subsection{O-014.b: a genuine positive-permanent example beyond the earlier gates}

Here is an exact test case demonstrating that the bounded
coefficient refinement can add an exclusion, including when
the ambient multiplicity is larger than one. Set
\[
 m=256,\quad r=m^2=65536,\quad b=32,\quad
 c=(32,2^{65534}),\quad n=D(32,c),\quad d=n.
\]
The exact integer is
\[
n=N_{15}(c)+N_{16}(c)-1.
                                                               \tag{O-014.2}
\]
Take the partition
\[
                 \lambda=(n^2-131100,32,2^{65534}).     \tag{O-014.3}
\]
Then
\[
              4\leq q^m_\lambda(d[n])\leq K_n(\lambda)
                     =a_\lambda(d[n]).
                                                               \tag{O-014.4}
\]
The comparison of every copy follows from F-015, not from
the explicit one-copy construction below.

To reproduce (O-014.2), put $R=65534$. Since $u=15$, $v=16$
and the first cap is inactive in these degrees,
\[
 N_j(c)=\sum_{k=0}^j(j-k+1)
          \sum_{t=0}^{\lfloor k/2\rfloor}
                         \binom Rt\binom{R-t}{k-2t}.
                                                               \tag{O-014.5}
\]
The inner sum selects the tail variables used twice and once;
the factor $j-k+1$ counts the possible first-coordinate
exponents in a monomial of total degree at most $j$.
Thus (O-014.5), followed by $n=N_{15}+N_{16}-1$, is a
finite integer certificate independent of decimal rounding.
The complete decimal integer is retained in
\texttt{boxed\_degree\_split\_certificate.json}.

\paragraph{Positive permanent multiplicity.}
For the internally padded permanent, choose the nondegenerate
Hessian point from O-011 and an orthogonal basis beginning
with its nonisotropic radial direction. Every leading Hessian
minor in that basis is nonzero, including $W_2$ and $W_r$;
the form itself is nonzero in the first direction. Therefore
\[
                    W_r(F)W_2(F)^{15}F(e_1)^{d-r-30}
                                                               \tag{O-014.6}
\]
is a nonzero degree-$d$ highest-weight function on the
permanent orbit. Its second row is $2+30=32$, its remaining
nonfirst rows are $2$, and its first row is determined by
total size $nd$. Its weight is exactly (O-014.3), proving
$q^m_\lambda>0$. The stronger bound $q^m_\lambda\geq4$
follows from M-012.3, whose hypotheses hold because $b=32$
and $d\geq r+30$. It retains four independent binary seed
classes after multiplication by one common nonzero Hessian
factor on the integral permanent coordinate ring.

For the ambient lower bound, use F-013's four-dimensional
highest-weight space at
\[
                    (r(n-2)-4,8,2^{r-2})
\]
in outer degree $r$. Multiplication by the nonzero polynomial
$W_2^{12}F(e_1)^{d-r-24}$ sends all four vectors to the
highest-weight space at (O-014.3). The ambient polynomial
ring is a domain, so multiplication by this fixed polynomial
is injective. Thus $a_\lambda\geq4$. No value for the actual
exact permanent multiplicity is inferred from this ambient calculation.

\paragraph{Comparison with previously proved sufficient conditions.}
Exact integer arithmetic in the accompanying certificate gives
\[
                   m^{25}\leq n<2^m-1,
\]
so the BIP range holds and the full Grenet containment gate
F-009 does not apply. The unbounded size is
\[
B(32,65536)=\binom{65550}{15}+\binom{65551}{16}-1>n,
\]
so the earlier second-row determinant criterion also fails.
For any split of the $r-1=65535$ tail variables, some part
has at least $32768$ variables. Since all caps are at least
two, its monomial count through degree $m$ is at least
$\binom{32768}{256}$, by choosing squarefree monomials of
degree $m$. Consequently
\[
        T_m(\lambda)\geq\binom{32768}{256}>n.
\]
The analogous degree-$b$ variable-splitting bound is at least
$\binom{32768}{32}>n$. The uniform row cutoff does not apply,
since $r>51$, and the second row is greater than $25$.
Since $d=n$, the quadratic Chow bound $n\geq d(d-1)$
and the adaptive bound $n\geq B_d(32)=32d-497$ also fail.
The coarse padding requirements hold: $r=m^2$ and
$|\bar\lambda|=131100\leq md$.

These failures describe the scope of the comparison: the
example adds to the explicitly listed sufficient criteria.
It does not prove that every other conceivable method would
fail. The conclusion is the exact full determinant multiplicity
in (O-014.4), with at least four actual permanent copies.
The exact permanent multiplicity remains undetermined.
The example certifies strict refinement of the listed explicit
conditions; it is not an exhibited surviving $q>K_n$ candidate.
The later nested-support criteria F-016 and M-011 also handle
this example. O-015.c gives a separate four-copy example
beyond all the stabilization criteria, using M-011 directly.

\paragraph{Geometric dictionary.}
The bounded coefficient space is an affine linear subscheme
of the scheme of degree-$b$ forms, cut out by coefficient
coordinates outside the prescribed exponent box. Its coordinates
are selected by torus weights, using SGA~3, I, Proposition~4.7.3,
pp.~35--36. The two trees define a morphism from that affine
coefficient scheme to the scheme of linear matrices over
$\mathbb Z$, and the determinant identity proves factorization
through the padded-form morphism. F-015 then interprets
restriction as a torus-equivariant map of coordinate modules;
representations as comodules are SGA~3, I, Proposition~4.7.2,
pp.~34--35, and Schur functor naturality is Macdonald,
I Appendix~A, (5.3)--(5.6), pp.~155--156. The integral
factorization does not itself split irreducibles integrally:
the full-copy comparison is made after base change to
$\mathbb C$, with the explicit factorial invertibility in
O-012.b. Multiplication of the Hessian sections takes place
in the actual ambient coordinate algebra, and the positivity
claim comes from their nonzero restriction to the permanent
orbit rather than a formal identity in its $K_0$ class.

\section{Determinant graphs for nested highest-weight supports}

\subsection{O-015.a: arbitrary downward-closed coefficient supports}
Fix a base degree $b\geq1$ and $r\geq2$. Let
$\Gamma\subseteq\mathbb Z_{\geq0}^{r-1}$ be a nonempty finite
coordinatewise downward-closed set contained in $|\gamma|\leq b$.
Consider all forms
\[
             f(L,z)=\sum_{\gamma\in\Gamma}
                         a_\gamma z^\gamma L^{b-|\gamma|}.
\]
Put $u=\lfloor(b-1)/2\rfloor$, $v=\lceil(b-1)/2\rceil$,
$\Gamma_j=\{\gamma\in\Gamma:|\gamma|\leq j\}$, and
\[
                    D_\Gamma=|\Gamma_u|+|\Gamma_v|-1.
                                                               \tag{O-015.1}
\]
There is a homogeneous linear matrix of size $D_\Gamma$
whose determinant is the polynomial
\[
                 \sum_{\gamma\in\Gamma}
                         a_\gamma z^\gamma L^{D_\Gamma-|\gamma|}.
                                                               \tag{O-015.2}
\]
For a split $A\sqcup B=\{2,\ldots,r\}$, let $\Gamma_A$
and $\Gamma_B$ be the coordinate projections of $\Gamma$,
including the zero index. There is also such a matrix of size
\[
                         S_{\Gamma,A}=|\Gamma_A|+|\Gamma_B|-1.
                                                               \tag{O-015.3}
\]
In (O-015.2) replace $D_\Gamma$ by this size. In each case
all displayed exponents are nonnegative. Thus every
$L^{n-b}f$ in the prescribed coefficient space belongs to
$\Omega_n$ when $n\geq b$ and $n$ is at least one of the
displayed sizes, provided $r\leq n^2$.

\begin{proof}
For (O-015.1), use the two degree-split trees on
$\Gamma_u$ and $\Gamma_v$. Deleting a variable from a
multi-index stays in $\Gamma$, so all root paths exist.
Every split $\gamma=\alpha+\beta$ of a retained monomial
also has $\alpha,\beta\in\Gamma$, since each is coordinatewise
bounded by $\gamma$. For $|\gamma|<b$, use the coefficient
edge $a_\gamma L$; for $|\gamma|=b$, first remove one
variable $z_i$, split the remaining degree $b-1$, and use
$a_\gamma z_i$. These are precisely the cycles of O-014.a.
When $|\gamma|<b$ the cycle has $|\gamma|+1$ vertices;
when $|\gamma|=b$ it has $b$ vertices. Since these vertices
are distinct in the constructed graph, the corresponding
exponent in (O-015.2) is nonnegative even if the total size
is less than $b$. Permutation signs cancel all tree-edge
signs and yield (O-015.2).

For (O-015.3), the projections are downward closed, and
the projections of each $\gamma\in\Gamma$ give the two
required vertices. Use the variable-split trees on these
two sets and a coefficient edge $a_\gamma L$ between them.
Each cycle has $|\gamma|+1$ vertices, so its contribution
is $a_\gamma z^\gamma L^{S_{\Gamma,A}-|\gamma|}$, again
with nonnegative exponent. Padding either matrix by an
$L$-diagonal block of the required size gives
$L^{n-b}f$. Formally the unpadded identity can be written
$L^b\det M=L^{\operatorname{size}(M)}f$; the termwise
argument proves polynomiality even when $\operatorname{size}(M)<b$.
All identities are integral in the coefficient parameters.
\end{proof}

\subsection{O-015.b: nested support counts}
For a partition $\lambda$ with $r\geq2$ rows and
$b=\lambda_2$, take
\[
 \Gamma(\lambda)=\left\{\gamma\in\mathbb Z_{\geq0}^{r-1}:
              \sum_{i=j}^r\gamma_i\leq\lambda_j
                         \quad(2\leq j\leq r)\right\}.
                                                               \tag{O-015.4}
\]
The $j=2$ inequality bounds the total degree by $b$;
all inequalities are preserved by coordinatewise decrease.
Thus O-015.a applies. The pure second-variable powers through
degree $b$ belong to this set. Consequently $D_\Gamma\geq b$
and every $S_{\Gamma,A}\geq b+1$.
F-016 proves that this support contains the whole stabilized
highest-weight space, so either size bound preserves all
ambient determinant multiplicities at the corresponding
weight. The suffix conditions strengthen coordinatewise caps.

For $r\geq3$ write $c=\lambda_3$ and $R=r-2$.
For $0\leq j\leq b$, ignoring the later suffix inequalities
gives the useful bound
\[
 |\Gamma_j(\lambda)|\leq
 \sum_{h=0}^{\min(c,j)}(j-h+1)\binom{R+h-1}{h}.
                                                               \tag{O-015.5}
\]
Here $h$ is the total degree in variables $3,\ldots,r$;
the second variable then has $j-h+1$ possible exponents.
Equality holds if $\lambda_3=\cdots=\lambda_r=c$.
At the variable split $A=\{2\}$, $B=\{3,\ldots,r\}$,
\[
                   S_{\Gamma,A}\leq
                          b+\binom{r-2+c}{c}.          \tag{O-015.6}
\]
Again equality holds for the constant later tail. In particular,
for tails $(b,2^{r-2})$ the sufficient full-multiplicity bound
from F-016 is simply
\[
                       n\geq b+\binom r2.             \tag{O-015.7}
\]
For comparison the degree-split count in this case, for $2\leq j\leq b$,
is
\[
              |\Gamma_j|=(j+1)+j(r-2)+(j-1)\binom{r-1}{2}.
\]
The variable split can be substantially smaller, so both
constructions are retained.

\subsection{O-015.c: an actual four-copy exclusion beyond inner-degree stabilization}

Set
\[
 m=256,\quad r=65536,\quad n=m^{25}=2^{200},\quad
 d=3n,\quad b=2n,
\]
and take
\[
           \lambda=(3n^2-2n-131068,2n,2^{65534}).       \tag{O-015.8}
\]
Then the actual multiplicities satisfy
\[
                         4\leq q^m_\lambda(d[n])
                                      \leq K_n(\lambda).        \tag{O-015.9}
\]
Neither exact multiplicity is claimed here.

For the upper comparison, use M-011 at the cut $k=3$.
Although $\lambda_2=2n>n$, the retained permanent-degree
forms have total degree at most two in variables $3,\ldots,r$.
The prefix variable has degree at most $m$, so the required
matrix size is
\[
               m+\binom r2=2147451136<2^{200}=n.
\]
This proves $q\leq K_n$ without any inner-degree stabilization.
For the lower bound, M-012.3 applies with the same even $b$:
$d\geq r+b-2$ because $n\geq r-2$. It gives four
independent highest-weight classes on the permanent closure,
proving $q\geq4$. More explicitly, multiply the four binary
seed classes by
\[
              W_r\,W_2^{n-4}F(e_1)^{n-r+2}.
\]
The common factor is nonzero on the reduced irreducible
permanent closure, and the degree and weight are those in
(O-015.8). The seed classes each have degree six; the
displayed factor has degree $d-6$.

This example satisfies the BIP threshold with equality, while
$n<2^{256}-1$, so the full Grenet containment criterion does
not apply. All the stabilization-based full-multiplicity
criteria F-014--F-016 require $n\geq\lambda_2$ and fail.
For every variable split in the older M-009 box criterion,
one part has at least $32768$ variables, all with caps at
least two. Therefore
\[
                 T_m(\lambda)\geq\binom{32768}{256}>2^{200}.
\]
The quadratic Chow criterion fails since $d(d-1)>n$.
The adaptive one fails since, with $b<d$,
\[
                    B_d(b)=db-\binom b2-1
                                =4n^2+n-1>n.
\]
The coarse GCT shape conditions hold: $r=m^2$ and
$|\bar\lambda|=2n+131068\leq md$. In particular the
necessary positive-multiplicity condition $q\geq2$ is
verified rather than left unevaluated.

The accompanying exact integer certificate checks every
comparison. This certifies strict refinement of the listed
explicit sufficient conditions by M-011. It exhibits an
excluded shape with genuine multiplicity at least four; it
does not exhibit an actual obstruction or establish a value
for either of its two multiplicities.

\paragraph{Geometric dictionary.}
The downward-closed support defines an affine coefficient
subscheme selected by torus weights, in the grading/comodule
dictionary of SGA~3, I, Proposition~4.7.3, pp.~35--36.
The graph morphisms to linear matrices and their determinant
identities are defined over $\mathbb Z$. The supports need
not be stable under the entire general linear group: the
highest-weight and ideal arguments are performed in the
ambient representation, using SGA~3, I, Proposition~4.7.2,
pp.~34--35, and functorial Schur decomposition as in
Macdonald, I Appendix~A, (5.3)--(5.6), pp.~155--156.
The characteristic-zero all-copy restrictions are F-016 and
M-011. The four surviving permanent sections in the example
come from the explicit integral binary substitution of M-012,
followed by multiplication in the integral coordinate algebra
of the reduced irreducible complex orbit closure.

\section{A boundary for exact universal padded determinant constructions}

\subsection{O-016: the degree-split exponent is optimal for this exact method}
Fix $b\geq2$, and let $R\geq5$ be the number of variables.
Write
\[
 h=\lfloor b/2\rfloor,\qquad
 D_j(R)=\binom{R+j-1}{j}.
\]
For a Zariski-general homogeneous degree-$b$ form $f$ in
$R$ variables over $\mathbb C$, its exact affine-linear
determinantal complexity satisfies
\[
 \operatorname{dc}(f)\geq
 1+\left\lceil\frac{D_b(R)}{D_h(R)+D_{b-h}(R)}\right\rceil
                         =\Omega_b(R^{\lfloor b/2\rfloor}).
                                                               \tag{O-016.1}
\]
Consequently any construction representing every padded
degree-$b$ form in $r=R+1$ variables by an exact homogeneous
linear determinant needs size
$\Omega_b(r^{\lfloor b/2\rfloor})$. The size
$B(b,r)=O_b(r^{\lfloor b/2\rfloor})$ in O-012.a therefore
has the optimal exponent in $r$ for fixed $b$, within this
exact universal problem.

The known input is Chatterjee--Kumar--Volk,
\emph{Determinants vs. Algebraic Branching Programs}, ITCS~2024,
article~27: Observation~2.6, p.~27:6, and Theorem~3.1,
p.~27:8. They imply that if a homogeneous form has singular
locus of dimension less than $R-4$ and an exact determinant
representation of size $s$, it has a homogeneous algebraic
branching program of width at most $s-1$. Their Definition~2.7,
p.~27:7, counts vertices rather than edges for program size;
Theorem~4.1, p.~27:9, more generally gives size $O(b^5s)$
for all homogeneous forms. The sharper width statement is
the only complexity input needed below. These are known
theorems, not new results of this investigation.

\begin{proof}
There is a nonempty Zariski-open set of degree-$b$ forms
whose projective hypersurface is smooth. To see both openness
and nonemptiness, the incidence defined by all partial
derivatives in the product of the coefficient space with
$\mathbb P^{R-1}$ has closed projection, since the second
factor is projective (EGA~II, Theorem~5.5.3(i), p.~104,
with properness Definition~5.4.1, p.~100), while the Fermat form
$z_1^b+\cdots+z_R^b$ is outside that projection.
For such a form the affine singular locus is exactly the
origin, of dimension zero, which is less than $R-4$.

By the cited regular-representation result, any determinant
representation of such an $f$ of size $s$ gives a homogeneous
branching program of width at most $s-1$. Cut the program
at layer $h$. Its path sum then has the form
\[
                f=\sum_{i=1}^{s-1}P_iQ_i,
       \qquad \deg P_i=h,\quad\deg Q_i=b-h,             \tag{O-016.2}
\]
inserting zero summands if the actual width is smaller.
This is simply the partition of all source-to-sink paths by
their vertex at that layer.

For a fixed integer $s$, the forms in (O-016.2) lie in the
image of the polynomial morphism
\[
 \left(\operatorname{Sym}^h\mathbb C^R
             \times\operatorname{Sym}^{b-h}\mathbb C^R\right)^{s-1}
 \longrightarrow\operatorname{Sym}^b\mathbb C^R,
                   \quad (P_i,Q_i)_i\longmapsto\sum_iP_iQ_i.
\]
The closure of its image has dimension at most
$(s-1)(D_h(R)+D_{b-h}(R))$. If this is strictly smaller
than $D_b(R)$, that image closure is proper. Its union with
the proper closed nonsmooth locus is still proper, since
the ambient affine coefficient space is irreducible.
Outside the union no form has determinant complexity at most
$s$. Taking the largest integer $s$ satisfying the strict
dimension inequality proves (O-016.1).

For fixed $b$, the numerator has degree $b$ in $R$ and
the denominator has degree $\lceil b/2\rceil$, with
positive leading coefficients. This gives the stated
asymptotic exponent. Finally, an exact identity
\[
                    \det M(L,z)=L^{s-b}f(z)
\]
with $M$ homogeneous linear and $f$ independent of $L$
specializes at $L=1$ to an affine-linear determinant for
$f(z)$ of size $s$. A universal padded construction includes
these forms, so it inherits the lower bound. O-012.a has
the matching exponent because
$\lceil(b-1)/2\rceil=\lfloor b/2\rfloor$.
\end{proof}

\paragraph{The base-degree twenty-six threshold.}
Take $r=m^2$ and $b=26$, so $R=m^2-1$. Then
\[
 \frac{D_{26}(R)}{2D_{13}(R)}
   =\frac{13!}{2\cdot26!}\prod_{j=13}^{25}(R+j)
                         =\Theta(m^{26}).              \tag{O-016.3}
\]
For all sufficiently large $m$, this exceeds $m^{25}$.
Thus not every degree-$26$ form in these variables can be
handled by an exact padded determinant of size $m^{25}$.
The uniform second-row constant twenty-five obtained from
an exact universal base-form construction cannot be raised
to twenty-six by merely improving the constant factor in
that universal exact construction. Restricted supports and
permanent-specific families remain available, as already
used in F-015--F-017 and M-011.

\paragraph{Exact complexity versus orbit-closure containment.}
The last conclusion is deliberately about exact determinant
identities. It does not prove that the corresponding padded
linear space is absent from $\Omega_{m^{25}}$, nor that a
multiplicity obstruction exists. A point of an orbit closure
can be a limit of determinants whose specialization at $L=1$
has higher-degree terms that disappear only in the limit.
The homogeneous fixed-degree conversion theorem does not
automatically apply uniformly to those approximating forms.
The source itself discusses the unresolved size-preserving
homogeneous-extraction issue on p.~27:5. No passage from
exact to border determinantal complexity is assumed here.
The generic lower bound is also not an explicit lower bound
for the permanent or for any identified GCT coefficient.

\paragraph{Geometric dictionary.}
The coefficient and product parameter spaces above are affine
schemes, and the dimension argument bounds the dimension of
a scheme-theoretic image after base change to $\mathbb C$.
The smooth-form open set is the complement of the proper
projection of a projective incidence scheme. The padded
matrix identity is an integral morphism, while specialization
at the nonzero section $L=1$ is a fiber operation.
These geometric operations do not identify an image with
its closure or assert lifting through a limiting family.
The distinction is the same coherent-quotient versus explicit
parametrization distinction used throughout the GCT comparison.

\section{A border-valid dimension boundary for the mixed construction}

\subsection{O-017: the mixed staircase reaches the universal cut-family exponent}
Fix integers $k\geq2$, $s\geq1$, and a real exponent
$c\geq1$. Put $r=m^2$, and suppose $n=n(m)\geq m$
with $n(m)=O(m^c)$. Take one padding variable $L$,
$k-2$ early variables, and $R=r-k+1$ later variables.
For sufficiently large $m$, consider the whole linear space
of degree-$m$ forms having total later-variable degree at
most $s$. Its dimension is
\[
 D_m(k,s)=\sum_{h=0}^s
       \binom{R+h-1}{h}\binom{m-h+k-2}{k-2}
        =\frac{m^{2s+k-2}}{s!(k-2)!}
                                +O_{k,s}(m^{2s+k-3}).    \tag{O-017.1}
\]
If
\[
                            k+2s-4>2c,                 \tag{O-017.2}
\]
then not every form in that space, after multiplication by
$L^{n-m}$, belongs to $\Omega_n$ for all sufficiently large
$m$. This is a dimension statement about the actual orbit
closure, and therefore allows border representations.

\begin{proof}
For later degree $h$, the first binomial coefficient in
(O-017.1) counts the later monomials. The remaining degree
$m-h$ is distributed among the $k-2$ early variables and
$L$, giving the second binomial coefficient. Distinct
monomials are linearly independent, so the sum is the exact
dimension. Here $m\geq s$ is harmless because the conclusion
is asymptotic with fixed $s$. Since $R=m^2-k+1$,
the term $h=s$ has the displayed positive leading coefficient
and degree $2s+k-2$ in $m$; all smaller $h$ have lower
degree. Multiplication by the fixed nonzero monomial
$L^{n-m}$ is injective, so it preserves that dimension.

O-010 gives the dimension bound
\[
 \dim\bigl(\Omega_n\cap\operatorname{Sym}^n((\mathbb C^r)^*)\bigr)
                              \leq (r-2)n^2+2.         \tag{O-017.3}
\]
Its use here does not replace the orbit closure by the exact
image of substitutions. Indeed the intersection is the
\emph{closure} of the $r$-variable determinant-pencil image:
the projection to the first $r$ variables sends the full
determinant substitution image to that pencil image, and
continuity sends its closure into the closure of the latter.
A point of the intersection is fixed by this projection.
The reverse inclusion follows because each pencil is an
endomorphism substitution of the determinant. O-010 bounds
the dimension of this image closure by its explicit generic
gauge parametrization.

If the whole padded linear family were contained in
$\Omega_n$, its dimension would be at most the right-hand
side of (O-017.3). Since $r=m^2$ and $n=O(m^c)$, that
right-hand side is $O(m^{2c+2})$. Condition (O-017.2)
is exactly $2s+k-2>2c+2$. The positive leading term in
(O-017.1) then contradicts the required dimension inequality
for all sufficiently large $m$.
\end{proof}

\paragraph{Meaning for the staircase.}
The BIP staircase in F-017 uses $c=25$ and
$k=54-2s$, so $k+2s-4=50=2c$. The IP staircase
uses $c=4$, $k=12-2s$, again giving equality; a fixed
factor such as $n=3m^4+1$ does not change the exponent.
Thus those constructions reach the allowed exponent for
their universal cut families. Increasing the permitted later
degree by one at the same cut crosses (O-017.2).
No construction proving containment of that \emph{whole}
larger truncated family in the determinant orbit closure can
work at the same power range for all large $m$.

For a finite illustration in the IP range, take $m=100$,
$r=10000$, $n=3m^4+1$, $k=2$, and the raised cap $s=6$.
Then
\[
 D_m(2,6)=\binom{10005}{6}>(9998)(300000001)^2+2.
\]
The exact comparison is in the accompanying certificate.
For a BIP illustration, take $m=10^{15}$, $n=m^{25}$,
$k=4$, and the raised cap $s=26$; the same certificate
checks $D_m(4,26)>(m^2-2)m^{50}+2$ by integer arithmetic.

This boundary does not apply to a smaller family selected
by additional nested row bounds, nor to the actual family
of projected permanents. The dimension of an unrestricted
coefficient family is not the dimension or evaluation rank
of those smaller families. In particular, (O-017.2) does
not exhibit an actual multiplicity obstruction, and it does
not say that the F-017 staircase is a complete classification.
It identifies the limitation of extending the same universal
cut-family containment argument by one unit in the suffix cap.

\paragraph{Geometric dictionary.}
The truncated family is a finite torus-weight affine linear
subscheme, with the grading interpretation of SGA~3, I,
Proposition~4.7.3, pp.~35--36. Padding by a fixed section
is a closed linear embedding. The determinant side is a
scheme-theoretic image closure after base change to
$\mathbb C$, and its dimension is bounded by the explicit
parametrization in O-010. This comparison already includes
limits and makes no normality, cohomological purity, or
saturation assumption. It is separate from O-016's stronger
fixed-base-degree exponent statement, whose method requires
exact determinant representations.

\section{Using two later-row bounds simultaneously}

\subsection{O-018.a: a joint-row gate}
In the BIP numerical range $n\geq m^{25}$, every actual
multiplicity obstruction satisfies
\[
                         \lambda_3\geq139
                           \quad\hbox{or}\quad\lambda_{52}\geq3.
                                                               \tag{O-018.1}
\]
The numerical constant is a sufficient bound obtained from
the count below; it is not claimed optimal.
A-025 subsequently strengthens $139$ to $557$ by using
two support trees and retaining the total-degree constraint.
The elementary count here remains useful for the explicit
full-profile example below.

\begin{proof}
The earlier reductions already give $m\geq190$ and
$52\leq r=\ell(\lambda)\leq m^2$ for an actual obstruction.
Suppose instead that $\lambda_3\leq138$ and
$\lambda_{52}\leq2$. In a projected degree-$m$ form,
the nested support F-016 relevant to a highest-weight
function has total degree at most $138$ in variables
$3,\ldots,r$, and at most two in variables $52,\ldots,r$.
There is one earlier variable, $x_2$, forty-nine middle
variables, and $r-51$ last variables.

Let $\Gamma$ be the set of all nonfirst-variable multi-indices
of degree at most $m$ satisfying the full nested bounds.
It is nonempty and downward closed. Ignoring coupled degree
conditions only increases its size, so
\[
 |\Gamma|\leq(m+1)\binom{187}{49}\binom{r-49}{2}
                           \leq\binom{187}{49}m^5.     \tag{O-018.2}
\]
The second inequality uses $m+1\leq2m$ and
$\binom{r-49}{2}\leq r^2/2\leq m^4/2$.
Exact integer arithmetic gives
\[
                         \binom{187}{49}\leq190^{20}.
\]
Since $m\geq190$, (O-018.2) is at most $m^{25}\leq n$.

Place all of $\Gamma$ in one monomial tree, and take the
other tree to consist of its root. O-015.a gives a determinant
of size $|\Gamma|$ for every retained padded form.
For every determinant-ideal highest-weight vector, first align
the padding with $x_1$ as in M-009, then delete the coefficient
monomials outside $\Gamma$. F-016 preserves its evaluation,
and the retained form is the determinant just constructed.
Thus every such vector vanishes on the permanent closure.
The inclusion of all highest-weight ideal multiplicity spaces
gives $q\leq K_n$, a contradiction. This proves (O-018.1).
This is also a special case of the general support-factorization
criterion F-018, with one factor equal to the whole support
and the other consisting only of zero.
\end{proof}

\subsection{O-018.b: a four-copy example beyond every single-cut containment test}
Set
\[
 m=2^{200},\quad r=m^2,\quad n=m^{25}=2^{5000},\quad d=3n,
\]
and define
\[
 \lambda=(3n^2-2n-2r-1172,\ 2n,\ 26^{49},\ 2^{r-51}).
                                                               \tag{O-018.3}
\]
It has $r$ rows and size $nd$. Then
\[
                              4\leq q^m_\lambda(d[n])
                                              \leq K_n(\lambda).
                                                               \tag{O-018.4}
\]
The exact values of these multiplicities are not asserted.
Every single-cut universal coefficient family associated to
this partition has dimension greater than the corresponding
determinant-pencil closure. Thus the comparison in (O-018.4)
uses information beyond every one-cut containment test, not
only beyond its uniform staircase corollaries.

\paragraph{The complete support fits.}
For this partition the full nested degree-$m$ support says
that the middle forty-nine variables and the last $r-51$
variables together have degree at most $26$, while the
last variables alone have degree at most two. The exponent
of $x_2$ is limited only by the total degree $m$, because
$\lambda_2=2n>m$.
The exact number of retained indices is
\[
 G=\sum_{t=0}^2\binom{r-52+t}{t}
       \sum_{j=0}^{26-t}\binom{48+j}{j}(m-j-t+1).
                                                               \tag{O-018.5}
\]
Here $t$ is the last-variable degree, $j$ the middle degree,
and the final factor counts powers of $x_2$.
In particular,
\[
 G\leq(m+1)\binom{75}{26}\binom{r-49}{2}
          \leq\binom{75}{26}m^5<2^{75}m^5<m^{25}=n.
\]
The one-tree construction in O-018.a therefore proves
$q\leq K_n$ for this partition as well.

\paragraph{At least four actual permanent copies.}
Take the four independent highest-weight classes on the
permanent closure of degree six and weight $(6n-6,6)$,
supplied by M-012.2. Multiply all four by the common factor
\[
                   W_2^{n-16}W_{51}^{12}W_r
                                       F(e_1)^{n-r-586}.        \tag{O-018.6}
\]
All exponents are nonnegative. Its degree is
$2n-32+612+r+n-r-586=3n-6=d-6$.
After adding the seed weight, its second row is
$6+2(n-16)+24+2=2n$; rows three through fifty-one
are $24+2=26$; the remaining nonfirst rows are two.
The first row is determined by the total degree, giving
exactly (O-018.3).

O-011 supplies a nondegenerate permanent Hessian at a point
where the value of the form and its radial Hessian entry
are nonzero. Extend that radial direction to an orthogonal
basis. Every leading Hessian minor, including $W_2,W_{51},W_r$,
is then nonzero. Thus (O-018.6) is a nonzero class on the
permanent closure. Its reduced coordinate ring is a domain,
so multiplication by this one class preserves the independence
of all four seed classes. This proves the lower bound in
(O-018.4); it does not require all seed classes to be nonzero
at the particular Hessian point.

\paragraph{Every single cut is too large as a universal family.}
For $2\leq k\leq r$, let $\mathcal U_k$ be the whole
space of degree-$m$ forms whose degree in variables
$k,\ldots,r$ is at most $\min(m,\lambda_k)$, with no
other later-row restriction. Also allow $k=r+1$, with an
empty last group and no restriction. Its padded copy has
the same dimension. We show
\[
                         \dim\mathcal U_k>m^{52}
                                    >(r-2)n^2+2         \tag{O-018.7}
\]
for every cut. By O-010, none of these whole padded families
can be contained in $\Omega_n$, even allowing limits.

If $k=2$, all degrees through $m$ are allowed in the last
group. If $k>r/2$, at least $r/2-2$ early variables are
unrestricted. In either case squarefree monomials of degree
twenty-seven in an available group, with the remaining degree
on $L$, give
\[
 \dim\mathcal U_k\geq\binom{r/2-2}{27}
       \geq\frac{(r/4)^{27}}{27!}
       =\frac{m^{54}}{2^{54}27!}>m^{52}.                \tag{O-018.8}
\]
The factorial inequality follows from
$2^{54}27!<2^{54+135}<2^{400}=m^2$.

For $3\leq k\leq51$, the last-variable cap is twenty-six
and at least one early variable is available. Choose
twenty-six distinct last variables and any power through
$m-26$ of one early variable. This gives
\[
 \dim\mathcal U_k\geq(m-25)\binom{r-50}{26}
       \geq\frac{m^{53}}{2^{27}26!}>m^{52},            \tag{O-018.9}
\]
using $m-25\geq m/2$, $r-75\geq r/2$, and
$2^{27}26!<2^{157}<m$.

For $52\leq k\leq r/2$, use fifty early variables and
last-variable degree exactly two. The dimension is at least
\[
 \binom{m+48}{50}\binom{r/2+1}{2}
       \geq\frac{m^{54}}{2^{53}50!}>m^{52},            \tag{O-018.10}
\]
because $2^{53}50!<2^{353}<m^2$.
These ranges exhaust all cuts, including the empty-tail
case in (O-018.8). Finally,
$(r-2)n^2+2=(m^2-2)m^{50}+2<m^{52}$.
This proves (O-018.7).

\paragraph{Other explicit gates and scope.}
The BIP threshold holds with equality. All F-017 staircase
inequalities hold: the rows through fifty-one in question
are twenty-six and row fifty-two is two. The full-multiplicity
stabilization criteria requiring $n\geq\lambda_2$ fail
since $\lambda_2=2n$. The universal quadratic and grouped
Chow degree criteria fail: their exact thresholds at $L=dm$
are at least $\binom d2>n$.
Grenet's full containment criterion fails because
$m=2^{200}>5001$ and hence $2^m-1>2^{5000}=n$.
The older coordinate-box test fails as well: some part in
every split has at least $r/2$ variables, so its count is
at least $\binom{r/2}{27}>n$.
The coarse padding inequalities hold since
$|\bar\lambda|=2n+2r+1172\leq md$ and $r=m^2$.

The finite certificate checks the exact count (O-018.5),
the three uniform lower bounds covering every cut, the
joint-row constant in (O-018.2), and every stated numerical
comparison. No enumeration of the enormous number of cuts
is needed. This is an excluded shape with actual multiplicity
at least four, and a strict reduction of the earlier explicit
conditions. It is not an actual $q>K_n$ obstruction.

\paragraph{Geometric dictionary.}
The entire nested profile selects a downward-closed finite
torus-weight coefficient subscheme. Passing from individual
row caps to this profile is a calculation with exterior
matrix coefficients and their Cauchy span, as in F-016;
the exact background locators are Macdonald I(4.3), p.~63,
and I Appendix~A(5.3)--(5.6), pp.~155--156, and SGA~3 I,
Proposition~4.7.3, pp.~35--36. The one-tree substitution is
an integral morphism to the linear-matrix scheme. The actual
comparison is obtained only after retaining every
highest-weight ideal copy in characteristic zero. The
dimension obstruction uses the image closure, while the
four permanent copies use multiplication in its integral
coordinate ring; these are distinct arguments on the actual
schemes, not deductions from a common virtual $K_0$ class.

\section{The dimension boundary for universal jet dominance}

\subsection{O-019: the next full jet order cannot be dominant at the same power}
Let $n\geq2$, $3\leq r\leq n^2$, and $0\leq b\leq n$.
Write a degree-$n$ form in variables $x_1,z_1,\ldots,z_{r-1}$
as
\[
 F=\sum_{|\alpha|\leq n}
                  c_\alpha x_1^{n-|\alpha|}z^\alpha.
\]
Its jet of order $b$ at the first coordinate direction means
the coefficient vector $(c_\alpha)_{|\alpha|\leq b}$,
equivalently the terms of total $z$-degree at most $b$ in
$F(1,z)$. The whole jet space has dimension
\[
                    D_b^{\mathrm{jet}}(r)=\binom{r-1+b}{b}.    \tag{O-019.1}
\]
If the jets of $r$-variable size-$n$ linear determinants
are Zariski dense in this entire space, necessarily
\[
                   D_b^{\mathrm{jet}}(r)\leq(r-2)n^2+2.       \tag{O-019.2}
\]
The same necessary condition holds when one first takes
the determinant orbit closure and then projects to jets.

\begin{proof}
The jet operation is a linear projection on the affine
coefficient space. Its restriction to the closure of the
$r$-variable determinant-pencil image has image closure of
dimension at most the dimension of that source. By O-010,
the latter is at most $(r-2)n^2+2$.
The identification of this pencil closure with
$\Omega_n\cap\operatorname{Sym}^n((\mathbb C^r)^*)$
is recalled in O-017 and already includes border limits.
A dense jet image would have dimension $D_b^{\mathrm{jet}}(r)$, proving
(O-019.2). The stars-and-bars count of monomials in $r-1$
variables through degree $b$ gives (O-019.1).
\end{proof}

Now fix an integer $c\geq1$, take $r=m^2$, and let
$n=m^c$. For a fixed $b$ the two sides have orders
\[
                   D_b^{\mathrm{jet}}(m^2)=\Theta_b(m^{2b}),
                \qquad (m^2-2)m^{2c}+2=\Theta(m^{2c+2}).
\]
Thus universal jet dominance is impossible for all sufficiently
large $m$ when $b\geq c+2$. In particular, the order
$b=c+1$ reached in M-013 for $c\geq6$ is the largest
fixed integer order compatible with dimension at this power
threshold. Raising that whole jet order by one cannot be
achieved by a different matrix construction, even one that
uses limits, without increasing the power range.

For a finite BIP-scale check, take $m=10^{15}$,
$r=m^2$, $n=m^{25}$, and $b=27$. Exact arithmetic gives
\[
                \binom{m^2+26}{27}>(m^2-2)m^{50}+2.
\]
The accompanying certificate records the two integers and
their difference. This is a dimension check of the complete
jet target, not a computation of a Kronecker coefficient or
of an actual multiplicity.

\paragraph{Why this does not conflict with the positive jet theorem.}
M-013 gives full determinant multiplicity for the indicated
highest weights by proving dominance only after discarding
the higher coefficient terms. Those higher coefficients are
free in the matrices producing a prescribed low-order jet.
Consequently jet dominance does not assert an exact determinant
for each truncated padded form, and it does not imply
containment of the whole truncated coefficient linear space
in the determinant orbit closure. O-016's exact-representation
boundary can therefore coexist with order-twenty-six jet
dominance in the BIP range.

The new dimension boundary limits dominance onto the
\emph{entire} next jet space. It leaves open projections
onto smaller nested supports, interpolation on a particular
permanent coefficient image, and direct comparisons of the
two evaluation ranks. In particular it is not a proof of
an actual multiplicity obstruction beyond the stated order.

\paragraph{Geometric dictionary.}
The finite jet coordinates form a quotient of the torus-graded
coefficient module, under SGA~3, I, Proposition~4.7.3,
pp.~35--36. The jet map is a morphism of affine coefficient
schemes; dominance is the injectivity of the associated
map on coordinate algebras. Its dimension cannot exceed
that of the determinant image closure. This argument uses
neither exact lifting of all jet points nor normality of
that image. The highest-weight comparison from such a
dominant projection is the separate characteristic-zero
all-copy argument in M-013.

\section{Full multiplicity from transverse jets along a binary plane}

\subsection{O-020: arbitrary second rows with third row at most two}
Let $r\geq3$ and $n\geq r-1$. For every outer degree $d$
and every partition $\lambda\vdash nd$ with
$\ell(\lambda)\leq r$ and $\lambda_3\leq2$,
\[
                         K_n(\lambda)=a_\lambda(d[n]).
                                                               \tag{O-020.1}
\]
Missing third parts are zero. There is no bound on
$\lambda_2$ and no upper bound on $d$.
The geometric input is dominance of the determinant projection
onto the entire coefficient space of degree-$n$ forms whose
total degree outside a fixed binary plane is at most two.
Higher transverse coefficients remain unrestricted.

Write the variables as $x,y,z_1,\ldots,z_R$, where $R=r-2$.
The target coefficient space is
\[
 \mathcal T_{n,2}=
 \operatorname{Sym}^n\mathbb C^2
 \oplus\left(\operatorname{Sym}^{n-1}\mathbb C^2
                               \otimes\mathbb C^R\right)
 \oplus\left(\operatorname{Sym}^{n-2}\mathbb C^2
                               \otimes\operatorname{Sym}^2\mathbb C^R\right).
                                                               \tag{O-020.2}
\]
Its dimension is $(n+1)+Rn+\binom{R+1}{2}(n-1)$.
The map from all $n$-by-$n$ linear pencils in these variables
to $\mathcal T_{n,2}$, obtained by deleting terms of
total $z$-degree at least three from the determinant,
has a surjective differential at an explicit integral pencil.

\paragraph{A difference-labeled upper-triangular pencil.}
Since $n\geq R+1$, put
\[
 \zeta_i=z_i\quad(1\leq i\leq R),\qquad
 \zeta_i=0\quad(R+1\leq i\leq n),\qquad
 u_{ij}=\zeta_i-\zeta_j\quad(i<j).
\]
Let $U=(u_{ij})$ be strictly upper triangular, choose distinct integers
$a_1,\ldots,a_n$, and define
\[
 \ell_i=y-a_ix,\qquad \Delta=\prod_{i=1}^n\ell_i,
             A_0=\operatorname{diag}(\ell_1,\ldots,\ell_n)-U.
                                                               \tag{O-020.3}
\]
Then $\det A_0=\Delta$.

\paragraph{Binary coefficients and residue coordinates.}
The $n$ polynomials $\Delta/\ell_i$ form a basis of the
binary forms of degree $n-1$. After setting $x=1$, their
evaluation matrix at $y=a_j$ is diagonal with nonzero
diagonal entries.
Send a polynomial $p(y)$ of degree
at most $n-2$ to the vector
\[
                \left(\frac{p(a_i)}{\Delta'(a_i)}\right)_{i=1}^n.
\]
Lagrange interpolation shows that its coordinates sum to
zero; the map is injective and both spaces have dimension
$n-1$, so it identifies this polynomial space with the
zero-sum hyperplane $H_n\subseteq\mathbb C^n$. The image of
$\Delta/((y-a_i)(y-a_j))$ is
$(e_i-e_j)/(a_i-a_j)$. Thus, after rescaling each independent
edge parameter by a nonzero scalar, the binary coefficients
$\Delta/(\ell_i\ell_j)$ can be treated as the vectors
$e_i-e_j$ in $H_n$.

\paragraph{A quadratic incidence isomorphism.}
First take $n=R+1$, and let $Z_1$ denote the linear forms
in $z_1,\ldots,z_R$. Consider
\[
 \Phi:\bigoplus_{1\leq i<j\leq R+1} Z_1
       \longrightarrow H_{R+1}\otimes\operatorname{Sym}^2Z_1,
 \qquad
 (f_{ij})\longmapsto
   \sum_{i<j}(e_i-e_j)(\zeta_i-\zeta_j)f_{ij}.
                                                               \tag{O-020.4}
\]
Its source and target both have dimension
$R\binom{R+1}{2}$. The map is an isomorphism.

To prove injectivity, write $f_{ji}=f_{ij}$ and
$f_{ij}=\sum_{k=1}^R c_{ij,k}z_k$. A kernel element satisfies
\[
          \sum_{j\neq i}(\zeta_i-\zeta_j)f_{ij}=0
                                  \quad(1\leq i\leq R+1).
\]
For distinct $i,j\leq R$, the coefficient of $z_j^2$
in the $i$-th equation gives $c_{ij,j}=0$.
The $j$-th equation likewise gives $c_{ij,i}=0$.
For three distinct $i,j,k\leq R$, the coefficient of
$z_jz_k$ in the $i$-th equation gives
$c_{ij,k}+c_{ik,j}=0$. Applying the same relation at
$j$ and $k$, and using $c_{ji,k}=c_{ij,k}$, shows that
$c_{ij,k}=-c_{ij,k}$. Hence it is zero in characteristic
zero. This also covers $R=1,2$, where there are no such
remaining triple coefficients. Every $f_{ij}$ with
$i,j\leq R$ therefore vanishes. The $i$-th equation now
reduces to $z_i f_{i,R+1}=0$, which forces the remaining
linear forms to vanish. This proves injectivity and the
dimension count proves surjectivity.

For $n>R+1$ add the vertices one at a time. The previous
image is $H_{n-1}\otimes\operatorname{Sym}^2Z_1$, embedded
with last coordinate zero. For a new vertex $v$, the
edges $\{i,v\}$, $1\leq i\leq R$, have labels $z_i$.
Their new residue coordinates are $-z_i f_{iv}$, which
span every quadratic polynomial, since every quadratic
monomial is a product of a variable with a linear form.
Choose these parameters to obtain any prescribed new
coordinate. The remaining error has new coordinate zero
and total coordinate sum zero, so it belongs to the old
image and can be removed by the old edge parameters.
This proves surjectivity onto
$H_n\otimes\operatorname{Sym}^2Z_1$ for every $n\geq R+1$.
It checks every additional residue coordinate, rather than
merely comparing the numbers of parameters.

\begin{proof}[Proof of the surjective differential]
The cofactor derivative formula is
\[
             d(\det)_{A_0}(H E_{ij})
                         =H\,\Delta\,(A_0^{-1})_{ji}.
                                                               \tag{O-020.5}
\]
The expression is polynomial although it is convenient to
compute the inverse over $\mathbb C(x,y,z)$.
Since $U$ is strictly upper triangular, an inverse entry
is a finite sum over increasing paths. In particular,
\[
 (A_0^{-1})_{ii}=\ell_i^{-1},\qquad
 (A_0^{-1})_{ij}
       =\frac{u_{ij}}{\ell_i\ell_j}
                    +\text{terms of total $z$-degree at least two}
                       \quad(i<j).                    \tag{O-020.6}
\]
Every extra path edge contributes another $z$ factor.
Multiplication by $\Delta$ clears the denominators,
because a strictly increasing path visits each vertex once.

Vary diagonal entries first. Taking $H=x$ or $H=y$
in (O-020.5) spans every binary degree-$n$ form:
the polynomials $\Delta/\ell_i$ span degree $n-1$,
and multiplication by $x,y$ spans the next degree.
Taking $H=z_a$ on diagonal entries spans independently
every term linear in the $z$ variables with an arbitrary
binary coefficient of degree $n-1$.

For any edge $\{i,j\}$, $i<j$, vary the opposite,
lower-triangular entry by an arbitrary $H\in Z_1$.
Equations (O-020.5)--(O-020.6), after deleting transverse
degree at least three, give exactly
\[
                  (\zeta_i-\zeta_j)H
                          \frac{\Delta}{\ell_i\ell_j}.
                                                               \tag{O-020.7}
\]
The terms from longer paths have degree at least three
after multiplication by $H$ and disappear from the
target projection. The residue identification and the
surjectivity of (O-020.4), including its added-vertex
extension, show that these directions span every quadratic
transverse polynomial with arbitrary binary coefficients
of degree $n-2$.

These directions therefore span every summand of
(O-020.2). Diagonal directions control transverse orders
zero and one, while the lower-triangular directions just
used have no contributions in those orders. All matrix-entry coefficients are
independent parameters. Thus the differential is surjective,
with no requirement to fix higher transverse terms.
The algebraic lowest-homogeneous-term argument in M-013
now proves dominance of this coefficient projection.
\end{proof}

\begin{proof}[Proof of the multiplicity equality]
Let $w$ be any degree-$d$ highest-weight vector of weight
$\lambda$ in the determinant ideal. It depends only on
the first $r$ coordinates. By the $j=3$ support bound
of F-016, every coefficient generator used by $w$ has
total degree at most $\lambda_3\leq2$ in variables
$3,\ldots,r$. Hence $w$ is a polynomial in the
coordinates of $\mathcal T_{n,2}$.
Every linear pencil determinant belongs to $\Omega_n$,
so this polynomial vanishes after the dominant projection
just proved. Its pullback is injective, so $w=0$.
This holds for every highest-weight copy, and therefore
the ideal multiplicity space vanishes entirely. The
ambient quotient then gives (O-020.1).
\end{proof}

For permanent shapes $r\leq m^2$, the condition
$n\geq m^2-1$ is enough for this full equality whenever
$\lambda_3\leq2$, with the one- and two-row cases already
covered by O-006. This conclusion is stronger than a
comparison with a particular permanent multiplicity: every
ambient copy survives, even with an arbitrarily large second
row. It is a different sufficient condition from M-013's
bounded second-row jet theorem.

\paragraph{An independently checkable tree construction.}
The earlier proof with the weaker sufficient bound
$n\geq2R=2r-4$ is retained in
\texttt{binary\_transverse\_jets\_tree.tex}.
It colors $R$ edge-disjoint spanning trees by the individual
variables $z_a$; each tree's edge cofactors form a binary
degree-$n-2$ basis. The paths
$j,(j-i,j+i)_{1\leq i\leq R-1},j-R$ on $\mathbb Z/(2R)$
give an explicit packing, including $R=1$.
Its independent cofactor calculations pass eight exact
Jacobian rank tests, including the smallest case $R=1,n=2$.
The difference-label argument above improves the threshold
by using the quadratic incidence isomorphism instead of
disjoint color classes.

\paragraph{The exact integral incidence cokernel.}
The particular map $\Phi$ in (O-020.4), with $n=R+1$,
is defined between free abelian groups of the same rank.
Its integral cokernel is
\[
          \operatorname{coker}(\Phi_{\mathbb Z})
                      \cong(\mathbb Z/2\mathbb Z)^{\binom R3}.
                                                               \tag{O-020.8}
\]
Here the target uses the ordinary monomial basis of
$\operatorname{Sym}^2\mathbb Z^R$. Delete the last residue
coordinate to identify $H_{R+1,\mathbb Z}$ with $\mathbb Z^R$.
The $R^2$ coefficients of the edges from $i$ to the zero
vertex occur with coefficient one in the rows given by
vertex $i$ and a monomial $z_i z_k$. They occur in no
other rows after the last residue coordinate is deleted.
Unimodular elimination removes these columns and rows.

Next, the two endpoint coefficients of each finite edge
$\{i,j\}$ occur with coefficient minus one in the remaining
rows $(i,z_j^2)$ and $(j,z_i^2)$, respectively.
They are unit pivots and can also be eliminated.
The remaining coordinates split by unordered triples
$\{i,j,k\}$. Each triple has the three variables
$c_{ij,k},c_{ik,j},c_{jk,i}$ and the three relations
\[
             a+b,\qquad a+c,\qquad b+c,
\]
up to signs. This $3$-by-$3$ block has Smith form
$\operatorname{diag}(1,1,2)$: a unit entry and a unit
$2$-by-$2$ minor give its first two invariant factors,
and its determinant has absolute value two. This proves
(O-020.8), including $R=1,2$, where there are no triples.
Thus this selected quadratic incidence submap is an
isomorphism after inverting two, and its rank drops by
$\binom R3$ after reduction modulo two.
No failure of the full determinant jet map in characteristic
two is inferred; other coefficient directions are not
included in $\Phi$.

\paragraph{Geometric dictionary.}
The target is the degree-$n$ part of the quotient by
$(z_1,\ldots,z_R)^3$, the second infinitesimal neighborhood
transverse to the binary coordinate plane. It is a finite
locally free coefficient module. The difference-labeled graph, diagonal
linear forms, and cofactor derivatives define integral
morphisms. Distinctness of the $a_i$ and the nonzero
Vandermonde factors and the precisely computed two-torsion
in (O-020.8) explain the passage to
$\mathbb Q$ or $\mathbb C$; no equal-rank assertion in
every residue characteristic is inferred.
The suffix support is a torus/comodule calculation under
SGA~3 I, Propositions~4.7.2--4.7.3, pp.~34--36,
and the Cauchy and Schur functors of Macdonald I(4.3),
p.~63, and Appendix~A(5.3)--(5.6), pp.~155--156.
Dominance injects the actual coordinate algebra of the
detected coefficients, which makes the whole highest-weight
ideal component vanish. Neither exact lifting of each
truncated form nor containment of that entire linear space
in the determinant closure is asserted.

\section{Full permanent multiplicities for a thin third row}

\subsection{O-021: a transverse permanent jet and padding theorem}
Let $r\geq3$, $n\geq m$, and let $\lambda\vdash nd$ satisfy
\[
       \ell(\lambda)\leq r,\qquad \lambda_3\leq2,\qquad
                    m\geq\max\{r,\lambda_2+1\}.
\]
Then, with the same internal padding as the supplied IP and
BIP papers,
\[
                    q^m_\lambda(d[n])
                         =K_n(\lambda)=a_\lambda(d[n]).
                                                               \tag{O-021.1}
\]
There is no upper bound on the outer degree $d$.
The statement determines all copies of these types on the
permanent side, rather than only their occurrence.

\paragraph{The unpadded permanent has the same transverse differential.}
For any $s\geq r-1$, use the variables
$x,y,z_1,\ldots,z_R$, $R=r-2$, and the difference-labeled
matrix of O-020, but change the signs on its upper-triangular
entries:
\[
        A_0=\operatorname{diag}(\ell_1,\ldots,\ell_s)+U,
        \qquad \ell_i=y-a_i x,\qquad
        U_{ij}=\zeta_i-\zeta_j\quad(i<j).
\]
Here $\zeta_i=z_i$ for $i\leq R$ and zero otherwise,
and the $a_i$ are distinct. The underlying directed graph
is acyclic. Consequently
$\operatorname{per}A_0=\Delta=\prod_i\ell_i$.

The permanent minor obtained by deleting row $j$ and
column $i$, with $i<j$, is the sum over directed paths
from $i$ to $j$ of the path's edge product multiplied
by the diagonal entries at its unused vertices.
Indeed every term of that minor consists of this path
and loops; an additional nontrivial cycle is impossible
in an acyclic graph. Conversely each path and its unused
loops gives exactly one minor term. In particular the
minor equals
\[
        (\zeta_i-\zeta_j)\frac{\Delta}{\ell_i\ell_j}
               +\text{terms of transverse degree at least two}.
                                                               \tag{O-021.2}
\]
The diagonal minors are $\Delta/\ell_i$.
This is an identity with weighted paths; negative coefficients
inside the difference labels cause no difficulty.

Vary diagonal entries by $x,y,z_a$, and lower-triangular
entries by linear forms in $z$. After deleting transverse
degree at least three, their permanent differentials are
exactly the selected determinant differentials in O-020.
The quadratic incidence isomorphism and the added-vertex
argument therefore apply unchanged. The permanent map
from all $s$-by-$s$ linear pencils dominates the full space
\[
 \operatorname{Sym}^s\mathbb C^2
 \ \oplus\ \bigl(\operatorname{Sym}^{s-1}\mathbb C^2
                                      \otimes\mathbb C^R\bigr)
 \ \oplus\ \bigl(\operatorname{Sym}^{s-2}\mathbb C^2
                               \otimes\operatorname{Sym}^2\mathbb C^R\bigr).
                                                               \tag{O-021.3}
\]
Its higher transverse coefficients are unrestricted.

\begin{proof}[Proof of the full multiplicity equality]
If $a_\lambda=0$, both quotient multiplicities vanish,
so assume it is positive. One-row types have
$q=K=a=1$, since the highest-weight function $F(e_1)^d$
is nonzero on both closures. Thus assume $\lambda_2>0$.
Then $\ell(\lambda)\leq d$.
Set $s=m-1$ and $b=\lambda_2$. The hypotheses give
$s\geq r-1$ and $s\geq b$. By F-014, the complete
ambient highest-weight multiplicity spaces at inner degrees
$s$ and $n$ have the same dimension:
\[
 a_\mu(d[s])=a_\lambda(d[n]),\qquad
             \mu=(sd-|\bar\lambda|,\bar\lambda).
                                                               \tag{O-021.4}
\]
Here $\bar\lambda=(\lambda_2,\lambda_3,\ldots)$.
The first part of $\mu$ is at least its second part,
because $|\bar\lambda|\leq(d-1)b$ and $s\geq b$.

The padding lift $T$ of O-012.b is injective from the
entire highest-weight space of weight $\mu$ into that
of weight $\lambda$. Equation (O-021.4) makes it an
isomorphism of these multiplicity spaces.
For the actual padding map
$P(F)=x^{n-s}F$, its pullback on $T$ is the diagonal
operator from the differential pairing in O-012.b:
\[
 P^*T=\operatorname{Sym}^d(\mathcal D),\qquad
 \mathcal D(e^\alpha)
        =\frac{(\alpha_1+n-s)!}{\alpha_1!}\,e^\alpha,
                 \qquad |\alpha|=s.                   \tag{O-021.5}
\]
Every diagonal entry is nonzero in characteristic zero.
The operator also preserves every monomial support condition.
It need not preserve the highest-weight subspace; that
additional property is not used.

By F-016 at its third-row cut, every polynomial in the
highest-weight space of weight $\mu$ uses only coefficient
generators of transverse degree at most two, since
$\mu_3=\lambda_3\leq2$. Thus the same is true after
applying the invertible operator in (O-021.5).
Any nonzero linear combination of the lifted copies
therefore pulls back to a nonzero polynomial in the
coordinates of (O-021.3). The dominant permanent map
proved above cannot annihilate that polynomial.
This proves independence of every lifted copy on the
padded permanent-pencil family.

Finally that entire family belongs to the required
internally padded permanent closure. For an arbitrary
$s$-by-$s$ pencil $A$, use the $m$-by-$m$ block matrix
\[
                         \widetilde A=\operatorname{diag}(x,A).
\]
Its distinguished entry is exactly $\widetilde A_{11}=x$,
and
\[
 \widetilde A_{11}^{\,n-m}\operatorname{per}\widetilde A
                      =x^{n-s}\operatorname{per}A.
                                                               \tag{O-021.6}
\]
These are linear endomorphism substitutions in
$X_{11}^{n-m}\operatorname{per}_m$. They belong to its
orbit closure because invertible endomorphisms are dense.
The variable space fits: $r\leq m\leq m^2\leq n^2$.
The independence just proved forces the entire permanent
ideal highest-weight space of type $\lambda$ to vanish.
Hence $q^m_\lambda=a_\lambda$.
O-020 also gives $K_n=a_\lambda$, since $n\geq m\geq r$.
This proves (O-021.1), including $n=m$.
\end{proof}

\paragraph{A strict four-copy example below every exact support-cover size.}
For a downward-closed support $\Gamma$ containing
$x_2^b$ and $x_i^2$ for every $3\leq i\leq r$, with
$b\geq2$, the two-tree invariant of F-018 satisfies
\[
                            \kappa(\Gamma)\geq b+r-2.
                                                               \tag{O-021.7}
\]
To prove this, let $P,Q$ be any admissible pair.
A decomposition of the pure power $b e_2$ forces
$a e_2\in P$, $a'e_2\in Q$ with $a+a'\geq b-1$.
Downward closure supplies both full axis chains, so these
alone contribute at least $a+a'+1\geq b$ to
$|P|+|Q|-1$. For $2e_i$, at least one of the two
endpoints has a positive $i$-th coordinate: the coefficient
edge can contribute at most one such unit. Downward closure
then forces $e_i\in P\cup Q$. These $r-2$ distinct
off-axis vertices are additional to the two chains.
This proves (O-021.7).

Take the concrete values
\[
       m=9,\quad n=10,\quad d=r=6,\qquad
                    \lambda=(44,8,2,2,2,2).
                                                               \tag{O-021.8}
\]
O-021.1 applies because $m=\max\{6,8+1\}$.
The exact near-rectangle formula F-013 gives
\[
      a_\lambda(6[10])=p_6(6,8)-p_5(6,8)=11-7=4.
\]
Thus
\[
                        q^9_\lambda(6[10])
                                =K_{10}(\lambda)=4.   \tag{O-021.9}
\]
Its nested support is
$\Gamma=\{\gamma:|\gamma|\leq8,\
                  \gamma_3+\cdots+\gamma_6\leq2\}$.
It contains the powers required by (O-021.7), so
$\kappa(\Gamma)\geq12>n>m$. Since the total-degree
cap is already eight, this is the same support at base
degrees $8,9,10$. Every exact two-tree support-cover
criterion of F-018 therefore misses this example, while
the transverse dominant projection and the all-copy lift
determine its actual multiplicities. This is a limitation
of that exact-cover architecture, not a lower bound for
arbitrary determinant or permanent representations.

\paragraph{Geometric and integral scope.}
The weighted path minors and the block substitution
(O-021.6) are integral polynomial identities. The
transverse differential uses the nonzero Vandermonde
factors and the selected incidence map whose exact
two-torsion is computed in O-020.8. The padding pullback
also uses invertible factorial scalars in characteristic
zero. No equality of integral Schur decompositions after
arbitrary base change is asserted.
The torus support is the comodule grading of SGA~3 I,
Propositions~4.7.2--4.7.3, pp.~34--36; F-014 and F-016
use the Cauchy and exterior functors of Macdonald I(4.3),
p.~63, Appendix~A(5.3)--(5.6), pp.~155--156, and
SGA~6 VI, Theorem~3.3, pp.~392--393.
The argument retains the injective coordinate pullback
through the dominant permanent projection and the actual
padding map. It supplies multiplicity information beyond
the virtual lambda class, with no saturation inference.

\section{A positive BIP example detected beyond exact support covers}

\subsection{O-022: the prefix jet criterion makes a strict further reduction}
Set
\[
 m=2^{200},\qquad r=m^2,\qquad n=m^{25}=2^{5000},
 \qquad d=3n,
\]
and put
\[
           \lambda=(3n^2-26r+20,\ 32,\ 26^{r-2}).
                                                               \tag{O-022.1}
\]
This is a partition of $nd$ with $r$ parts. It satisfies
\[
                  4\leq q^m_\lambda(d[n])
                                \leq K_n(\lambda)=a_\lambda(d[n]).
                                                               \tag{O-022.2}
\]
The exact values of these multiplicities are not asserted.
The full multiplicity conclusion follows from M-017.
Every exact two-tree support cover of F-018 has size greater
than $n$ for this shape. Moreover every entire one-cut
universal family is too large, even in border dimension,
and the entire point-jet target is also too large.
Thus (O-022.2) is a strict reduction beyond those proved
criteria, with actual permanent multiplicity at least four.

\paragraph{The prefix extension fits.}
Use M-017 with $b=26$, $\ell=12$, $t=6$, and $k=3$.
Here $p=1$, $\lceil\sqrt{r-1}\rceil=m$, and the size is
\[
                       N=8m\binom{m^2+11}{12}.
                                                               \tag{O-022.3}
\]
The exact integer constant in M-017 is
\[
 C_*=\left\lfloor\frac{190^{24}}{\binom{36111}{12}}\right\rfloor
                  =478126755.
                                                               \tag{O-022.4}
\]
The saved certificate gives the quotient and positive
remainder and verifies the defining two integer inequalities.
The normalized binomial ratio is nonincreasing in $m$.
Since $8<C_*$ and $m\geq190$, (O-022.3) is strictly
less than $m^{25}=n$. The target has total nonfirst
degree at most $32$ and degree at most $26$ in variables
$3,\ldots,r$. These are exactly the relevant two nested
bounds of $\lambda$. M-017 therefore gives $K_n=a_\lambda$.

\paragraph{Four permanent copies survive.}
M-012 supplies four independent permanent-closure
highest-weight classes of degree six and weight
$(6n-6,6)$. Write $c_0(F)=F(e_1)$ and let $W_r$
be the leading Hessian determinant from O-011.
Multiply each class by the common highest-weight function
\[
                          W_r^{13}c_0^{\,d-6-13r}.
                                                               \tag{O-022.5}
\]
Its exponent is nonnegative. O-011 gives a permanent
point and basis where $W_r$ and $c_0$ are simultaneously
nonzero, since $r=m^2$. The coordinate algebra of the
irreducible permanent orbit closure is a domain, so
multiplication by (O-022.5) preserves all four independent
classes. Their total degree is $d$; their first weight is
$nd-26r+20$, their second is $6+26=32$, and every
remaining weight is $26$. Hence $q^m_\lambda\geq4$.
The ambient quotient bound and the preceding equality
give (O-022.2).

\paragraph{Every exact support cover is too large.}
Let $\Gamma$ denote the full nested support at base degree
$32$ or at base degree $m$. Both supports are the same:
\[
       \Gamma=\{\gamma:|\gamma|\leq32,\
                            \gamma_3+\cdots+\gamma_r\leq26\}.
                                                               \tag{O-022.6}
\]
It contains every homogeneous degree-$26$ monomial in
the $R'=r-2$ last variables. If a support cover had
size $S\leq n$, its integral determinant identity would
represent every such form $f$ as
$x_1^{S-26}f$. Here $S\geq26$ by F-018.
Specializing $x_1=1$ and the unused early variable $x_2=0$
would give an affine linear determinant of size at most
$n$ for every homogeneous degree-$26$ form in $R'$ variables.

This contradicts the exact generic lower bound O-016.
Indeed its middle-layer ratio is
\[
 \frac{D_{26}(R')}{2D_{13}(R')}
     =\frac{13!}{2\cdot26!}\prod_{j=13}^{25}(R'+j)
     \geq\frac{13!\,m^{26}}{2^{14}\,26!}>m^{25}=n.
                                                               \tag{O-022.7}
\]
We used $R'=m^2-2\geq m^2/2$ and
$2^{14}26!/13!<2^{79}<m$. Thus a generic form has
exact determinantal complexity greater than $n$, and
$\kappa(\Gamma)>n$. This argument concerns exact
representations. It does not infer a border-complexity
lower bound for this restricted family.

\paragraph{Every one-cut enlargement fails even by dimension.}
The determinant-pencil closure in $r$ variables has
dimension at most
\[
                      (r-2)n^2+2<m^{52}.
                                                               \tag{O-022.8}
\]
For a cut $k$, the entire universal degree-$m$ family
allows unrestricted early degree in variables
$2,\ldots,k-1$ and suffix degree at most $\lambda_k$,
with total degree at most $m$.

At $k=2$, it contains all squarefree degree-$27$
monomials in the last variables, because $\lambda_2=32$.
At every cut $k>r/2$, including $k=r+1$, there are at
least $r/2-1$ unrestricted early variables. Choosing
$27$ of them gives in both cases the common lower bound
\[
     \binom{r/2-1}{27}
          \geq\frac{m^{54}}{2^{54}27!}>m^{52}.
                                                               \tag{O-022.9}
\]
The last inequality follows from
$2^{54}27!<2^{189}<m^2$.

For $3\leq k\leq r/2$, there are at least $r/2$
suffix variables, and their allowed total degree is
$26$. Choose $26$ distinct suffix variables and an
arbitrary exponent from $0$ through $m-26$ on $x_2$.
The family therefore has dimension at least
\[
 (m-25)\binom{r/2}{26}
           \geq\frac{m^{53}}{2^{53}26!}>m^{52},
                                                               \tag{O-022.10}
\]
using $m-25\geq m/2$ and
$2^{53}26!<2^{183}<m$.
These three ranges exhaust all cuts. Multiplication by
the padding power is injective, so the same dimensions
hold for the padded linear spaces. Equations
(O-022.8)--(O-022.10) exclude containment in the
determinant closure, including border containment.
The full point-jet target through degree $32$ contains
the monomials counted in (O-022.9), so universal dominance
onto that entire target is likewise impossible by O-019.
The restricted projection of M-017 avoids this dimension
obstruction by retaining both nested bounds.

\paragraph{The preceding explicit candidate conditions are met.}
The shape satisfies $\ell(\lambda)=m^2$,
$|\bar\lambda|=26m^2-20\leq md$, and
$\ell(\lambda)\leq d$. It meets the complete BIP
staircase and the joint A-025 condition, since
$\lambda_2=32$, all subsequent parts are $26$, and
$\lambda_{52}=26$. It is outside O-020's
third-row-at-most-two range.
The uniform exact Chow degree threshold $P_d(dm)$ of
A-023 is at least $\binom d2>n$ here.
Its stronger shape-specific threshold also fails:
$P_2(L)=1$ and every recurrence step increases by at
least one for $L\geq1$, so
$P_d(32)\geq d-1=3n-1>n$.
The whole-permanent Grenet threshold also fails:
$2^m-1>2^{5000}=n$.

For completeness, every coordinate-box split of M-009
has a group containing at least $r/2$ variables.
All individual caps are at least $26$, so it contains
all squarefree monomials of degree $m$. Its size is at least
\[
                 \binom{r/2}{m}\geq(m/2)^m>m^{25}.
                                                               \tag{O-022.11}
\]
The binomial inequality follows termwise from
$(r/2-i)/(m-i)\geq(r/2)/m$.
The final comparison follows from $m=2^{200}$.
All exact support-cover variants are already excluded
more generally by (O-022.7).
These checks establish strict refinement of the listed
proved conditions. They do not exhibit an actual
multiplicity obstruction: (O-022.2) proves the opposite
comparison for this particular shape.

\paragraph{Geometric dictionary and exact versus border scope.}
The target (O-022.6) is a finite torus-weight coefficient
module under SGA~3 I, Proposition~4.7.3, pp.~35--36.
M-017 gives a universally injective pullback on its
polynomial coordinate algebra, and F-016 identifies every
relevant Schur-copy coefficient using the exterior functors
of SGA~6 VI, Theorem~3.3, pp.~392--393, and Macdonald
I(4.3), p.~63, Appendix~A(5.3)--(5.6), pp.~155--156.
The four copies in (O-022.5) are multiplied in the actual
integral coordinate algebra of an irreducible orbit closure,
not only in its virtual representation class.
Equation (O-022.7) uses an exact-complexity bound and is
kept separate from the border-valid dimension statements
(O-022.8)--(O-022.10). The latter use the actual reduced
image closure and its projection, as in O-010.
Higher coefficients remain free in the dominant prefix
projection; no exact lift of every truncated form is asserted.

\section{A further positive example requiring a nonsimplex extension}

\subsection{O-023: the order-ideal extension improves every prefix-simplex test}
Keep
$m=2^{200}$, $r=m^2$, $n=m^{25}$, and $d=3n$,
and now define
\[
 \lambda=(nd-26r-1000000+50,\ 1000000,\ 28,\ 26^{r-3}).
                                                               \tag{O-023.1}
\]
Then
\[
                 4\leq q^m_\lambda(d[n])
                               \leq K_n(\lambda)=a_\lambda(d[n]).
                                                               \tag{O-023.2}
\]
Every choice of parameters in the prefix-simplex criterion
M-017 fails to reach this shape at size $n$.
The exact support-cover and one-cut failures of O-022
also remain. Thus M-018 makes a further strict reduction
of the proved candidate conditions.
No exact numerical value of the multiplicities in
(O-023.2) is asserted.

\paragraph{The nonsimplex extension fits.}
For the order ideal
\[
 A=\{x_2^a x_3^c:\ a+c\leq999974,\quad c\leq28\}
\]
one has
\[
       |A|=29\cdot999975-\frac{28\cdot29}{2}
                          =28998869.
                                                               \tag{O-023.3}
\]
The full nested support of (O-023.1) has total degree
at most one million, suffix degree from coordinate three
at most twenty-eight, and suffix degree from coordinate
four at most twenty-six.
If its monomial degree exceeds twenty-six, choose an
early submonomial in $x_2,x_3$ of the excess degree.
There are enough early factors because the later suffix
has degree at most twenty-six. Its degree is at most
$999974$, and its $x_3$ exponent is at most twenty-eight,
so it lies in $A$. The remaining monomial has degree
twenty-six. Smaller monomials use the zero element of $A$.
Thus the full support is contained in $A+D_{26}$.

Use M-018 with $b=26$ and
$T=m\binom{m^2+11}{12}$.
The size is $28998870T<n$, because
$28998870<C_*=478126755$ and the uniform normalized
estimate of M-017 applies. This proves $K_n=a_\lambda$.

\paragraph{Four positive permanent copies.}
Multiply the four degree-six binary seed classes of M-012
by
\[
             W_r^{13}\,W_3\,W_2^{499983}\,
                         c_0^{\,d-13r-999975}.          \tag{O-023.4}
\]
The exponent is nonnegative. O-011 supplies one basis
at a permanent point where every displayed leading
Hessian minor and $c_0$ is nonzero. Therefore the common
factor is nonzero in the integral coordinate algebra of
the irreducible permanent closure, and multiplication
preserves the four independent seed classes.
Their degree is
$6+13r+3+2\cdot499983+(d-13r-999975)=d$.
Their second weight is
$6+26+2+999966=1000000$, their third weight is $28$,
and all later weights are $26$. The first weight is
the remaining part $nd-26r-999950$, as required.
This proves $q^m_\lambda\geq4$ and hence (O-023.2).

\paragraph{All prefix-simplex choices fail.}
Consider any M-017 parameter choice that could cover this
highest-weight support. Its even base order is $b=2\ell+2$,
and it must satisfy $\lambda_2\leq b+t$ and
$\lambda_k\leq b$. Taking more ambient variables than
the actual length $r$ only increases its size; it suffices
to give lower bounds using this length.

If $b\geq28$, then $\ell\geq13$, and even the unextended
core has size at least
\[
     2m\binom{r+12}{13}
           \geq\frac{2m^{27}}{13!}>m^{25}=n.             \tag{O-023.5}
\]
If $b=26$, the first possible cut is $k=4$, since
$\lambda_3=28$. Therefore $p=k-2\geq2$ and
$t\geq999974$. Its size is at least
\[
 \frac{m^{25}}{12!}
       \left(1+\binom{999976}{2}\right)>m^{25},
                                                               \tag{O-023.6}
\]
because $T\geq m^{25}/12!$ and the parenthesized integer
is larger than $12!$.
Finally, if $b\leq24$, every actual part is larger than
$b$, so the cut must be at or after $r+1$.
Then $p\geq r-1$, $t\geq13$, and $T\geq m$.
Its size is at least
\[
                    m\binom{r+12}{13}
                          \geq\frac{m^{27}}{13!}>n.    \tag{O-023.7}
\]
These cases exhaust all allowed even base orders.
The strict inequalities use $m^2>13!$, and are verified
by the exact certificate. Thus no M-017 parameter choice
satisfies its size condition.

\paragraph{The earlier failures persist.}
The nested support (O-022.6) is contained in the present
support: its total degree is at most $32$, its suffix
from the third coordinate is at most $26$, and therefore
its suffix from the fourth coordinate is at most $26$.
In particular the present support contains all degree-$26$
forms in the last $r-2$ variables. The exact-complexity
argument (O-022.7) again forces every F-018 exact support
cover to have size greater than $n$.

Each part of (O-023.1), except the first, is at least
the corresponding part of (O-022.1).
Every one-cut family and coordinate box is consequently
an enlargement of its O-022 counterpart. The border-valid
dimension failures (O-022.8)--(O-022.10) and the box
size failure (O-022.11) remain valid.
The full point-jet target is still too large, since it
contains the degree-$27$ monomials used there.
The shape meets the BIP staircase, A-025's joint gate,
and the M-017 coupled necessary conditions; the complete
M-017 size failure was proved above.
Its tail is $26r+999950\leq md$ and its length is $m^2$.
The shape-specific Chow threshold is at least
$d-1=3n-1>n$; the uniform threshold and Grenet failure
are unchanged. O-020 does not apply because $\lambda_3=28$.

These checks concern exact proved exclusions, including
ones with actual multiplicity at least four. They do not
claim that (O-023.1) is an obstruction: (O-023.2)
proves that it is not.

\paragraph{Geometric dictionary.}
The set $A$ is a finite monomial order ideal in an integral
torus-graded coefficient module, as in SGA~3 I,
Proposition~4.7.3, pp.~35--36. M-018 gives its actual
universally injective coordinate pullback. The full
highest-weight support is the Schur/exterior calculation
of F-016, with SGA~6 VI, Theorem~3.3, pp.~392--393,
and Macdonald I(4.3), p.~63, Appendix~A(5.3)--(5.6),
pp.~155--156. The domain multiplication retains four
sections of the actual orbit-closure coordinate algebra.
The exact-complexity obstruction and the separate
border dimension arguments retain the distinct scopes
spelled out in O-022. No saturation statement follows
from a virtual class alone.

\section{A canonical positive count for order-ideal jet extensions}

\subsection{O-024: subtracting the jet order from every suffix bound}
Let $\lambda$ have length at most $r\geq2$, and let
$b\geq0$ be an integer. Set
\[
           \mu_j=(\lambda_j-b)_+\quad(2\leq j\leq r),
            \qquad u_+=\max(u,0),
\]
and write $\Gamma_\lambda$ for the full nested suffix
support of M-018. Then
\[
                         \Gamma_\lambda\subseteq\Gamma_\mu+D_b.
                                                               \tag{O-024.1}
\]
Moreover
\[
        |\Gamma_\mu|
            =\#\{\nu:\nu\text{ is a partition and }
                         0\leq\nu_j\leq\mu_j\ (2\leq j\leq r)\}.
                                                               \tag{O-024.2}
\]
Thus the needed extension size has a direct positive count:
the number of subpartitions of the tail obtained by
subtracting $b$ and replacing negative entries by zero.
This is an explicit consequence of the general M-018
criterion, not a stronger existential criterion.

\begin{proof}
Take $\gamma\in\Gamma_\lambda$. Remove at most $b$
variable factors, always starting with the largest index,
and call their multiset $\eta$. If $|\gamma|\leq b$,
remove all factors. Set $\alpha=\gamma-\eta$.
For every suffix beginning at $j$, this procedure removes
exactly the smaller of $b$ and that suffix's original
size before removing any earlier variable. Consequently
\[
         \sum_{i=j}^r\alpha_i
                =\left(\sum_{i=j}^r\gamma_i-b\right)_+
                \leq(\lambda_j-b)_+=\mu_j.
                                                               \tag{O-024.3}
\]
Hence $\alpha\in\Gamma_\mu$ and $\eta\in D_b$,
proving (O-024.1).

For (O-024.2), send $\alpha\in\Gamma_\mu$ to its
suffix cumulative sums
$\nu_j=\sum_{i=j}^r\alpha_i$.
They form a weakly decreasing sequence of nonnegative
integers bounded componentwise by $\mu$.
Conversely any such subpartition gives
$\alpha_j=\nu_j-\nu_{j+1}$, with $\nu_{r+1}=0$.
These two operations are inverse, proving the positive count.
\end{proof}

For an even jet order $b=2\ell+2$, put
\[
 T=\lceil\sqrt{r-1}\rceil\binom{r-1+\ell}{\ell},
 \qquad c(\mu)=\#\{\nu:\nu\subseteq\mu\}.
\]
M-018 now gives the explicit sufficient condition
\[
 n\geq T(1+c(\mu))
       \quad\Longrightarrow\quad K_n(\lambda)=a_\lambda(d[n]).
                                                               \tag{O-024.4}
\]
If $n\geq m\geq T(1+c(\mu))$, it also gives
$q^m_\lambda=K_n=a_\lambda$ with internal padding.
These are full-copy conclusions at arbitrary outer degree.

\paragraph{An exact positive recurrence.}
The count requires no cancellation. For $0\leq t\leq\mu_r$
put $F_r(t)=1$, and put it zero otherwise.
Recursively define
\[
 F_j(t)=\mathbf1_{\{0\leq t\leq\mu_j\}}
                         \sum_{s=0}^t F_{j+1}(s)
                    \quad(j=r-1,\ldots,2).
                                                               \tag{O-024.5}
\]
Here $F_j(t)$ counts allowed suffixes with cumulative
sum $t$ at position $j$; choosing the next cumulative
sum $s\leq t$ proves the recurrence.
Finally $c(\mu)=\sum_{t=0}^{\mu_2}F_2(t)$.
This is a finite positive computation, not a claim of
polynomial complexity in a binary encoding of the row lengths.

\paragraph{A sharper explicit three-row condition in BIP range.}
Every actual obstruction at $n\geq m^{25}$ satisfies
\[
 \boxed{\lambda_2\geq159375611\quad\text{or}\quad
          \lambda_3\geq29\quad\text{or}\quad\lambda_4\geq27.}
                                                               \tag{O-024.6}
\]
Indeed suppose the three inequalities fail. With $b=26$,
the shaved tail has at most two nonzero parts, bounded
by $(159375584,2)$. If a two-part partition is $(A,C)$,
with $A\geq C$, its number of subpartitions is
\[
                  \sum_{j=0}^C(A-j+1)
                 =(C+1)(A+1)-\frac{C(C+1)}2.
                                                               \tag{O-024.7}
\]
Thus here
$c(\mu)\leq478126752$ and
$1+c(\mu)\leq478126753<C_*=478126755$.
The uniform normalized estimate of M-017 makes
$T(1+c(\mu))<m^{25}\leq n$, so (O-024.4)
gives $K_n=a_\lambda$ and excludes an actual obstruction.
The exact arithmetic also gives
$3\cdot159375611-77=478126756>C_*$;
the displayed constant is the largest obtained by this
particular two-part count with the stated other caps.
It is not optimal among all order-ideal choices.

\paragraph{The canonical order ideal need not be smallest.}
For the full degree-four simplex in two variables, with
$b=2$, the canonical shaved ideal is the degree-two
simplex and has six members. The five-member order ideal
\[
                     \{1,x,x^2,y,y^2\}
\]
also covers the degree-four simplex after adding $D_2$:
a monomial of degree at most two uses the member $1$;
one of degree three or four has some exponent at least
two, so remove $x^2$ or $y^2$.
Thus (O-024.2) is an explicit sufficient count, not
the solution of the finite optimization in M-018.

\paragraph{Grothendieck dictionary.}
The exponent sets are finite order ideals in a torus
character monoid, and their free coefficient modules
are comodules as in SGA~3 I, Proposition~4.7.3,
pp.~35--36. The positive count (O-024.2) indexes
an explicit integral monomial basis. M-018 supplies
the actual universally injective coordinate pullback;
F-016's Schur/exterior support calculation uses SGA~6
VI, Theorem~3.3, pp.~392--393, and Macdonald I(4.3),
p.~63, Appendix~A(5.3)--(5.6), pp.~155--156.
The partition bijection gives the size of a coefficient
module. It does not by itself give positivity of a new
Kronecker coefficient or saturation of its support;
the multiplicity conclusion still uses the proved
dominant coefficient morphism.

\section{A positive example beyond every one-sided jet extension}

\subsection{O-025: strictness of the two-sided order-ideal criterion}
Set
\[
 m=2^{200},\qquad r=m^2,\qquad n=m^{25},\qquad d=3n,
             \qquad B_0=1000000,
\]
and define
\[
 \lambda=(nd-26r-1999928,\ B_0+6,\ B_0,\ 26^{r-3}).
                                                               \tag{O-025.1}
\]
Then
\[
                  4\leq q^m_\lambda(d[n])
                                \leq K_n(\lambda)=a_\lambda(d[n]).
                                                               \tag{O-025.2}
\]
The two-sided construction M-019 proves the ambient
determinant equality. No choice of base jet order and
finite one-sided order ideal in M-018 can prove that
equality at size $n$ for this support. This compares
all of M-018's admissible choices, not just a particular
canonical ideal or a prefix simplex.
The exact values of the multiplicities remain unspecified.

\paragraph{The two-sided extension fits.}
Use the canonical shaving O-024 with $b=26$.
The shaved support has only two nonzero bounds:
\[
                       (B_0-20,\ B_0-26).
\]
It is contained in the sum of the two axis ideals
\[
 \mathcal A=\{1,x_2,\ldots,x_2^{B_0-20}\},\qquad
 \mathcal B=\{1,x_3,\ldots,x_3^{B_0-26}\}.
\]
Therefore $\Gamma_\lambda\subseteq\mathcal A+\mathcal B+D_{26}$.
At this perfect-square length, both
$\lceil\sqrt r\rceil$ and $\lceil\sqrt{r-1}\rceil$ equal $m$.
With
$T=m\binom{r+11}{12}$, the two-sided size is
\[
          N=(|\mathcal A|+|\mathcal B|)T
                =(2B_0-44)T=1999956T<n.                \tag{O-025.3}
\]
The final inequality follows from $1999956<C_*=478126755$
and M-017's normalized estimate. M-019 gives $K_n=a_\lambda$.
No uniform dominance relation between the two criteria
at other, nonsquare ambient lengths is needed here.

\paragraph{Four positive permanent copies.}
Multiply the four degree-six binary classes of M-012 by
\[
                      W_r^{13}\,W_3^{499987}\,
                             c_0^{\,d-13r-1499967}.     \tag{O-025.4}
\]
The last exponent is nonnegative, and the Hessian argument
O-011 makes all factors simultaneously nonzero on the
permanent closure. Its coordinate algebra is an integral
domain, so all four classes remain independent.
The total degree is
$6+13r+3\cdot499987+(d-13r-1499967)=d$.
The second weight is
$6+26+2\cdot499987=1000006$, the third is
$26+2\cdot499987=1000000$, and the later weights are $26$.
Their first weight is the remaining part in (O-025.1).
This proves $q^m_\lambda\geq4$ and hence (O-025.2).

\paragraph{Excluding every one-sided choice.}
Suppose, for a contradiction, that M-018 applies with
some even base order $b=2\ell+2$ and a nonempty finite
order ideal $A$:
\[
 \Gamma_\lambda\subseteq A+D_b,\qquad
             T_b(1+|A|)\leq n,\qquad
             T_b=m\binom{r-1+\ell}{\ell}.
                                                               \tag{O-025.5}
\]
Taking more ambient variables can only increase this
core size. Monomials using those additional variables
cannot participate in a nonnegative decomposition of
an exponent supported in the first $r$ variables, so
discarding them from $A$ retains every decomposition
needed below. Thus (O-025.5) is the weakest relevant
size bound even if a larger ambient length was chosen.

If $b\geq28$, then $\ell\geq13$ and
\[
 T_b(1+|A|)\geq2m\binom{r+12}{13}
                       \geq\frac{2m^{27}}{13!}>n.
                                                               \tag{O-025.6}
\]
This already contradicts feasibility.

If $b=26$, then $T_b\geq n/12!$, so
$|A|<12!=479001600$. There are $r-3>12!$ suffix
variables $x_4,\ldots,x_r$. At least one of their unit
monomials, say $x_j$, is missing from $A$.
Downward closure implies that no member of $A$ uses $x_j$.
For every two-variable monomial $\nu(x_2,x_3)$ of total
degree at most $B_0-26=999974$, the monomial
$x_j^{26}\nu$ lies in $\Gamma_\lambda$.
In a decomposition of its exponent as $\alpha+\eta$,
with $\alpha\in A$ and $|\eta|\leq26$,
the entire $x_j^{26}$ must belong to $\eta$.
It exhausts that degree budget, so $\eta=26e_j$ and
$\alpha=\nu$. Consequently $A$ contains the entire
two-variable simplex of degree $999974$, and
\[
                    |A|\geq\binom{999976}{2}
                        =499975500300>12!,
                                                               \tag{O-025.7}
\]
a contradiction.

It remains to treat $2\leq b\leq24$. Then
$0\leq\ell\leq11$ and
$k=26-b=24-2\ell\geq2$.
Put $R'=r-3$. The support $\Gamma_\lambda$ contains
every squarefree degree-$26$ monomial in those last
$R'$ variables. In its decomposition $\gamma=\alpha+\eta$,
the exponent $\alpha$ is squarefree and has degree at
least $k$. By downward closure, $A$ contains a
squarefree $k$-submonomial of $\gamma$.
Thus the squarefree degree-$k$ members of $A$ meet every
$26$-subset of the suffix variables. Double-counting
the incidences gives
\[
 |A|\geq
 \frac{\binom{R'}{26}}{\binom{R'-k}{26-k}}
       =\frac{\binom{R'}k}{\binom{26}k}
       \geq\frac{r^k}{2^k\,26!}.                       \tag{O-025.8}
\]
The final inequality follows because every numerator
factor in $\binom{r-3}k$ is at least $r/2$;
the additional factor $(26-k)!$ was harmlessly omitted.
Also
$T_b\geq m r^\ell/\ell!$ and $n=m r^{12}$.
It follows that
\[
 \frac{T_b|A|}{n}
       \geq\frac{r^{12-\ell}}{2^k\ell!\,26!}>1.
                                                               \tag{O-025.9}
\]
Indeed $12-\ell\geq1$ and
$2^k\ell!\,26!\leq2^{24}11!\,26!<2^{139}<r=2^{400}$.
This contradicts (O-025.5) for every remaining base order.
All cases are exhausted, proving the asserted strict
separation from every one-sided M-018 choice.

\paragraph{Other established failures remain.}
The full support (O-022.6) is contained in the present
support, and every present row cap is at least its
O-022 counterpart. Thus O-022's exact generic
degree-$26$ complexity argument still excludes every
F-018 exact support cover, while its separate dimension
arguments exclude every entire one-cut family and the
whole point-jet target.
The coordinate-box size failure also persists.
The shape meets the BIP staircase and A-025's joint gate.
The shape-specific Chow threshold is at least
$d-1=3n-1>n$, and the uniform Chow and Grenet thresholds
fail as before. Its length is $m^2$, its tail is
$26r+1999928\leq md$, and its third row is larger than two.
It remains outside every one-sided jet exclusion by the
complete argument (O-025.5)--(O-025.9).
The new two-sided comparison therefore reduces the
actual remaining gate on a shape with at least four
permanent copies. It does not produce a multiplicity
obstruction; (O-025.2) rules one out here.

\paragraph{Geometric dictionary and scope.}
The two axis ideals are finite free integral torus-weight
modules under SGA~3 I, Proposition~4.7.3, pp.~35--36.
The constant bridge and the two attached families in
M-019 give their actual universally injective coordinate
pullback. F-016 supplies the full Schur-copy support
through SGA~6 VI, Theorem~3.3, pp.~392--393, and
Macdonald I(4.3), p.~63, Appendix~A(5.3)--(5.6),
pp.~155--156. The counting argument (O-025.8)
is a positive incidence count in the suffix monomial
basis, not an alternating virtual identity.
The exact representation barrier inherited from O-022
and its border dimension barriers retain their distinct
scopes. No general saturation or unresolved evaluation-rank
comparison is inferred.

\section{Exact positive cover counts for a two-prefix profile}

\subsection{O-026: a fixed-base optimization within the existing criteria}
Let $L\geq M\geq b\geq1$, let $r\geq4$, and put
\[
 q=r-3,\qquad u=L-b,\qquad v=M-b,\qquad
 c=(v+1)(u+1)-\frac{v(v+1)}2.
\]
Consider the support
\[
 \Gamma=\left\{\gamma:
       |\gamma|\leq L,\quad
       \sum_{i=3}^r\gamma_i\leq M,\quad
       \sum_{i=4}^r\gamma_i\leq b\right\}.
                                                               \tag{O-026.1}
\]
Assume
\[
                             q>\max\{c,u+v+2\}.
                                                               \tag{O-026.2}
\]
Among nonempty finite monomial order ideals, the following
two exact minima hold:
\[
 \begin{split}
 \min_{\Gamma\subseteq A+D_b}|A|&=c,\\
 \min_{\Gamma\subseteq A+B+D_b}(|A|+|B|)&=u+v+2.
 \end{split}                                                  \tag{O-026.3}
\]
The unique one-sided optimum is
\[
                  A_0=\{x_2^a x_3^t:a+t\leq u,\ t\leq v\}.
                                                               \tag{O-026.4}
\]
Two axis ideals of degrees $u$ and $v$ attain the
two-sided minimum. These are exact positive counts for
the fixed-base support optimization in M-018 and M-019.
They are consequences and optimality statements within
those already proved criteria, not new existential
multiplicity criteria.

\begin{proof}[The one-sided optimum]
O-024 gives $\Gamma\subseteq A_0+D_b$, and its
subpartition count gives $|A_0|=c$.
Take any proposed order ideal $A$ with $|A|\leq c$.
Since there are $q>c$ suffix variables, some unit
$x_j$, $j\geq4$, is absent from $A$.
Downward closure implies that every member of $A$
has zero $j$-th exponent.

For every prefix monomial $\nu\in A_0$,
$x_j^b\nu$ belongs to $\Gamma$:
its total degree is at most $b+u=L$, its suffix from
coordinate three is at most $b+v=M$, and its final
suffix degree is $b$. In a decomposition of its exponent
as $\alpha+\eta$, with $\alpha\in A$ and $|\eta|\leq b$,
the $j$-th exponent forces $\eta=be_j$.
Thus $\alpha=\nu$, proving $A_0\subseteq A$.
Therefore $|A|\geq c$; equality forces $A=A_0$.
Any ideal with more than $c$ members cannot improve
the known optimum. This proves both the minimum and
its uniqueness.
\end{proof}

\begin{proof}[The two-sided optimum]
The two axis ideals
$A=\{1,x_2,\ldots,x_2^u\}$ and
$B=\{1,x_3,\ldots,x_3^v\}$ contain a sum covering
$A_0$, so they attain the upper bound $u+v+2$.
Now take any proposed pair with
$|A|+|B|\leq u+v+2$.
By (O-026.2) there is a suffix unit $x_j$ in neither
ideal, and no member of either ideal uses that variable.

Apply the cover to $x_j^b x_2^u$.
Its $D_b$ part must be precisely $x_j^b$, so its
remaining decomposition uses pure $x_2$ powers in
$A$ and $B$. If $a_2,b_2$ are their largest $x_2$
axis degrees, this forces $a_2+b_2\geq u$.
Likewise $x_j^b x_3^v$ gives $a_3+b_3\geq v$.
Every ideal contains its two full axis chains, sharing
only the constant monomial. Consequently
\[
 |A|+|B|\geq(a_2+a_3+1)+(b_2+b_3+1)
                              \geq u+v+2.
\]
Larger pairs cannot improve the already attained upper
bound, proving the second minimum.
\end{proof}

\paragraph{The exact counts in the positive example.}
For O-025, take
$L=1000006$, $M=1000000$, and $b=26$.
Then
\[
 \begin{split}
 u&=999980,\qquad v=999974,\\
 c&=499981500150,\qquad u+v+2=1999956.
 \end{split}                                                  \tag{O-026.5}
\]
The suffix count $r-3=2^{400}-3$ satisfies (O-026.2).
At the common fixed core size
$T=m\binom{r+11}{12}$, the optimal one-sided
matrix size is therefore $499981500151T$, whereas
the optimal two-sided size at this same base order is
$1999956T$.
O-025 has already proved that the latter fits inside
$n$ and the former does not, with actual permanent
multiplicity at least four.

For this particular example the one-sided count can
also be optimized over every even base order of M-018.
The optimum still occurs at $b=26$ with the unique
ideal (O-026.4). Indeed its size is less than
$2^{39}n$, since $T<n$ and
$499981500151<2^{39}$.
For $b\geq28$, (O-025.6) gives a lower ratio to $n$
of at least $2r/13!>2^{39}$.
For $b\leq24$, (O-025.9) gives a lower ratio greater
than $r/2^{139}=2^{261}>2^{39}$.
These cases exhaust the other orders. Enlarging the
ambient length cannot lower the core size or help
decompose monomials supported on the original variables.
This establishes the stated optimum for the M-018
architecture, without an optimality assertion for
arbitrary determinant maps or for all M-019 base orders.

\paragraph{Geometric dictionary.}
The counts in (O-026.3) are ranks of explicitly indexed
free coefficient modules over $\mathbb Z$, with their
torus-comodule structures from SGA~3 I,
Proposition~4.7.3, pp.~35--36.
Nonnegative character exponents make the missing-variable
argument exact: a monomial containing that character
cannot be cancelled inside a support decomposition.
The already proved morphisms M-018 and M-019 convert
the counts into universally injective coordinate
pullbacks. Their characteristic-zero Schur-copy step
uses SGA~6 VI, Theorem~3.3, pp.~392--393, and the
Cauchy/Schur formalism of Macdonald I(4.3), p.~63,
Appendix~A(5.3)--(5.6), pp.~155--156.
The optimization concerns these explicit coefficient
schemes and maps; neither a virtual lambda-ring identity
nor purity alone provides the comparison.

\section{Two transverse labels per matrix edge}

\subsection*{Finding O-027: the exact determinant row threshold when \(\lambda_3\leq2\)}
Let \(n\geq2\). For every outer degree \(d\) and every
\(\lambda\vdash nd\) with
\[
                    \ell(\lambda)\leq2n,\qquad \lambda_3\leq2,
                                                               \tag{O-027.1}
\]
the determinant orbit closure retains all ambient copies:
\[
                         K_n(\lambda)=a_\lambda(d[n]).
                                                               \tag{O-027.2}
\]
For the unpadded permanent of size \(n\), the same conditions give
\[
                         q^n_\lambda(d[n])=a_\lambda(d[n]).
                                                               \tag{O-027.3}
\]
These statements include all highest-weight multiplicity copies.
They follow from a single explicit construction which makes the complete
projection to transverse order two along a binary plane dominant for
both determinant and permanent in \(2n\) variables.
The construction improves O-020's sufficient determinant size
\(n\geq r-1\) to \(n\geq\lceil r/2\rceil\).
O-020 and its integral incidence calculation remain valid independently.

For \(n\geq3\), the row threshold \(2n\) in the universal assertion
(O-027.2) is sharp. The known LMR module has
\[
 \Lambda_n=\bigl(2n^3-4n^2+1,\ 2n^2-4n+1,\ 2^{\,2n-1}\bigr),
 \qquad d=2n(n-1),
                                                               \tag{O-027.4}
\]
so that \(\ell(\Lambda_n)=2n+1\) and \((\Lambda_n)_3=2\).
Theorem 2.3.1 and Section 3.2 of Landsberg--Manivel--Ressayre
give a nonzero copy of this module in the determinant ideal.
Consequently \(K_n(\Lambda_n)<a_{\Lambda_n}(d[n])\).
This sharpness statement concerns ambient determinant multiplicity;
it does not assert a permanent comparison for all \(n\).
The finite permanent endpoint certificates in A-031 supply that
additional comparison for their stated sizes.

\paragraph{The binary pencil and its quadratic coefficient space.}
Put \(N=n-1\), and use the variables
\[
                 x,y,\ z_{i1},z_{i2}\quad(1\leq i\leq N).
                                                               \tag{O-027.5}
\]
Let \(V=\bigoplus_{i=1}^{N}V_i\), where
\(V_i=\langle z_{i1},z_{i2}\rangle\), and choose distinct scalars
\(a_1,\ldots,a_n\). Write
\[
                \ell_i=y-a_i x,\qquad \Delta=\prod_{i=1}^n\ell_i.
\]
The desired coefficient space is
\[
 \mathcal J=
 \operatorname{Sym}^n\langle x,y\rangle
 \ \oplus\
 \operatorname{Sym}^{n-1}\langle x,y\rangle\otimes V
 \ \oplus\
 \operatorname{Sym}^{n-2}\langle x,y\rangle\otimes\operatorname{Sym}^2V.
                                                               \tag{O-027.6}
\]
Projection to \(\mathcal J\) discards only terms of total \(z\)-degree
at least three. Higher discarded coefficients are unrestricted.

For \(i<j<n\), set
\[
 A^0_{ij}=z_{i1}+2z_{j1},\qquad
 A^0_{ji}=z_{i2}+2z_{j2}.
\]
For \(i<n\), set
\[
                  A^0_{in}=z_{i1},\qquad A^0_{ni}=z_{i2},
\]
and set \(A^0_{ii}=\ell_i\).
Thus \(A^0\) is a completely specified matrix of linear forms.
The quadratic transverse part of its permanent is
\[
       \sum_{i<j}\frac{\Delta}{\ell_i\ell_j}A^0_{ij}A^0_{ji};
                                                               \tag{O-027.7}
\]
the corresponding determinant part is its negative.
Indeed, a permutation with exactly two off-diagonal entries is a
transposition. All cycles of length at least three disappear in
the projection (O-027.6).

\paragraph{The two-label incidence lemma.}
Let
\[
 H_n=\{(c_1,\ldots,c_n):\textstyle\sum_i c_i=0\}.
\]
For \(i<j<n\), put
\[
 U_{ij}=\langle z_{i1}+2z_{j1},\,z_{i2}+2z_{j2}\rangle,
 \qquad U_{in}=V_i.
\]
The linear map
\[
 \Psi:\bigoplus_{i<j}(U_{ij}\otimes V)
       \longrightarrow H_n\otimes\operatorname{Sym}^2V,
 \qquad
 (u\otimes v)_{ij}\longmapsto(e_i-e_j)\otimes uv
                                                               \tag{O-027.8}
\]
is surjective over every field of characteristic different from
two and three.

\begin{proof}
A functional on the target can be represented by symmetric bilinear
forms \(Q_1,\ldots,Q_n\) on \(V\), modulo adding the same form to every
\(Q_i\). Set \(Q_n=0\). A functional annihilates the image of \(\Psi\)
exactly when
\[
             (Q_i-Q_j)(u,v)=0
       \quad\text{for all }u\in U_{ij},\ v\in V.
                                                               \tag{O-027.9}
\]
The reference edges \(\{i,n\}\) first imply that \(Q_i\) kills \(V_i\)
in either argument.

For \(i<j<n\), apply (O-027.9) to
\(u=z_i+2z_j\), where the coordinates of \(z_i\in V_i\) and
\(z_j\in V_j\) agree under the displayed bases.
Taking the other argument in \(V_i\) shows that the
\(V_i\)-diagonal block of \(Q_j\) is zero.
Taking it in \(V_j\), and dividing by two, shows that the
\(V_j\)-diagonal block of \(Q_i\) is zero.
Together with the reference edges, this kills every diagonal
two-by-two block of every \(Q_i\).

Every block still possibly nonzero is indexed by three distinct
nonreference vertices. Fix \(i<j<k<n\), and denote by
\[
            A=(Q_i)_{jk},\qquad B=(Q_j)_{ik},\qquad C=(Q_k)_{ij}
\]
their two-by-two matrices. The edges \(\{i,j\}\), \(\{i,k\}\),
and \(\{j,k\}\), respectively, give
\[
                         B=2A,\qquad C=2A^{\mathsf T},
                         \qquad C=2B.
                                                               \tag{O-027.10}
\]
Hence \(A^{\mathsf T}=2A\), and transposition gives \(A=4A\).
Characteristic different from three forces \(A=0\), and then
\(B=C=0\). This treats every remaining block.
There is no triple to treat when \(N<3\); the preceding diagonal
argument already gives the same conclusion.
The dual kernel is zero, proving surjectivity.
\end{proof}

\paragraph{From binary residues to the full coefficient differential.}
The map
\[
 \operatorname{Sym}^{n-2}\langle x,y\rangle\longrightarrow H_n,
 \qquad
 f\longmapsto
 \left(\frac{f(1,a_i)}{\prod_{j\ne i}(a_i-a_j)}\right)_{i=1}^n
                                                               \tag{O-027.11}
\]
is an isomorphism. Its image is in \(H_n\) by the coefficient
of \(y^{-1}\) in the partial-fraction expansion of
\(f(1,y)/\Delta(1,y)\). The map is injective because a polynomial
of degree at most \(n-2\) vanishing at all \(n\) distinct \(a_i\)
is zero; source and target both have dimension \(n-1\).
In particular,
\[
 \frac{\Delta}{\ell_i\ell_j}\longmapsto
                    \frac{e_i-e_j}{a_i-a_j}.
                                                               \tag{O-027.12}
\]

Vary a single off-diagonal entry \(A_{ij}\) by an arbitrary
linear form \(H\in V\). The transverse-order-two part of the
derivative is exactly
\[
              \frac{\Delta}{\ell_i\ell_j}\,A^0_{ji}H
                                                               \tag{O-027.13}
\]
for the permanent and its negative for the determinant.
It has no transverse degree zero or one. Varying both
orientations of every edge, using (O-027.12), gives precisely
the image in (O-027.8), up to nonzero scalar factors.
The incidence lemma therefore supplies the whole third
summand of (O-027.6).

Next vary \(A_{ii}\) by \(H\in V\). Up to transverse order two,
the derivative is
\[
                         \frac{\Delta}{\ell_i}H.
                                                               \tag{O-027.14}
\]
There is no transverse-degree-one term in the remaining
principal minor, since a nonidentity permutation cannot have
exactly one off-diagonal entry. The \(n\) binary forms
\(\Delta/\ell_i\) form a basis in degree \(n-1\), as evaluation
at \(a_i\) shows. These variations supply the entire middle
summand, without affecting the other two.

Finally vary each diagonal entry by \(x\) or \(y\).
The transverse-degree-zero images span
\(\operatorname{Sym}^n\langle x,y\rangle\), because they are
\(x\Delta/\ell_i\) and \(y\Delta/\ell_i\).
Their possible transverse-degree-two components can be
subtracted using (O-027.13); their transverse-degree-one
components are zero. The complete differential onto
\(\mathcal J\) is consequently surjective, for both determinant
and permanent.

\paragraph{Dominance and all-copy consequence.}
The coefficients of a matrix of linear forms are affine
parameters, and projection of its determinant or permanent
to (O-027.6) is a polynomial map.
A surjective differential at \(A^0\) gives an injective
coordinate pullback: restrict source directions to a linear
right inverse of the differential, translate source and
target to their base values, and use the lowest nonzero
homogeneous part of a polynomial under a map
\(t\mapsto t+O(t^2)\).
Thus the coefficient projection is dominant.
Setting unnecessary variables to zero gives the same
conclusion for every \(r\leq2n\): coefficients using only
the retained variables are independent of the discarded
matrix-entry coefficients, so variable restriction commutes
with projection to the retained target coordinates.

For any \(\lambda\) satisfying (O-027.1), every highest-weight
vector in the degree-\(d\) ambient coordinate ring uses only
the first \(\ell(\lambda)\) variable directions.
The flag-minor support theorem F-016 says that every coefficient
generator appearing in such a vector has total exponent at
most \(\lambda_3\leq2\) in coordinates three and later.
It is therefore a polynomial in the coordinates of
\(\mathcal J\).
If the vector lies in the determinant ideal, it vanishes
on every matrix-pencil substitution into the determinant,
since such substitutions belong to the orbit closure.
The injective pullback just proved forces it to be zero.
The same reasoning applies to the unpadded permanent.
Every vector in the full highest-weight ideal multiplicity
space is covered by this argument. Consequently that space
is zero, proving (O-027.2)--(O-027.3).

\paragraph{Internal padding.}
The argument above varies all matrix-entry coefficients.
In particular, it does not claim that \(A_{11}=x\) can be
fixed in the same construction.
For the internally padded permanent, the verified O-021
padding argument instead gives the following consequence:
if \(n\geq m\), \(\ell(\lambda)\leq r\), \(\lambda_3\leq2\),
and
\[
                  m\geq
                \max\{\lceil r/2\rceil+1,\ \lambda_2+1\},
                                                               \tag{O-027.15}
\]
then
\[
                           q^m_\lambda(d[n])
                         =K_n(\lambda)=a_\lambda(d[n]).
                                                               \tag{O-027.16}
\]
Indeed, put \(s=m-1\). The new construction supplies the
full permanent coefficient projection in size \(s\), because
\(r\leq2s\). The other inequality gives \(s\geq\lambda_2\),
so F-014 stabilizes the complete highest-weight space at
base degree \(s\). The diagonal, invertible differential
transpose in O-012 preserves its transverse-order-two
coefficient support. Finally the block matrix
\(\operatorname{diag}(x,A_s)\) has fixed internal entry \(x\)
and permanent \(x\operatorname{per}(A_s)\), exactly as in
O-021. This proves all-copy equality after padding.
For one-row partitions, the nonzero pure-power coordinate
gives the same equality directly.

\paragraph{Sharp boundary for the determinant coefficient projection.}
For \(r\) variables, the target in transverse order two has dimension
\[
 D_n(r)=(n+1)+n(r-2)+(n-1)\binom{r-1}{2}.
                                                               \tag{O-027.17}
\]
The closure of the \(r\)-variable determinant-pencil image
has dimension at most \((r-2)n^2+2\), by O-010's
left-right gauge argument. At \(r=2n+1\), subtraction gives
\[
             D_n(2n+1)-\bigl((2n-1)n^2+2\bigr)=n-1>0.
                                                               \tag{O-027.18}
\]
Hence the full coefficient projection cannot be dominant
there. For larger \(r\), the difference increases.
Together with the construction, this makes \(r=2n\) the
exact maximal number of variables for this determinant
coefficient projection, within the ambient variable range.
This is a geometric statement about the whole projection;
the separate representation-theoretic sharpness in
(O-027.4) uses the known LMR ideal copy.

\paragraph{The selected integral incidence cokernel.}
The incidence map (O-027.8) is defined over \(\mathbb Z\).
With \(N=n-1\), its finite cokernel is
\[
 \operatorname{coker}\Psi_{\mathbb Z}
 \simeq
 (\mathbb Z/2)^{\,3\binom N2+3\binom N3}
       \ \oplus\
 (\mathbb Z/6)^{\,\binom N3}.
                                                               \tag{O-027.19}
\]
Here binomial coefficients are zero outside their usual range.
To see this, use the dual matrix solely to compute Smith
invariants, which are unchanged by transposition.
The reference-edge equations eliminate, by unit pivots,
every block of \(Q_i\) involving \(V_i\).
For each pair \(i<j\), the \(V_i\)-diagonal block of \(Q_j\)
is eliminated by unit pivots, while the three independent
entries of the \(V_j\)-diagonal block of \(Q_i\) have the
relations \(2q=0\).
For each triple \(i<j<k\), use the unit relations
\(B=2A\) and \(C=2B\) in (O-027.10).
The remaining four-entry block is
\[
                          2(A^{\mathsf T}-2A)=0.
\]
Its two diagonal entries contribute \(2,2\); on its two
off-diagonal entries the matrix is
\[
                    \begin{pmatrix}-4&2\\2&-4\end{pmatrix},
\]
whose Smith invariants are \(2,6\).
All pairs and triples involve disjoint remaining variables,
giving (O-027.19).
This is the cokernel of the selected incidence map;
it is not a claim that every full determinant or permanent
jet map must fail in characteristic two or three.

\paragraph{Grothendieck dictionary and base-change scope.}
The matrix pencil, determinant, permanent, coefficient
projection and incidence homomorphism are polynomial or
linear constructions over \(\mathbb Z\).
Equation (O-027.19) makes the integral limitation explicit.
After inverting six and the differences \(a_i-a_j\),
the incidence and binary interpolation maps are split
surjective, and the full differential has a linear
right inverse. The resulting lowest-term coordinate
injection survives arbitrary commutative base change
over this localized ring.
The passage to irreducible highest-weight multiplicities
is made over \(\mathbb C\), using F-016 and
characteristic-zero semisimplicity.
In Grothendieck's language, the graded coordinate algebras
and their ideals are comodules as in SGA~3, I,
Propositions~4.7.2--4.7.3, pp.~34--36.
Their Schur operations belong to the special
\(\lambda\)-structures of SGA~6, VI, Theorem~3.3,
pp.~392--393. The conclusion is certified by the actual
coordinate injection and its explicit integral localization,
not by positivity of a virtual \(K\)-class.


\section{What is, and is not, supplied by positivity}

Brion, \emph{Positivity in the Grothendieck group of complex flag
varieties}, Theorem~1, asserts that for a subvariety $Y$ with rational
singularities the coefficient $b_w$ in
$[\mathcal O_Y]=\sum_w b_w[\mathcal O_{X_w}]$ satisfies
$(-1)^{\operatorname{codim}X_w-\operatorname{codim}Y}b_w\geq0$.
Its application to products of Schubert structure sheaves is explained
immediately after that theorem (and uses Lemma~2). These are the
alternating positive structure constants represented in type A by
Grothendieck polynomials.

There are three separate operations here. Induction multiplication on
$\bigoplus_N R(S_N)$ has Littlewood--Richardson constants. Tensor
multiplication inside a fixed $R(S_N)$ has Kronecker constants.
Multiplication of coherent-sheaf classes on a flag variety has
Grothendieck-polynomial structure constants. Equality of these
products, or a map preserving their distinguished bases and their
stretching, has not been established by a $\lambda$-ring identity.
In particular a proposed general saturation implication from special
$\lambda$ identities is refuted by O-002: those identities hold in
the representation rings that exhibit the holes.  A proposed
implication from normality or positive multiplication alone is refuted
by O-003. Neither observation rules out an additional geometric
construction with substantially stronger hypotheses in a restricted
family. The missing extra statement must be supplied and proved.

The positive basis theorem actually established by the supplied GCT IV
is Theorem~15.21(vii), with the explicit rule in Theorem~18.2, for
two factors with at most two rows. The full compatible basis programme
is still stated there as Conjecture~19.1, and Section~19.2 explains
the dimension failure of the direct construction beyond that case.
It cannot be promoted to a proved general geometric positivity theorem
merely by identifying exterior powers in $K_0$.

The exact known occurrence conclusions in the GCT range are:
Ikenmeyer--Panova Theorem~1.6, $n>3m^4$, gives
$g_{n,d}(\lambda)=0\Rightarrow q^m_\lambda(d[n])=0$;
B\"urgisser--Ikenmeyer--Panova Theorem~1.4, $n\geq m^{25}$, gives
$q^m_\lambda(d[n])>0\Rightarrow K_n(\lambda)>0$.
The latter does not assert $q^m_\lambda\leq K_n$.
O-005--O-007 prove the latter inequality in the requested hook and
two-row families (and extend it to three-row outputs), but supply
no general multiplicity obstruction and no separation of complexity
classes.

\section{Geometric placement of the findings}

Finite symmetric groups are constant finite \'etale group schemes
over $\mathbb Z$ (SGA~3, Expos\'e~I, Section~4.1), and their
representations are comodules by SGA~3, I, Proposition~4.7.2.
Passing to $K_0$ over a characteristic-zero field turns the symmetric,
exterior, and Schur functors used above into the usual operations on
locally free representations. The special $\lambda$ identities live
in SGA~6, Expos\'e~V, Definition~2.4 and identities~(2.4.2), and
the realization for locally free sheaves is Expos\'e~VI,
Theorems~3.2--3.3 (pp.~391--393). The exterior filtration behind the
direct-sum and exact-sequence formulas is V, Lemma~2.2.1
(pp.~315--316). The $\gamma$-filtration is V, 3.10,
equation~(3.10.1), p.~331. None of these operations makes an arbitrary
virtual Burnside class effective; the counterexample in cote~133,
scan~019 must be corrected before making the translation.

O-001 and O-002 state multiplicities in
$K_0(\operatorname{Rep}_{\mathbb Q}S_N)$ and, through Schur--Weyl
duality, in the restriction of $S_\lambda(E\otimes F)$ along the
group-scheme homomorphism
$\mathrm{GL}(E)\times\mathrm{GL}(F)\to\mathrm{GL}(E\otimes F)$.
This is the Cauchy/Schur dictionary relevant to cote~133, scan~028.
O-003 has the more explicit integral graded-affine-scheme realization
already given above. O-004 is a statement about invariant summands
of Schur functors of the adjoint bundle on the characteristic-zero
fiber. O-005 is the exterior-power calculation in that exact tensor
category.

O-006.a is a statement about equivariant coherent quotients of the
coordinate algebra of an affine space of sections. O-006.b and
O-007 apply it to the scheme-theoretic orbit closure of an integral
form: the determinant and its substitution morphism are defined
over $\mathbb Z$, and their characteristic-zero fibers give the
orbit closures above. Semisimplicity and the displayed irreducible
decompositions are used only after base change to $\mathbb Q$ or
$\mathbb C$, not integrally and not in all residue characteristics.
O-007 also has a concrete Picard/Quot interpretation: Beauville's
matrix gives a resolution
$0\to\mathcal O_{\mathbb P^2}(-2)^n\to
\mathcal O_{\mathbb P^2}(-1)^n\to L\to0$ on a plane curve,
with $L$ a point of its Picard scheme outside the theta locus and
the final arrow a coherent quotient. The Hilbert/Picard/Quot
constructions provide the parameter spaces, while the determinantal
and multiplicity conclusions are the arguments of this run and the
explicit cited theorem, not statements attributed retrospectively
to Grothendieck. For the exact Quot representability statement, see
Grothendieck, FGA, Bourbaki expos\'e~221, \emph{Techniques de
construction en g\'eom\'etrie analytique, IV: Les sch\'emas de Hilbert},
Theorem~3.1, p.~259, and Theorem~3.2, p.~260, for a fixed Hilbert
polynomial; the identification of the Hilbert functor with the Quot
functor of $\mathcal O_X$ follows Proposition~3.8 on p.~265.
Remark~3.9, p.~266, explains the reduction to projective schemes over
$\mathbb Z$. O-008 is a statement about the homogeneous ideal of the
scheme-theoretic image of the same integral substitution morphism
after base change to $\mathbb C$, with its dimension bound proved
directly by the displayed parametrization.
O-009 is the invariant-multiplicity statement on the characteristic-zero
fiber obtained from the stated positivity and plethysm theorems.
O-010 is a dimension constraint on containment of a linear affine
subscheme in a scheme-theoretic image; its proof is an explicit
coordinate-ring transcendence-degree calculation, so requires no
unproved normality or purity property of that image.

% END agents/obstructions/obstructions_findings.tex

\part{Comparisons of actual orbit-closure multiplicities}

% BEGIN agents/archive/HESSIAN_OBSTRUCTION.tex
% Run GCT_KRON_20261120; GPT-6 Astra Pro; effort not exposed.
\section{An actual multiplicity obstruction for the unpadded permanent}

\subsection*{A-030. The Hessian bracket and its unique ambient copy}
Work over \(\mathbb C\). Let \(n\geq3\), \(R=n^2\),
\(W=\operatorname{Mat}_n(\mathbb C)\), and put
\[
 q=R(n-2),\qquad d=R+1,\qquad
 \lambda=(n+q+1,3,\underbrace{2,\ldots,2}_{R-2}).       \tag{A-030.1}
\]
Thus \(|\lambda|=nd\) and \(\ell(\lambda)=R\).
In the notation of the main report, with the permanent unpadded
and of the same size as the determinant, one has
\[
 a_\lambda(d[n])=1,\qquad
 K_n(\lambda)=0,\qquad
 q_\lambda^{\,n}(d)=1.                               \tag{A-030.2}
\]
In particular this is an actual multiplicity obstruction
\(q_\lambda^{\,n}(d)>K_n(\lambda)\), for every \(n\geq3\).
For example, \(n=3\) gives \(d=10\) and
\(\lambda=(13,3,2,2,2,2,2,2,2)\).
It is outside the large-padding range of the B\"urgisser--Ikenmeyer--Panova
theorem. No assertion of historical priority is made.
In fact this is an ordinary rectangular Kronecker occurrence
obstruction as well:
\[
               g(\lambda,(d^n),(d^n))=0.                       \tag{A-030.2a}
\]
Accordingly the symmetric Kronecker multiplicity also vanishes;
there is no orbit-to-closure multiplicity loss in this family.

\paragraph{Known inputs and the additional calculation.}
The permanent Hessian used below is exactly the classical
Mignon--Ressayre point, not a newly discovered full-rank point:
\emph{A quadratic bound for the determinant and permanent problem},
author preprint dated July 20, 2004, Section 3.2.2,
Proposition 3.1, printed p.~8. Their Lemma 3.4 on the same page
gives its first and second derivatives, and pp.~9--10 prove
invertibility. Their Theorem 2, p.~10, is the resulting quadratic
lower bound on determinantal complexity. We use Proposition 3.1
as a known result. The determinant's Hessian power identity is
classical as well: Thomas Muir, \emph{The Hessian of a General
Determinant}, \emph{Transactions of the Royal Society of Edinburgh}
40 (1905), 203--207, DOI \texttt{10.1017/S008045680003427X};
the publisher's extract explicitly states the power
\(n(n-2)\). The short computation below fixes the nonzero
constant for ordinary second derivatives and a specified basis.

The new work in this report is the explicit assembly of the Hessian
bracket, its Schur type, and F-024's ambient multiplicity-one
calculation into (A-030.2), together with an exact determinant
normalization check. This description does not claim that this
assembly is absent from the literature. Hessian equations for
determinant orbit closures and a stronger quadratic border bound
are already established by Landsberg--Manivel--Ressayre,
\emph{Hypersurfaces with degenerate duals and the Geometric
Complexity Theory Program},
\emph{Commentarii Mathematici Helvetici} 88 (2013), 469--484,
Theorem 1.1.1, printed p.~470; Theorem 2.3.1, p.~474;
and Theorem 3.1.1, p.~476. Their Section 3.2, p.~476,
explicitly distinguishes one ideal copy from an entire isotypic
component. Here F-024 supplies precisely the ambient uniqueness
that turns the explicit ideal copy into the full isotypic component.

\paragraph{The Hessian power identity, with normalization.}
Use the ordered matrix-entry basis on \(W\). For a homogeneous
degree-\(n\) polynomial \(F\) set
\[
 \mathcal H(F)(x)=\det\!\left(
           \frac{\partial^2 F}{\partial x_i\partial x_j}(x)
                         \right)_{1\leq i,j\leq R}.
\]
This is a homogeneous form of degree \(q\) in \(x\), and a
polynomial of degree \(R\) in the coefficients of \(F\).
For \(D=\det_n\), the classical identity, in this normalization, is
\[
 \mathcal H(D)=c_nD^{n(n-2)},\qquad
 c_n=(n-1)(-1)^{n(n+1)/2-1}\ne0.                     \tag{A-030.3}
\]
Indeed, the second differential at the identity is
\(\operatorname{tr}(X)\operatorname{tr}(Y)-\operatorname{tr}(XY)\).
Relative to the entrywise pairing its operator is
\(X\mapsto\operatorname{tr}(X)I-X^{\mathsf T}\).
Its eigenvalues are \(n-1\) on the scalar line, \(-1\) on the
symmetric trace-zero matrices, and \(1\) on the skew-symmetric
matrices. This gives \(c_n\). For an invertible matrix \(A\),
the change \(X\mapsto A^{-1}X\) has determinant
\(\det(A)^{-n}\), while the second differential gains the scalar
\(\det(A)\). Taking its \(R\)-by-\(R\) determinant gives
\(c_n\det(A)^{R-2n}\). Polynomial continuation proves (A-030.3).
This verifies the constant rather than imposing a convention from
a differently normalized polarization.

\paragraph{A polynomial covariant in one irreducible module.}
The action is \((g\cdot F)(x)=F(g^{-1}x)\), so the chain rule gives
\[
 \mathcal H(g\cdot F)(x)
       =\det(g)^{-2}\mathcal H(F)(g^{-1}x).             \tag{A-030.4}
\]
Consequently \(\mathcal H\) is intrinsically valued in
\((\det W^*)^{\otimes2}\otimes\operatorname{Sym}^{q}W^*\).
For \(a=n+q-1\), form the one-form covariant
\[
 \Theta(F)=q\mathcal H(F)\,\mathrm dF
                         -nF\,\mathrm d\mathcal H(F).          \tag{A-030.5}
\]
Its degree in the coefficients of \(F\) is \(R+1=d\).
Its polynomial coefficients have degree \(a\) in \(x\).
Euler's identity gives
\[
                       \Theta(F)_x(x)=0.                      \tag{A-030.6}
\]
The kernel of multiplication
\(\operatorname{Sym}^aW^*\otimes W^*
\longrightarrow\operatorname{Sym}^{a+1}W^*\)
is \(\mathbb S_{(a,1)}W^*\), by the characteristic-zero
Pieri decomposition. Thus (A-030.5) is a polynomial
\(\operatorname{GL}(W)\)-equivariant map to the irreducible module
\[
 U=(\det W^*)^{\otimes2}\otimes\mathbb S_{(a,1)}W^*
                 =\mathbb S_\lambda W^*.                       \tag{A-030.7}
\]
The two determinant factors add two boxes to every one of the
\(R\) rows. They do not produce a two-row representation:
the resulting highest-weight partition is exactly (A-030.1).

Equation (A-030.3), together with \(q=n\,n(n-2)\), implies
\(\Theta(D)=0\). Equivariance then gives zero on the determinant
orbit. Since every coefficient of \(\Theta\) is a polynomial
in the coefficients of \(F\), it gives zero on the full affine
orbit closure \(\Omega_n\) as well. This is a closure assertion,
not only an orbit-stabilizer calculation.

\paragraph{Nonvanishing on the permanent.}
Let \(P=\operatorname{perm}_n\) and let \(A\) have all entries
equal to one except \(A_{11}=1-n\). Mignon--Ressayre,
Proposition 3.1 and Lemma 3.4 at the locators above, give
\[
 P(A)=0,\qquad
 \frac{\partial P}{\partial x_{11}}(A)=(n-1)!,\qquad
 \operatorname{rank}\operatorname{Hess}(P)(A)=n^2.               \tag{A-030.8}
\]
Thus \(\mathcal H(P)(A)\ne0\), and evaluating (A-030.5) yields
\[
 \Theta(P)_A=q\mathcal H(P)(A)\,\mathrm dP_A\ne0.                \tag{A-030.9}
\]
In particular the coefficient map
\[
 \Theta^*:U^*\longrightarrow
       \mathbb C[\operatorname{Sym}^nW^*]_d
       =\operatorname{Sym}^{d}(\operatorname{Sym}^nW)            \tag{A-030.10}
\]
is nonzero. Its source \(U^*=\mathbb S_\lambda W\) is
irreducible, so it is injective. Its image is contained in the
ideal of \(\Omega_n\), whereas its restriction to the permanent
orbit closure is nonzero, hence injective on this simple module.

Now F-024 gives, for every outer degree \(d\geq3\) and inner
degree \(n\geq3\),
\[
 a_{(d(n-2)+3,\,3,\,2^{d-3})}(d[n])=1
 \quad\hbox{when the ambient dimension is at least }d-1.         \tag{A-030.11}
\]
Here \(d=R+1\) and \(\dim W=R=d-1\), and the first part
\(d(n-2)+3\) is \(n+q+1\). Therefore the copy in
(A-030.10) is the entire ambient isotypic component. It dies
on the determinant closure and survives on the permanent closure.
Complete reducibility in each finite degree proves (A-030.2).
In particular it is the ambient uniqueness, in addition to the
separating equation, that proves the exact multiplicity values by
this argument. The following independent vanishing calculation
gives a second route to the determinant side.

\paragraph{The ordinary Kronecker coefficient already vanishes.}
Write \(W=A_0\otimes B_0\) with
\(\dim A_0=\dim B_0=n\), and put
\[
 E=n+q=n(n-1)^2,\qquad \nu=(E-1,1),\qquad
                         t=d-2n=(n-1)^2.
\]
As a \(\operatorname{GL}(W)\)-module,
\(\mathbb S_\lambda W=(\det W)^2\otimes\mathbb S_\nu W\).
On restriction to \(\operatorname{GL}(A_0)\times
\operatorname{GL}(B_0)\),
\(\det(A_0\otimes B_0)^2=(\det A_0)^{2n}(\det B_0)^{2n}\).
The Schur--Weyl definition of Kronecker multiplicities therefore gives
\[
 g(\lambda,(d^n),(d^n))
                         =g((E-1,1),(t^n),(t^n)).              \tag{A-030.14}
\]
The standard hook \((E-1,1)\) has multiplicity equal to the
number of removable corners of \((t^n)\) minus one in the
tensor square of its Specht module. This is zero for a rectangle;
equivalently it is the degree-one coefficient of F-003's
rectangular-hook product. This proves (A-030.2a).
Since \(K_n(\lambda)\) is bounded by the symmetric Kronecker
multiplicity, which in turn is bounded by the ordinary one,
it independently gives \(K_n(\lambda)=0\). Thus the example
does not rely on substituting an orbit coefficient for a larger
or smaller closure coefficient.

\paragraph{A saturation hole with actual permanent positivity.}
For every integer \(s\geq1\), the same determinant twist gives
\begin{align}
 g(s\lambda,((sd)^n),((sd)^n))
   &=g((s(E-1),s),((st)^n),((st)^n))\notag\\
   &=[z^s]\prod_{j=2}^{n}(1-z^j)^{-1}.                \tag{A-030.15}
\end{align}
The second equality is the known rectangular two-row formula
F-002: it is \(p_s(n,st)-p_{s-1}(n,st)\), where \(p_u(n,st)\)
counts partitions of \(u\) in an \(n\)-by-\(st\) rectangle.
Since \(st=s(n-1)^2\geq s\), the width bound has no effect at
these two coefficients, giving the displayed product.
Consequently the unstretched coefficient is zero, the
twice-stretched coefficient is one, and the coefficient is positive
for every \(s\geq2\): all such integers are sums of twos and threes,
and \(n\geq3\).

Moreover,
\[
                         q_{s\lambda}^{\,n}(sd)>0
                         \quad\text{for every }s\geq1.          \tag{A-030.16}
\]
Indeed choose a nonzero highest-weight vector in the copy
(A-030.10) after restriction to the permanent orbit closure.
That closure is irreducible, so its coordinate algebra is a
domain. Its \(s\)-th power is a nonzero highest-weight vector
of weight \(s\lambda\) and degree \(sd\). This proves
(A-030.16). Thus the ordinary Kronecker saturation hole here
has genuine permanent positivity, rather than just a positive
ambient plethysm coefficient. Equations (A-030.15)--(A-030.16)
do not determine the determinant-closure multiplicities
\(K_n(s\lambda)\) for \(s\geq2\); positive ordinary Kronecker
multiplicity is not a boundary-extension theorem.

\paragraph{Optional exact normalization of the permanent Hessian.}
The following explicit scalar sharpens the nonsingularity input;
it is not needed in order to invoke the known Proposition 3.1:
\[
 \mathcal H(P)(A)=
 ((n-3)!)^{n^2}(-2)^{(n-2)^2}
             (n-2)^{4(n-1)}(n-1)^{2n}.                        \tag{A-030.12}
\]
Here is a check that also specifies all normalization choices.
More generally take \(X(t)\) with entry \(11\) equal to \(t\)
and every other entry equal to one. Put
\(k=n-1\), \(b=(n-2)!\),
\(c=(n-3)!(t+n-3)\). Two Hessian entries in the same row or
column of the underlying matrix have value zero. Other entries
have value \(b\) if one of their rows or columns is numbered
one, and value \(c\) otherwise. This follows either from
the permanent minor formula, or directly from Lemma 3.2 of
Mignon--Ressayre, p.~7.

Decompose \(\mathbb C^k=\mathbb C\mathbf1\oplus E\), where
\(E\) is the sum-zero subspace, using the permutations of the
last \(k\) rows and columns. On the bottom-block summand
\(E\otimes E\), of dimension \((k-1)^2\), the Hessian has
eigenvalue \(c\). The first-row and matching column-constant
sum-zero modes, and the analogous first-column modes, have
two-by-two blocks
\[
 \begin{pmatrix}0&-b\sqrt{k}\\-b\sqrt{k}&-c(k-1)\end{pmatrix},
\]
of determinant \(-b^2k\), repeated \(2(k-1)\) times in total.
On the remaining four-dimensional constant subspace, use the
unit vectors of the \(11\) coordinate, first row, first column,
and bottom block. The matrix is
\[
 \begin{pmatrix}
 0&0&0&bk\\
 0&0&bk&b\sqrt{k}(k-1)\\
 0&bk&0&b\sqrt{k}(k-1)\\
 bk&b\sqrt{k}(k-1)&b\sqrt{k}(k-1)&c(k-1)^2
 \end{pmatrix},
\]
whose determinant is \(b^4k^4\). These are an orthonormal
decomposition of the entrywise bilinear form, so their product is
\[
          \mathcal H(P)(X(t))
               =b^{4k}k^{2k+2}c^{(k-1)^2}.                     \tag{A-030.13}
\]
At \(t=-k=1-n\), \(c=-2(n-3)!\), and (A-030.12) follows.
The factorial exponent is
\(4(n-1)+(n-2)^2=n^2\). No choice of square root affects the
block determinants. Exact permutation-sum differentiation and
integer determinants for \(n=3,4,5,6\) are recorded separately;
these finite checks supplement the displayed general derivation.

\paragraph{Padding limitation.}
For the same determinant size \(n\), the report's internally
padded permanent \(x_{11}^{n-m}\operatorname{perm}_m\),
with \(m<n\), uses at most \(m^2<n^2\) essential variables.
The present partition has
\(n^2\) nonzero parts. Therefore its entire ambient isotypic
component vanishes on this subspace variety: equivalently,
\(q_\lambda^{\,m}(d)=0\), by the standard Schur length bound
for the polynomial coordinate ring on forms in fewer variables.
One can also see the failure of this particular equation directly:
the full \(R\)-variable Hessian determinant, and hence
\(\Theta\), is identically zero on such a form. Thus (A-030.2)
does not transfer to \(n>m\), and supplies no contradiction to
the large-padding no-occurrence theorem or a superpolynomial
determinantal-complexity lower bound.
For the necessary second-row rule established in F-022 and
M-016, namely \(n\geq m^c\Rightarrow\lambda_2\geq c+2\)
for an actual obstruction, the present family attains
\(\lambda_2=3\) at \(c=1\) and \(n=m\). It proves sharpness
of that endpoint only, with no conclusion about the endpoints
at \(c\geq2\).

\paragraph{Grothendieck dictionary and exact scope.}
The construction is a polynomial map between finite-dimensional
representations, and its coefficients define a homogeneous
\(\operatorname{GL}(W)\)-stable ideal in the coordinate algebra
of the affine space \(\operatorname{Sym}^nW^*\). The determinant
closure is the scheme-theoretic image closure; vanishing of the
coefficient map on its dense orbit gives a factorization through
its ideal. Schur operations here are actual characteristic-zero
representations, the setting of the SGA 6, VI, Theorem 3.3
dictionary already used in the report. After taking the degree-\(d\)
piece, (A-030.10) is an explicit kernel inclusion for the
restriction map to the determinant closure, and (A-030.9)
proves its nonzero image under restriction to the permanent
closure. F-024 computes the ambient rank of this isotypic part.
Thus this example performs the ideal-kernel and quotient-rank
comparison that positivity of a virtual class alone cannot give.
No claim of invariants commuting with arbitrary integral base
change, or of a statement in positive characteristic, is made.

% END agents/archive/HESSIAN_OBSTRUCTION.tex

% BEGIN agents/jet_review/per3_full_transverse_jets.tex
\section{An actual multiplicity obstruction with both sides positive}

\subsection*{Finding F-025: full permanent multiplicity on a seven-variable cubic projection}
For the unpadded \(3\)-by-\(3\) permanent, every partition
\(\lambda\vdash3d\) with \(\ell(\lambda)\leq7\)
and \(\lambda_3\leq2\) satisfies
\begin{equation}\label{eq:F025full}
                         q^3_\lambda(d[3])=a_\lambda(d[3]).
\end{equation}
Missing parts are zero. The statement holds for every outer
degree \(d\), including every highest-weight multiplicity
copy. In particular, for
\[
                        \lambda=(19,7,2^5),\qquad d=12,
                                                               \tag{F-025.1}
\]
one has the actual orbit-closure comparison
\begin{equation}\label{eq:F025obstruction}
                  q^3_\lambda(12[3])=6>5=K_3(\lambda).
\end{equation}
The determinant-side value in \eqref{eq:F025obstruction}
is a published input specified below. The new permanent-side
comparison is proved by an explicit integer coefficient map
and a completely specified nonzero Jacobian minor.
This example has \(m=n=3\); it lies outside the IP and
BIP large-padding ranges. It does not settle the remaining
comparison in those ranges.

\paragraph{The coefficient projection.}
Let \(x=(x_1,\ldots,x_7)\), and set
\[
 \Gamma=\{\alpha\in\mathbb Z_{\geq0}^7:
           |\alpha|=3,\quad\alpha_3+\cdots+\alpha_7\leq2\}.
                                                               \tag{F-025.2}
\]
There are
\[
               |\Gamma|=4+5\cdot3+\binom62\cdot2=49
\]
monomials: the suffix degree is respectively zero, one or
two, leaving a binary monomial of degree three, two or one.
For a \(3\)-by-\(3\) matrix \(A(x)\) of linear forms,
consider the polynomial map
\[
 \Phi:\mathbb A^{63}\longrightarrow\mathbb A^{49},
 \qquad A\longmapsto
        ([x^\alpha]\operatorname{per}_3 A(x))_{\alpha\in\Gamma}.
                                                               \tag{F-025.3}
\]
All coordinates and equations here are over \(\mathbb Z\).
We prove that \(\Phi\) is dominant over \(\mathbb C\).

\paragraph{A fully specified integer witness.}
At the following pencil, row \((i,j)\) gives the seven
coefficients of \(A_{ij}(x)\) in the order
\(x_1,\ldots,x_7\):
\begin{center}
\begin{tabular}{@{}c|rrrrrrr@{}}
\toprule
\((i,j)\)&1&2&3&4&5&6&7\\\midrule
\((1,1)\)&-3&1&3&-3&3&-2&1\\
\((1,2)\)&4&3&-4&-2&1&-2&0\\
\((1,3)\)&4&0&4&-3&3&2&0\\
\((2,1)\)&3&3&3&-3&2&-4&1\\
\((2,2)\)&3&0&-4&3&4&3&2\\
\((2,3)\)&4&0&-1&3&-2&1&0\\
\((3,1)\)&-3&4&2&0&0&-3&-3\\
\((3,2)\)&1&-4&2&2&2&-1&-4\\
\((3,3)\)&4&3&1&2&1&-2&4\\\bottomrule
\end{tabular}
\end{center}
Call this pencil \(A^0\). Order \(\Gamma\) in increasing
lexicographic order of its seven exponent coordinates.
Number source coordinates from zero by
\[
                 c(i,j,k)=7(3(i-1)+(j-1))+(k-1)
 \quad(1\leq i,j\leq3,\ 1\leq k\leq7).
\]
If \(\{u,v\}=\{1,2,3\}\setminus\{i\}\) and
\(\{s,t\}=\{1,2,3\}\setminus\{j\}\), with both
pairs in increasing order, the Jacobian entries are exactly
\[
 J_{\alpha,c(i,j,k)}
   =[x^\alpha]\bigl(x_k
       (A^0_{us}A^0_{vt}+A^0_{ut}A^0_{vs})\bigr).
                                                               \tag{F-025.4}
\]
This follows by differentiating the six-term permanent.
It is also a direct formula specifying every entry of
an integer \(49\)-by-\(63\) matrix.

Retain, in increasing order, the following 49 columns:
\[
              C=\{0,1,\ldots,39\}\cup\{41,42,\ldots,47\}
                                  \cup\{49,50\}.
                                                               \tag{F-025.5}
\]
Exact integer elimination gives
\[
 \begin{gathered}
 \Delta:=\det J_{\Gamma,C}\\
 =-333464241200762901244704190454683312552421068381594185434187935121408.
 \end{gathered}
                                                               \tag{F-025.6}
\]
In particular
\[
                         \Delta\equiv163652\pmod{1000003},
                                                               \tag{F-025.7}
\]
which already certifies \(\Delta\ne0\).
All entries of this finite certificate are specified by the
table, (F-025.2), (F-025.4) and (F-025.5). The included
\path{agents/jet_review/check_per3_binary_transverse_jacobian.py}
constructs the cofactor matrix and verifies the modular minor.
The separate
\path{agents/jet_review/verify_per3_binary_transverse_jacobian.py}
expands the permanent as six products, changes each one of
the 63 source coefficients by one, and compares the resulting
polynomial differences with every column of \(J\).
Those differences are the exact derivatives because the
permanent is linear in each individual matrix entry.
It then computes (F-025.6) over the integers. Thus the
cofactor construction and the full-polynomial construction
are independently specified and agree exactly; no floating-point
rank or numerical tolerance is used.

\begin{proof}[Dominance and all ambient permanent copies]
Fix the fourteen source coordinates omitted from \(C\)
at their values in \(A^0\). On the remaining affine
49-dimensional space, the coefficient map has derivative
with determinant \(\Delta\ne0\). Translate source and
target to their base values and multiply the target by
the inverse of this derivative over \(\mathbb Q\).
The restricted map has the form \(F(t)=t+O(t^2)\).
For a nonzero polynomial in the target coordinates, its
lowest nonzero homogeneous part is unchanged by substitution
into \(F\). Therefore polynomial pullback is injective,
which proves dominance over \(\mathbb C\).
Equivalently, the integral argument works after inverting
\(\Delta\). A unit minor over all of \(\mathbb Z\)
is not asserted.

Let \(w\) be any highest-weight vector of type \(\lambda\)
in the degree-\(d\) ideal of the permanent orbit closure
in its nine-dimensional variable space. Since
\(\ell(\lambda)\leq7\), nonnegative polynomial weights
force every coefficient generator in \(w\) to involve
only \(x_1,\ldots,x_7\). By F-016, its suffix degree
from coordinate three on is at most \(\lambda_3\leq2\).
Thus \(w\) is a polynomial in exactly the coefficient
coordinates (F-025.2), or a subset of them.

Every \(\operatorname{per}_3 A(x)\) is a linear
endomorphism substitution into the original permanent,
and hence belongs to its orbit closure: invertible
endomorphisms are dense in the full endomorphism space.
Consequently \(w\) pulls back to zero along \(\Phi\).
The proved injection forces \(w=0\). This applies to
every vector in the highest-weight ideal multiplicity space.
The entire such space is therefore zero, and the quotient
coordinate ring retains all \(a_\lambda(d[3])\) copies.
This proves \eqref{eq:F025full}. No restriction on \(d\)
was used. Since \(n=m=3\), the internal padding exponent
is zero, so all 63 pencil coefficients are allowed parameters.
\end{proof}

\paragraph{The published determinant-side value and the actual comparison.}
Landsberg--Manivel--Ressayre,
\emph{Hypersurfaces with degenerate duals and the Geometric
Complexity Theory Program}, Comment. Math. Helv. 88 (2013),
469--484, Section~3.2, printed p.~476, state that the module
of highest weight \(12\omega_1+5\omega_2+2\omega_7\)
has ambient multiplicity six in
\(\operatorname{Sym}^{12}\operatorname{Sym}^3\mathbb C^9\)
and ideal multiplicity one for the determinant orbit closure.
The primary publication is
\url{https://ems.press/content/serial-article-files/43321},
DOI \(10.4171/\mathrm{CMH}/292\).
These are known input values, not a new calculation of the
determinant ideal in this run.

The unique polynomial \(\mathrm{GL}_9\) lift in total degree 36 is
\((19,7,2^5,0,0)\). Its adjacent differences are
\(12,5,0,0,0,0,2,0\), and its size is 36.
Characteristic-zero semisimplicity therefore gives
\(K_3=6-1=5\). Formula \eqref{eq:F025full} applies
because its length is seven and its third part is two,
giving \(q^3_\lambda=6\). This proves
\eqref{eq:F025obstruction}, with both actual multiplicities
positive. The independent exact cycle-index computation
in \path{agents/jet_review/lmr_weight_dimensions_certificate.json}
also confirms the ambient value six; it does not claim
to recompute the cited determinant ideal multiplicity.

\paragraph{The orbit and ordinary Kronecker upper bounds are larger still.}
For \(\mu=(12,12,12)\), exact character sums give
\[
 g(\lambda,\mu,\mu)=11,\qquad
             \operatorname{sk}(\lambda,\mu)=10.
                                                               \tag{F-025.8}
\]
Here is a fully finite specification of that computation.
For a cycle partition \(\rho\vdash36\), set
\(z_\rho=\prod_i i^{m_i(\rho)}m_i(\rho)!\), and let
\(\rho^{[2]}\) replace each even part \(2j\) by two
parts \(j\), leaving odd parts unchanged. Then
\[
\begin{split}
 \sum_{\rho\vdash36}
  \frac{\chi^\lambda(\rho)\chi^\mu(\rho)^2}{z_\rho}&=11,\\
 \sum_{\rho\vdash36}
  \frac{\chi^\lambda(\rho)\chi^\mu(\rho^{[2]})}{z_\rho}&=9.
\end{split}                                                   \tag{F-025.9}
\]
The character of a symmetric square at a permutation is
one half of its squared character plus its character at
the squared permutation. Thus the two sums give
\((11+9)/2=10\).

The included
\path{agents/jet_review/check_lmr_rectangular_kronecker.py}
evaluates all \(17977\) cycle partitions with exact rational
arithmetic. Its character recurrence can be specified without
a character table: for a partition \(\nu\) of length \(s\),
form the distinct nonnegative integers
\(\beta_i=\nu_i+s-i\). To remove a cycle part \(k\),
replace each eligible bead \(b\) by \(b-k\geq0\) when
\(b-k\) is not already present. Sort the resulting beads
decreasingly, subtract \((s-1,\ldots,0)\), and call the
resulting partition \(\tau\). That term has sign
\((-1)^{\#\{c:b-k<c<b\}}\). Sum this sign times the
character of \(\tau\) at the remaining cycles, terminating
with the empty partition. This is the border-strip character
recurrence. It also follows immediately by extracting the
coefficient of \(x^{\nu+(s-1,\ldots,0)}\) in the
Vandermonde alternant times the product of power sums:
subtracting \(k\) from one exponent and reordering gives
exactly the stated signs. The Frobenius characteristic
identity is Macdonald, I(7.3)--(7.6), pp.~113--114.
The companion character certificate independently checks
the recurrence against connected border strips in all sizes
through eight.

Consequently the exact comparison in this example is
\[
             K_3=5<q^3_\lambda=a_\lambda=6
                          <\operatorname{sk}=10<g=11.
                                                               \tag{F-025.10}
\]
Thus this actual closure multiplicity obstruction is detected
by neither the stronger test \(q>\operatorname{sk}\) nor
the stronger test \(q>g\). The values \(g\) and
\(\operatorname{sk}\) in (F-025.8) are exact finite
computations; the value \(K_3=5\) remains the separately
cited determinant-closure input.

\paragraph{Grothendieck dictionary and limits.}
The coefficient projection and the explicit pencil define
morphisms of affine schemes over \(\mathbb Z\).
Over \(\mathbb Z[\Delta^{-1}]\), the displayed square
minor is invertible, and the translated source restriction
has universally injective coordinate pullback by the
lowest-term argument. Over \(\mathbb C\), F-016's
Schur-functor support places the complete ideal highest-weight
space in this coordinate algebra. These are the torus
comodules and ideal comodules of SGA~3, I,
Propositions~4.7.2--4.7.3, pp.~34--36, and the Schur
operations of SGA~6, VI, Theorem~3.3, pp.~392--393.
The proof uses an actual coordinate injection to determine
all permanent copies. Positivity of a virtual class or
nonvanishing of an ordinary Kronecker coefficient would
not give the inequality between six and five.
The known determinant-side value and the explicit permanent
morphism have separate check events. This concrete
multiplicity obstruction does not extend the cited IP/BIP
thresholds and does not imply an asymptotic complexity
separation.

% END agents/jet_review/per3_full_transverse_jets.tex

% BEGIN agents/archive/FINITE_PERMANENT_JETS.tex
% Run GCT_KRON_20261120; GPT-6 Astra Pro; effort not exposed.
\section{Finite permanent jet dominance and further actual obstructions}

\subsection*{A-031. Full multiplicities through the maximal possible variable count}
For each \(m\in\{3,4,5,6\}\), every partition
\(\lambda\vdash md\) satisfying
\[
                  \ell(\lambda)\leq2m+2,\qquad\lambda_3\leq2
\]
has full unpadded permanent multiplicity:
\[
                       q^m_\lambda(d[m])=a_\lambda(d[m])
                       \quad\text{for every }d\geq1.           \tag{A-031.1}
\]
Missing parts are zero. This includes every ambient copy, with no
bound on the second row or outer degree. The proof consists of
specified nonzero integer Jacobian minors and the all-copy support
theorem F-016. The value \(2m+2\) is the largest possible variable
count for dominance of the complete coefficient projection used here,
by the dimension argument below. That dimension argument is a
boundary of the projection method, not a converse to (A-031.1).

There are resulting actual obstructions with both multiplicities
positive. Put
\[
 d_m=2m(m-1),\qquad
 \Lambda_m=(2m^2(m-2)+1,\ 2m^2-4m+1,\ 2^{\,2m-1}),              \tag{A-031.2}
\]
and set
\[
 b_m=[z^{\,2m^2-4m-1}]\prod_{j=2}^m(1-z^j)^{-1}>0.
\]
For every \(m\in\{3,4,5,6\}\),
\[
 q^m_{\Lambda_m}(d_m[m])=a_{\Lambda_m}(d_m[m])
                    >K_m(\Lambda_m)\geq b_m.                   \tag{A-031.3}
\]
More concretely:
\[
\begin{array}{c|c|c|c}
 m&d_m&\Lambda_m&\text{proved comparison}\\ \hline
 3&12&(19,7,2^5)&q=6>K=5\\
 4&24&(65,17,2^7)&q=a>K\geq7\\
 5&40&(151,31,2^9)&q=a>K\geq63\\
 6&60&(289,49,2^{11})&q=a>K\geq594
\end{array}                                                   \tag{A-031.4}
\]
Here \(a=a_{\Lambda_m}(d_m[m])\) in each row. In particular the
last three permanent multiplicities are at least \(8,64,595\).
Their exact ambient values and exact determinant multiplicities
are not asserted. The \(m=3\) numerical values are the separately
verified F-025 comparison. These examples have \(n=m\) and are
outside the large-padding IP/BIP ranges. No historical priority
or asymptotic complexity separation is claimed.

\paragraph{The coefficient map and finite witnesses.}
For integers \(m,r\), write
\[
 \Gamma_{m,r}=\{\alpha\in\mathbb Z_{\geq0}^r:
              |\alpha|=m,\ \alpha_3+\cdots+\alpha_r\leq2\}.
\]
Its size is
\[
 D_m(r)=(m+1)+m(r-2)+(m-1)\binom{r-1}{2}.                       \tag{A-031.5}
\]
Let \(A(x)=(\sum_{k=1}^r a_{ijk}x_k)_{1\leq i,j\leq m}\).
The polynomial map over \(\mathbb Z\) under consideration is
\[
 \Phi_{m,r}:\mathbb A^{m^2r}\longrightarrow\mathbb A^{D_m(r)},
 \qquad
 (a_{ijk})\longmapsto
 ([x^\alpha]\operatorname{per}_m A(x))_{\alpha\in\Gamma_{m,r}}.
                                                               \tag{A-031.6}
\]
At an integer pencil \(A^0\), its Jacobian is specified exactly by
\[
 J_{\alpha,(i,j,k)}
     =[x^\alpha]\left(x_k\,
          \operatorname{per}_{m-1}A^0_{\widehat i,\widehat j}(x)
                    \right).                                  \tag{A-031.7}
\]
Rows are in increasing lexicographic order of the \(r\) exponents.
Source columns, numbered from zero, are
\[
                  r\bigl(m(i-1)+(j-1)\bigr)+(k-1).
\]

Here is a fully specified mathematical recipe for all four pencils.
Start with \(u_0=20261121\), and define integers recursively by
\[
 u_{t+1}=(1103515245u_t+12345)\bmod 2^{31}.
\]
For source column \(c\), the coefficient of the pencil is
\[
                        a^0_c=(u_{c+1}\bmod9)-4.
\]
The sequence is restarted at \(u_0\) for each \(m\), and the first
\(m^2r\) coefficients are used in the source order specified above.
Thus this recipe defines every entry without a random choice or
an external data dependency.

The following table gives nonzero determinants of selected
\(D_m(r)\)-column square minors, reduced modulo \(p=1000003\).
In each row, all \(D_m(r)\) output rows are retained.
\[
\begin{array}{c|c|r|r|r}
 m&r&D_m(r)&m^2r&\det J_{\Gamma,C}\pmod p\\ \hline
 3&8&64&72&88699\\
 4&10&145&160&998223\\
 5&12&276&300&434335\\
 6&14&469&504&601963
\end{array}                                                    \tag{A-031.8}
\]
The retained columns \(C\), in increasing order, are the complements
of the following omitted sets in \(\{0,\ldots,m^2r-1\}\).
Here \([a,b]_{\mathbb Z}\) means all integers from \(a\) through \(b\):
\[
\begin{aligned}
 O_3&=\{47,55,63\}\cup[67,71]_{\mathbb Z},\\
 O_4&=\{77,118,129,139\}\cup[149,159]_{\mathbb Z},\\
 O_5&=\{118,179,238,251,263,275\}\cup[282,299]_{\mathbb Z},\\
 O_6&=\{167,251,334,419,433,447,461,474\}
                                      \cup[477,503]_{\mathbb Z}.
\end{aligned}                                                  \tag{A-031.9}
\]
The omitted counts are \(8,15,24,35\), leaving exactly the required
square minors. The coefficient recipe, exponent order,
(A-031.7), and (A-031.9) completely specify these integer matrices.
The displayed nonzero residues prove their integer determinants
are nonzero.

For convenient reproduction, the report bundle additionally stores
the complete pencils, ordered target exponents, selected columns,
and every integer Jacobian entry:
\begin{itemize}
\item \path{agents/archive/per3_8vars_lcg_certificate.json};
\item \path{agents/archive/per4_10vars_lcg_certificate.json};
\item \path{agents/archive/per5_12vars_lcg_certificate.json};
\item \path{agents/archive/per6_14vars_lcg_certificate.json}.
\end{itemize}
The exact generator is
\path{agents/archive/check_permanent_binary_jets_lcg.py},
using the polynomial routines in
\path{agents/archive/check_permanent_binary_jets.py}.
Those routines compute cofactors by a subset dynamic program:
a state records the used columns, and adding the next row
multiplies by its selected linear form and sums over unused columns.
Terms of suffix degree greater than two may be deleted immediately,
because later multiplication cannot lower that degree.
Multiplication of each final cofactor by \(x_k\) gives (A-031.7).
All coefficients are integers, followed by exact modular elimination.
The modular products are below \(p^2<2^{63}\), so the
integer-array operations do not introduce floating-point arithmetic
or overflow. A separate direct-permutation construction reproduced
all \(346984\) integer Jacobian entries for these four pencils and
all four selected minor residues. It also checked the recurrence
and the omitted sets against (A-031.9). These independent records are
\path{agents/obstructions/finite_permanent_jets_independent_certificate.json}
and \path{agents/obstructions/finite_permanent_recipe_certificate.json}.

\begin{proof}[Dominance and retention of all copies]
For any one of the four rows in (A-031.8), fix the source
coordinates not among its selected columns at their values in
the stored integer pencil. The resulting square polynomial map
has an invertible derivative over \(\mathbb Q\) at that point.
Translate source and target and compose with the inverse of the
derivative. The resulting map has the form \(t\mapsto t+O(t^2)\).
Substitution preserves the lowest nonzero homogeneous part of
any nonzero polynomial, so its pullback on polynomial algebras is
injective. This proves dominance of \(\Phi_{m,2m+2}\).

Put \(R=m^2\), the original number of permanent variables.
One has \(2m+2\leq R\) for each \(m\) here.
Let \(w\) be any highest-weight vector of type \(\lambda\)
in the degree-\(d\) ideal of the permanent orbit closure in
\(\operatorname{Sym}^m\mathbb C^{R*}\).
Since \(\ell(\lambda)\leq2m+2\), its nonnegative coordinate
weights force every coefficient generator occurring in \(w\)
to use only the first \(2m+2\) variables.
If the ambient multiplicity is zero there is nothing to prove;
otherwise ordinary Pieri gives \(\ell(\lambda)\leq d\).
F-016 therefore applies, and at its third-row cut says that every
coefficient generator in every such highest-weight vector has
suffix degree at most \(\lambda_3\leq2\). For shapes with fewer
than three rows, their zero suffix weights give the same assertion
directly. Thus \(w\) is a
polynomial in the target coordinates of \(\Phi_{m,2m+2}\).

Every permanent of an \(m\)-by-\(m\) matrix of linear forms in
these variables is obtained from the original permanent by a
linear endomorphism substitution. Such a substitution belongs to
its orbit closure, since invertible endomorphisms are dense and
the substitution map is polynomial. Hence \(w\) pulls back to zero
under \(\Phi_{m,2m+2}\). The proved injection forces \(w=0\).
This kills the entire ideal highest-weight multiplicity space,
not just a chosen vector. Characteristic-zero complete
reducibility then gives (A-031.1).
\end{proof}

\paragraph{The known determinant equations and the positive lower bound.}
The nonzero determinant-ideal copy used here is a known input:
Landsberg--Manivel--Ressayre,
\emph{Hypersurfaces with degenerate duals and the Geometric
Complexity Theory Program}, Comment. Math. Helv. 88 (2013),
Theorem 2.3.1, printed p.~474, together with the determinant
specialization in Section 3.1, pp.~475--476.
The exact primary PDF is
\url{https://ems.press/content/serial-article-files/43321}.
Their weight at degree \(2m(m-1)\) is
\[
 2m(m-1)(m-2)\omega_1+(2m^2-4m-1)\omega_2+2\omega_{2m+1}.
                                                               \tag{A-031.10}
\]
Theorem 2.3.1 supplies this formula, including its factor two
in the first coefficient, and the published Section 3.2 agrees.
The unique polynomial \(\mathrm{GL}_{m^2}\) lift of (A-031.10)
in total degree \(m\,d_m\) is exactly \(\Lambda_m\).
The source gives a nonzero copy in the determinant-closure ideal,
so
\[
             K_m(\Lambda_m)\leq a_{\Lambda_m}(d_m[m])-1.         \tag{A-031.11}
\]
This input only requires \(2m+1\leq m^2\), valid for \(m\geq3\).
By (A-031.1), \(q^m_{\Lambda_m}=a_{\Lambda_m}\) for each of
the four tested sizes. This proves the strict comparison.

For positivity of its determinant side, let
\[
 r_0=2m+1,\quad k=d_m-r_0=2m^2-4m-1,\quad
 \eta=(r_0(m-2)+2,2^{r_0-1}),\quad
 \beta=(k(m-1),k).
\]
The determinant part of O-011 gives a nonzero highest-weight
section of weight \(\eta\), degree \(r_0\), and multiplicity one:
it uses a nondegenerate \(r_0\)-dimensional Hessian restriction
at the identity matrix. This determinant assertion does not
require a padding exponent. O-006.b and the binary coefficient
formula give, in degree \(k\),
\[
 K_m(\beta)=a_\beta(k[m])
   =p_k(m,k)-p_{k-1}(m,k)
   =[z^k]\prod_{j=2}^m(1-z^j)^{-1}=b_m.                        \tag{A-031.12}
\]
The width bound is inactive at these coefficients.
Here \(k\geq5\), and \(m\geq3\), so \(b_m>0\):
an odd \(k\geq5\) is \(3\) plus an even nonnegative integer.

The determinant orbit closure is irreducible, so its coordinate
ring is a domain. Multiplication by the fixed nonzero section
of weight \(\eta\) injects the \(b_m\)-dimensional highest-weight
space of weight \(\beta\) into the highest-weight space of
weight \(\eta+\beta\), in degree \(r_0+k=d_m\).
Direct addition gives \(\eta+\beta=\Lambda_m\), proving
\(K_m(\Lambda_m)\geq b_m\). The coefficient values in
(A-031.4) are obtained by the finite product in (A-031.12).
For \(m=3\), the sharper exact values \(a=6,K=5\) are the
published Section 3.2 input audited in
\path{agents/archive/LMR_SOURCE_SCOPE.md} and used in F-025.
No corresponding exact numerical ideal multiplicity is assumed
for \(m=4,5,6\).

\paragraph{Why the full projection cannot use more variables.}
This boundary statement holds for every \(m\geq2\).
On the open set of pencils for which every entry linear form
is nonzero, the row/column scaling group
\[
 \{(u_1,\ldots,u_m,v_1,\ldots,v_m)\in(\mathbb C^\times)^{2m}:
                        \prod_i u_i\prod_j v_j=1\}/\mathbb C^\times
\]
acts freely and leaves the permanent unchanged. The quotient
removes the kernel \(u_i=c,\ v_j=c^{-1}\), so the effective
dimension is \(2m-2\). Consequently
\[
             \dim\overline{\operatorname{im}\Phi_{m,r}}
                         \leq m^2r-2m+2.                       \tag{A-031.13}
\]
But \(D_m(2m+3)=m^2(2m+3)\), and
\(D_m(r)-m^2r\) strictly increases thereafter.
Thus \(\Phi_{m,r}\) cannot be dominant when \(r\geq2m+3\).
The finite witnesses attain the last possible integer value
\(r=2m+2\). Failure of this full projection at larger \(r\)
does not imply failure of \(q=a\) for any specified Schur type.

\paragraph{Remaining general gate and geometric dictionary.}
An arbitrary-\(m\) construction proving dominance of
\(\Phi_{m,2m+1}\), or of the stronger endpoint
\(\Phi_{m,2m+2}\), is not proved here. The finite witnesses
establish exactly the four sizes stated. A generic binary
diagonal pencil by itself leaves coupled scaling directions
and does not supply an all-\(m\) right inverse. Replacing the
specific Hessian family of a permanent by generic independent
quadratic forms would also require a new justification.

The displayed maps are polynomial morphisms of affine schemes
over \(\mathbb Z\). Their specified minor becomes invertible
after inverting a nonzero integer, and the lowest-term argument
then proves an injective coordinate pullback, compatible with
base extension. The passage to full Schur multiplicity spaces
uses F-016 and characteristic-zero semisimplicity, with the
SGA 3 comodule and SGA 6 Schur-functor dictionary already
specified in that result. The actual comparison uses three
separate pieces: a coordinate injection on the permanent side,
a known nonzero ideal copy on the determinant side, and
domain multiplication for the positive determinant lower bound.
No unproved general padding or base-change statement is used.

% END agents/archive/FINITE_PERMANENT_JETS.tex

% BEGIN agents/jet_review/per3_fixed_padding_copies.tex
\section{Preserving six permanent copies under internal padding}

\subsection*{Finding F-026: a fixed-entry coefficient projection and an all-copy injection}
Let \(n\geq3\), and let \(\nu\vdash3d\) satisfy
\(\ell(\nu)\leq7\) and \(\nu_3\leq2\). Define
\[
                    \nu[n]=(\nu_1+d(n-3),\nu_2,\ldots).
\]
For the internally padded \(3\)-by-\(3\) permanent,
\begin{equation}\label{eq:F026copies}
                   q^3_{\nu[n]}(d[n])\geq a_\nu(d[3]).
\end{equation}
In particular, for every \(n\geq3\),
\begin{equation}\label{eq:F026six}
              q^3_{(12n-17,7,2^5)}(12[n])\geq6.
\end{equation}
The right side is a count of independent actual permanent
copies. An equality with the full ambient multiplicity
at inner degree \(n\) is not asserted when \(\nu_2>3\).
This distinction applies in particular to \eqref{eq:F026six}.

\paragraph{The coefficient map with its distinguished entry fixed.}
Use the same seven variables, target \(\Gamma\),
coefficient ordering and source numbering as in F-025.
Require the first matrix entry to equal \(x_1\), so
only the 56 source coordinates numbered 7 through 62
may vary. The following new integer pencil meets that
condition:
\begin{center}
\begin{tabular}{@{}c|rrrrrrr@{}}
\toprule
\((i,j)\)&1&2&3&4&5&6&7\\\midrule
\((1,1)\)&1&0&0&0&0&0&0\\
\((1,2)\)&0&4&2&-3&-1&0&-3\\
\((1,3)\)&2&1&3&4&1&3&-3\\
\((2,1)\)&0&3&4&-1&2&-2&2\\
\((2,2)\)&-3&-4&-2&-1&1&-2&3\\
\((2,3)\)&3&3&0&2&-3&4&2\\
\((3,1)\)&-2&-3&-3&3&1&-1&3\\
\((3,2)\)&0&-1&2&2&1&4&-1\\
\((3,3)\)&4&-3&-1&-3&-1&4&-1\\\bottomrule
\end{tabular}
\end{center}
The Jacobian is again specified by the coefficient formula
(F-025.4), now at this pencil. Retain the 49 allowed columns
\[
                 C'=\{7,8,\ldots,47\}
                        \cup\{49,50,\ldots,54\}\cup\{56,57\}.
                                                               \tag{F-026.1}
\]
Their square determinant is
\[
 \begin{gathered}
       \Delta'=31568108869448599905478876197553019094163200564336768,
          \\
          \Delta'\equiv641253\pmod{1000003}.
 \end{gathered}                                      \tag{F-026.2}
\]
Thus the projection onto all 49 coordinates in \(\Gamma\)
is dominant even with \(A_{11}=x_1\) fixed.
The algebraic argument is exactly the invertible-differential
argument in F-025: restrict to these 49 directions, translate,
and normalize the linear term to obtain \(t+O(t^2)\).
It works over \(\mathbb Z[(\Delta')^{-1}]\).

The complete finite witness is the displayed table together
with \(\Gamma\), (F-025.4), and (F-026.1).
The scripts
\path{agents/jet_review/check_per3_fixed_entry_transverse.py}
and
\path{agents/jet_review/verify_per3_fixed_entry_transverse.py}
construct the cofactor matrix and independently expand the
full permanent. They verify all coefficient columns and
compute (F-026.2) exactly. The second calculation includes
the seven columns which are subsequently held fixed, but
the selected minor uses only the 56 permitted columns.

\paragraph{The highest-weight injection and the actual padding pullback.}
Let \(U\) be the coordinate-function space on the first
seven variables. For a nonzero ambient type, let \(H_\nu\)
be its entire highest-weight multiplicity space in
\(\operatorname{Sym}^d\operatorname{Sym}^3 U\).
Multiplication in each column of the Cauchy matrix model
by \(e_1^{n-3}\) gives an injective map
\[
 T_n:H_\nu\longrightarrow
           \bigl(\operatorname{Sym}^d\operatorname{Sym}^n U\bigr)
                          ^{\mathrm{HW},\nu[n]}.
                                                               \tag{F-026.3}
\]
Indeed the upper unipotent group fixes \(e_1\), the
weight changes by \(d(n-3)e_1\), and column permutations
commute with this multiplication. The map is injective on
monomial tensors and therefore on their invariant subspace.
Only an injection is used here; F-014's extra conditions
for surjectivity are not assumed.

Let \(R_n\) be the actual pullback induced by padding a
cubic form \(F\) as \(x_1^{n-3}F\).
Under the natural differential, or apolar, pairing, on a
cubic monomial generator \(e^\alpha\) one has
\[
       (R_nT_n)(e^\alpha)
          =\frac{(\alpha_1+n-3)!}{\alpha_1!}\,e^\alpha.
                                                               \tag{F-026.4}
\]
The factors in the other coordinates cancel. For products
of \(d\) generators this formula extends multiplicatively.
Thus \(R_nT_n\) is an invertible diagonal operator on
the relevant polynomial coefficient algebra over
\(\mathbb C\). This is the same precise padding
calculation as O-012 and F-015; it is not an identification
of the highest-weight lift with the padding morphism itself.

\begin{proof}[Every original ambient copy remains independent]
By F-016, every polynomial in \(H_\nu\) uses only cubic
coefficient generators whose suffix degree from the third
variable is at most \(\nu_3\leq2\). Hence
\(H_\nu\subset\mathbb C[\Gamma]\). The diagonal
operator (F-026.4) preserves this coefficient algebra and
is invertible on it. It need not preserve its highest-weight
subspace; that is unnecessary.

Suppose a linear combination of the highest-weight vectors
\(T_n(H_\nu)\) vanishes on the internally padded permanent
closure. Evaluate it on the substitutions furnished by the
fixed-entry pencils above. Because \(A_{11}=x_1\), these
are exactly
\[
                         x_1^{n-3}\operatorname{per}_3 A(x),
\]
with the padding form equal to the substituted distinguished
matrix entry. They are points of the required orbit closure
by endomorphism substitution in its \(n^2\)-variable space.
Pulling the alleged linear relation back gives the
corresponding combination of \((R_nT_n)(H_\nu)\)
vanishing on the dominant \(\Gamma\)-projection of the
fixed-entry permanent family. Injectivity of that projection's
coordinate pullback makes this combination zero as a polynomial
in \(\mathbb C[\Gamma]\). The invertibility of (F-026.4)
then makes the original combination in \(H_\nu\) zero.

Therefore the \(a_\nu(d[3])\) highest-weight vectors in
the image of (F-026.3) remain independent on the padded
permanent closure. This proves \eqref{eq:F026copies}.
For \(\nu=(19,7,2^5)\), the known and independently
checked ambient coefficient is six by F-025. Its first
part shifts by \(12(n-3)\), giving \eqref{eq:F026six}.
\end{proof}

If additionally \(\nu_2\leq3\), the all-copy stabilization
F-014 identifies the new ambient multiplicity with the old
one, so \eqref{eq:F026copies} is an equality with both.
For the second row seven in \eqref{eq:F026six}, no such
conclusion follows from this argument.

\paragraph{The shifted obstruction stops at inner degree three.}
For the particular family \(\nu[n]=(12n-17,7,2^5)\),
O-027 gives full determinant multiplicity for every \(n\geq4\):
its length is seven, \(7\leq2n\), and its third row is two.
Consequently
\[
       6\leq q^3_{\nu[n]}(12[n])
                 \leq K_n(\nu[n])=a_{\nu[n]}(12[n])
                           \qquad(n\geq4).
                                                               \tag{F-026.5}
\]
Thus the exact \(6>5\) obstruction at \(n=3\) from F-025
does not persist under these first-row shifts. Six actual
permanent copies do persist, and the determinant closure
has enough copies to accommodate them at every larger
inner degree. The finite independent determinant-Jacobian
certificates in this folder also verify the seven-variable
projection at \(n=4,5\), with ranks 70 and 91 respectively;
the complete arbitrary-size proof of (F-026.5) is O-027.

\paragraph{Grothendieck dictionary and scope.}
The fixed-entry pencil parameter space is an affine linear
subscheme defined over \(\mathbb Z\), and the coefficient
projection is dominant after base change to \(\mathbb C\),
with an integral coordinate injection after inverting
\(\Delta'\). The highest-weight map is multiplication
by the first basis line in each column; the Schur-functor
and Cauchy interpretation is SGA~6, VI, Theorem~3.3,
pp.~392--393, together with Macdonald I(4.3), p.~63,
Appendix~A(5.3)--(5.6), pp.~155--156.
The precise torus grading and ideal-comodule statements
are SGA~3, I, Propositions~4.7.2--4.7.3, pp.~34--36.
The factorial comparison is used in characteristic zero;
it is not silently asserted to be invertible in every
residue characteristic. The theorem counts actual surviving
copies and does not infer an actual obstruction from their
positivity alone.

% END agents/jet_review/per3_fixed_padding_copies.tex

% BEGIN proofs/projection_bound.tex
\section{A stronger multiplicity exclusion from projections}

The following two findings give a comparison of actual coordinate-ring
multiplicities, rather than just nonvanishing of an upper bound.
They are derived here in an explicit form; no assertion of historical
priority is intended. They use the polynomial-functor dictionary of
cote~133, scan~028, but require an additional determinantal construction.

\subsection*{Finding M-006: projection containment compares all copies}
Let \(V\) be a complex vector space, let \(U\) be an \(r\)-dimensional
quotient of \(V\), and choose a linear section \(U\hookrightarrow V\).
Write
\[
 F_d(V)=\operatorname{Sym}^d(\operatorname{Sym}^n V),\qquad
 \pi:\operatorname{Sym}^n V^*\longrightarrow
      \operatorname{Sym}^n U^*
\]
for restriction followed, when necessary, by the inclusion of forms
depending only on \(U\).
Let \(X,Y\subseteq\operatorname{Sym}^n V^*\) be reduced closed
\(\mathrm{GL}(V)\)-stable cones. Suppose
\begin{equation}\label{eq:M006hyp}
                         \pi(Y)\subseteq X.
\end{equation}
Then for every \(d\) and every partition \(\lambda\vdash nd\)
with \(\ell(\lambda)\leq r\),
\begin{equation}\label{eq:M006}
 \operatorname{mult}_{S_\lambda V}\mathbb C[Y]_d
 \ \leq\
 \operatorname{mult}_{S_\lambda V}\mathbb C[X]_d.
\end{equation}
Only the stated containment of the projected variety is required;
containment of \(Y\) itself has not been assumed.

\begin{proof}
Choose coordinates in which \(U\) is the span of the first \(r\)
basis vectors. The natural Schur--Weyl decomposition gives
\[
 F_d(V)=\bigoplus_{\lambda\vdash nd} S_\lambda V\otimes M_\lambda,
 \qquad
 M_\lambda=[\lambda]^{S_n\wr S_d}.
\]
The ideals \(I(X)_d\) and \(I(Y)_d\) are
\(\mathrm{GL}(V)\)-submodules; hence, in characteristic zero, they
have the forms
\[
 I(X)_d=\bigoplus_\lambda S_\lambda V\otimes J_{X,\lambda},
 \qquad
 I(Y)_d=\bigoplus_\lambda S_\lambda V\otimes J_{Y,\lambda}
\]
for subspaces of the same \(M_\lambda\).
The natural inclusion \(F_d(U)\hookrightarrow F_d(V)\) identifies
\(S_\lambda U\otimes M_\lambda\) with the corresponding summand
of \(F_d(U)\). If \(\ell(\lambda)\leq r\), it is nonzero.
Any \(f\in S_\lambda U\otimes J_{X,\lambda}\) belongs to \(I(X)_d\)
and depends only on the projection to \(U\). For \(y\in Y\),
\[
                              f(y)=f(\pi(y))=0
\]
by \eqref{eq:M006hyp}. Since \(Y\) is reduced over the algebraically
closed field \(\mathbb C\), this means \(f\in I(Y)_d\).
The \(\mathrm{GL}(V)\)-span of the nonzero
\(S_\lambda U\subseteq S_\lambda V\) is the irreducible
\(S_\lambda V\). Thus
\(S_\lambda V\otimes J_{X,\lambda}\subseteq I(Y)_d\), giving
\(J_{X,\lambda}\subseteq J_{Y,\lambda}\).
Subtracting their dimensions from \(\dim M_\lambda\) proves
\eqref{eq:M006}. This proves the comparison for all copies of the
representation, not just its presence.
\end{proof}

\paragraph{Equivalent highest-weight formulation.}
A polynomial highest-weight vector of weight
\((\lambda_1,\ldots,\lambda_r,0,\ldots,0)\) lies in \(F_d(U)\):
the polynomial weights have nonnegative coordinates, so trailing
zero weights allow no coordinate outside \(U\).
The proof therefore says that every equation of type \(\lambda\)
for \(X\) is also an equation for \(Y\) once all its projected
arguments already lie in \(X\).

\subsection*{Finding M-007: an explicit determinant for every padded form}
Let \(m\geq2\), \(r\geq2\), and put
\[
 p=\left\lfloor\frac{r-1}{2}\right\rfloor,\quad
 q=\left\lceil\frac{r-1}{2}\right\rceil,\qquad
 S(m,r)=\binom{m+p}{p}+\binom{m+q}{q}-1.
\]
For any degree-\(m\) homogeneous form \(f\) in at most \(r\)
variables and any nonzero linear form \(L\) in the same variable
space, the form \(L^{n-m}f\) is the determinant of an \(n\times n\)
matrix of homogeneous linear forms whenever \(n\geq S(m,r)\).
The entries may use the coefficients of \(f\) as constants.
The underlying determinant identity is an identity over
\(\mathbb Z\) in universal coefficient variables.

\begin{proof}
Take \(L\) as one coordinate and split the other \(r-1\) coordinates
into \(x=(x_1,\ldots,x_p)\) and \(y=(y_1,\ldots,y_q)\). Write
\begin{equation}\label{eq:M007expansion}
 f=\sum_{|\alpha|+|\beta|\leq m}
       c_{\alpha,\beta}x^\alpha y^\beta
                      L^{m-|\alpha|-|\beta|}.
\end{equation}
For \(p=0\) the only \(x\)-multi-index is the empty one.

Construct a rooted tree whose vertices are all multi-indices
\(\alpha\in\mathbb Z_{\geq0}^{p}\) with \(|\alpha|\leq m\).
Its root is \(0\). For every nonzero \(\alpha\), let \(i\) be
the largest index for which \(\alpha_i>0\), and draw the edge
\(\alpha-e_i\longrightarrow\alpha\), with label \(-x_i\).
There is a unique path from the root to \(\alpha\); its product
of labels is \((-1)^{|\alpha|}x^\alpha\).
Make a second tree on the analogous \(y\)-multi-indices
\(\beta\), but reverse all its edges, using label \(-y_j\).
Its unique path from \(\beta\) to the root has product
\((-1)^{|\beta|}y^\beta\).
Identify the two roots into one central vertex, and keep all
other vertices distinct.

For each pair with \(|\alpha|+|\beta|\leq m\), insert a directed
edge from the prefix vertex \(\alpha\) to the suffix vertex
\(\beta\), labelled \(c_{\alpha,\beta}L\).
The pair \((0,0)\) gives the central loop
\(c_{0,0}L\). Put a loop labelled \(L\) at every noncentral
vertex. Let \(M\) be the matrix whose entry \(M_{uv}\) is the
sum of the labels of all edges \(u\to v\).
All its entries are homogeneous linear forms, and its size is
the number of vertices,
\[
                    S=\binom{m+p}{p}+\binom{m+q}{q}-1.
\]

We now evaluate the determinant by its permutation expansion.
Aside from loops, every directed cycle must pass through the
central vertex. Indeed, the prefix tree only points outwards,
the suffix tree only points inwards, and coefficient edges
only point from a prefix vertex to a suffix vertex.
A cycle through the central vertex uses exactly one coefficient
edge, from some \(\alpha\) to some \(\beta\).
Consequently a cycle cover either uses the central loop, or
chooses exactly one such nontrivial cycle and uses loops at
all unused vertices.

Put \(k=|\alpha|+|\beta|\). The chosen cycle has \(k+1\) edges:
\(k\) variable edges and one coefficient edge.
Its permutation sign is \((-1)^k\); its variable-edge labels
have the same sign. The signs cancel. The coefficient edge
and the \(S-k-1\) unused loops together contribute \(L^{S-k}\).
This cycle cover therefore contributes exactly
\[
                   c_{\alpha,\beta}x^\alpha y^\beta L^{S-k}.
\]
The central-loop case gives the same expression for \(k=0\).
Summing all cycle covers and using \eqref{eq:M007expansion} gives
\[
                                \det M=L^{S-m}f.
\]
For \(r\geq2\), \(S\geq m+1\), so the exponent is nonnegative.
If \(n>S\), adjoining a diagonal block \(L I_{n-S}\) gives the
required matrix of size \(n\).
All cancellations and coefficient signs in this argument are
integral polynomial identities, so the construction with
indeterminate \(c_{\alpha,\beta}\) is defined over \(\mathbb Z\).
\end{proof}

\paragraph{Boundary cases and complexity meaning.}
For \(r=1\), \(f\) is a scalar multiple of \(L^m\), and a
diagonal determinant suffices. A degree-one form is similarly
immediate. If \(L=0\) and \(n>m\), the padded form is zero and
has a zero determinantal representation.
The result is a size bound with arbitrary scalar coefficients;
it does not assert that the coefficients of a projected
permanent can be computed in polynomial time.

\subsection*{Finding M-008: row bounds for actual GCT multiplicity obstructions}
Let \(m\geq2\), \(n>m\), \(V=\mathbb C^{n^2}\), and use
\(q^m_\lambda(d[n])\) and \(K_n(\lambda)\) for the actual
padded-permanent and determinant-closure multiplicities.
If \(n\geq S(m,r)\) and \(r\leq n^2\), then
\begin{equation}\label{eq:M008}
        q^m_\lambda(d[n])\leq K_n(\lambda)
   \qquad\text{for every \(d\) and every \(\ell(\lambda)\leq r\)}.
\end{equation}
In particular:
\begin{enumerate}
\item If \(n>3m^4\), no actual multiplicity obstruction has
      at most nine rows.
\item If \(n\geq m^{25}\), no actual multiplicity obstruction
      has at most fifty-one rows.
\item More generally, for every fixed integer \(c\geq6\),
      \(n\geq m^c\) excludes all outputs with at most \(2c+1\)
      rows.
\end{enumerate}
These are statements about the closure multiplicity \(K_n\),
so they also exclude the stronger sufficient test
\(q^m_\lambda>g_{n,d}(\lambda)\) in these ranges.

\begin{proof}
An \(r\)-variable restriction of any translate of the padded
permanent has the form \(L^{n-m}f\), where \(f\) has degree \(m\)
in those variables. Finding M-007 makes it a determinant when
\(L\ne0\); the zero case was handled above.
Every linear substitution into \(\det_n\) belongs to
\(\Omega_n\), since \(\mathrm{GL}(V)\) is dense in
\(\operatorname{End}(V)\). Thus the projection of the entire
padded-permanent orbit lies in \(\Omega_n\).
The preimage of this closed set under the restriction morphism
is closed, so the same containment holds for the orbit closure.
Apply M-006 with \(Y=Z_{n,m}\) and \(X=\Omega_n\).

For the numerical bounds, when \(r=2c+1\) one has
\(p=q=c\) and \(S(m,r)=2\binom{m+c}{c}-1\).
For fixed \(c\), the real function
\[
 \frac{2\binom{m+c}{c}}{m^c}
       =\frac{2}{c!}\prod_{j=1}^c\left(1+\frac jm\right)
\]
is decreasing for \(m>0\).
For \(m\geq2\) it is at most
\[
                  \frac{(c+1)(c+2)}{2^c}.
\]
For \(c\geq6\), this is at most one: it is \(56/64\)
at \(c=6\), and the ratio at the next integer is
\((c+3)/(2(c+1))<1\).
Therefore \(S(m,2c+1)<m^c\), proving the third assertion
and, with \(c=25\), the second.
For \(c=4\), the same bound is \(30/16=15/8<3\);
thus \(S(m,9)<3m^4\), proving the first.
All these row counts fit inside \(n^2\) in the indicated ranges.
\end{proof}

\paragraph{Exact scope of the new gate.}
At the Ikenmeyer--Panova threshold \(n>3m^4\), the remaining
actual-multiplicity candidates have at least ten rows; at the
B\"urgisser--Ikenmeyer--Panova threshold \(n\geq m^{25}\), they
have at least fifty-two. They must also have positive
permanent-side multiplicity and satisfy the usual padded-shape
conditions. No assertion that these remaining candidates are
obstructions has been proved.
Every ambient constituent of \(\operatorname{Sym}^d\operatorname{Sym}^nV\)
has at most \(d\) rows: it is a constituent of \(h_n^d\) under the
character map, and Pieri's rule adds at most one new row at each step
(Macdonald, I, (5.16)). Thus the respective remaining actual-obstruction
degrees also satisfy \(d\geq10\) and \(d\geq52\).
This is independent of the stronger inequality \(d>n/m\), which is
known to be necessary for the sufficient test \(q>g\), not for every
possible \(q>K_n\).
The four-row ambient loss in O-008 is consistent with M-008:
retaining full ambient multiplicity and retaining at least the
permanent multiplicity are different assertions.

\paragraph{Grothendieck dictionary.}
The restriction map and the Schur decomposition in M-006 are
polynomial-functor operations on locally free sheaves, as in
cote~133, scan~028, and Macdonald, I, Appendix~A,
(5.3)--(5.6), pp.~155--156.
The matrix in M-007 defines a morphism over
\(\mathbb Z[c_{\alpha,\beta}]\) from an affine space of
coefficient data to the affine scheme of linear matrix pencils.
Its determinant is the displayed universal section of
\(\mathcal O(n)\). The projection containment is a containment
of reduced characteristic-zero orbit-closure schemes, and its
ideal inclusion on isotypic summands is the proof of M-008.
SGA~3, I, 4.7.2 supplies the equivariant-module language and
SGA~6, VI, 3.3 the lambda operations on their \(K_0\).
No semisimple splitting over all of \(\mathbb Z\) is inferred
from the characteristic-zero multiplicity comparison.

% END proofs/projection_bound.tex

% BEGIN proofs/weight_box_bound.tex
\section{A shape-sensitive comparison of actual multiplicities}

\subsection*{Finding M-009: bounded weight boxes improve the row gate}
Let \(m\geq2\), \(n>m\), \(d\geq1\), and let
\(\lambda\vdash nd\), with \(r=\ell(\lambda)\leq n^2\).
For a subset \(A\subseteq\{2,\ldots,r\}\), define
\[
 N_A(m,\lambda)=
 \#\{\alpha\in\mathbb Z_{\geq0}^{A}:
       \alpha_i\leq\lambda_i\ (i\in A),\quad |\alpha|\leq m\}.
\]
In particular \(N_\varnothing=1\). Put
\[
 T_m(\lambda)=
 \min_{A\sqcup B=\{2,\ldots,r\}}
       \bigl(N_A(m,\lambda)+N_B(m,\lambda)-1\bigr).
\]
Then
\begin{equation}\label{eq:M009}
 n\geq T_m(\lambda)
 \quad\Longrightarrow\quad
 q^m_\lambda(d[n])\leq K_n(\lambda).
\end{equation}
This compares the actual coordinate-ring multiplicities. It refines
M-008 because deleting the individual bounds \(\alpha_i\leq\lambda_i\)
recovers its two-trie bound.

\begin{proof}
Use \(V=\mathbb C^{n^2}\), with basis \(e_1,\ldots,e_{n^2}\),
and let \(x_i\) denote its dual coordinates. Let \(w\) be a highest-weight
vector of weight \(\lambda\) in \(I(\Omega_n)_d\), for the upper
triangular Borel in \(\mathrm{GL}(V)\). A degree-\(d\) polynomial
function on forms is a polynomial in the coefficient functions
\(a_\eta\), where \(|\eta|=n\); the weight of \(a_\eta\) is
\(\eta\). Every coefficient monomial occurring in \(w\) has total
weight \(\lambda\). All these weights have nonnegative coordinates.
It follows that \(w\) depends only on the first \(r\) coordinates,
and that no coefficient \(a_\eta\) with
\(\eta_i>\lambda_i\) can occur in \(w\).

Restrict a general translate of the padded permanent to those
\(r\) coordinates. It is \(L^{n-m}f\), with \(f\) homogeneous
of degree \(m\). On the nonempty open set on which the coefficient
\(L_1=L(e_1)\) is nonzero, an upper triangular change of basis
aligns \(L\) with \(x_1\). Explicitly, choose an upper unipotent
matrix whose first row is \((1,L_2/L_1,\ldots,L_r/L_1)\);
then \(L\circ g^{-1}=L_1x_1\).
The vector \(w\) is a Borel eigenvector, so this change of basis
rescales its value by a nonzero scalar (it is unchanged for the
unipotent choice). Absorbing \(L_1^{n-m}\) into \(f\), it suffices
to evaluate \(w\) on a form \(x_1^{n-m}f\).

Expand
\[
 f=\sum_{|\alpha|\leq m}
            c_\alpha x_2^{\alpha_2}\cdots x_r^{\alpha_r}
                         x_1^{m-|\alpha|}.
\]
Delete all terms for which some \(\alpha_i>\lambda_i\),
\(i\geq2\), and call the resulting form \(f_\lambda\).
The preceding weight observation proves exactly
\[
       w(x_1^{n-m}f)=w(x_1^{n-m}f_\lambda).
\]
There is no approximation in this equality: every deleted form
coefficient is absent from the polynomial \(w\).

Choose a partition \(A\sqcup B\) attaining \(T_m(\lambda)\).
Make the two trees of M-007, but retain only the multi-indices
bounded by \(\lambda_i\) in their respective coordinates and
of total degree at most \(m\). These vertex sets are closed
under decreasing any coordinate, so the parent rule remains valid.
Their total number of vertices, after identifying their roots, is
\(S=N_A+N_B-1\). Add the coefficient edges for all retained
pairs of total degree at most \(m\), and put \(x_1\) on the
unused-vertex loops. The same cycle-cover argument gives, term by term,
\[
 \det M=\sum_{\substack{|\alpha|\leq m\\
                       \alpha_i\leq\lambda_i\ (i\geq2)}}
                c_\alpha x_2^{\alpha_2}\cdots x_r^{\alpha_r}
                                x_1^{S-|\alpha|}.
\]
Every exponent is nonnegative: the relevant cycle has
\(|\alpha|+1\) distinct vertices in a graph with \(S\) vertices.
This termwise argument is necessary because \(S\) can be smaller
than \(m\). Since \(n\geq S\), adjoining
\(x_1 I_{n-S}\) gives an \(n\times n\) homogeneous linear
determinant equal to \(x_1^{n-m}f_\lambda\).
It belongs to \(\Omega_n\), and \(w\) vanishes on it.

We have proved that \(w\) vanishes on a dense open subset of the
padded-permanent orbit. The orbit is the image of the irreducible
variety \(\mathrm{GL}(V)\), and the condition \(L(e_1)\ne0\)
defines a nonempty open subset of that parameter space. Polynomial
vanishing therefore extends to the entire orbit and its closure.
Hence every highest-weight vector of weight \(\lambda\) in
\(I(\Omega_n)_d\) belongs to \(I(Z_{n,m})_d\).

In the natural decomposition
\(F_d(V)=\bigoplus_\mu S_\mu V\otimes M_\mu\),
the highest-weight space of the \(\lambda\)-isotypic ideal
is its one-dimensional highest-weight line tensored with
\(J_{\Omega_n,\lambda}\). The inclusion just proved gives
\(J_{\Omega_n,\lambda}\subseteq J_{Z_{n,m},\lambda}\).
Subtracting their dimensions from \(\dim M_\lambda\) proves
\eqref{eq:M009} for all copies.
\end{proof}

\paragraph{Explicit thin-tail consequence.}
If \(\lambda_2\leq k\), then every tail row is at most \(k\),
and a balanced split gives
\[
 T_m(\lambda)\leq
  (k+1)^{\lfloor(r-1)/2\rfloor}
       +(k+1)^{\lceil(r-1)/2\rceil}-1.
\]
Thus the right-hand side being at most \(n\) excludes an actual
multiplicity obstruction, even for shapes with more rows than
M-008 covers. The bound counts all points of the box, so discarding
points of total degree greater than \(m\) can improve it further.
For unequal tail rows one may minimize the exact bounded-monomial
counts; no efficient optimization algorithm is claimed.

\paragraph{What this changes and what remains.}
The residual candidate gate now includes \(n<T_m(\lambda)\),
in addition to the independent occurrence, padding, and degree
conditions. Unlike the universal-form method limited by O-010,
this argument only preserves evaluation of a specified highest-weight
space. It does not assert containment of all projected forms in
the determinant closure. It supplies no conclusion for shapes
failing its inequality.

\paragraph{Grothendieck dictionary.}
The coefficient functions are weight sections for the diagonal torus
of \(\mathrm{GL}(V)\); the weight decomposition is the grading/comodule
correspondence of SGA~3, I, Proposition~4.7.3. The upper Borel
acts by its highest-weight character on the line being evaluated.
The deleted coefficients lie outside the finite multigraded support
of that section. The bounded-tree matrices are morphisms of affine
coefficient schemes defined over \(\mathbb Z\), with the displayed
determinant identities integral. The passage from a highest-weight
ideal inclusion to a multiplicity inequality uses semisimplicity
only on the characteristic-zero fiber, as in M-006 and the
Schur-functor dictionary of cote~133, scan~028. No equality of
integral multiplicity modules or of bad-characteristic fibers
is asserted.

% END proofs/weight_box_bound.tex

% BEGIN proofs/chow_degree_bound.tex
\section{A degree exclusion from the canonical Foulkes--Howe map}

\subsection*{Finding M-010: full ambient determinant multiplicity in small degree}
Let \(n\geq1\), let \(V=\mathbb C^{n^2}\), and let
\(\Omega_n\subseteq\operatorname{Sym}^nV^*\) be the reduced
determinant orbit closure. If \(d\geq2\) and
\[
                              n\geq d(d-1),
\]
then \(I(\Omega_n)_d=0\). Equivalently, for every partition
\(\lambda\vdash nd\) that fits in \(V\),
\begin{equation}\label{eq:M010}
                         K_n(\lambda)=a_\lambda(d[n]).
\end{equation}
The same equality holds in degree one. Consequently an actual
padded-permanent multiplicity obstruction must have
\[
                              d(d-1)>n.
\]
This is a statement about the closure multiplicity itself. It does
not replace the different, stronger-test degree restriction
\(d>n/m\) for \(q^m_\lambda>g_{n,d}(\lambda)\).

\paragraph{Known input and exact locator.}
OpenAI, \emph{Quadratic stabilization of the canonical
Foulkes--Howe map}, dated 25 September 2026,
Theorem~1 and, directly, Proposition~3 in Section~4, proves that
the Chow cone of products of \(n\) linear forms has no degree-\(d\)
equations when \(d\geq2\) and \(n\geq d(d-1)\).
The primary source is the authors' repository:
\begin{quote}\small
\url{https://github.com/openai/math/blob/main/preprints/Quadratic-Stabilization-of-the-Canonical-Foulkes-Howe-Map-September-25-2026/build/paper.tex}.
\end{quote}
The exact retrieved source has SHA-256
\begin{quote}\small\path{6b4eb1977533b9b1503927a55687c72afd9dc68b05b27b94cf404413195dd2c2}.
\end{quote}
Theorem~1 is at source lines 99--103, the differential injectivity
lemma at lines 260--291, and Proposition~3 and its proof at
lines 304--430. The complete written proof was read in this run
by root and independently by the archive worker. The detailed
audit is \path{agents/archive/FOULKES_INDEPENDENT_AUDIT.md}.
No formal verification is claimed, and this known theorem is
not presented as a new finding of this run.

\begin{proof}[Proof of the determinant consequence]
For linear forms \(L_1,\ldots,L_n\in V^*\),
\[
                     L_1\cdots L_n
                     =\det\operatorname{diag}(L_1,\ldots,L_n).
\]
This is a linear endomorphism substitution into \(\det_n\).
Every such substitution belongs to its orbit closure: the
invertible endomorphisms are dense in \(\operatorname{End}(V)\),
and substitution is a morphism. Thus the Chow cone is contained
in \(\Omega_n\), and the ideal inclusion is
\[
                     I(\Omega_n)_d\subseteq I(\operatorname{Ch}_n)_d.
\]
The cited Proposition~3 makes the right side zero in the stated
range. Hence the quotient defining \(\mathbb C[\Omega_n]_d\)
has no kernel, and every ambient irreducible copy survives.
For \(d=1\), pure powers lie in the Chow cone and span
\(\operatorname{Sym}^n V^*\), so there are no linear equations.
Finally the permanent-closure coordinate ring is also an ambient
quotient, so \(q^m_\lambda\leq a_\lambda=K_n(\lambda)\)
whenever the displayed equality holds. This proves the gate.
\end{proof}

\paragraph{Grothendieck dictionary.}
The multiplication map \((V^*)^n\to\operatorname{Sym}^n V^*\)
and the diagonal-matrix map are morphisms of affine schemes over
\(\mathbb Z\). Taking their scheme-theoretic images gives the
integral incidence underlying the containment. The equality of
degree pieces above is proved on the reduced characteristic-zero
fiber, where the cited differential argument and irreducible
decomposition apply. The homogeneous ideals are graded
\(\mathrm{GL}(V)\)-comodules, using SGA~3, Expos\'e~I,
Proposition~4.7.2, pp.~34--35, and the torus grading is
Proposition~4.7.3, pp.~35--36. The symmetric and Schur functors
are the operations in SGA~6, Expos\'e~VI, Theorem~3.3,
pp.~392--393. The ideal and quotient define coherent graded
families of the type parameterized by FGA, Bourbaki~221,
Theorems~3.1--3.2, pp.~259--260. None of these integral
definitions asserts that the characteristic-zero kernel
vanishing holds in every residue characteristic.

For \(n\geq m^{25}\) this leaves only degrees satisfying
\(d(d-1)>n\geq m^{25}\). Together with the other shape gates,
it strictly reduces the still unresolved evaluation-rank
comparison; it does not settle that comparison in higher degree.

% END proofs/chow_degree_bound.tex

% BEGIN agents/archive/SHAPE_CHOW_BOUND.tex
% Run GCT_KRON_20261120; GPT-6 Astra Pro; effort not exposed.
\section{A second-row-sensitive vanishing theorem for Chow equations}

\subsection*{A-020. A shape-dependent improvement of the degree gate}
Let $V$ be a finite-dimensional complex vector space, and let
$\operatorname{Ch}_b(V)\subseteq\operatorname{Sym}^bV$ be the
Zariski closure of products of $b$ vectors, for $b\geq1$.
For a partition $\lambda\vdash ab$, write $\lambda_2=0$ if it
has one row. Then
\[
 b\geq\lambda_2+\binom a2
 \quad\Longrightarrow\quad
 I(\operatorname{Ch}_b(V))_a[\lambda]=0.                 \tag{A-020.1}
\]
Here $[\lambda]$ denotes the complete $S_\lambda V^*$ isotypic
component, including every multiplicity copy. Consequently the
degree-$d$ coordinate ring of the determinant orbit closure in
$n^2$ variables has full ambient multiplicity at $\lambda$ if
\[
                         n\geq\lambda_2+\binom d2.    \tag{A-020.2}
\]
In this range $K_n(\lambda)=a_\lambda(d[n])$, so no other
orbit closure in the same representation can have a larger
degree-$d$ multiplicity at $\lambda$.

The argument refines the two differential sequences in OpenAI,
\emph{Quadratic stabilization of the canonical Foulkes--Howe map}
(September 25, 2026), Proposition 3, by keeping track of the
diagonal $GL(V)$ type. The uniform bound $b\geq a(a-1)$ from
that source remains a separate, complementary sufficient
condition. No historical-priority claim is made for the present
shape bound.

\paragraph{A differential kernel lemma.}
Let $U,Y$ be independent vector-variable blocks, and set
$D=\sum_iY_i\partial/\partial U_i$. On the polynomial space
of bidegree $(p,q)$, as a diagonal $GL(V)$ module,
\[
 \ker D=
 \begin{cases}
 S_{(q,p)}V^*,&p\leq q,\\
 0,&p>q.
 \end{cases}                                          \tag{A-020.3}
\]
The usual convention makes the first module zero if its
partition has too many rows. If additional independent vector
blocks of fixed nonnegative degrees are present, every
irreducible constituent $S_\eta V^*$ of the kernel satisfies
$\eta_2\geq p$. In particular the operator is injective on
the $\lambda$ isotypic component whenever $\lambda_2<p$,
irrespective of $q$.

\begin{proof}
The horizontal Pieri rule gives the multiplicity-free
decomposition
\[
 \operatorname{Sym}^pV^*\otimes\operatorname{Sym}^qV^*
 =\bigoplus_{i=0}^{\min(p,q)} S_{(p+q-i,i)}V^*.
\]
For dimension at least two, a highest-weight vector for the
$i$-th summand is
\[
 h_i=(U_1Y_2-U_2Y_1)^i U_1^{p-i}Y_1^{q-i}.
\]
The operator kills the displayed determinant and sends $h_i$
to
\[
 (p-i)(U_1Y_2-U_2Y_1)^i
                  U_1^{p-i-1}Y_1^{q-i+1}.
\]
This is nonzero for $i<p$, so equivariance and irreducibility
make $D$ injective on each such summand. The remaining summand
$i=p$, present exactly when $p\leq q$, is killed. This proves
(A-020.3). In dimension one the same conclusion follows from
$D(U^pY^q)=pU^{p-1}Y^{q+1}$; dimension zero is immediate.
The case $p=0$ says that the whole source is killed and agrees
with $S_{(q,0)}=\operatorname{Sym}^q$.

With spectators the kernel is the tensor product of this kernel
and their polynomial space, since tensoring vector spaces is
exact. That spectator space is a tensor product of symmetric
powers of $V^*$. Each successive horizontal Pieri step only
adds boxes to the partition $(q,p)$, so every final second row
has length at least $p$. This proves the claimed restricted
injectivity. The classical rule used here is Macdonald,
Chapter I, (5.16), printed p.~73. \end{proof}

\paragraph{Proof of (A-020.1).}
We induct on $a$. For $a=1$, pure powers span
$\operatorname{Sym}^bV$, so there are no linear equations.
For $a=2$, there are no quadratic equations when $b\geq2$
by the cited Proposition 3; when $b=1$ the Chow cone is all
of $V$. These base cases are stronger than required.

Let $a\geq3$, put $r=a-1$, $m=b-r$, and assume
$b\geq\lambda_2+\binom a2$. Then
\[
 m\geq\lambda_2+\binom r2\geq1.                       \tag{A-020.4}
\]
Suppose $T$ is in the $\lambda$ isotypic component of the
space of symmetric $a$-linear forms on $\operatorname{Sym}^bV$
whose diagonals vanish on products of $b$ vectors. Such forms
are precisely the polarizations of the degree-$a$ equations.
All constructions below are equivariant under the diagonal
$GL(V)$ action, so their images remain in type $\lambda$.

For a product $Q=y_1\cdots y_m$, the binary polynomial
$x^r+ct^r$ splits into $r$ linear factors over $\mathbb C$.
The coefficient of $c$ in the diagonal vanishing identity for
$(x^r+ct^r)Q$ is therefore
\[
                       T(x^rQ,\ldots,x^rQ,t^rQ)=0.     \tag{A-020.5}
\]
Now set $H=T(x^rQ,\ldots,x^rQ,t^b)$, keeping $x,t,y_i$
as vector-variable blocks. Apply
$D_{y_m,t}\cdots D_{y_1,t}$, where
$D_{Y,U}=\sum Y_i\partial_{U_i}$. The result is
$(b!/r!)$ times (A-020.5), hence zero. Before the $i$-th
operator the degree in $t$ is $b-i+1$, while that in the new
$y_i$ is $r$. Their difference is $m-i+1\geq1$.
The $p>q$ case of (A-020.3), with all other blocks as
spectators, makes the composition injective. Thus $H=0$.

Consider the equivariant linear map sending $T$ to the
polynomial family of symmetric $r$-linear forms
\[
 S_{x,t}(Q_1,\ldots,Q_r)
             =T(x^rQ_1,\ldots,x^rQ_r,t^b).
\]
Its codomain is
\[
 \operatorname{Sym}^{r^2}V^*\otimes\operatorname{Sym}^bV^*
             \otimes\operatorname{Sym}^r((\operatorname{Sym}^mV)^*).
\]
The identity $H=0$ says that every coefficient in $x,t$
of its last factor vanishes diagonally on products of $m$
vectors. Consequently the image lies in
\[
 \operatorname{Sym}^{r^2}V^*\otimes\operatorname{Sym}^bV^*
                         \otimes I(\operatorname{Ch}_m(V))_r.
                                                               \tag{A-020.6}
\]
By induction and (A-020.4), every constituent $S_\nu V^*$
of the last factor has
$\nu_2>m-\binom r2=b-\binom a2\geq\lambda_2$.
Multiplication by the two symmetric-power factors in (A-020.6)
is governed by two horizontal Pieri steps. It can only add boxes
to $\nu$, so cannot produce any constituent of type $\lambda$.
Since the map is equivariant and its domain is isotypic of type
$\lambda$, its image is zero. Therefore
\[
 T(x^rQ_1,\ldots,x^rQ_r,t^b)=0
       \quad\hbox{for all }Q_i\in\operatorname{Sym}^mV.       \tag{A-020.7}
\]

Finally introduce independent $v_1,\ldots,v_r,t,x$ and put
$f=T(v_1^b,\ldots,v_r^b,t^b)$. Apply $D_{x,v_i}^r$
successively for $i=1,\ldots,r$. The result is
\[
 \left(\frac{b!}{m!}\right)^r
                 T(x^rv_1^m,\ldots,x^rv_r^m,t^b)=0
\]
by (A-020.7). Before each individual derivative the source
$v_i$ has degree $p\geq b-r+1$. Our hypothesis implies
\[
 \lambda_2\leq b-\binom a2\leq b-r<p.
\]
Every intermediate polynomial remains in type $\lambda$;
the restricted injectivity part of (A-020.3) therefore applies
at every step, even when the common $x$-degree is larger than
the source degree. The composition is injective on that type,
so $f=0$. Pure powers span each input space, hence $T=0$.
This completes the induction and proves (A-020.1).

For (A-020.2), products of $n$ linear forms are determinants
of diagonal matrices of those forms. They are linear
endomorphism substitutions into $\det_n$, and belong to its
orbit closure because invertible endomorphisms are dense.
The determinant-closure ideal is therefore contained in the
Chow ideal. Its $\lambda$ component vanishes by (A-020.1).
Semisimplicity in characteristic zero then shows that every
ambient copy of $S_\lambda V^*$ survives in the quotient
coordinate ring. This is precisely the stated equality of
multiplicities. \hfill$\square$

\paragraph{Grothendieck dictionary and remaining gate.}
The Chow cone can be defined over $\mathbb Z$ as the scheme
closure of the multiplication morphism from $V^b$ to
$\operatorname{Sym}^bV$. Its homogeneous ideal is a
$GL(V)$-comodule; the operators $D_{Y,U}$ are integral,
equivariant differential maps between multigraded polynomial
modules. Representations as comodules are SGA 3, I,
Proposition 4.7.2, pp.~34--35; torus gradings are I,
Proposition 4.7.3, pp.~35--36. The symmetric and Schur
functors live in the exact $K_0$ and special $\lambda$-ring
framework of SGA 6, VI, Theorem 3.3, pp.~392--393.
The isotypic decomposition, nonzero differential coefficients,
and binary factorization in the proof are used over
$\mathbb C$; no assertion of identical kernel ranks in every
characteristic follows from the integral definitions.

Together with M-010, an actual multiplicity obstruction must
satisfy both $d(d-1)>n$ and
$\lambda_2>n-\binom d2$, in addition to the other proved
shape restrictions. The second condition is informative for
degrees with $\binom d2\leq n$, and retains all multiplicity
copies. It concerns the degree of equations and the second row,
not a saturation assertion for Kronecker support.

% END agents/archive/SHAPE_CHOW_BOUND.tex

% BEGIN agents/archive/ADAPTIVE_CHOW_BOUND.tex
% Run GCT_KRON_20261120; GPT-6 Astra Pro; effort not exposed.
\section{Adaptive differential orders for thin partition shapes}

\subsection*{A-021. A linear degree gate for a bounded second row}
For $a\geq2$ and an integer $L\geq0$ define
\[
 B_a(L)=L+\sum_{j=2}^{a-1}\min(j,L).                    \tag{A-021.1}
\]
In particular $B_a(0)=0$; for $L\geq1$ the closed form is
\[
 B_a(L)=
 \begin{cases}
 aL-\binom L2-1,&1\leq L<a,\\
 L+\binom a2-1,&L\geq a.
 \end{cases}                                          \tag{A-021.2}
\]
For every finite-dimensional complex $V$, every $a\geq2$,
$b\geq1$, and $\lambda\vdash ab$, one has
\[
 b\geq B_a(\lambda_2)
 \quad\Longrightarrow\quad
 I(\operatorname{Ch}_b(V))_a[\lambda]=0.                \tag{A-021.3}
\]
As in A-020, missing second parts are zero and $[\lambda]$
means the entire isotypic component. Therefore for the actual
degree-$d$ determinant-closure multiplicity,
\[
 n\geq B_d(\lambda_2)
 \quad\Longrightarrow\quad
                K_n(\lambda)=a_\lambda(d[n])            \tag{A-021.4}
\]
for $d\geq2$. Since $B_d(L)\leq(d-1)L$, the simpler condition
$n\geq(d-1)\lambda_2$ is sufficient. These inequalities retain
all copies, and exclude an actual multiplicity obstruction for
every competing orbit closure in the same ambient space.

This refines A-020 by varying the exponent of the split binary
form. The adaptive choice was proposed by the root worker in
this run, and its proof was completed and checked here.
The two differential sequences originate in the cited 2026
Foulkes--Howe proof; no historical-priority claim is intended.

\begin{proof}
First handle one-row types. The degree-$a$ polynomial module
$\operatorname{Sym}^a((\operatorname{Sym}^b V)^*)$ contains
exactly one copy of $S_{(ab)}V^*$: the top-weight monomial is
the $a$-th power of the coefficient of the pure $b$-th power.
Restriction to pure powers $v^b$ is nonzero on this copy, since
that monomial evaluates to $\ell(v)^{ab}$ for a nonzero linear
form $\ell$. Equivariance and irreducibility imply that the
one-row copy has zero intersection with the Chow ideal. This
proves (A-021.3) for $\lambda_2=0$, in every dimension where
the type occurs.

Assume henceforth $L=\lambda_2\geq1$. The functions $B_a(L)$
are nondecreasing in $L$, satisfy $B_2(L)=L$, and obey
\[
 B_a(L)=B_{a-1}(L)+\min(a-1,L)\quad(a\geq3),
 \qquad B_a(L)\geq L.                                 \tag{A-021.5}
\]
We induct on $a$. At $a=2$ there are no quadratic Chow
equations for any $b\geq1$: for $b\geq2$ this is the cited
Foulkes--Howe Proposition 3, and for $b=1$ the Chow cone is
the whole vector space. Also any partition of $2b$ satisfies
$\lambda_2\leq b$, consistently with the claimed base bound.

Let $a\geq3$, put
\[
 r=a-1,\qquad s=\min(r,L),\qquad m=b-s.
\]
Then $s\geq1$ and the hypothesis gives
\[
                   m\geq B_r(L)\geq L\geq1.            \tag{A-021.6}
\]
Take $T$ in the full $\lambda$ isotypic component of the
symmetric $a$-linear forms polarizing the Chow ideal.
For $Q=y_1\cdots y_m$, binary factorization of
$x^s+ct^s$ and extraction of the coefficient of $c$ give
\[
                     T(x^sQ,\ldots,x^sQ,t^sQ)=0.       \tag{A-021.7}
\]
Keep all $x,t,y_i$ as independent vector-variable blocks, and
set $H=T(x^sQ,\ldots,x^sQ,t^b)$. Apply successively
$D_{y_i,t}=\sum_j(y_i)_j\partial/\partial t_j$ for
$i=1,\ldots,m$. Their composition is $b!/s!$ times
(A-021.7), so is zero. Before the $i$-th derivative, the
source $t$ has degree $p=b-i+1\geq s+1$, whereas the new
destination $y_i$ has degree $q=r$.

If $s=r$, then $p>q$, giving injectivity by the differential
kernel lemma (A-020.3). If $s=L$, then $p>L$; the same lemma
gives injectivity on type $\lambda$ even if $p\leq q$.
These cases cover $s=\min(r,L)$, including equality. Each
intermediate polynomial has only type $\lambda$, since both
evaluation and differentiation are equivariant. Thus every
derivative is injective on the image under consideration, and
$H=0$.

The equivariant map
\[
 T\longmapsto\bigl[(x,t,Q_1,\ldots,Q_r)\longmapsto
                         T(x^sQ_1,\ldots,x^sQ_r,t^b)\bigr]
\]
therefore has image in
\[
 \operatorname{Sym}^{rs}V^*\otimes\operatorname{Sym}^bV^*
                         \otimes I(\operatorname{Ch}_m(V))_r.
                                                               \tag{A-021.8}
\]
Indeed its diagonal in the $Q_i$ vanishes on every product,
and expansion in the independent $x,t$ coordinates puts every
coefficient in that homogeneous Chow ideal. By induction,
(A-021.6), and monotonicity of $B_r$, the last factor contains
no type $\nu$ with $\nu_2\leq L$: for $\nu_2>0$ one has
$m\geq B_r(L)\geq B_r(\nu_2)$, while the one-row case
was already proved. Two horizontal Pieri steps for the two
symmetric-power factors only add boxes. Hence (A-021.8) has
no type $\lambda$, whose second row has length $L$.
Equivariance forces the image of our $T$ to be zero. We have
obtained
\[
 T(x^sQ_1,\ldots,x^sQ_r,t^b)=0
                  \quad(Q_i\in\operatorname{Sym}^mV). \tag{A-021.9}
\]

Now put $f=T(v_1^b,\ldots,v_r^b,t^b)$ and introduce the
independent block $x$. Applying $D_{x,v_i}^s$ successively
for $i=1,\ldots,r$ gives
\[
 \left(\frac{b!}{m!}\right)^r
                  T(x^sv_1^m,\ldots,x^sv_r^m,t^b)=0
\]
by (A-021.9). Before any individual derivative the source
$v_i$ has degree at least $b-s+1=m+1>L$ by (A-021.6).
Thus (A-020.3) makes each derivative injective on type
$\lambda$, regardless of the growing degree in $x$.
The composition is injective on the image of $T$, so $f=0$.
As pure powers span every input space, this implies $T=0$.
The induction proves (A-021.3).

The closed form (A-021.2) is the sum in (A-021.1): when
$L<a$, sum $1+\cdots+L$ followed by $a-1-L$ copies of $L$,
remove the $j=1$ term and then add the leading $L$.
When $L\geq a$, all $j$ in the sum are at most $L$.
The bound $B_a(L)\leq(a-1)L$ follows because its $a-1$
terms are each at most $L$ (and also holds for $L=0$).
Finally, Chow products are diagonal determinants, so the
ideal containment and semisimple quotient argument in A-020
give (A-021.4), retaining every ambient multiplicity copy.
\end{proof}

\paragraph{Meaning of the strengthened gate.}
For $1\leq\lambda_2=L<d$, the sufficient bound is
\[
                          n\geq dL-\binom L2-1.
\]
For a fixed second row this grows linearly in $d$, whereas
the uniform Foulkes--Howe condition grows quadratically.
Combining the two verified bounds, an actual multiplicity
obstruction must satisfy
\[
        d(d-1)>n,\qquad B_d(\lambda_2)>n,
        \quad\hbox{and hence}\quad (d-1)\lambda_2>n.    \tag{A-021.10}
\]
All previous padding, row-number, and ambient-multiplicity
conditions remain applicable. This is a restriction on actual
orbit-closure multiplicities, without asserting support
saturation for Kronecker coefficients.

The scheme and comodule interpretation is exactly A-020's:
the maps in (A-021.8) are equivariant maps of Schur and
symmetric functors on locally free modules, the differential
operators are integral, and the proof is carried out after
base change to $\mathbb C$. The exact Grothendieck locators
are SGA 3, I, Propositions 4.7.2--4.7.3, pp.~34--36, and
SGA 6, VI, Theorem 3.3, pp.~392--393. The use of polynomial
spectators, rather than fixed choices of $x,t$, is what
preserves the full comodule type throughout the argument.

% END agents/archive/ADAPTIVE_CHOW_BOUND.tex

% BEGIN agents/archive/OPTIMIZED_CHOW_BOUND.tex
% Run GCT_KRON_20261120; GPT-6 Astra Pro; effort not exposed.
\section{An optimized differential recurrence and a general padded-degree gate}

\subsection*{A-022. A uniform degree restriction for actual multiplicity obstructions}
For $L\geq1$ define integers $N_a(L)$, $a\geq2$, by
\[
 \begin{gathered}
 N_2(L)=1,\qquad s=\min(a-1,L),\\
 N_a(L)=s+\max\{N_{a-1}(L),\min((a-1)s-1,L)\}
                       \quad(a\geq3).
 \end{gathered}                                      \tag{A-022.1}
\]
They are nondecreasing in $L$. For every complex vector space
$V$, every $b\geq1$, and every $\lambda\vdash ab$ with
$\lambda_2\leq L$,
\[
 b\geq N_a(L)
 \quad\Longrightarrow\quad
 I(\operatorname{Ch}_b(V))_a[\lambda]=0.                \tag{A-022.2}
\]
The bounds satisfy
\[
 N_a(L)\leq B_a(L),\qquad
 N_a(L)\leq a(a-1)-1,                                 \tag{A-022.3}
\]
where $B_a$ is A-021's bound. This recurrence uses both
possible injectivity tests at the last differential step. It
is an optimization within those degree and second-row tests,
not an assertion of the least possible Chow stabilization bound.

For the GCT application let $n\geq m\geq1$, $d\geq2$.
If
\[
                              n\geq N_d(dm),           \tag{A-022.4}
\]
then $q^m_\lambda(d[n])\leq K_n(\lambda)$ for every partition
$\lambda\vdash nd$. More precisely, every type with
$q^m_\lambda(d[n])>0$ has full ambient determinant multiplicity.
Consequently an actual multiplicity obstruction requires
\[
 n<N_d(dm)\leq
 \min\left\{d(d-1)-1,\;dm+\binom d2-1\right\}.        \tag{A-022.5}
\]
In particular the broad condition
$n\geq dm+\binom d2-1$ already excludes all actual
multiplicity obstructions in that degree, by A-021 alone.

There is a closed finite maximum for the sharper bound needed
in (A-022.4): since $L=dm\geq d$,
\[
 \boxed{\quad N_d(L)=\binom d2+
 \max\left(0,\max_{2\leq j\leq d-1}
      \left\{\min(j^2-1,L)-\binom j2\right\}\right).\quad}
                                                               \tag{A-022.6}
\]
The empty inner maximum gives $N_2(L)=1$. For $d\geq3$
only the two indices
$\min(d-1,\lfloor\sqrt{L+1}\rfloor)$ and
$\min(d-1,\lfloor\sqrt{L+1}\rfloor+1)$ need be inspected.
An explicit sufficient bound, slightly weaker than the exact
maximum, is
\[
 n\geq \binom d2+
 \left\lfloor\frac{dm-1+\lfloor\sqrt{dm+1}\rfloor}{2}
                                                    \right\rfloor.
                                                               \tag{A-022.7}
\]
Thus the exact degree threshold can be computed using only
integer arithmetic and two candidates.

\begin{proof}[Proof of the recurrence]
The two applicable cases of A-020's differential kernel lemma
say that $D_{Y,U}$ is injective on type $\lambda$ whenever
the source degree $p$ satisfies
\[
                            p>\min(q,\lambda_2),       \tag{A-022.8}
\]
where $q$ is the destination degree. Indeed $p>q$ gives
injectivity on the entire space, while $p>\lambda_2$ gives
injectivity on the specified type with all spectators.

We prove (A-022.2) for each fixed upper bound $L$, by induction
on $a$. One-row types have no Chow equations by A-021's
one-row argument. At $a=2$ there are no quadratic equations
for any $b\geq1$, again as in A-021. Now set $r=a-1$,
$s=\min(r,L)$ and $m=b-s$. The recurrence ensures
\[
       m\geq N_r(L)\geq1,
       \qquad m\geq\min(rs-1,L).                      \tag{A-022.9}
\]
For $T$ of type $\lambda$, use the same first sequence as
in A-021, starting from
$H=T(x^sQ,\ldots,x^sQ,t^b)$, $Q=y_1\cdots y_m$.
Its iterated derivative is $b!/s!$ times
$T(x^sQ,\ldots,x^sQ,t^sQ)=0$, from binary factorization.
At every individual step the source degree is at least
$s+1$, and the untouched destination has degree $r$.
If $s=r$, it is injective by $p>q$. If $s=L$, it is
injective by $p>L\geq\lambda_2$. The images retain type
$\lambda$, so $H=0$.

Keeping $x,t$ as polynomial spectators gives the equivariant
image in
\[
 \operatorname{Sym}^{rs}V^*\otimes\operatorname{Sym}^bV^*
                              \otimes I(\operatorname{Ch}_m(V))_r.
\]
The coefficientwise ideal argument is the one proved in A-021.
By $m\geq N_r(L)$ and the induction hypothesis, the last
factor contains no type $\nu$ with $\nu_2\leq L$.
Two horizontal Pieri steps cannot decrease its second row.
Hence the image has no type $\lambda$, and
$T(x^sQ_1,\ldots,x^sQ_r,t^b)=0$ for independent $Q_i$.

Apply the second sequence $D_{x,v_i}^s$, $i=1,\ldots,r$,
to $T(v_1^b,\ldots,v_r^b,t^b)$. The resulting polynomial
is zero by the preceding identity. Before every individual
derivative, $p\geq m+1$ and $q\leq rs-1$.
By (A-022.9) and $\lambda_2\leq L$,
\[
 p>\min(rs-1,L)\geq\min(q,\lambda_2).
\]
Thus (A-022.8) makes every step injective on the relevant
type, and $T$ vanishes on all tuples of pure powers. Spanning
and multilinearity give $T=0$, proving (A-022.2).

For monotonicity, $s=\min(a-1,L)$ and
$\min((a-1)s-1,L)$ are nondecreasing in $L$; induction
in (A-022.1) proves the assertion. Also
$\min((a-1)s-1,L)\leq L\leq B_{a-1}(L)$ and
$N_{a-1}(L)\leq B_{a-1}(L)$ by induction. Adding $s$
therefore gives $N_a(L)\leq B_a(L)$. Finally $s\leq r$
and $\min(rs-1,L)\leq r^2-1$, while
$N_r(L)\leq r(r-1)-1\leq r^2-1$ by induction.
Thus $N_a(L)\leq r+r^2-1=a(a-1)-1$, proving (A-022.3).
\end{proof}

\begin{proof}[Proof of the padded-degree consequence]
We include the needed padding weight restriction. The orbit
closure of $\ell^{n-m}\operatorname{per}_m$ is contained in the
closure of all forms $\ell^{n-m}f$, with $f$ any degree-$m$
form. Pulling back a degree-$d$ polynomial along this
multiplication morphism embeds the degree-$d$ coordinate
module of that larger family into
\[
          \operatorname{Sym}^{d(n-m)}V\otimes
                         \operatorname{Sym}^d(\operatorname{Sym}^mV).
\]
The first factor records the coefficients of $\ell$, and the
second those of $f$. If a type $\lambda$ occurs on the
permanent orbit closure, it must occur in this module.
Decompose the second factor into $S_\nu V$ with $|\nu|=dm$.
Horizontal Pieri says that $\lambda/\nu$ is a horizontal
$d(n-m)$ strip. In particular $\lambda_2\leq\nu_1\leq dm$.
Also, a horizontal strip has at most one added box per column,
so $\lambda_1\geq d(n-m)$ and $|\bar\lambda|\leq dm$.
This is a statement about every
positive permanent-side multiplicity, not only a particular
highest-weight vector.

Now assume $n\geq N_d(dm)$. Apply (A-022.2) with $L=dm$
to every type with $q^m_\lambda>0$. Its Chow-ideal component
is zero. Since products of $n$ linear forms are diagonal
determinants, their cone lies in the determinant orbit closure;
therefore every such determinant-ideal component is zero as
well. This gives $K_n(\lambda)=a_\lambda(d[n])\geq q^m_\lambda$.
For $q^m_\lambda=0$ the comparison is automatic. Thus no
actual multiplicity obstruction occurs in this degree.
For $L=dm\geq d$, (A-021.2) reads
$B_d(dm)=dm+\binom d2-1$, and (A-022.3) gives (A-022.5).

The simpler degree restriction can be written without $B,N$:
every actual obstruction has
\[
 d>
 \frac{-(2m-1)+\sqrt{(2m-1)^2+8(n+1)}}{2}.              \tag{A-022.10}
\]
Indeed $n<dm+d(d-1)/2-1$ is exactly the strict quadratic
inequality whose positive root is displayed. This is a
necessary condition; the exact recurrence is stronger.
\end{proof}

\begin{proof}[Proof of the finite maximum]
When $L\geq d$, every $s$ occurring for $a\leq d$ is $a-1$.
Unroll (A-022.1) from $N_2=1$. The sum of the additive
increments $2+\cdots+(d-1)$ is $\binom d2-1$.
An intermediate term $\min(j^2-1,L)$ is inserted just
before adding $j$, so the increments before its insertion
sum to $\binom j2-1$. Taking the maximum of the initial
term and all inserted terms yields precisely (A-022.6).

Put $k=\lfloor\sqrt{L+1}\rfloor$. For integer $j\leq k$,
the summand in braces is $\binom{j+1}{2}-1$, increasing
with $j$. For $j\geq k+1$ it is $L-\binom j2$, decreasing
with $j$. Hence the maximum on $2\leq j\leq d-1$ occurs
at the two stated, clipped indices. The case $d=2$ was
already specified.

For completeness, the corresponding real function of $j$
has its maximum at $j=\sqrt{L+1}$, with value
$(L-1+\sqrt{L+1})/2$. Its integer maximum is thus at most
$\lfloor(L-1+\sqrt{L+1})/2\rfloor$, which equals
$\lfloor(L-1+\lfloor\sqrt{L+1}\rfloor)/2\rfloor$.
The equality follows by separating the integer and fractional
parts of the square root: its fractional part divided by two
cannot cross the next integer from an integer or half-integer.
Adding $\binom d2$ proves (A-022.7).
\end{proof}

\paragraph{Scope and geometric dictionary.}
This is a restriction on actual orbit-closure evaluation ranks.
For $m$ small compared with $d$, the sufficient expression
(A-022.7) has leading terms $\binom d2+dm/2$, with its
displayed square-root correction. The exact expressions (A-022.6)--(A-022.7)
are used for conclusions. No conclusion about all higher
degrees, global Kronecker saturation, or $\mathsf P$ versus
$\mathsf{NP}$ follows from this reduction alone.

The Chow and padding multiplication maps are morphisms of
affine schemes over $\mathbb Z$, and their homogeneous
coordinate modules are comodules as in SGA 3, I,
Propositions 4.7.2--4.7.3, pp.~34--36. Pieri is Macdonald,
I, (5.16), p.~73. The symmetric and Schur functors are
interpreted in the exact $K_0$ framework of SGA 6, VI,
Theorem 3.3, pp.~392--393. The proof uses the reduced
characteristic-zero fibers and preserves all isotypic
multiplicity copies; integral definitions alone do not
assert the same kernel statement in every characteristic.

% END agents/archive/OPTIMIZED_CHOW_BOUND.tex

% BEGIN agents/archive/GROUPED_CHOW_BOUND.tex
% Run GCT_KRON_20261120; GPT-6 Astra Pro; effort not exposed.
\section{Grouped derivatives give a cleaner universal degree gate}

\subsection*{A-023. Improving the recurrence by treating each derivative power at once}
For $L\geq1$ and $a\geq2$, define
\[
 P_2(L)=1,\qquad s=\min(a-1,L),\qquad
 P_a(L)=s+\max\{P_{a-1}(L),\min((a-2)s,L)\}
                        \quad(a\geq3).                \tag{A-023.1}
\]
Then for every $\lambda\vdash ab$ with $\lambda_2\leq L$,
\[
 b\geq P_a(L)
 \quad\Longrightarrow\quad
 I(\operatorname{Ch}_b(V))_a[\lambda]=0                 \tag{A-023.2}
\]
over every finite-dimensional complex $V$. Moreover
\[
           P_a(L)\leq N_a(L),\qquad P_a(L)\leq(a-1)^2. \tag{A-023.3}
\]
In particular $b\geq(a-1)^2$ gives the uniform vanishing
in every type. These are sufficient bounds, not claims of
optimal stabilization.

For the padded permanent comparison with $n\geq m\geq1$
and outer degree $d\geq2$, there is no actual degree-$d$
multiplicity obstruction whenever
\[
                              n\geq P_d(dm).           \tag{A-023.4}
\]
For $L\geq d$, the exact threshold is
\[
 \boxed{\quad P_d(L)=\binom d2+
       \max\left(0,\max_{2\leq j\leq d-1}
           \min\left\{\binom j2,L-\binom j2\right\}\right).\quad}
                                                               \tag{A-023.5}
\]
At $d=2$ the maximum is zero. Since the inner minimum is at
most $\lfloor L/2\rfloor$, a simple sufficient condition is
\[
       \boxed{\quad n\geq
              \left\lfloor\frac{d(d+m-1)}2\right\rfloor.\quad}
                                                               \tag{A-023.6}
\]
The exact maximum can improve this further. Thus an actual
multiplicity obstruction necessarily satisfies
\[
 n<P_d(dm)\leq
 \min\left\{(d-1)^2,
             \left\lfloor\frac{d(d+m-1)}2\right\rfloor\right\}.
                                                               \tag{A-023.7}
\]

\paragraph{The power of a differential operator.}
Let $D=D_{Y,U}$ on polynomial functions homogeneous of degrees
$p,q$ in independent vector blocks $U,Y$, and let $e$ be an
integer with $1\leq e\leq p$. With arbitrary polynomial
spectators of fixed degrees, $D^e$ is injective on type
$\lambda$ whenever
\[
                         p-e\geq\min(q,\lambda_2).     \tag{A-023.8}
\]
This is stronger than requiring every individual derivative to
be injective on its entire intermediate homogeneous space.

\begin{proof}
Use A-020's multiplicity-free Pieri decomposition of the two
blocks. On the summand $S_{(p+q-i,i)}V^*$, its highest-weight
vector is
$h_i=(U_1Y_2-U_2Y_1)^iU_1^{p-i}Y_1^{q-i}$.
Applying $D^e$ multiplies by the falling factorial
$(p-i)(p-i-1)\cdots(p-i-e+1)$ and changes the exponents
of $U_1,Y_1$ by $-e,+e$. It is nonzero exactly for
$i\leq p-e$. Thus the kernel is the sum of the summands
with $p-e<i\leq\min(p,q)$.

If $q\leq p-e$, this kernel is zero. Otherwise every
constituent of the kernel has second row at least $p-e+1$;
tensoring the spectator symmetric powers can only add boxes
by horizontal Pieri. It therefore has no type $\lambda$
if $\lambda_2\leq p-e$. These two cases prove (A-023.8).
The one-dimensional case follows from the same monomial
derivative, and types not fitting the dimension are zero.
\end{proof}

\begin{proof}[Proof of the refined recurrence]
The base $a=2$ and one-row types are the same as in A-021.
Fix an upper bound $L$ on the second row, let $r=a-1$,
$s=\min(r,L)$, and $m=b-s$. By (A-023.1),
\[
                m\geq P_r(L)\geq1,
                \qquad m\geq\min((r-1)s,L).            \tag{A-023.9}
\]
Let $T$ be in the $\lambda$ isotypic component of symmetric
$a$-linear forms polarizing the degree-$a$ Chow ideal.
The binary-factorization identity and the first differential
sequence of A-021 apply without change: at each step the
source degree is at least $s+1$, while the untouched
destination degree is $r$; either $s=r$ gives $p>q$, or
$s=L$ gives $p>L\geq\lambda_2$. Hence
$T(x^sQ,\ldots,x^sQ,t^b)=0$.

Retaining the $x,t$ polynomial variables gives the image in
$\operatorname{Sym}^{rs}V^*\otimes\operatorname{Sym}^bV^*
\otimes I(\operatorname{Ch}_m(V))_r$.
The induction hypothesis kills all its last-factor types
whose second row is at most $L$. The two symmetric-power
spectators cannot lower the second row, so the entire
$\lambda$ isotypic image is zero. Consequently
\[
 T(x^sQ_1,\ldots,x^sQ_r,t^b)=0
                   \quad(Q_i\in\operatorname{Sym}^mV). \tag{A-023.10}
\]

Now start from $f=T(v_1^b,\ldots,v_r^b,t^b)$, and apply
the complete powers $D_{x,v_i}^{\,s}$ for $i=1,\ldots,r$.
Before the $i$-th power, its source block has degree $p=b$
and the common destination $x$ has degree $q=(i-1)s$.
The earlier powers affected only other source blocks.
By (A-023.9),
\[
 p-s=m\geq\min((r-1)s,L)\geq\min(q,\lambda_2).
\]
The grouped injectivity criterion (A-023.8) therefore applies
to each complete power on type $\lambda$. Their composition
is injective, while its value is
\[
 \left(\frac{b!}{m!}\right)^r
                   T(x^sv_1^m,\ldots,x^sv_r^m,t^b)=0
\]
by (A-023.10). Thus $f=0$, and pure-power spanning yields
$T=0$. This proves (A-023.2), retaining all multiplicity copies.

The new second term $\min((r-1)s,L)$ is at most
$\min(rs-1,L)$, since $s\geq1$. Induction therefore
gives $P_a(L)\leq N_a(L)$. For the uniform bound, use
$s\leq r$, $P_r(L)\leq(r-1)^2$, and
$\min((r-1)s,L)\leq r(r-1)$ in (A-023.1). This gives
$P_a(L)\leq r+r(r-1)=r^2=(a-1)^2$.
The base $P_2=1$ verifies both inductions.
\end{proof}

\begin{proof}[Proof of the explicit degree gate]
The padding morphism and Pieri argument in A-022 show that
every type occurring on the padded permanent has
$\lambda_2\leq dm$. Apply (A-023.2) with $a=d$, $b=n$,
and $L=dm$. If $n\geq P_d(dm)$ its Chow-ideal component
vanishes. Chow products are diagonal determinants, so the
determinant-ideal component also vanishes, and its actual
multiplicity equals the ambient multiplicity. This proves
(A-023.4) and the resulting comparison with every
permanent-side copy.

If $L\geq d$, every $s$ in the recursion is $a-1$.
The same unrolling as in A-022 gives
\[
 P_d(L)=\binom d2+
 \max\left(0,\max_{2\leq j\leq d-1}
       \left\{\min(j(j-1),L)-\binom j2\right\}\right).
\]
Writing $t=\binom j2$ turns the expression in braces
into $\min(2t,L)-t=\min(t,L-t)$, proving (A-023.5).
Its integer value is at most $\lfloor L/2\rfloor$;
the outer maximum with zero obeys the same bound for $L\geq1$.
Thus
\[
 P_d(dm)\leq\binom d2+\left\lfloor\frac{dm}{2}\right\rfloor
              =\left\lfloor\frac{d(d+m-1)}2\right\rfloor.
\]
This proves (A-023.6)--(A-023.7).

For exact evaluation at $d\geq3$, let
\[
 k=\left\lfloor\frac{1+\sqrt{1+4L}}2\right\rfloor,
\]
the largest integer satisfying $k(k-1)\leq L$.
For $j\leq k$ the summand $\min(\binom j2,L-\binom j2)$
is $\binom j2$ and increases. For $j\geq k+1$ it is
$L-\binom j2$ and decreases. Hence it is enough to inspect
$\min(d-1,k)$ and $\min(d-1,k+1)$. As $L\geq d\geq3$,
both clipped indices are at least two. The integer $k$ can
equivalently be computed as
$(1+\operatorname{isqrt}(1+4L))\mathbin{//}2$, where
\(\operatorname{isqrt}\) and \(//\) denote integer square
root and integer division. This gives an exact computation,
with no numerical root approximation. \end{proof}

\paragraph{Why grouping changes the bound.}
The power $D^s$ need only be injective on its actual source.
Some intermediate homogeneous spaces contain extra irreducible
summands that are never reached from that source. Requiring
injectivity on all of those larger spaces introduced the
extra $s-1$ in A-022's final derivative test. The explicit
Pieri calculation removes that loss. The proof still uses
additional structure beyond an effective $K_0$ class: the
equivariant evaluation and differential maps, their kernels,
and preservation of all multiplicity copies.

These are morphisms of polynomial comodules in the language
of SGA 3, I, Propositions 4.7.2--4.7.3, pp.~34--36, and
symmetric/Schur operations in SGA 6, VI, Theorem 3.3,
pp.~392--393. The kernel calculation uses horizontal Pieri,
Macdonald, I, (5.16), p.~73. As in A-020--A-022, the
underlying maps are integral while the irreducible and
nonzero-coefficient argument is made in characteristic zero.
The improvement concerns the actual multiplicity comparison;
it asserts no Kronecker saturation or conclusion about
$\mathsf P$ versus $\mathsf{NP}$.

% END agents/archive/GROUPED_CHOW_BOUND.tex

% BEGIN agents/formulas/further_gates.tex
\section{Permanent-specific containment and the remaining rank comparison}

\subsection*{Finding F-009: a permanent-size gate valid for every partition}
The following uses a known upper bound.  Grenet,
\emph{An upper bound for the permanent versus determinant problem},
25 May 2012 version, unnumbered Lemma on p.~2 and unnumbered
Theorem on p.~3, constructs an affine determinantal representation
of size \(2^m-1\).  We record the homogeneous identity and its
coordinate-ring consequence explicitly.

Let \(m\geq2\), put \(s=2^m-1\), and first let \(L\) be an
auxiliary indeterminate.  There is an \(s\times s\) matrix of homogeneous
linear forms over \(\mathbb Z\) such that
\begin{equation}\label{eq:F009}
                    \det M=L^{s-m}\operatorname{perm}_m(X).
\end{equation}
Consequently, for every \(n\geq2^m-1\), the entire padded-permanent
orbit closure lies in the determinant orbit closure:
\[
                         Z_{n,m}\subseteq\Omega_n.
\]
In particular \(q^m_\lambda(d[n])\leq K_n(\lambda)\) for every
\(d\) and every partition \(\lambda\vdash nd\), without a row bound.

\begin{proof}[Proof of the homogeneous form and consequence]
Use one vertex for each subset of \([m]\), identifying the empty
and full subsets as a central vertex.  There are \(s\) vertices.
For \(S\subsetneq[m]\) and \(j\notin S\), put an edge
\(S\to S\cup\{j\}\), of weight \(x_{|S|+1,j}\).
Multiply the weights of all edges leaving the central vertex
by \((-1)^{m-1}\).  Put a loop of weight \(L\) at every
noncentral vertex and no loop at the center.

Every non-loop cycle passes through the center.  A cycle cover
therefore chooses one path from the empty set to the full set
in the original subset lattice, and loops at all unused vertices.
Such a path adds the elements in the order of a unique permutation
\(\sigma\in S_m\), and its variable product is
\(\prod_{i=1}^m x_{i,\sigma(i)}\).
The merged path is a cycle of length \(m\), whose permutation
sign is \((-1)^{m-1}\).  The sign placed on its first edge
cancels that sign.  The \(s-m\) unused loops supply
\(L^{s-m}\).  The permutation expansion of the adjacency-matrix
determinant thus gives \eqref{eq:F009}, including its signs.
This is an integral polynomial identity.

Appending the block \(L I_{n-s}\) gives a matrix of size \(n\)
with determinant \(L^{n-m}\operatorname{perm}_m(X)\).
Thus the padded permanent itself is an endomorphism substitution
into \(\det_n\), and lies in \(\Omega_n\).  The latter is closed
and \(\mathrm{GL}(V)\)-stable, so it contains the whole
padded-permanent orbit closure.  Restriction of coordinate rings
is a surjection, and characteristic-zero semisimplicity gives
the inequality of every isotypic multiplicity.
\end{proof}

For \(m=1\) a diagonal matrix gives the assertion directly.
The identity remains valid after every substitution for \(L\),
including \(L=0\); when \(m\geq2\), \(s>m\), so both sides
then vanish. In particular one may put \(L=x_{11}\), the internal
padding convention of the supplied IP and BIP papers and of
\(q^m_\lambda(d[n])\) here. The construction is a morphism over the integers
from coefficient data to linear matrix pencils.  Its characteristic-zero
image containment and the induced coherent-algebra quotient are
the scheme and \(K_0\) meanings of the multiplicity statement.
The finite subset lattice describes the vertices, but its
permanent-specific edge labels are additional data; the containment
is not inferred from a formal lambda identity alone.

\paragraph{Exact small-size consequences in the two stated GCT ranges.}
For integer \(m\geq2\),
\[
 2^m-1\leq3m^4+1\quad(2\leq m\leq18),\qquad
 2^m-1\leq m^{25}\quad(2\leq m\leq189).
\]
Therefore an actual multiplicity obstruction in the range
\(n>3m^4\) must have \(m\geq19\); one in the range
\(n\geq m^{25}\) must have \(m\geq190\).
These bounds are combined with, and independent of, the row and
degree exclusions in M-008.  They are uses of a known determinantal
upper bound, not lower bounds on the complexity of the permanent.

Here is a direct verification that the cutoffs are exact for this
particular comparison.  The function
\(x\log2-c\log x\) is convex for \(x>0\).
At the endpoints \(m=2,18\), one has \(2^m\leq3m^4\);
convexity gives the inequality throughout that interval.
Likewise \(2^2\leq2^{25}\) and the exact integer comparison
\(2^{189}\leq189^{25}\) imply the second inequality
throughout \([2,189]\).
At the next integers, exact arithmetic gives
\[
 2^{19}-1=524287>390964=3\cdot19^4+1,
 \qquad 2^{190}-1>190^{25}.
\]
The ratios \((1+1/m)^4\) and \((1+1/m)^{25}\) decrease with
\(m\), and exact integer checks give
\(20^4<2\cdot19^4\) and \(191^{25}<2\cdot190^{25}\).
It follows inductively that the respective failures persist for
all larger integers, even with the displayed additive constants.
The complete integer comparisons and small determinant checks
are recorded in \texttt{permanent\_gate\_certificate.json}.

\subsection*{Finding F-010: the exact multiplicity-space datum left to compare}
Let \(V\) be a complex vector space, \(f\in\operatorname{Sym}^nV^*\),
and \(X_f=\overline{\mathrm{GL}(V)\cdot f}\).
For degree \(d\), use the natural decomposition
\[
 F_d(V)=\operatorname{Sym}^d\operatorname{Sym}^nV
       =\bigoplus_{\lambda\vdash nd}S_\lambda V\otimes M_\lambda,
 \qquad M_\lambda=[\lambda]^{S_n\wr S_d}.
\]
The \(\lambda\)-component of the evaluation tensor \(f^{\otimes d}\)
is an element of
\(S_\lambda V^*\otimes M_\lambda^*\), and therefore a linear map
\[
                    T_{f,\lambda}:M_\lambda\longrightarrow S_\lambda V^*.
\]
Then
\begin{equation}\label{eq:F010}
 \operatorname{mult}_{S_\lambda V}\mathbb C[X_f]_d
                  =\operatorname{rank}T_{f,\lambda}.
\end{equation}
In particular the remaining actual GCT comparison is exactly
\[
 \operatorname{rank}T_{L^{n-m}\operatorname{perm}_m,\lambda}
                 \ \mathop{\leq}^{?}\
 \operatorname{rank}T_{\det_n,\lambda}.
\]
It is a comparison of positive ranks, not merely the question whether
either map is zero. Here \(L=x_{11}\) for the stated GCT convention;
the rank characterization itself applies to every form and every
choice of padding.

\begin{proof}
Write
\(I(X_f)_d=\bigoplus_\lambda S_\lambda V\otimes J_{f,\lambda}\).
For \(a\in M_\lambda\), the condition
\(T_{f,\lambda}(a)=0\) says exactly that every polynomial
\(s\otimes a\), \(s\in S_\lambda V\), vanishes at \(f\).
Since this entire copy of \(S_\lambda V\) is stable under the
general linear group, it then vanishes on the entire orbit:
evaluation at \(g f\) is evaluation of \(g^{-1}(s\otimes a)\)
at \(f\).  It vanishes on the closure as well.
Conversely, a copy contained in the ideal vanishes at \(f\),
so its multiplicity vector is in \(\ker T_{f,\lambda}\).
Thus \(J_{f,\lambda}=\ker T_{f,\lambda}\).  Subtraction from
\(\dim M_\lambda\) proves \eqref{eq:F010}.
\end{proof}

The Young-set and wreath-invariant constructions determine the
ambient space \(M_\lambda\) and its dimension.  They do not
determine the map \(T_{f,\lambda}\), which depends on the
specific tensor \(f^{\otimes d}\).  One sufficient additional
statement would be a factorization
\[
 T_{L^{n-m}\operatorname{perm}_m,\lambda}
                   =A_\lambda T_{\det_n,\lambda}.
\]
This is equivalent to inclusion of the corresponding kernels
and is stronger than a rank inequality alone.  M-006 supplies
precisely such kernel inclusion in its stated projection range;
no such inclusion is established in general.

\paragraph{A small model of the missing information.}
The distinction is already visible for the standard representation
\(S=\mathbb C^2\) of \(\mathrm{GL}_2\), with
\(W=S\otimes\mathbb C^2\).  Regard elements of \(W\) as
two-column matrices.  The orbit closure of
\(v_1=(e_1,0)\) is the linear subspace \(S\otimes\mathbb C e_1\),
whereas the orbit of \(v_2=(e_1,e_2)\) is the invertible matrices
and its closure is all of \(W\).  In their degree-one coordinate
rings, the multiplicity of \(S^*\) is respectively one and two.
It is positive in both.  The ambient representation and all its
lambda operations are identical in the two comparisons.
The asserted orbit closures follow directly from transitivity
on nonzero vectors and density of invertible matrices, so this
example needs no external orbit classification.

This example refutes an inference from the ambient lambda object
and occurrence alone to equality, or to a chosen inequality, of
orbit-closure multiplicities.  It is not a counterexample to any
specific permanent-versus-determinant inequality.
Geometrically \(T_{f,\lambda}\) is a fiberwise linear map on a
Schur isotypic multiplicity bundle, whose minors cut out its
rank loci.  Projectors and the Schur decomposition are used in
characteristic zero, as in the polynomial-functor dictionary of
Macdonald, I, Appendix~A, (5.3)--(5.6), pp.~155--156.
The specific evaluation tensor, or equivalent equations for
these rank loci, is the additional datum not supplied by the
class in \(K_0\).

\subsection*{Finding F-011: a complete scalar-multiplicity family}
Let \(d\geq2\), \(n\geq1\), and \(0\leq k\leq n\). Delete zero
parts when forming the partition
\[
 \lambda_{d,k}=
       \bigl(d(n-k)+k,\underbrace{k,\ldots,k}_{d-1}\bigr)
                  \vdash nd.
\]
Then the ambient plethysm coefficient is
\begin{equation}\label{eq:F011}
 a_{\lambda_{d,k}}(d[n])=
 \begin{cases}1&k\text{ even},\\0&k\text{ odd}.
 \end{cases}
\end{equation}
As usual this is the stable coefficient of the Schur function;
when applying it to a fixed vector space, the partition must fit
its dimension. For \(d=1\), the partition is just \((n)\) and the
coefficient is one, independently of the redundant parameter \(k\).
The proof is a direct Cauchy-formula calculation; no historical
priority is claimed for the formula.

\begin{proof}
Put \(E=\mathbb C^d\) and take \(V\) of dimension at least \(d\).
In \(\operatorname{Sym}(V\otimes E)\), the subspace of column
multidegree \((n,\ldots,n)\) is naturally
\((\operatorname{Sym}^n V)^{\otimes d}\).
The permutation matrices in \(\mathrm{GL}(E)\) permute its
\(d\) tensor factors. Their invariant space is
\(\operatorname{Sym}^d\operatorname{Sym}^n V\).
The Cauchy decomposition
\[
 \operatorname{Sym}(V\otimes E)
      =\bigoplus_\mu S_\mu V\otimes S_\mu E
\]
therefore identifies the required multiplicity with
\[
 \dim\left((S_{\lambda_{d,k}}E)_{(n,\ldots,n)}\right)^{S_d}.
\]
The precise classical identity is Macdonald, Chapter~I,
equation~(4.3), printed p.~63, together with the polynomial-functor
interpretation in Appendix~A, (5.3)--(5.6), pp.~155--156.

Subtracting \(k\) from each of the \(d\) highest-weight coordinates
gives
\[
 S_{\lambda_{d,k}}E
       =\det(E)^k\otimes\operatorname{Sym}^{d(n-k)}E.
\]
Its weight \((n,\ldots,n)\) space is one-dimensional:
after the determinant twist it is spanned by
\(e_1^{n-k}\cdots e_d^{n-k}\).
Every permutation matrix fixes this monomial and acts on
\(\det(E)^k\) by \(\operatorname{sgn}^k\).
For \(d\geq2\), this one-dimensional representation has an
invariant exactly when \(k\) is even. This proves \eqref{eq:F011}.
The cases \(k=0\), \(k=n\), and \(d=1\) are included or handled
as stated, with no negative exponent.
\end{proof}

\paragraph{Consequences for actual and sufficient multiplicity tests.}
For \(d\geq2\), odd \(k\) gives zero permanent-side multiplicity
for every \(m,n\), since the ambient multiplicity is zero.
For even \(k\), there is no actual multiplicity obstruction in
the B\"urgisser--Ikenmeyer--Panova range \(n\geq m^{25}\).
Indeed \(q\leq a=1\), and their Theorem~1.4 says that
\(q>0\) implies \(K_n>0\); in that case \(q=1\leq K_n\).
For the Ikenmeyer--Panova range \(n>3m^4\), their Theorem~1.6
gives the analogous exclusion of the sufficient test \(q>g\),
since \(q>0\) forces \(g>0\). The latter argument alone does
not exclude an actual inequality \(q>K_n\).

More generally, throughout the BIP range every actual multiplicity
obstruction must have
\[
                     q^m_\lambda(d[n])\geq2,
                \qquad a_\lambda(d[n])\geq2.
\]
Thus all ambient multiplicity-one shapes are excluded there,
and \eqref{eq:F011} provides an explicit infinite family to which
this applies. No general determinant-side nonvanishing is asserted
for even \(k\) outside the stated no-occurrence range.
The special case \(k=2\) has a stronger all-\(n\) determinant
and permanent interpretation through the Hessian calculation in O-011.

\paragraph{Grothendieck dictionary.}
The multidegree condition is a character-weight submodule for the
diagonal torus in \(\mathrm{GL}(E)\); the permutation matrices
form its constant finite subgroup \(S_d\).
The calculation concerns the normalizer's action on a one-dimensional
weight submodule of a Schur functor. It uses locally free symmetric,
determinant, and Schur operations, then finite-group invariants in
characteristic zero. Thus it is a calculation in
\(K_0(\operatorname{Rep}(\mathrm{GL}(V)\times(T\rtimes S_d)))\),
with the grading/comodule dictionary of SGA~3, I, Proposition~4.7.3,
and the lambda operations of SGA~6, VI, Theorem~3.3.
The sign is the orientation character of the determinant line;
forgetting it would incorrectly predict the odd cases to occur.

\subsection*{Finding F-012: mixed near-rectangular tails have zero ambient multiplicity}
Let \(d\geq2\), \(n\geq1\), \(0\leq k<n\), and
\(0\leq s\leq d-1\). Put \(a=n-k\geq1\), \(D=da\), and
delete zero parts from
\[
 \lambda_{d,k,s}
  =\bigl(D+k-s,(k+1)^s,k^{d-s-1}\bigr)\vdash nd.
\]
This is a partition: \(D-s\geq1\), and for \(s>0\) its
first part is at least \(k+1\). Its ambient coefficient is
\begin{equation}\label{eq:F012}
 a_{\lambda_{d,k,s}}(d[n])=
 \begin{cases}
 1& s=0\text{ and }k\text{ even},\\
 1& s=d-1\text{ and }k\text{ odd},\\
 0&\text{otherwise}.
 \end{cases}
\end{equation}
The second positive endpoint is exactly F-011 with \(k\) replaced
by \(k+1\). In particular every mixed tail, with
\(1\leq s\leq d-2\), has zero permanent-side multiplicity,
for all sizes: its ambient multiplicity already vanishes.
At \(k=0\), the assertion recovers hook vanishing in the
indicated range.

\begin{proof}
Use the Cauchy weight-space description from F-011, with
\(E=\mathbb C^d\). Subtracting the rectangle \((k^d)\) gives
\[
 S_{\lambda_{d,k,s}}E
       =\det(E)^k\otimes S_{(D-s,1^s)}E.
\]
It remains to understand the balanced weight \((a,\ldots,a)\)
of the hook Schur module as a representation of the column
permutation group \(S_d\).

The hook form of Jacobi--Trudi gives, in the representation ring,
\[
 [S_{(D-s,1^s)}E]
    =\sum_{t=0}^{s}(-1)^{s-t}
       [\operatorname{Sym}^{D-t}E\otimes\Lambda^t E].
\]
This is the usual Jacobi--Trudi identity of Macdonald, I, (3.4),
expanded along the hook; no positivity is being inferred from
its signs. Taking the balanced weight commutes with restriction
to its normalizer, so it preserves this equality as an equality
of \(S_d\)-representation classes.

In the summand indexed by \(t\), a basis of that weight space is
indexed by the \(t\)-subsets \(J\subseteq[d]\): its basis vector is
\[
  \left(\prod_{i=1}^d e_i^{a-\mathbf1_{i\in J}}\right)
                 \otimes\bigwedge_{j\in J}e_j,
\]
with increasing wedge order. Every exponent is nonnegative
because \(a\geq1\). A column permutation permutes the first
factor and has exactly the oriented-subset sign in the second.
Consequently this weight space is the \(S_d\)-module
\(\Lambda^t P_d\), where \(P_d\) is the permutation module.
Write \(P_d=\mathbf1\oplus U_d\), with \(U_d\) the augmentation
representation. The exterior-series identity gives
\[
 \sum_{t=0}^s(-1)^{s-t}[\Lambda^t P_d]
                              =[\Lambda^s U_d].
\]
Thus the balanced-weight hook module has exactly that character.
Characteristic-zero semisimplicity identifies its isomorphism
class with \(\Lambda^s U_d\).

Including the determinant twist, the desired coefficient is
\[
                    \dim(\operatorname{sgn}^k\otimes
                                      \Lambda^s U_d)^{S_d}.
\]
The only invariant in \(\Lambda^\bullet U_d\) is in degree zero,
as proved by the orientation calculation in O-005.
Moreover \(U_d\) is self-dual and
\(\det(U_d)=\operatorname{sgn}\), so exterior duality gives
\[
  \operatorname{sgn}\otimes\Lambda^s U_d
                      \simeq\Lambda^{d-1-s}U_d.
\]
For even \(k\) the invariant occurs exactly at \(s=0\);
for odd \(k\), exactly at \(s=d-1\). Its dimension is one
in each case. This proves \eqref{eq:F012} and all stated
vanishing consequences.
\end{proof}

\paragraph{Dictionary and implication.}
This proof takes the torus-weight submodule of a Schur operation
on a locally free rank-\(d\) sheaf, then the invariant module
of the constant group scheme \(S_d\). Its exact computation
uses the exterior lambda identity for
\(P_d=\mathbf1\oplus U_d\), with the orientation character
retained. These are the grading and representation dictionaries
of SGA~3, I, 4.7.2--4.7.3, and the exact-sequence exterior
filtration of SGA~6, V, Lemma~2.2.1 and VI, Theorem~3.3.
It provides a genuine all-size exclusion of specified GCT
candidate shapes. It compares no positive multiplicities for
the shapes outside this vanishing family.

\subsection*{Finding F-013: exact ambient multiplicities after a rectangular shift}
Let \(d\geq2\), \(n\geq1\), \(0\leq k\leq n\), put
\(a=n-k\), \(D=da\), and let \(0\leq t\leq D/2\) be an integer.
Delete zero parts from
\[
                 \lambda=(D-t+k,t+k,k^{d-2})\vdash nd.
\]
Let \(p_j(d,b)\) count partitions of \(j\) in a rectangle with
\(d\) rows and width \(b\); set it equal to zero for \(j<0\)
or \(b<0\), and set \(p_0(d,0)=1\).
With \(\delta=d(d-1)/2\), one has
\begin{equation}\label{eq:F013}
 a_\lambda(d[n])=
 \begin{cases}
 p_t(d,a)-p_{t-1}(d,a),&k\text{ even},\\
 p_{t-\delta}(d,a-d+1)-p_{t-1-\delta}(d,a-d+1),
       &k\text{ odd and }a\geq d-1,\\
 0,&k\text{ odd and }a<d-1.
 \end{cases}
\end{equation}
Thus this ambient multiplicity is reduced to the same partition
differences as the rectangular/two-row Kronecker family F-002.
The formula is a direct consequence of the Cauchy and two-row
Jacobi--Trudi identities, with the sign action recorded explicitly.
It is not a claim that the two orbit closures have these full
ambient multiplicities.

\begin{proof}
Again put \(E=\mathbb C^d\) and use the column weight
\((n,\ldots,n)\) description in F-011. Removing the determinant
twist gives
\[
 S_\lambda E=\det(E)^k\otimes S_{(D-t,t)}E.
\]
Jacobi--Trudi, equivalently the difference of the two adjacent
horizontal Pieri decompositions, is
\[
 [S_{(D-t,t)}E]
 = [\operatorname{Sym}^{D-t}E\otimes\operatorname{Sym}^{t}E]
  -[\operatorname{Sym}^{D-t+1}E\otimes
                                      \operatorname{Sym}^{t-1}E],
\]
with the second summand zero for \(t=0\).
The classical locators are Macdonald, I, (3.4) and (5.16).

The balanced weight \((a,\ldots,a)\) in the first tensor product
has a monomial basis indexed by
\[
             B_t=\{\beta\in\{0,\ldots,a\}^d:|\beta|=t\}:
 \qquad e^{a\mathbf1-\beta}\otimes e^\beta.
\]
The column permutation group permutes this basis. The second
tensor product gives \(\mathbb C[B_{t-1}]\) in exactly the same
way. Thus the required coefficient is the difference of the
\(\operatorname{sgn}^k\)-multiplicities in those two permutation
modules.

If \(k\) is even, each orbit contributes one invariant.
Sorting the entries of \(\beta\) identifies its orbit with a
partition of \(t\) having at most \(d\) parts, each at most \(a\).
This gives \(p_t(d,a)-p_{t-1}(d,a)\).

If \(k\) is odd, an orbit contributes to the sign isotypic
space precisely when its stabilizer has no odd permutation.
The stabilizer is the product of the symmetric groups on the
equal-entry positions. It contains an odd transposition if any
two entries are equal, and is trivial otherwise. Hence the
contributing orbits are exactly the strictly increasing lists
\[
                0\leq b_1<b_2<\cdots<b_d\leq a.
\]
There are none when \(a<d-1\). Otherwise subtract
\((0,1,\ldots,d-1)\). The resulting list is weakly increasing,
lies in \([0,a-d+1]\), and has sum \(t-\delta\).
Reversing its order is a partition in the stated rectangle.
Applying the same bijection at \(t-1\) proves the second line
of \eqref{eq:F013}.
All finite-group invariant and sign-isotypic functors are exact
in characteristic zero, so the subtraction is legitimate.
\end{proof}

In the odd case, the nonzero shifted degree \(u=t-\delta\)
satisfies \(u\leq d(a-d+1)/2\). Thus the displayed differences
stay in the first half of their rectangles, consistently with
their representation-theoretic nonnegativity. The conventions
also cover \(a=0,t=0\), and the single-point rectangle in the
odd case \(a=d-1,t=\delta\).

\paragraph{A concrete four-copy family beyond occurrence.}
If \(d\geq6\) and \(n\geq8\), take \(k=2,t=6\). Then
\[
 \lambda=(d(n-2)-4,8,2^{d-2}),\qquad
                         a_\lambda(d[n])=4.
\]
Indeed the rectangle bounds do not affect partitions of six or
five. Removing the part of size one leaves the four partitions
\((6),(4,2),(3,3),(2,2,2)\).
For \(d\leq m^2\), \(m\geq3\), and \(d\geq6\), this family
also meets the coarse internal-padding conditions:
its tail has size \(2d+4\leq md\) and its length is \(d\).
The ambient coefficient four does not prove positive permanent
multiplicity or compare it with the determinant closure.
It gives a precise small multiplicity space on which F-010's
remaining rank comparison can be studied, including at arbitrarily
large row counts.

\paragraph{Dictionary.}
This is a balanced-torus-weight calculation of a Schur functor
followed by the trivial or orientation isotypic summand for
the constant group scheme \(S_d\). The two permutation modules
are linearizations of the actual finite sets \(B_t\) and
\(B_{t-1}\). The sign contribution is detected by the stabilizer,
so repeated entries are excluded for a representation-theoretic
reason, not by an assumption that a virtual Burnside class is
effective. The character-weight and comodule language is SGA~3,
I, 4.7.2--4.7.3; symmetric, determinant and Schur operations
are those of SGA~6, VI, Theorem~3.3 and the polynomial-functor
dictionary already used above.

% END agents/formulas/further_gates.tex

% BEGIN agents/formulas/second_row_stabilization.tex
\section{Second-row stabilization and actual multiplicities}

\subsection*{Finding F-014: stabilization controlled by the second row}
Let \(d\geq1\), \(n\geq1\), and let \(\lambda\vdash nd\)
have \(r\leq d\) rows. If \(r=1\), its ambient multiplicity is
one and it occurs in the determinant coordinate ring. Suppose
\(r\geq2\), put \(b=\lambda_2\), and assume \(n\geq b\).
Then
\[
       \mu=(db-|\bar\lambda|,\lambda_2,\ldots,\lambda_r)
                                                     \vdash db
\]
is a partition and
\begin{equation}\label{eq:F014stab}
                         a_\lambda(d[n])=a_\mu(d[b]).
\end{equation}
More precisely, multiplication by the fixed first basis vector
in each symmetric-power factor identifies the entire highest-weight
spaces of these two weights. The statement is established here
directly from Cauchy and flag minors; no assertion of historical
priority for plethysm stabilization is intended.

\begin{proof}
Let \(U\) have dimension \(r\), let \(E=\mathbb C^d\), and use
variables \(z_{ij}\), \(1\leq i\leq r\), \(1\leq j\leq d\),
for the symmetric algebra \(\operatorname{Sym}(U\otimes E)\).
Rows carry the \(\mathrm{GL}(U)\) action and columns the commuting
\(\mathrm{GL}(E)\) action. The Cauchy decomposition says that
the space \(H_\lambda\) of row highest-weight vectors of weight
\(\lambda\) is the irreducible column module \(S_\lambda E\).
The condition \(r\leq d\) ensures that this module is nonzero.

For a column subset \(J\) of size \(h\), call
\[
             \Delta_{h,J}=\det(z_{ij})_{1\leq i\leq h,\ j\in J}
\]
a flag minor. It is a row highest-weight vector of weight
\((1^h)\). Consider all products of flag minors whose heights,
with multiplicity, are the column lengths \(\lambda'_j\).
Each product is a row highest-weight vector of weight \(\lambda\).
Their linear span is stable under \(\mathrm{GL}(E)\), because
the column change-of-basis formula for a minor expresses it
as a linear combination of minors of the same height.
This span is nonzero: choose the first \(h\) columns for every
minor of height \(h\); their product is a nonzero polynomial.
Irreducibility of \(H_\lambda\simeq S_\lambda E\) now proves
that these products span all of \(H_\lambda\).
This proves the required spanning assertion without presuming
an unproved basis theorem for the products.

Exactly \(\lambda_2=b\) of the minor factors have height at
least two. In the expansion of any one minor, each column is
used at most once. Height-one minors use only first-row variables.
It follows that every monomial in every spanning product, and
therefore in every element of \(H_\lambda\), has degree at most
\(b\) in the variables below the first row of any fixed column.

Take the column multidegree \((n,\ldots,n)\) part of
\(H_\lambda\). Every monomial there consequently has exponent
at least \(n-b\) on each \(z_{1j}\). Hence every polynomial in
that space is divisible by
\[
                         Q^{n-b},\qquad Q=\prod_{j=1}^d z_{1j}.
\]
The polynomial \(Q\) is fixed by the upper row unipotent group,
has row weight \(de_1\), and is fixed by column permutations.
Division by \(Q^{n-b}\) therefore produces a row highest-weight
vector of weight \(\mu=\lambda-d(n-b)e_1\), of column
multidegree \((b,\ldots,b)\). Conversely multiplication by
\(Q^{n-b}\) sends every such vector back to the original space.
These operations are inverse and commute with column permutations.

The tuple \(\mu\) is a partition because
\[
 |\bar\lambda|\leq(r-1)b\leq(d-1)b,
                  \qquad db-|\bar\lambda|\geq b=\lambda_2.
\]
Finally, column-permutation invariants in the balanced column
space are \(\operatorname{Sym}^d\operatorname{Sym}^n U\)
and \(\operatorname{Sym}^d\operatorname{Sym}^b U\), respectively.
Taking invariants in the displayed isomorphism proves
\eqref{eq:F014stab}, including all multiplicity copies.
The one-row case consists of a single highest-weight monomial
and follows directly.
\end{proof}

\paragraph{Full determinant multiplicities.}
For \(b\geq1\), \(r\geq2\), write
\[
 u=\left\lfloor\frac{b-1}{2}\right\rfloor,\qquad
 v=\left\lceil\frac{b-1}{2}\right\rceil,\qquad
 B(b,r)=\binom{r-1+u}{u}+\binom{r-1+v}{v}-1.
\]
If \(r\leq n^2\) and \(n\geq B(b,r)\), then
\begin{equation}\label{eq:F014det}
                         K_n(\lambda)=a_\lambda(d[n]).
\end{equation}
In particular there is no actual permanent-versus-determinant
multiplicity obstruction of this shape, at any outer degree.

\begin{proof}
The degree-splitting construction O-012.a puts every form
\(x_1^{n-b}f\), \(f\in\operatorname{Sym}^b U^*\), in
\(\Omega_n\) under the stated inequality.
It also gives \(B(b,r)\geq b\), so the stabilization applies.
By the arbitrary-rank padding inequality O-012.b,
\[
 K_n(\mu+d(n-b)e_1)\geq a_\mu(d[b])=a_\lambda(d[n]).
\]
Since the partition on the left is \(\lambda\), and an orbit-closure
coordinate ring is an ambient quotient, the reverse inequality
also holds. This proves \eqref{eq:F014det}.

The distinction between the two maps used here is essential.
Multiplication by \(e_1^{n-b}\) on the coordinate tensor space
is Borel equivariant up to a character and gives the highest-weight
space isomorphism. Actual pullback to padded forms is its
differential transpose. Their composition is diagonal with
nonzero factorial entries, as proved in O-012.b; it preserves
linear independence of every lifted copy. No false assertion
that those two maps are the same has been used.
\end{proof}

\paragraph{Explicit GCT second-row exclusions.}
For \(m\geq2\), a positive permanent-side multiplicity satisfies
\(r\leq m^2\) in the internal-padding convention.
The number \(\binom{r-1+j}{j}\) counts degree-\(j\) monomials
in \(r\) variables, so it is at most \(r^j\), by mapping
words of length \(j\) to their commutative monomials.
Thus
\[
                         B(b,r)\leq r^u+r^v-1.
\]
If \(b\leq5\), then \(u,v\leq2\), and
\[
               B(b,r)\leq2r^2-1\leq2m^4-1<3m^4.
\]
Consequently in the range \(n>3m^4\), every actual multiplicity
obstruction must satisfy
\begin{equation}\label{eq:F014IP}
                              \lambda_2\geq6.
\end{equation}
If \(b\leq25\), then \(u,v\leq12\), and
\[
 B(b,r)\leq2r^{12}-1\leq2m^{24}-1<m^{25}
                                           \qquad(m\geq2).
\]
Therefore in the range \(n\geq m^{25}\), every actual
multiplicity obstruction must satisfy
\begin{equation}\label{eq:F014BIP}
                              \lambda_2\geq26.
\end{equation}
Shapes with more than \(d\) rows have zero ambient multiplicity;
one-row shapes have full determinant multiplicity, so neither
exception leaves a gap in these corollaries.
The exclusions hold even when \(d\) is much larger than the
number of rows. They therefore remain additional restrictions
after any independent degree exclusion for \(d(d-1)\leq n\).

\paragraph{Grothendieck dictionary.}
The matrix polynomial ring is the coordinate algebra of the
affine scheme of maps between two locally free sheaves. The
flag minors are exterior-power matrix coefficients; their row
characters and commuting column action are the Cauchy/Schur
functor operations of cote~133, scan~028, and Macdonald, I,
(4.3), p.~63, and Appendix~A, (5.3)--(5.6), pp.~155--156.
The balanced multidegree is a torus-character submodule, with
the grading/comodule correspondence of SGA~3, I, 4.7.3.
Multiplication and division by the displayed common section
give an explicit isomorphism on those characteristic-zero
highest-weight modules. The determinant construction is an
integral identity, whereas irreducible splitting and the
factorial-invertibility argument are used only in characteristic
zero. The induced coherent quotient on the orbit-closure scheme
is what proves equality of the actual multiplicities.

% END agents/formulas/second_row_stabilization.tex

% BEGIN agents/formulas/restricted_padding.tex
\section{Full multiplicity from a restricted coefficient box}

\subsection*{Finding F-015: shape-sensitive padding preserves every copy}
Let \(\lambda\vdash nd\) have \(2\leq r\leq\min\{d,n^2\}\)
rows, and put \(b=\lambda_2\). For \(j\geq0\), define
\[
 C_j(\lambda)=\#\{\alpha\in\mathbb Z_{\geq0}^{r-1}:
       \alpha_i\leq\lambda_i\ (2\leq i\leq r),\quad
       |\alpha|\leq j\}.
\]
Set
\[
 \begin{split}
 u&=\left\lfloor\frac{b-1}{2}\right\rfloor,\qquad
 v=\left\lceil\frac{b-1}{2}\right\rceil,\\
 D_{\mathrm{box}}(\lambda)&=C_u(\lambda)+C_v(\lambda)-1.
 \end{split}
\]
Also let \(T_b(\lambda)\) denote the bounded variable-splitting
count of M-009, with its degree parameter replaced by \(b\):
\[
 \begin{split}
 N_A(b,\lambda)&=
 \#\{\alpha\in\mathbb Z_{\geq0}^{A}:\alpha_i\leq\lambda_i,
                                            |\alpha|\leq b\},\\
 T_b(\lambda)&=
 \min_{A\sqcup B=\{2,\ldots,r\}}
                  \bigl(N_A(b,\lambda)+N_B(b,\lambda)-1\bigr).
 \end{split}
\]
Then
\begin{equation}\label{eq:F015}
 n\geq\min\{D_{\mathrm{box}}(\lambda),T_b(\lambda)\}
 \quad\Longrightarrow\quad
                       K_n(\lambda)=a_\lambda(d[n]).
\end{equation}
The assertion holds at every outer degree \(d\). It excludes
actual multiplicity obstructions for every orbit closure in the
same ambient form space, including the internally padded permanent.
It is a full-multiplicity conclusion, stronger than mere occurrence.

The count \(D_{\mathrm{box}}\) is at most \(B(b,r)\) from F-014,
because imposing the coordinate bounds only deletes monomials.
When \(b\leq m\), the other count satisfies
\(T_b(\lambda)\leq T_m(\lambda)\). Thus the criterion can
strengthen both the uncapped second-row bound and M-009.

\begin{proof}
The coordinate bound on \(\alpha_2\) is \(b\). Hence
\(C_j(\lambda)\geq j+1\) for \(0\leq j\leq b\), by
using powers of the second variable alone. It follows that
\(D_{\mathrm{box}}\geq u+v+1=b\). In every variable split,
the part containing the second variable contributes at least
\(b+1\) monomials and the other part contributes at least one;
therefore \(T_b\geq b+1\). The hypothesis in
\eqref{eq:F015} consequently implies \(n\geq b\).

Let \(U\) have basis \(e_1,\ldots,e_r\), with dual variables
\(x_1,\ldots,x_r\), and put
\[
       \mu=(db-|\bar\lambda|,\lambda_2,\ldots,\lambda_r).
\]
By F-014 this is a partition and
\(a_\mu(d[b])=a_\lambda(d[n])\).
Let \(H_\mu\) be the complete highest-weight space of weight
\(\mu\) in \(\operatorname{Sym}^d\operatorname{Sym}^b U\).
Define the monomial subspace
\[
 W=\operatorname{span}\{e^\eta:|\eta|=b,\quad
                    \eta_i\leq\lambda_i\ (2\leq i\leq r)\}
                  \subseteq\operatorname{Sym}^b U.
\]
Every monomial in a polynomial of weight \(\mu\) has total
nonnegative weight \(\mu\). A generator \(e^\eta\) with
\(\eta_i>\mu_i=\lambda_i\) for some \(i\geq2\) cannot
appear in such a monomial. Therefore
\begin{equation}\label{eq:F015support}
                         H_\mu\subseteq\operatorname{Sym}^d W.
\end{equation}
This statement concerns every highest-weight vector, including
all multiplicity copies. It does not require \(W\) to be a
Borel-stable subspace.

Let \(\mathcal C\subseteq\operatorname{Sym}^b U^*\) be the
span of the dual monomials \(x^\eta\) with the same support
conditions. The differential pairing restricts to a nondegenerate
diagonal pairing between \(W\) and \(\mathcal C\), with
\(\langle e^\eta,x^\eta\rangle=\eta!\).
Consequently restriction of polynomial functions to
\(\mathcal C\) is injective on \(\operatorname{Sym}^d W\),
and in particular on \(H_\mu\).

Put \(p=n-b\). Multiplication
\(T:e^\eta\mapsto e_1^p e^\eta\) from
\(\operatorname{Sym}^b U\) to \(\operatorname{Sym}^n U\)
is upper-Borel equivariant up to the character \(g_{11}^p\).
Thus \(\operatorname{Sym}^d T\) takes \(H_\mu\) injectively
to highest-weight vectors of weight
\(\mu+dp e_1=\lambda\).
With \(U\) the span of the first \(r\) basis vectors of the
ambient space, these are also highest-weight vectors for the
ambient general linear group: raising operators involving a
later source index act by zero.
Consider actual form padding
\[
                 P:\mathcal C\longrightarrow
                       \operatorname{Sym}^n U^*,\qquad
                         f\longmapsto x_1^p f.
\]
As in O-012.b, the transpose \(R\) of padding is
\(\partial_{e_1}^p\), so
\[
                 RT(e^\eta)
                   =\frac{(\eta_1+p)!}{\eta_1!}e^\eta.
\]
This invertible diagonal operator preserves \(W\). Pulling
the lifted vectors back to \(P(\mathcal C)\) is consequently
the restriction of the invertible operator
\(\operatorname{Sym}^d(RT)\) on \(\operatorname{Sym}^d W\).
By \eqref{eq:F015support}, a basis of \(H_\mu\) remains
linearly independent after lifting and restriction.

It remains to put that restricted padded family in \(\Omega_n\).
The coefficient-box determinant construction O-014.a, obtained
by restricting both degree-split trees of O-012.a to the
downward-closed index sets above, gives
\[
           x_1^{n-b}\mathcal C\subseteq\Omega_n
                    \qquad\text{when }n\geq D_{\mathrm{box}}.
\]
Alternatively the bounded variable-split construction in M-009,
applied to degree \(b\), gives the same containment when
\(n\geq T_b\). In either case the displayed basis remains
linearly independent in the determinant coordinate ring.
Thus
\[
 K_n(\lambda)\geq\dim H_\mu
       =a_\mu(d[b])=a_\lambda(d[n]).
\]
The reverse inequality holds because that coordinate ring is
a quotient of the ambient polynomial ring. This proves
\eqref{eq:F015}.
\end{proof}

\paragraph{Exact scope of the improvement.}
The argument preserves all ambient highest-weight copies by
restricting to a linear space adapted to their support. It does
not replace the desired comparison by a positive occurrence
statement, and it does not infer containment of every padded
degree-\(b\) form. The rank criterion F-010 is now decided
whenever \eqref{eq:F015} holds: the determinant evaluation map
has its maximum possible rank. When the displayed inequality
fails, the construction supplies no rank inequality.
The permanent degree \(m\) enters only when applying separate
permanent-side shape constraints such as \(r\leq m^2\).

\paragraph{Grothendieck dictionary.}
The box is a finite torus-weight support in a locally free
symmetric-power module, with the grading/comodule interpretation
of SGA~3, I, 4.7.3. Its dual is an affine coefficient scheme.
The graph identities define integral morphisms from that scheme
to the determinant locus. The restriction and padding maps are
symmetric-power maps, and their differential transpose has
explicit factorial diagonal entries. Cauchy and the Schur
functor decomposition, Macdonald I(4.3) and Appendix~A(5.3)--(5.6),
identify the characteristic-zero highest-weight multiplicity
spaces. Invertibility of those factorials and semisimplicity
are asserted only in characteristic zero. This is a coherent
quotient calculation on the orbit-closure scheme; it does not
claim an effective Burnside lift or positivity of virtual
classes in unrelated fibers.

% END agents/formulas/restricted_padding.tex

% BEGIN agents/formulas/nested_support.tex
\section{Nested support bounds from exterior minors}

\subsection*{Finding F-016: tail-sum support and full determinant multiplicity}
Let \(\lambda\vdash nd\), with
\(r=\ell(\lambda)\leq d\), and let \(U\) have basis
\(e_1,\ldots,e_r\). Every highest-weight vector of weight
\(\lambda\) in \(\operatorname{Sym}^d\operatorname{Sym}^n U\)
belongs to \(\operatorname{Sym}^d W_{n,\lambda}\), where
\begin{equation}\label{eq:F016support}
 W_{n,\lambda}=
 \operatorname{span}\left\{e^\eta:|\eta|=n,\quad
             \sum_{i=j}^r\eta_i\leq\lambda_j
                             \text{ for every }2\leq j\leq r\right\}.
\end{equation}
This assertion does not require \(n\geq\lambda_2\).
It constrains the sum of the exponents in each terminal set
of rows, and is stronger than the individual coordinate bounds.

\begin{proof}
Use the matrix polynomial model
\(\operatorname{Sym}(U\otimes\mathbb C^d)\), with row
variables indexed by \(i\) and column variables by \(h\).
As proved in F-014, the entire row highest-weight space of
weight \(\lambda\) is spanned by products of flag minors
whose heights are the column lengths of \(\lambda\).
For completeness, their span is a nonzero column-group-stable
subspace of the irreducible Cauchy module
\(S_\lambda\mathbb C^d\), and hence is the whole space.

Exactly \(\lambda_j\) minor factors in each such product
have height at least \(j\). A minor of smaller height uses
no row \(j,\ldots,r\). In a monomial of a minor of height
at least \(j\), any fixed column contributes at most one
variable. Consequently every monomial in every spanning product,
and therefore every polynomial in the entire highest-weight
space, has degree at most \(\lambda_j\) in rows
\(j,\ldots,r\) of any one fixed column.

Now take the column multidegree \((n,\ldots,n)\) subspace.
Its monomial basis is naturally the tensor-product monomial
basis of \((\operatorname{Sym}^n U)^{\otimes d}\).
The preceding bound holds separately in every column, so
the relevant highest-weight subspace lies in
\(W_{n,\lambda}^{\otimes d}\). Taking the column-permutation
invariants gives
\[
       (W_{n,\lambda}^{\otimes d})^{S_d}
                                  =\operatorname{Sym}^d W_{n,\lambda}.
\]
Under this identification, an unordered product of monomial
generators is the sum of its distinct column assignments,
up to a nonzero scalar. Different unordered products have
disjoint sets of assignments. Thus no cancellation in this
passage could conceal a forbidden coefficient generator.
This proves \eqref{eq:F016support} for every copy.
\end{proof}

Suppose now \(r\geq2\), put \(b=\lambda_2\), and define
the finite, coordinatewise downward-closed set
\[
 \Gamma_\lambda=
 \left\{\alpha\in\mathbb Z_{\geq0}^{r-1}:
          \sum_{i=j}^r\alpha_i\leq\lambda_j
                              \text{ for all }2\leq j\leq r\right\}.
\]
In particular \(|\alpha|\leq b\). For \(s\geq0\), write
\[
 C_s^{\mathrm{flag}}(\lambda)=
       \#\{\alpha\in\Gamma_\lambda:|\alpha|\leq s\}.
\]
For a subset \(A\subseteq\{2,\ldots,r\}\), write
\(\alpha|_{A^c}=0\) to mean that its coordinates outside
\(A\) vanish, and put
\[
 \begin{split}
 N_A^{\mathrm{flag}}(\lambda)&=
       \#\{\alpha\in\Gamma_\lambda:\alpha|_{A^c}=0\},\\
 D_{\mathrm{flag}}(\lambda)&=
 C_{\lfloor(b-1)/2\rfloor}^{\mathrm{flag}}(\lambda)
 +C_{\lceil(b-1)/2\rceil}^{\mathrm{flag}}(\lambda)-1,\\
 T_{\mathrm{flag}}(\lambda)&=
 \min_{A\sqcup B=\{2,\ldots,r\}}
       \bigl(N_A^{\mathrm{flag}}(\lambda)
                         +N_B^{\mathrm{flag}}(\lambda)-1\bigr).
 \end{split}
\]
If also \(r\leq n^2\), then
\begin{equation}\label{eq:F016K}
 n\geq\min\{D_{\mathrm{flag}}(\lambda),T_{\mathrm{flag}}(\lambda)\}
 \quad\Longrightarrow\quad
                     K_n(\lambda)=a_\lambda(d[n]).
\end{equation}

\begin{proof}
All powers of the second variable up to degree \(b\) lie
in \(\Gamma_\lambda\). Exactly as in F-015,
\(D_{\mathrm{flag}}\geq b\) and
\(T_{\mathrm{flag}}\geq b+1\). Thus \(n\geq b\).
Apply F-014 to obtain
\[
 \mu=(db-|\bar\lambda|,\lambda_2,\ldots,\lambda_r),
                 \qquad a_\mu(d[b])=a_\lambda(d[n]).
\]
The first part of this finding, now applied at degree \(b\)
and weight \(\mu\), puts its entire highest-weight space
in \(\operatorname{Sym}^d W\), where
\[
 W=\operatorname{span}
      \{e_1^{b-|\alpha|}e_2^{\alpha_2}\cdots e_r^{\alpha_r}:
                                        \alpha\in\Gamma_\lambda\}.
\]
The tail bounds for \(\mu\) and \(\lambda\) are identical.

Let \(\mathcal C\) be the corresponding dual monomial
coefficient space of degree-\(b\) forms. The graph construction
O-015, or the constructions O-014.a and M-009 restricted to
these downward-closed vertex sets, puts
\(x_1^{n-b}\mathcal C\) in \(\Omega_n\) under the respective
degree-split or variable-split inequality. Indeed, every
coordinatewise smaller multi-index satisfies all terminal-sum
bounds. Splitting a monomial between the trees, or first
removing the variable carried by its coefficient edge, keeps
both endpoints and every tree path in the permitted sets.
The cycle-cover signs and exponents are unchanged.

The restriction-and-padding proof of F-015 applies to this
monomial space \(W\) without any assumption that it is
Borel-stable. Restriction to \(\mathcal C\) is injective
on \(\operatorname{Sym}^d W\), and the composition of the
highest-weight lift with actual padding pullback acts on each
generator by the nonzero scalar
\[
                       \frac{(b-|\alpha|+n-b)!}
                            {(b-|\alpha|)!}.
\]
Thus all \(a_\mu(d[b])=a_\lambda(d[n])\) copies remain
independent in the determinant coordinate ring. The ambient
upper bound gives \eqref{eq:F016K}.
\end{proof}

\paragraph{A simple bound from a cut in the rows.}
For \(2\leq k\leq r\), split the variables into rows
\(2,\ldots,k-1\) and \(k,\ldots,r\). Ignoring other
terminal-sum bounds can only increase the counts. The first
part has \(k-2\) variables of total degree at most \(b\),
and the second has \(r-k+1\) variables of total degree at
most \(\lambda_k\). Hence
\begin{equation}\label{eq:F016cut}
 T_{\mathrm{flag}}(\lambda)
 \leq \binom{b+k-2}{k-2}
           +\binom{r-k+1+\lambda_k}{\lambda_k}-1.
\end{equation}
If the right-hand side is at most \(n\), every ambient
copy of \(\lambda\) occurs on the determinant closure.
At \(k=3\), this reads
\begin{equation}\label{eq:F016thirdrow}
 n\geq\lambda_2+\binom{r-2+\lambda_3}{\lambda_3}
             \quad\Longrightarrow\quad K_n(\lambda)=a_\lambda(d[n]).
\end{equation}
In particular, for \(r\geq3\) and \(\lambda_3\leq2\),
the readily checked sufficient condition is
\[
                       n\geq\lambda_2+\binom r2.
\]
This allows second rows almost as large as \(n\), even at
arbitrarily large outer degree. It is much stronger on these
tails than the uniform constant second-row exclusion alone.

For the exact tail \((b,2^{r-2})\), \(b\geq2\), the nested
condition is simply that the variables after the second have
total degree at most two. Its degree-split counts, for \(s\geq1\),
are
\[
 C_s^{\mathrm{flag}}
  =(s+1)+s(r-2)+(s-1)\binom{r-1}{2}
  =r+(s-1)\binom r2,
\]
where the displayed formula is used only for \(s\leq b\).
Thus \(D_{\mathrm{flag}}=2r+(b-3)\binom r2-1\) for
\(b\geq3\), and \(D_{\mathrm{flag}}=r\) for \(b=2\).
The variable-split bound \(b+\binom r2\) can be smaller still.

\paragraph{Relation to the actual GCT comparison.}
The support statement \eqref{eq:F016support} applies directly
to determinant-ideal highest-weight functions before any
stabilization. After aligning the padding form, it justifies
deleting degree-\(m\) coefficients outside the same terminal-sum
constraints while preserving evaluation. The root finding
M-011 develops that variant without requiring \(n\geq\lambda_2\).
By contrast, \eqref{eq:F016K} proves the stronger equality
with the full ambient multiplicity, using the stabilized base
degree \(b=\lambda_2\). These two consequences have different
hypotheses and neither is inferred from occurrence alone.

\paragraph{Grothendieck dictionary.}
The additional information comes from exterior-power matrix
coefficients: a flag minor cannot use the same column twice.
The Cauchy decomposition identifies their column-group span
with the complete Schur multiplicity module, as in Macdonald,
I(4.3), p.~63, and Appendix~A(5.3)--(5.6), pp.~155--156.
Passing to a terminal row set is a support calculation in the
torus grading of the matrix scheme, under SGA~3, I, 4.7.3.
The permitted coefficient support defines an affine linear
subscheme; it need only be stable under the diagonal torus.
The associated determinant morphisms are integral identities.
The all-copy conclusion uses characteristic-zero highest-weight
modules and the factorial pairing, with exactly the same
base-change qualifications as F-015. It supplies an explicit
condition under which the actual coherent quotient loses no
copies of the indicated Schur functor.

% END agents/formulas/nested_support.tex

% BEGIN agents/formulas/mixed_tail_graph.tex
\section{A mixed determinant construction and a staircase gate}

\subsection*{Finding F-017: split early variables and later degree}
Let \(m\geq2\), \(n>m\), and \(\lambda\vdash nd\), with
\(r=\ell(\lambda)\). Fix \(2\leq k\leq r\) and an integer
\(s\geq1\) such that \(\lambda_k\leq s\). Write
\(R=r-k+1\). Choose nonnegative integers
\[
                   p+q=k-2,\qquad u+v=s-1,
\]
and set
\begin{equation}\label{eq:F017size}
 \mathcal S_m(p,q,u,v;R)=
 \binom{m+p}{p}\binom{R+u}{u}
 +\binom{m+q}{q}\binom{R+v}{v}-1.
\end{equation}
Then
\begin{equation}\label{eq:F017compare}
 n\geq\mathcal S_m(p,q,u,v;R)
       \quad\Longrightarrow\quad
                  q^m_\lambda(d[n])\leq K_n(\lambda).
\end{equation}
The graph establishing the bound mixes a split of the early
variables with a split of the degree in the later variables.
Its coefficient edge carries one later variable when there
is one. This saves a full degree from the two stored parts.

\paragraph{The universal determinant identity.}
Let \(L\), \(z_1,\ldots,z_{k-2}\), and
\(y_1,\ldots,y_R\) be variables. Suppose that a homogeneous
degree-\(m\) form has degree at most \(s\) in the \(y\)'s:
\[
 f=\sum_{|\alpha|+|\gamma|\leq m,\ |\gamma|\leq s}
              c_{\alpha,\gamma}z^\alpha y^\gamma
                                      L^{m-|\alpha|-|\gamma|}.
\]
There is a homogeneous linear matrix with
\(S=\mathcal S_m(p,q,u,v;R)\) vertices and determinant
\begin{equation}\label{eq:F017identity}
 \det M=
 \sum_{|\alpha|+|\gamma|\leq m,\ |\gamma|\leq s}
      c_{\alpha,\gamma}z^\alpha y^\gamma
                                      L^{S-|\alpha|-|\gamma|}.
\end{equation}
Every exponent actually occurring on the right is nonnegative.
Thus, for \(n\geq\max\{m,S\}\), adjoining
\(L I_{n-S}\) represents \(L^{n-m}f\) as an
\(n\)-by-\(n\) homogeneous linear determinant.
The coefficients of this identity are integral polynomials
in the parameters \(c_{\alpha,\gamma}\).

\begin{proof}
Split the early variables \(z\) into disjoint sets of sizes
\(p\) and \(q\). The first tree has as vertices pairs
\((\alpha_A,\gamma_A)\), where \(\alpha_A\) is a
multi-index in the first set with \(|\alpha_A|\leq m\)
and \(\gamma_A\) is a multi-index in \emph{all} \(R\)
later variables with \(|\gamma_A|\leq u\).
The second tree has the analogous pairs
\((\alpha_B,\gamma_B)\), with the other early variables
and \(|\gamma_B|\leq v\).
Use a fixed parent rule which decreases the last nonzero
coordinate of a pair. The sets are coordinatewise downward
closed, so every root path remains present and is unique.
Orient the first tree away from its root and the second
toward its root. Merge the roots and no other vertices.
In particular, the two copies of a later-variable multi-index
are distinct vertices, despite having the same variable labels.
Their total number is exactly \(S\) in \eqref{eq:F017size}.

All tree edges have labels equal to the negative of the
variable they add or remove. For a monomial with \(\gamma=0\),
insert an edge from \((\alpha_A,0)\) to
\((\alpha_B,0)\), labeled \(c_{\alpha,0}L\).
For \(\gamma\ne0\), choose one \(y_j\) with
\(\gamma_j>0\) by a fixed rule. Split
\[
              \gamma-e_j=\gamma_A+\gamma_B,
                     \qquad |\gamma_A|\leq u,
                            \quad |\gamma_B|\leq v.
\]
This is possible because \(|\gamma|-1\leq s-1=u+v\);
taking a fixed initial segment of its sorted variable word
gives a unique choice. Add the edge from
\((\alpha_A,\gamma_A)\) to
\((\alpha_B,\gamma_B)\), labeled
\(c_{\alpha,\gamma}y_j\).
Parallel labels are added in the relevant matrix entry.
Put an \(L\)-loop at every noncentral vertex. The central
diagonal consists of its coefficient-loop labels only.

Every non-loop cycle passes through the central vertex and
uses exactly one coefficient edge. Hence a cycle cover chooses
one coefficient-edge summand and its two unique tree paths,
with loops at all unused vertices.
For \(\gamma=0\), that cycle has \(|\alpha|+1\) edges,
of which \(|\alpha|\) are negative tree edges. The
permutation sign cancels those signs and gives
\(c_{\alpha,0}z^\alpha L^{S-|\alpha|}\).
For \(\gamma\ne0\), the paths have
\(|\alpha|+|\gamma|-1\) tree edges and the cycle has
\(|\alpha|+|\gamma|\) edges. Its sign again cancels
the tree signs; the coefficient edge supplies the removed
\(y_j\). The contribution is
\(c_{\alpha,\gamma}z^\alpha y^\gamma
L^{S-|\alpha|-|\gamma|}\).
This includes a central coefficient loop when the paths
are empty. Since each cycle uses distinct vertices, all
displayed exponents are nonnegative. Summing proves
\eqref{eq:F017identity} and the claimed padding identity.
Writing the result term by term avoids an unjustified
negative power when \(S<m\).
\end{proof}

\begin{proof}[Proof of the actual multiplicity comparison]
If \(q^m_\lambda=0\), the assertion is immediate. Otherwise
\(r\leq\min\{d,m^2\}\), so the natural highest-weight
space is supported on an \(r\)-dimensional coordinate subspace.
Let \(w\) be any highest-weight vector of weight \(\lambda\)
in the determinant ideal. A restricted generic translate
of the internally padded permanent is \(L^{n-m}f\),
where \(f\) has degree \(m\). On the nonempty open set
\(L(e_1)\ne0\), the upper-unipotent change of variables
from M-009 aligns \(L\) with \(x_1\) without changing
the vanishing of \(w\).

The support theorem F-016 says that no coefficient generator
used by \(w\) has total exponent greater than \(\lambda_k\)
in \(x_k,\ldots,x_r\). Delete all monomials of \(f\)
with larger degree in those variables and call the result
\(f_{\leq\lambda_k}\). Padding changes only the exponent
of \(x_1\), so
\[
       w(x_1^{n-m}f)=w(x_1^{n-m}f_{\leq\lambda_k}).
\]
The retained form has later-variable degree at most
\(\lambda_k\leq s\). Apply \eqref{eq:F017identity}
with \(z=(x_2,\ldots,x_{k-1})\) and
\(y=(x_k,\ldots,x_r)\). Under the size hypothesis, its
padded form is a determinant and \(w\) vanishes on it.
Thus every such \(w\) vanishes on a dense open subset
of the permanent orbit and hence on its reduced closure.
The highest-weight ideal multiplicity spaces satisfy
\(J_{\Omega_n,\lambda}\subseteq J_{Z_{n,m},\lambda}\).
This is an inclusion for every copy; taking quotient dimensions
proves \eqref{eq:F017compare}.
\end{proof}

\paragraph{A staircase in the BIP numerical range.}
Suppose \(n\geq m^{25}\). Every actual multiplicity
obstruction satisfies the twenty-five inequalities
\begin{equation}\label{eq:F017BIP}
                \lambda_{54-2s}\geq s+1
                               \qquad(1\leq s\leq25).
\end{equation}
In particular,
\[
       \lambda_4\geq26,\quad\lambda_6\geq25,\quad
       \lambda_8\geq24,\quad\ldots,\quad\lambda_{52}\geq2.
\]
Odd-indexed consequences follow from monotonicity of the rows.
For example \(\lambda_3\geq26\) and \(\lambda_{27}\geq14\).

\begin{proof}
By M-008 an actual obstruction already has \(r\geq52\),
so every index \(k=54-2s\) in \eqref{eq:F017BIP} exists.
Suppose \(\lambda_k\leq s\). Choose
\(u=\lfloor(s-1)/2\rfloor\),
\(v=\lceil(s-1)/2\rceil\), and
\[
                     p=25-2u,\qquad q=25-2v.
\]
Then \(0\leq u,v\leq12\), \(p,q\geq1\), and
\(p+q=k-2\). Each term of \eqref{eq:F017size}
has total power budget \(p+2u=q+2v=25\).

Here are uniform estimates including all small parameters.
For \(a\geq0\), \(m\geq2\), the product formula gives
\[
 \frac{\binom{m+a}{a}}{m^a}
       \leq A_a:=\frac{\binom{a+2}{2}}{2^a}.
\]
Since \(R\leq r-1\) and \(r\geq2\), similarly
\[
 \frac{\binom{R+j}{j}}{r^j}
       \leq\frac{\binom{r-1+j}{j}}{r^j}
       \leq B_j:=\frac{j+1}{2^j}.
\]
Both inequalities follow by observing that each factor
\(1+i/m\), or \(1+i/r\), decreases as its denominator
increases. The formulas include \(a=0\) and \(j=0\).
The sequence \(A_a\) is at most \(3/2\) and decreases
for \(a\geq1\); \(B_j\leq1\) and decreases for \(j\geq1\).
For \(0\leq j\leq12\), with \(a=25-2j\), this gives
\[
 A_a B_j\leq
 \begin{cases}
 55/512,&j\leq8,\quad a\geq9,\\
 15/512,&j\geq9,
 \end{cases}
\]
because \(A_9=55/512\) and \(B_9=10/512\).
Using \(r\leq m^2\), each term in the size formula is
therefore at most \((55/512)m^{25}\), and
\[
                 \mathcal S_m
                   \leq\frac{110}{512}m^{25}-1<m^{25}\leq n.
\]
The comparison theorem forbids the alleged obstruction,
which proves \eqref{eq:F017BIP}.
\end{proof}

\paragraph{A staircase in the IP numerical range.}
If \(n>3m^4\), every actual multiplicity obstruction satisfies
\begin{equation}\label{eq:F017IP}
              \lambda_{12-2s}\geq s+1
                                     \qquad(1\leq s\leq5),
\end{equation}
or explicitly
\[
       \lambda_2\geq6,\quad\lambda_4\geq5,\quad
       \lambda_6\geq4,\quad\lambda_8\geq3,\quad\lambda_{10}\geq2.
\]

\begin{proof}
M-008 gives \(r\geq10\) for an actual obstruction.
If \(k=12-2s\) and \(\lambda_k\leq s\), choose
\(u=\lfloor(s-1)/2\rfloor\),
\(v=\lceil(s-1)/2\rceil\),
\(p=4-2u\), and \(q=4-2v\).
These are nonnegative and satisfy all required sums,
with \(p+2u=q+2v=4\).
The estimates just proved give \(A_p,A_q\leq3/2\)
and \(B_u,B_v\leq1\). Thus
\[
                   \mathcal S_m\leq3m^4-1<n.
\]
Again \eqref{eq:F017compare} excludes the obstruction.
\end{proof}

\paragraph{The corresponding full ambient equality.}
Assume \(r\leq\min\{d,n^2\}\). For \(b=\lambda_2\), replace \(m\) by \(b\) in
\eqref{eq:F017size}, take \(s\geq\lambda_k\), and assume
\(n\geq\max\{b,\mathcal S_b\}\).
Then the same coefficient-space construction, F-014, and
the restriction-and-padding proof of F-015 give
\[
                         K_n(\lambda)=a_\lambda(d[n]).
\]
Indeed all stabilized highest-weight vectors use only the
degree-\(b\) coefficient space with later-variable degree
at most \(\lambda_k\). The differential padding composition
is diagonal and invertible on this monomial space, so every
ambient copy remains independent. This full-equality variant
has a different size parameter from the direct degree-\(m\)
comparison above. The staircase conclusions use the latter
and make no assumption that \(n\geq\lambda_2\).

\paragraph{Dictionary and remaining gate.}
The coefficient support is detected by exterior matrix
coefficients of Schur functors and the torus grading
(F-016; Macdonald I(4.3), Appendix~A(5.3)--(5.6);
SGA~3 I4.7.3). The two tree copies are auxiliary affine
matrix coordinates, and the coefficient edges encode the
universal linear section of the form bundle. Their determinant
identity is integral in the universal coefficients. Passing
from highest-weight vanishing to the coherent ideal inclusion
uses the characteristic-zero comodule and Schur decomposition.
The inequalities constrain the actual multiplicity test
\(q>K\), at every outer degree, and not only the stronger
sufficient ordinary-Kronecker test. Outside these size bounds,
no evaluation-rank comparison follows from this construction.

% END agents/formulas/mixed_tail_graph.tex

% BEGIN agents/formulas/support_cover.tex
\section{A finite support-cover criterion}

\subsection*{Finding F-018: two monomial order ideals and one linear edge}
Let \(\Gamma\subseteq\mathbb Z_{\geq0}^{r-1}\) be a
nonempty finite coordinatewise downward-closed set. Write
\(E=\{0,e_2,\ldots,e_r\}\). Call \((P,Q)\) admissible
if \(P,Q\subseteq\Gamma\) are nonempty downward-closed
sets and \(\Gamma\subseteq P+Q+E\). Define
\begin{equation}\label{eq:F018kappa}
 \kappa(\Gamma)=\min_{(P,Q)\ \mathrm{admissible}}
                                      \bigl(|P|+|Q|-1\bigr).
\end{equation}
Here the sum is the set of sums of multi-indices. The minimum
exists, since \(P=\Gamma,Q=\{0\}\) is admissible, and
\[
                 \max_{\gamma\in\Gamma}|\gamma|
                         \leq\kappa(\Gamma)\leq|\Gamma|.
\]
This is a finite optimization problem for this particular
graph construction. No polynomial-time algorithm or optimality
among all linear determinant representations is asserted.

\paragraph{Universal identity.}
For any integer \(h\geq\max_{\gamma\in\Gamma}|\gamma|\),
put
\[
               f(L,z)=\sum_{\gamma\in\Gamma}
                         c_\gamma z^\gamma L^{h-|\gamma|}.
\]
Every admissible pair \(P,Q\) yields a homogeneous linear
matrix with \(S=|P|+|Q|-1\) vertices and
\begin{equation}\label{eq:F018det}
                 \det M=\sum_{\gamma\in\Gamma}
                         c_\gamma z^\gamma L^{S-|\gamma|}.
\end{equation}
In particular, if \(n\geq\max\{h,\kappa(\Gamma)\}\),
every form \(L^{n-h}f\) in this coefficient space is an
\(n\)-by-\(n\) homogeneous linear determinant.

\begin{proof}
Choose one decomposition
\(\gamma=\alpha_\gamma+\beta_\gamma+\epsilon_\gamma\)
with \(\alpha_\gamma\in P\), \(\beta_\gamma\in Q\), and
\(\epsilon_\gamma\in E\), for every \(\gamma\).
Make the two monomial trees on \(P\) and \(Q\), using
the parent rule which deletes the last nonzero coordinate.
Downward closure guarantees all parents and root paths.
Orient the first tree away from its root and the second
toward its root, merge the two roots, and keep all other
vertices distinct. Label tree edges by the negative of the
corresponding variable.

Insert one coefficient-edge summand from \(\alpha_\gamma\)
to \(\beta_\gamma\), labeled \(c_\gamma L\) if
\(\epsilon_\gamma=0\), or \(c_\gamma z_i\) if
\(\epsilon_\gamma=e_i\). Add parallel summands in their
matrix entry. Put an \(L\)-loop at every noncentral
vertex; the central diagonal has only its coefficient labels.
Every non-loop cycle contains the central vertex and exactly
one coefficient edge. Its two paths are unique.

For \(\epsilon_\gamma=0\), the cycle has
\(|\gamma|+1\) vertices; for a unit vector it has
\(|\gamma|\) vertices. The determinant sign cancels the
negative tree-edge signs in both cases. The chosen coefficient
edge and the unused loops give precisely
\(c_\gamma z^\gamma L^{S-|\gamma|}\). This includes
central coefficient loops. The vertices of a cycle are
distinct, so all exponents are nonnegative. This proves
\eqref{eq:F018det} and the bound \(S\geq|\gamma|\)
for each \(\gamma\). One can also see the latter directly
from \(|P|\geq|\alpha_\gamma|+1\) and
\(|Q|\geq|\beta_\gamma|+1\).
Padding by \(L I_{n-S}\) gives the assertion at size \(n\).
All identities are polynomial identities over the integers
in the coefficient parameters.
\end{proof}

\paragraph{Two actual multiplicity criteria.}
For a partition \(\lambda\) of length \(r\geq2\), define
\[
 \Gamma_{h,\lambda}=
 \left\{\gamma\in\mathbb Z_{\geq0}^{r-1}:|\gamma|\leq h,
      \quad\sum_{i=j}^r\gamma_i\leq\lambda_j
                                      \ (2\leq j\leq r)\right\}.
\]
For internal padding, \(n>m\geq2\), and
\(\lambda\vdash nd\), one has
\begin{equation}\label{eq:F018q}
 n\geq\kappa(\Gamma_{m,\lambda})
       \quad\Longrightarrow\quad q^m_\lambda(d[n])\leq K_n(\lambda).
\end{equation}
If \(r\leq\min\{d,n^2\}\) and \(b=\lambda_2\), then
\begin{equation}\label{eq:F018K}
 n\geq\kappa(\Gamma_{b,\lambda})
       \quad\Longrightarrow\quad K_n(\lambda)=a_\lambda(d[n]).
\end{equation}
The first comparison uses the permanent degree; the second
uses the stabilized inner degree. They must not be interchanged.

\begin{proof}
For \eqref{eq:F018q}, discard the immediate case \(q=0\).
Positive \(q\) implies \(r\leq\min\{d,m^2\}\).
Take any determinant-ideal highest-weight vector of weight
\(\lambda\). Align the padding linear form with the first
coordinate on the dense open set used in M-009. The support
theorem F-016 permits deleting exactly the coefficient monomials
outside \(\Gamma_{m,\lambda}\) without changing its value.
The universal identity above makes the retained padded form
a determinant at size \(n\). Consequently that vector
vanishes on the permanent orbit and its reduced closure.
This holds for the whole highest-weight space, so it gives
the ideal multiplicity-space inclusion and the inequality
in \eqref{eq:F018q}, as in F-017.

For \eqref{eq:F018K}, the set \(\Gamma_{b,\lambda}\)
contains the pure second-variable power of degree \(b\).
Hence \(\kappa(\Gamma_{b,\lambda})\geq b\), and
the hypothesis implies \(n\geq b\).
F-014 stabilizes the ambient multiplicity at
\(\mu=(db-|\bar\lambda|,\bar\lambda)\vdash db\).
F-016 puts the entire highest-weight space of weight \(\mu\)
in the coefficient algebra of \(\Gamma_{b,\lambda}\).
Its dual padded coefficient space is contained in the
determinant closure by \eqref{eq:F018det}. The diagonal,
invertible padding pullback of F-015 preserves the independence
of all copies. Thus \(K_n(\lambda)\geq a_\lambda(d[n])\);
the reverse ambient quotient inequality proves equality.
\end{proof}

The preceding exact two-tree constructions give special choices of
the two order ideals and the coefficient-edge decomposition
after intersecting their vertex sets with \(\Gamma\), if
necessary. Every used endpoint and its parents are coordinatewise
bounded by their supported monomial, so this pruning retains
the required decompositions and can only reduce the size.
In applying \eqref{eq:F018q}, a pair may encode the entire
nested profile, instead of replacing it by just one row cut.
The elementary bound \(\kappa(\Gamma)\leq|\Gamma|\)
already becomes useful when joint terminal constraints make
the full support much smaller than every one-cut enlargement.

\paragraph{The permanent version preserves its distinguished padding entry.}
For \(r\leq\min\{d,m^2\}\), \(b=\lambda_2\), and
\(n>m\geq2\), there is also the exact statement
\begin{equation}\label{eq:F018per}
 m\geq\kappa(\Gamma_{b,\lambda})
       \quad\Longrightarrow\quad
          q^m_\lambda(d[n])=K_n(\lambda)=a_\lambda(d[n]).
\end{equation}
This determines full permanent-side multiplicities in additional
families. It is not a new determinant-side exclusion when
\eqref{eq:F018K} already applies.

\begin{proof}
In the graph for an admissible pair, replace each negative
tree-variable label by its positive counterpart. A permanent
cycle cover has no permutation sign, so its contribution is
now exactly the same monomial as in \eqref{eq:F018det}.
Thus the resulting matrix \(A\) satisfies
\[
                 \operatorname{perm}A
                    =\sum_{\gamma\in\Gamma}
                              c_\gamma z^\gamma L^{S-|\gamma|}.
\]
If \(S\geq2\), every noncentral diagonal entry is exactly
\(L\); coefficient edges never create a noncentral diagonal
because the two tree copies are distinct. A simultaneous
row-and-column permutation moves one such entry to position
\((1,1)\), without changing the permanent. Adjoin
\(L I_{m-S}\) to reach size \(m\).
If \(S=1\), first adjoin one \(L\)-diagonal entry and
move it to position \((1,1)\); this is possible because
\(m\geq2\). The remaining padding is the same.

Apply this to \(\Gamma=\Gamma_{b,\lambda}\), so that
\(b\leq S\leq m\). For every form in its degree-\(b\)
coefficient space the constructed matrix has
\[
                 A_{11}=L,\qquad
                 \operatorname{perm}A=L^{m-b}f.
\]
Substitution in the internally padded permanent gives
\(A_{11}^{n-m}\operatorname{perm}A=L^{n-b}f\).
Every such substitution lies in its orbit closure by density
of invertible endomorphisms. The same all-copy stabilization
and invertible padding restriction used above now give
\(q^m_\lambda\geq a_\lambda\). The ambient quotient
bound forces equality. Since \(n>m\geq\kappa\),
\eqref{eq:F018K} also gives \(K_n=a_\lambda\).
\end{proof}

\paragraph{Scheme and lambda-operation interpretation.}
The order ideals are finite monomial supports in a torus-graded
symmetric-power module, under SGA~3 I4.7.3. Their coefficient
schemes map to schemes of linear matrices over \(\mathbb Z\),
with determinant and permanent identities just proved.
The highest-weight support is supplied by exterior minors
and Cauchy/Schur functors (Macdonald I(4.3), Appendix~A(5.3)--(5.6);
SGA~6 VI3.3). The coordinate-ring conclusions retain the
specific quotient and evaluation maps: this is exactly the
additional data missing from a bare effective lambda class
in F-010. Characteristic-zero factorial invertibility and
semisimplicity are used only in the all-copy lift. The graph
identities themselves are integral. No saturation or general
positive effective Burnside lift is deduced.

% END agents/formulas/support_cover.tex

% BEGIN agents/archive/JOINT_SUPPORT_COVER.tex
% Run GCT_KRON_20261120; GPT-6 Astra Pro; effort not exposed.
\section{A two-tree improvement of the joint later-row gate}

\subsection*{A-025. Simultaneously using the third and fifty-second rows}
In the BIP numerical range $n\geq m^{25}$, every actual
internally padded permanent multiplicity obstruction satisfies
\[
                   \boxed{\quad\lambda_3\geq557
                                  \quad\text{or}\quad\lambda_{52}\geq3.\quad}
                                                         \tag{A-025.1}
\]
This strengthens the sufficient constant $139$ in O-018.a.
The new construction was proposed by the root worker, and
the covering proof and exact uniform estimate were completed
here. No optimality among determinant representations is claimed.

\paragraph{A general joint-support cover.}
Assume $r\geq52$, and write $R=r-51$. Separate the
nonfirst variables into $x_2$, the forty-nine variables
$x_3,\ldots,x_{51}$, and the $R$ variables
$x_{52},\ldots,x_r$, called the early, middle, and last
groups in this proof. Suppose a coordinatewise downward-closed
support $\Gamma$ has total degree at most $m$, middle-plus-last
degree at most $B$, and last degree at most two. Other nested
conditions may be present. Let $u,v\geq0$ satisfy
\[
                              u+v=B-1,\qquad B\geq1.
\]
Define the following two monomial order ideals, intersected
with $\Gamma$:
\[
 \begin{split}
 P&=\Gamma\cap\{\text{early degree}\leq m,
                     \ \text{middle degree}\leq u,
                     \ \text{last degree}=0\},\\
 Q&=\Gamma\cap\{\text{early degree}=0,
                     \ \text{middle degree}\leq v,
                     \ \text{last degree}\leq1\}.
 \end{split}                                           \tag{A-025.2}
\]
Then $(P,Q)$ is admissible for F-018: with
$E=\{0,e_2,\ldots,e_r\}$ one has $\Gamma\subseteq P+Q+E$.
In particular
\[
 \kappa(\Gamma)\leq
       (m+1)\binom{49+u}{49}
          +(r-50)\binom{49+v}{49}-1.                  \tag{A-025.3}
\]

\begin{proof}
Take $\gamma\in\Gamma$. If its last degree is one or two,
remove one last-variable unit vector and put it on the
coefficient edge. Put the remaining last degree, at most one,
entirely in $Q$. The middle degree is then at most $B-1$.
If the last degree is zero but the middle degree is positive,
remove one middle-variable unit vector for the coefficient edge;
again the remaining middle degree is at most $B-1=u+v$.
Split this remaining middle multiset into degrees at most
$u$ and at most $v$, and put the two parts in $P,Q$,
respectively. Put every early-variable factor in $P$.
If both middle and last degrees are zero, put the whole early
monomial in $P$ and use $Q=0$ and edge exponent zero.

Each chosen endpoint is coordinatewise at most $\gamma$.
Downward closure therefore keeps both endpoints in $\Gamma$;
it also keeps every parent on their monomial-tree paths.
All requirements in (A-025.2) hold. This proves the cover.

Ignoring the intersections with $\Gamma$ can only enlarge
the cardinalities. Stars and bars gives
$|P|\leq(m+1)\binom{49+u}{49}$ and
$|Q|\leq(R+1)\binom{49+v}{49}$, since the last part
of $Q$ is either zero or one of its $R$ variables.
F-018 uses $|P|+|Q|-1$ vertices, proving (A-025.3).
The coefficient edge carries precisely the removed variable,
or $x_1$ when its exponent is zero. Thus the determinant
identity is the integral identity of F-018, with no change
to its cycle signs or padding exponents.
\end{proof}

\begin{proof}[A first uniform numerical bound]
The earlier actual-obstruction reductions used in O-018.a
give $m\geq190$ and $52\leq r\leq m^2$. Suppose
$\lambda_3\leq419$ and $\lambda_{52}\leq2$.
For the full nested support $\Gamma_{m,\lambda}$ of F-018,
the preceding cover applies with
\[
                         B=419,\qquad u=222,\qquad v=196.
\]
Write $C=\binom{271}{49}$ and $D=\binom{245}{49}$.
By (A-025.3), its determinant size is strictly less than
\[
                            (m+1)C+m^2D.
\]
The normalized upper bound is
\[
              \frac{(m+1)C+m^2D}{m^{25}}
                        =C(m^{-24}+m^{-25})+Dm^{-23}.
\]
Every term decreases as $m$ increases through positive real
values. At $m=190$, exact integer arithmetic gives
\[
       50\left(191\binom{271}{49}
                    +36100\binom{245}{49}\right)<49\cdot190^{25}.
                                                         \tag{A-025.4}
\]
The saved certificate expands both integers and their positive
difference; the binomials are computed by exact factorial
products, not floating point. Therefore for every $m\geq190$
the cover size is less than $(49/50)m^{25}<n$.

F-018 now gives $q^m_\lambda(d[n])\leq K_n(\lambda)$,
contradicting an actual obstruction. More explicitly, F-016
puts every coefficient generator of a determinant-ideal
highest-weight vector inside the nested support after aligning
the padding form. The cover represents every retained padded
form as a determinant, so every such vector vanishes on the
permanent closure. This includes every ideal multiplicity copy.
Thus the contradiction concerns the actual quotient dimensions.
It proves the weaker bound $\lambda_3\geq420$ or
$\lambda_{52}\geq3$. Retaining the total-degree constraint
below improves this to (A-025.1).
\end{proof}

\paragraph{Exact counts after total-degree pruning.}
The intersections in (A-025.2) retain total degree at most $m$.
Let $t=\min(m,u)$ and define
\[
 \begin{split}
 A(m,u)&=\sum_{j=0}^t\binom{48+j}{j}(m-j+1)\\
       &=(m+1)\binom{49+t}{49}-49\binom{49+t}{50},\\
 Q(m,v)&=\binom{49+\min(v,m)}{49}
        +(m^2-51)\binom{49+\min(v,m-1)}{49}.
 \end{split}                                             \tag{A-025.5}
\]
For $m\geq190$ and $r\leq m^2$, these give
\[
               \kappa(\Gamma)\leq A(m,u)+Q(m,v)-1.        \tag{A-025.6}
\]
Indeed a middle monomial of degree $j$ in $P$ has
$m-j+1$ possible early exponents, and there are
$\binom{48+j}{j}$ such middle monomials. The binomial sum
uses the identities
\[
 \sum_{j=0}^t\binom{48+j}{48}=\binom{49+t}{49},\qquad
 \sum_{j=0}^t j\binom{48+j}{48}=49\binom{49+t}{50},
\]
which follow from Pascal's identity and
$j\binom{48+j}{48}=49\binom{48+j}{49}$.
For $Q$, last degree zero allows middle degree through
$\min(v,m)$; last degree one has $R\leq m^2-51$
possible variables and allows middle degree through
$\min(v,m-1)$. This proves (A-025.5)--(A-025.6).
Other nested restrictions can only decrease these counts.

\begin{proof}[A first total-degree-pruned bound]
Suppose $\lambda_3\leq538$ and $\lambda_{52}\leq2$.
Use the same covering proof with
\[
                       B=538,\qquad u=325,\qquad v=212.
\]
These satisfy $u+v=B-1$. Set $T(m)=A(m,325)+Q(m,212)$.
The saved exact certificate verifies all 142 integer cases
\[
                  100T(m)<99m^{25}
                               \qquad(190\leq m\leq331).       \tag{A-025.7}
\]
Each record contains $m$, $T(m)$, $m^{25}$, and the positive
integer $99m^{25}-100T(m)$. The maximum ratio in this
finite range occurs at $m=331$, also checked by exact
cross multiplication. Every acceptance test in the certificate
uses integer arithmetic.

For the infinite tail $m\geq331$, write
$C'=\binom{374}{49}$ and $D'=\binom{261}{49}$.
The caps have then stabilized in (A-025.5), giving
\[
              T(m)=C'(m-635/2)+D'(m^2-50).
\]
For every real $m\geq331$, differentiation yields
\[
 2m^{26}\frac{d}{dm}\left(\frac{T(m)}{m^{25}}\right)
       =C'(15875-48m)+D'(2500-46m^2)<0.                 \tag{A-025.8}
\]
Both coefficients on the right are strictly negative there.
Thus the normalized count decreases on the whole tail,
and (A-025.7) at $331$ proves $T(m)<(99/100)m^{25}$
for every integer $m\geq331$ as well.

Combining the finite certificate and tail estimate shows
that the admissible cover has size less than $m^{25}\leq n$
for every allowed $m$. The same F-018 ideal inclusion used
in the first bound gives $q^m_\lambda\leq K_n(\lambda)$.
This contradiction proves $\lambda_3\geq539$ or
$\lambda_{52}\geq3$.
\end{proof}

\paragraph{Allocating the early variable to the other tree.}
There is a better cover in the same family. Put all early
degree in $Q$ instead of $P$. Thus $P$ has only middle
variables of degree at most $u$, while $Q$ has unrestricted
early degree, middle degree at most $v$, and last degree
at most one; both sets are intersected with $\Gamma$.
The covering proof is identical, with the early monomial
assigned to $Q$. Its total-degree-pruned bound is
\[
 \begin{split}
 \kappa(\Gamma)&\leq T_*(m;u,v)-1,\\
 T_*(m;u,v)&=\binom{49+\min(u,m)}{49}
                 +A(m,v)+(m^2-51)A(m-1,v),
 \end{split}                                            \tag{A-025.9}
\]
where $A$ is defined in (A-025.5). Last degree zero gives
the middle-plus-early count $A(m,v)$; last degree one
uses one of at most $m^2-51$ last variables and leaves
total degree $m-1$ for the other two groups.

\begin{proof}[Proof of (A-025.1)]
Suppose $\lambda_3\leq556$ and $\lambda_{52}\leq2$.
Use (A-025.9) with $u=360,v=195$, so $u+v=556-1$.
The certificate verifies all 171 integer cases
\[
             25T_*(m;360,195)<24m^{25}
                              \qquad(190\leq m\leq360).        \tag{A-025.10}
\]
The finite maximum ratio occurs at $m=360$; every
comparison is made by exact integer cross multiplication.

For real $m\geq360$, write $C''=\binom{409}{49}$,
$D''=\binom{244}{49}$, and $c=1911/10=49\cdot195/50$.
The normalized count is the sum of
\[
 \frac{C''}{m^{25}},\qquad
 \frac{D''(m+1-c)}{m^{25}},\qquad
 \frac{D''(m^2-51)(m-c)}{m^{25}}.                       \tag{A-025.11}
\]
Each is positive and strictly decreasing in this range.
For the last term its logarithmic derivative is
\[
 \frac{2m}{m^2-51}+\frac1{m-c}-\frac{25}{m}
       \leq\frac3m+\frac3m-\frac{25}{m}<0.
\]
Here $m^2\geq153$ implies the first bound, and
$m\geq360>3c/2$ implies $m-c\geq m/3$.
For the middle term the same bound gives logarithmic
derivative $1/(m+1-c)-25/m\leq3/m-25/m<0$.
The first term is plainly decreasing. Thus (A-025.10)
at $360$ proves the required bound for all larger integers.

The cover size is consequently less than $m^{25}\leq n$
for every $m\geq190$. F-018 again preserves all highest-weight
ideal copies and gives $q^m_\lambda\leq K_n(\lambda)$,
contradicting an actual obstruction. This proves (A-025.1).
\end{proof}

\paragraph{Verification and scope.}
The accompanying exact-integer script
\texttt{check\_joint\_support.py} and certificate
\texttt{JOINT\_SUPPORT\_CERTIFICATE.json} verify (A-025.4)
and every exact integer case in (A-025.7) and (A-025.10).
They also check all splits for the coarse bound at $m=190$:
$B=419$ has successful splits, while $B=420$ has none.
This last finite calculation is only about (A-025.3);
it is not an optimality assertion for the pruned support,
the full F-018 optimization, or arbitrary determinants.

The support uses the torus grading and exterior-minor argument
of F-016, with SGA 3, I, Proposition 4.7.3, pp.~35--36,
and Schur/exterior functors from SGA 6, VI, Theorem 3.3,
pp.~392--393. The determinant morphism and coefficient
identities are defined over $\mathbb Z$; the actual
highest-weight ideal comparison is in characteristic zero.
The argument uses two row constraints simultaneously and
retains the evaluation maps. It does not deduce the comparison
from a virtual lambda-ring identity or a support-saturation claim.

% END agents/archive/JOINT_SUPPORT_COVER.tex

% BEGIN proofs/nested_tail_bound.tex
\section{Later rows control a stronger multiplicity comparison}

\subsection*{Finding M-011: a cut through the nested tail support}
Let \(m\geq2\), \(n>m\), and let \(\lambda\vdash nd\)
have \(r\) rows. For \(2\leq k\leq r+1\), put
\(\lambda_{r+1}=0\), \(t=\min\{m,\lambda_k\}\), and
\[
 C_{m,k}(\lambda)=
 \binom{m+k-2}{k-2}
 +\binom{r-k+1+t}{t}-1.
\]
Then
\begin{equation}\label{eq:M011}
 n\geq C_{m,k}(\lambda)
 \quad\Longrightarrow\quad
                  q^m_\lambda(d[n])\leq K_n(\lambda).
\end{equation}
The multiplicities are those of the actual reduced orbit
closures, with the internal-padding convention. If the ambient
multiplicity is zero, the conclusion is immediate; otherwise
\(r\leq d\). The one-row case also has no obstruction.

Two uniform consequences, at every outer degree, are
\begin{align}
 n>3m^4, q^m_\lambda>K_n(\lambda)
       &\quad\Longrightarrow\quad \lambda_6\geq3,
                                                    \label{eq:M011IP}\\
 n\geq m^{25},\ q^m_\lambda>K_n(\lambda)
       &\quad\Longrightarrow\quad \lambda_{27}\geq13.
                                                    \label{eq:M011BIP}
\end{align}
These constrain later rows even when the second row is larger
than \(n\). They therefore apply beyond the stabilization
criteria that require \(n\geq\lambda_2\).

\paragraph{Support lemma.}
Every highest-weight vector of weight \(\lambda\) in
\(\operatorname{Sym}^d\operatorname{Sym}^n\mathbb C^r\)
uses only coefficient generators indexed by \(\eta\) with
\[
        |\eta|=n,\qquad
        \sum_{i=j}^r\eta_i\leq\lambda_j
                       \quad(2\leq j\leq r).          \tag{M-011.1}
\]
Here ``uses'' refers to every generator in every coefficient
monomial of the polynomial, not only the sum of their weights.

To prove this, realize the symmetric power as the column-permutation
invariants in the balanced column-degree \((n,\ldots,n)\)
polynomials on an \(r\)-by-\(d\) matrix. By the complete
flag-minor spanning argument in F-014, its row highest-weight
space is spanned by products of flag minors whose heights are
the column lengths of \(\lambda\). Exactly \(\lambda_j\)
minor factors have height at least \(j\). In any fixed column,
a monomial of a single minor uses at most one variable; a minor
of height less than \(j\) uses none from rows \(j,\ldots,r\).
Thus each matrix monomial has total degree at most \(\lambda_j\)
in these rows of that column. Taking linear combinations cannot
introduce a monomial that violates the bound. Finally, the
averaging map from \(\operatorname{Sym}^d\operatorname{Sym}^n\)
sends a monomial basis to disjoint orbit sums of column monomials.
The condition holds in every column and is permutation invariant.
It therefore implies precisely (M-011.1) on each coefficient
generator. This also proves the support assertion for every
multiplicity copy, as in F-016.

\begin{proof}[Proof of the comparison]
If \(q^m_\lambda=0\), the conclusion is immediate. Otherwise
the permanent length bound gives \(r\leq m^2<n^2\), so
the coordinate subspace used below exists.
Take a highest-weight vector \(w\) of weight \(\lambda\) in
\(I(\Omega_n)_d\). It depends only on the first \(r\)
coordinates. A general restricted translate of the padded
permanent is \(L^{n-m}f\), with \(f\) of degree \(m\).
On the dense open set \(L(e_1)\ne0\), the upper-unipotent
change of basis in M-009 aligns \(L\) with \(x_1\) and
preserves the vanishing of \(w\). Hence it is enough to
consider \(x_1^{n-m}f\).

Expand \(f\) in monomials
\(x_1^{m-|\alpha|}x_2^{\alpha_2}\cdots x_r^{\alpha_r}\).
Delete every term violating one of the inequalities
\(\sum_{i=j}^r\alpha_i\leq\lambda_j\), \(j\geq2\),
and denote the resulting form by \(f_{\mathrm{tail}}\).
Padding changes only the exponent of \(x_1\), so (M-011.1)
gives the exact equality
\[
 w(x_1^{n-m}f)=w(x_1^{n-m}f_{\mathrm{tail}}).
\]
Split the other variables into the prefix
\(x_2,\ldots,x_{k-1}\) and the suffix \(x_k,\ldots,x_r\).
All retained monomials have prefix degree at most \(m\)
and suffix degree at most \(t=\min(m,\lambda_k)\).

Use the variable-split graph of M-007 with all prefix
multi-indices of degree at most \(m\) and all suffix
multi-indices of degree at most \(t\). These sets are
downward closed. Identify the two roots, insert coefficient
edges only for monomials actually retained in
\(f_{\mathrm{tail}}\), and put the \(x_1\)-loops and
negative variable-edge labels as before. The number of vertices
is \(S=C_{m,k}(\lambda)\). Every nontrivial cycle has one
coefficient edge and the two unique tree paths. The cycle
permutation sign cancels the signs on those variable edges.
A monomial with tail degree \(|\alpha|\) contributes
\(c_\alpha x^\alpha x_1^{S-|\alpha|}\).
The relevant cycle has \(|\alpha|+1\) distinct vertices,
so this exponent is nonnegative, including when \(S<m\).
Adjoining \(x_1 I_{n-S}\) gives a determinant equal to
\(x_1^{n-m}f_{\mathrm{tail}}\) whenever \(n\geq S\).

The vector \(w\) vanishes on this determinant and hence
on the dense open subset of the permanent orbit just used.
It vanishes on the entire reduced orbit closure. This holds
for every highest-weight vector in the determinant ideal,
so the ideal multiplicity spaces satisfy
\(J_{\Omega_n,\lambda}\subseteq J_{Z_{n,m},\lambda}\).
Subtracting their dimensions from the common ambient
multiplicity proves \eqref{eq:M011}.
\end{proof}

\begin{proof}[Proof of the two uniform consequences]
A positive permanent multiplicity has \(r\leq m^2\).
In the IP range, M-008 already excludes \(r\leq9\), so
we may use \(k=6\). If \(\lambda_6\leq2\), then
\[
 C_{m,6}(\lambda)
 \leq\binom{m+4}{4}+\binom{r-3}{2}-1
 \leq m^4+r^2-1<3m^4.
\]
The middle bound uses \(m\geq2\): the ratio
\(\binom{m+4}{4}/m^4=\prod_{i=1}^4(1+i/m)/4!\)
decreases with \(m\) and is \(15/16\) at \(m=2\).
The comparison theorem therefore forbids an obstruction.

In the BIP range, M-008 already excludes \(r\leq51\),
so use \(k=27\). If \(\lambda_{27}\leq12\),
\[
 \begin{split}
 C_{m,27}(\lambda)
 &\leq\binom{m+25}{25}+\binom{r-14}{12}-1\\
 &\leq\tfrac12m^{25}+m^{24}-1<m^{25}.
 \end{split}
\]
For the first term,
\(\binom{m+25}{25}/m^{25}\) decreases with \(m\),
and its value at \(m=2\) is \(351/2^{25}<1/2\).
For the second term, the monomials of degree at most twelve
in \(r-26\) variables inject into words of length twelve
on an alphabet of \(r-25\leq r\) letters, after inserting
a slack variable. Their number is at most \(r^{12}\leq m^{24}\).
Finally \(m^{24}\leq m^{25}/2\). This proves
\eqref{eq:M011BIP}.
\end{proof}

\paragraph{A general power-range consequence.}
For every integer \(c\geq6\), the same proof gives
\[
 n\geq m^c,\ q^m_\lambda>K_n(\lambda)
 \quad\Longrightarrow\quad
       \lambda_{c+2}\geq\lfloor c/2\rfloor+1.          \tag{M-011.2}
\]
Indeed M-008 leaves at least \(2c+2\) rows, so take
\(k=c+2\), \(t_0=\lfloor c/2\rfloor\). If
\(\lambda_k\leq t_0\), the prefix count is at most
\(m^c/2\). This follows by the same decreasing-ratio
argument and \(\binom{c+2}{2}<2^{c-1}\) for \(c\geq6\);
the last inequality holds at six and remains true on
incrementing \(c\), since its left ratio is
\((c+3)/(c+1)<2\). The suffix count is at most
\[
 \binom{r-c-1+t_0}{t_0}
       \leq\frac{r^{t_0}}{t_0!}
       \leq\frac{m^{2t_0}}{t_0!}\leq\tfrac12m^c,
\]
since \(t_0\geq3\) and \(t_0\leq c+1\).
Their sum minus one is less than \(m^c\), proving the claim.

\paragraph{Grothendieck dictionary.}
The support lemma uses exterior-power minors of the universal
map between locally free sheaves, followed by the Cauchy
decomposition of its symmetric algebra; these are the Schur
operations of cote~133, scan~028, and SGA~6, VI,
Theorem~3.3, pp.~392--393. The suffix inequalities describe
finite torus-weight submodules, using SGA~3, I,
Proposition~4.7.3, pp.~35--36. Their affine coefficient
schemes map to the determinant locus by integral polynomial
identities. Highest-weight vanishing then gives an inclusion
of homogeneous ideal comodules on the characteristic-zero
fiber, in the representation/comodule dictionary of SGA~3,
I, Proposition~4.7.2, pp.~34--35. The resulting inequality
concerns coherent coordinate-ring quotients. It uses neither
an effective lift of an irreducible finite-group representation
to a set nor a saturation assumption.

% END proofs/nested_tail_bound.tex

% BEGIN proofs/binary_permanent_seed.tex
\section{An explicit permanent family with all ambient copies}

\subsection*{Finding M-012: binary containment and a four-copy seed}
For \(n>m\geq2\), every binary form
\(L^{n-m}f(L,y)\), with \(f\) homogeneous of degree \(m\),
belongs to the internally padded permanent orbit closure
\(Z_{n,m}\). Consequently, for \(d\geq2\) and
\(0\leq t\leq m\),
\[
 \lambda=(nd-t,t)
 \quad\Longrightarrow\quad
           q^m_\lambda(d[n])=K_n(\lambda)=a_\lambda(d[n]).
                                                               \tag{M-012.1}
\]
Zero final parts are deleted. In particular, for \(m\geq6\)
and \(d\geq6\),
\[
 q^m_{(nd-6,6)}(d[n])=K_n((nd-6,6))
                    =a_{(nd-6,6)}(d[n])=4.             \tag{M-012.2}
\]
This supplies an actual permanent-side multiplicity larger than
one, established by containment and independence of all copies.

\begin{proof}[Proof of binary containment]
Write \(f=\sum_{j=0}^m c_j L^j y^{m-j}\). Define an
\(m\)-by-\(m\) matrix \(A\) by
\[
 A_{ii}=L,\quad A_{i,i+1}=y\quad(1\leq i<m),\qquad
 A_{m1}=c_0y+c_1L,\quad A_{mj}=c_jL\quad(2\leq j\leq m),
\]
with all remaining entries zero. In particular \(A_{11}=L\).
For a nonzero permanent term, let \(j\) be the column chosen
in the last row. The other choices are forced: rows \(i<j\)
use their diagonal entries, and rows \(j\leq i<m\)
use their superdiagonal entries. For example the forcing
starts at column one when \(j>1\), and then proceeds
successively to the right. Thus
\[
 \operatorname{perm}A
   =\sum_{j=1}^m L^{j-1}y^{m-j}A_{mj}
   =\sum_{j=0}^m c_j L^j y^{m-j}=f.
\]
Substitution into the internally padded polynomial gives
\(A_{11}^{n-m}\operatorname{perm}A=L^{n-m}f\).
This is a linear endomorphism substitution in the original
\(n^2\) variables, so it lies in the orbit closure by
density of the invertible endomorphisms. The identity is
integral in the coefficient variables \(c_j\).
\end{proof}

\begin{proof}[Proof of the multiplicity equalities]
The case \(t=0\) is the single pure-power coefficient.
For \(t\geq1\), put \(\mu=(dm-t,t)\).
The arbitrary-rank padding inequality O-012.b, now applied
to the binary family contained in \(Z_{n,m}\), yields
\[
                  q^m_\lambda(d[n])\geq a_\mu(d[m]).
\]
Because \(n\geq m\geq t\), F-014 stabilizes both sides
at inner degree \(t\), and gives
\(a_\mu(d[m])=a_\lambda(d[n])\). The ambient quotient
bound gives the reverse inequality. Every binary degree-\(n\)
form is a product of \(n\) linear forms over \(\mathbb C\),
hence a diagonal determinant. The two-row argument O-006
therefore gives \(K_n(\lambda)=a_\lambda(d[n])\), proving
(M-012.1).

For \(t=6\), \(d\geq6\), \(n>m\geq6\), the binary
plethysm formula in F-013 with determinant shift zero gives
\(a_{(nd-6,6)}=p_6(d,n)-p_5(d,n)=4\).
Equivalently the four partitions of six with no part one
are \((6),(4,2),(3,3),(2,2,2)\); both rectangle bounds
are inactive. This proves (M-012.2).
\end{proof}

\paragraph{A higher-row family with permanent multiplicity at least four.}
Suppose \(n>m\geq6\), \(2\leq r\leq m^2\), \(b\geq8\)
is even, and \(d\geq r+b-2\). Then
\[
 \lambda=(nd-b-2(r-2),b,2^{r-2})
 \quad\Longrightarrow\quad q^m_\lambda(d[n])\geq4.    \tag{M-012.3}
\]
Take a basis of the four highest-weight classes at degree
six and weight \((6n-6,6)\) in (M-012.2). Multiply every
class by the same function
\[
       W_r\, W_2^{(b-8)/2}\, F(e_1)^{d-r-b+2},
\]
where \(W_j\) is the leading Hessian minor from O-011.
The exponents are nonnegative. This common factor is nonzero
on \(Z_{n,m}\): O-011 supplies a nondegenerate Hessian
and a nonisotropic first direction with \(F(e_1)\ne0\);
extending that direction to an orthogonal basis makes all
leading Hessian minors nonzero simultaneously. The reduced
orbit closure is irreducible, so its coordinate ring is a
domain. Multiplication by this nonzero class preserves the
linear independence of all four seed classes.
The product has degree \(d\), second row \(b\), and
rows three through \(r\) equal to two; its first row is
therefore \(nd-b-2(r-2)\). Products of highest-weight
vectors add their weights, so these are four independent
highest-weight classes of the claimed weight. The tuple is
a partition because \(nd\geq2(r+b-2)\).
This proves (M-012.3), without claiming its lower bound is
the exact multiplicity in this higher-row family.

\paragraph{Grothendieck dictionary.}
The displayed matrix gives a morphism from an affine binary
coefficient scheme over \(\mathbb Z\) into the permanent
substitution scheme, preserving the distinguished padding
section. The padding comparison is a map of locally free
symmetric powers and its differential transpose, as in
O-012.b. The Cauchy and Schur operations are SGA~6, VI,
Theorem~3.3, pp.~392--393, and the highest-weight characters
are the torus comodules of SGA~3, I, Propositions~4.7.2--4.7.3,
pp.~34--36. On the reduced characteristic-zero orbit closure,
the four independent coherent sections remain independent
after multiplication by a nonzero section because its
coordinate ring is integral. This is an actual rank statement,
not an inference from support or from a virtual \(K_0\) class.

% END proofs/binary_permanent_seed.tex

% BEGIN proofs/jet_dominance.tex
\section{Taylor coefficients give a stronger all-copy comparison}

\subsection*{Finding M-013: a determinant jet map with surjective differential}
Let \(r\geq2\), put \(R=r-1\), and choose \(\ell\geq0\).
Define
\[
 b=2\ell+2,\qquad h=\lceil\sqrt R\rceil,\qquad
 J_b(r)=2h\binom{R+\ell}{\ell}.
\]
If \(n\geq J_b(r)\), the map from \(n\)-by-\(n\)
linear matrix pencils in \(x_1,z_1,\ldots,z_R\) to the
Taylor coefficients through degree \(b\) of
\(\det A(1,z)\) is dominant. More precisely, its differential
is surjective at an explicitly constructed pencil over \(\mathbb Z\).
Consequently, for every outer degree \(d\) and every
\(\lambda\vdash nd\) with \(\ell(\lambda)\leq r\)
and \(\lambda_2\leq b\),
\begin{equation}\label{eq:M013K}
                         K_n(\lambda)=a_\lambda(d[n]).
\end{equation}
As usual one-row shapes have second row zero. We only discuss
types fitting in the ambient \(n^2\) variables; the displayed
size bound ensures \(r\leq n^2\).
Also \(J_b(r)\geq b\), since
\(\binom{R+\ell}{\ell}\geq\ell+1\) and \(h\geq1\),
so the stated jet order never exceeds the determinant degree.

The construction permits other Taylor coefficients to vary.
It does not represent every polynomial in the truncated space
as an exact padded determinant. Dominance of its detected
coefficients is sufficient to preserve every highest-weight copy.

\paragraph{The acyclic pencil.}
Choose an injection from the \(R\) variable labels into pairs
\((a,c)\in\{1,\ldots,h\}^2\). Denote the corresponding
variable by \(z_{ac}\); unused pairs have label zero.
For each \(a\), take a copy of the monomial tree on
\(\alpha\in\mathbb Z_{\geq0}^R\), \(|\alpha|\leq\ell\),
oriented toward its root. For each \(c\), take a separate
copy oriented away from its root. The parent of a nonzero
multi-index removes its last nonzero coordinate; label that
edge by the corresponding variable, with positive sign.
Connect inward root \(a\) to outward root \(c\) by the
edge labeled \(z_{ac}\) when this is nonzero. No roots
are identified, and no other edges are present.

There are \(J_b(r)\) vertices. The graph is acyclic:
paths run inward, cross at most one bridge, and then run
outward. Let \(B(z)\) be its adjacency matrix. After a
topological ordering it is strictly triangular, hence nilpotent,
and
\[
 A_0(x_1,z)=x_1 I-B(z),\qquad
 \det A_0=x_1^{J_b(r)},\qquad
 A_0(1,z)^{-1}=I+B(z)+B(z)^2+\cdots.
\]
Only finitely many powers occur. Between inward vertex
\((\alpha,a)\) and outward vertex \((\beta,c)\) there
is exactly one path if the bridge is present. Its product is
\[
                         z^\alpha z_{ac}z^\beta.
                                                               \tag{M-013.1}
\]
Thus the corresponding inverse entry is exactly this monomial.
Every monomial of degree between one and \(2\ell+1\)
has this form: choose one of its variables for the bridge
and split its remaining multiset between \(\alpha,\beta\)
of degrees at most \(\ell\). The diagonal inverse entries
also give the constant monomial one. It follows that the
inverse entries span every polynomial of degree at most
\(b-1=2\ell+1\). For \(n>J_b(r)\), adjoin an
\(x_1 I\) block; all these assertions remain true on the
original block after setting \(x_1=1\).

\begin{proof}[Proof of jet dominance]
Let \(\mathcal P\) be the affine space of all linear matrix
pencils of size \(n\) in these \(r\) variables, and let
\(\mathcal J_b\) be the vector space with basis
\(z^\gamma\), \(|\gamma|\leq b\). The polynomial map
\[
 \Phi:\mathcal P\longrightarrow\mathcal J_b,
 \qquad A\longmapsto\det A(1,z)\pmod{(z_1,\ldots,z_R)^{b+1}}
\]
records exactly the desired coefficient jet. A variation of
one matrix entry by a linear form \(H\) has differential
\[
 d(\det)_{A_0}(H E_{ij})
       =H(1,z)\,\det A_0(1,z)\,(A_0(1,z)^{-1})_{ji}
       =H(1,z)(A_0(1,z)^{-1})_{ji}.
\]
This is also the cofactor derivative identity, so no analytic
matrix inverse is required. Taking \(H=x_1\) and using
(M-013.1) yields every monomial through degree \(b-1\).
Taking \(H=z_t\) yields every degree-\(b\) monomial as
well. The differential of \(\Phi\) at \(A_0\) is
therefore surjective onto \(\mathcal J_b\).

Here is an explicit algebraic justification of dominance.
Choose \(\dim\mathcal J_b\) source directions on which
this differential is an isomorphism and restrict \(\Phi\)
to the resulting affine subspace through \(A_0\). Translate
source and target origins and make an invertible linear change
of target coordinates. The restricted polynomial map then
has the form \(F(t)=t+\text{terms of degree at least two}\).
If a nonzero polynomial \(P\) vanished on its image, the
lowest nonzero homogeneous part of \(P(F(t))\) would be
the same as that of \(P(t)\), a contradiction. Thus
pullback on polynomial functions is injective. This proves
that \(\Phi\) is dominant, and completes the first claim.

In fact the chosen differential directions can be individual
matrix-entry coefficient coordinates, each mapping to one distinct
monomial with coefficient one. The corresponding Jacobian minor
is the identity over \(\mathbb Z\). The translated restricted map
already has \(F(t)=t+O(t^2)\) integrally. Its lowest-term argument
therefore proves injectivity of its polynomial pullback after
every commutative base change. This assertion does not require
an algebraic section or a flatness claim; the representation-theoretic
conclusion below continues to be asserted in characteristic zero.
\end{proof}

\begin{proof}[Proof that every ambient copy survives]
Take a highest-weight vector \(w\) of weight \(\lambda\)
in the degree-\(d\) determinant ideal. Nonnegative polynomial
weights show that it depends only on the first \(r\)
coordinates. The \(j=2\) case of F-016 says, more strongly,
that every coefficient generator used by \(w\) has total
degree at most \(\lambda_2\leq b\) outside the first
coordinate. Consequently \(w\) is a polynomial in precisely
the coefficients recorded by \(\mathcal J_b\).

Every determinant of a linear pencil in these variables is
a linear endomorphism substitution into \(\det_n\), and
belongs to \(\Omega_n\). Thus \(w\circ\Phi=0\).
The injectivity of pullback just proved gives \(w=0\).
This applies to every highest-weight vector, not merely to
one chosen copy. The complete ideal multiplicity space at
\(\lambda\) is zero, which is \eqref{eq:M013K}.
Types of zero ambient multiplicity cause no exception.
\end{proof}

\paragraph{Uniform consequence in every power range.}
For \(m\geq2\), every integer \(c\geq6\), and
\(n\geq m^c\), an actual internally padded permanent
multiplicity obstruction must satisfy
\begin{equation}\label{eq:M013uniform}
                              \lambda_2\geq c+2.
\end{equation}
In particular, the BIP range now requires \(\lambda_2\geq27\),
which strengthens the earlier constant twenty-six.

\begin{proof}
If \(q^m_\lambda>0\), then \(r=\ell(\lambda)\leq m^2\).
One-row shapes have full determinant multiplicity, so assume
\(r\geq2\). For every \(\ell\geq3\),
\[
 \frac{2\binom{m^2+\ell-1}{\ell}}{m^{2\ell}}
 \leq\frac{2\binom{\ell+3}{3}}{4^\ell}
 \leq\frac58<1.                                      \tag{M-013.2}
\]
The first inequality follows from the product formula:
each factor \(1+i/m^2\), \(0\leq i<\ell\), decreases
with \(m\), so its largest allowed value occurs at \(m=2\).
The second inequality is equality at \(\ell=3\); the
ratio on increasing \(\ell\) is
\((\ell+4)/(4(\ell+1))<1\).

If \(c=2\ell+1\) is odd, then \(\ell\geq3\),
\(R=r-1\leq m^2-1\), and \(h=\lceil\sqrt R\rceil\leq m\).
For the even jet order \(b=c+1=2\ell+2\),
\[
 J_b(r)\leq2m\binom{m^2+\ell-1}{\ell}
                      \leq\tfrac58m^{2\ell+1}<n.
\]
Thus every type with \(\lambda_2\leq c+1\) has
\(K_n=a_\lambda\) by the jet theorem.

If \(c=2\ell\) is even, use the exact degree-split
construction O-012.a at base degree \(c+1=2\ell+1\).
Its size, for \(r\) variables, is
\[
 2\binom{r-1+\ell}{\ell}-1
       \leq\tfrac58m^{2\ell}-1<n.
\]
Monotonicity in the base degree and F-014 therefore give
\(K_n=a_\lambda\) whenever \(\lambda_2\leq c+1\)
in this case too. In either parity \(q^m_\lambda\leq a_\lambda\)
excludes the alleged obstruction, proving \eqref{eq:M013uniform}.
\end{proof}

\paragraph{Relation to the staircase.}
Combining this result with F-017, the BIP range requires the
complete staircase
\[
                    \lambda_{2j}\geq28-j
                              \qquad(1\leq j\leq26).
\]
The new endpoint is \(\lambda_2\geq27\); the other
twenty-five inequalities are the mixed graph theorem.
The jet argument has changed the required geometric input:
a dominant map to the coefficients detected by a section is
enough. Exact containment of the entire truncated polynomial
family is a stronger condition and is not claimed here.

\paragraph{Grothendieck dictionary.}
The truncated algebra
\(\mathbb Z[z_1,\ldots,z_R]/(z_1,\ldots,z_R)^{b+1}\)
is the explicit finite infinitesimal neighborhood used to
record the coefficients. The graph, pencil and cofactor maps
are defined over \(\mathbb Z\). Their monomial gradings
are torus comodules as in SGA~3, I, Proposition~4.7.3,
pp.~35--36. The flag-minor support is the Schur/exterior
functor calculation of SGA~6, VI, Theorem~3.3, pp.~392--393.
The target jet module is a finite locally free coefficient
module, and the unit Jacobian minor makes its polynomial coordinate
ring inject into that of the pencil parameter scheme after arbitrary
base change. The consequent characteristic-zero vanishing of the ideal
comodule uses SGA~3, I, Proposition~4.7.2, pp.~34--35.
This is an injection of actual coordinate functions, which
contains more information than equality of virtual classes
in \(K_0\). The coherent coefficient quotients fit the
parameter-space framework of FGA, Bourbaki~221, following
Proposition~3.8, pp.~265--266. No unproved positivity of a
motivic trace or integral irreducible splitting is used.

% END proofs/jet_dominance.tex

% BEGIN proofs/general_staircase.tex
\section{A uniform staircase at every integral power threshold}

\subsection*{Finding M-014: the mixed construction for all exponents at least six}
Let \(m\geq2\), \(c\geq6\) be integers and \(n\geq m^c\).
Every actual internally padded permanent multiplicity obstruction
\(q^m_\lambda(d[n])>K_n(\lambda)\) satisfies
\begin{equation}\label{eq:M014stairs}
             \lambda_{2j}\geq c+3-j
                         \qquad(1\leq j\leq c+1).
\end{equation}
Parts beyond the length are zero, as usual. This generalizes
the BIP staircase at \(c=25\). The separate IP staircase
at \(n>3m^4\) is already F-017 and is not obtained by
substituting \(c=4\) into the hypotheses above.

\begin{proof}
First, an actual obstruction requires \(m>2c\).
If \(4\leq m\leq2c\), then
\[
          m^c\geq4^c=2^{2c}\geq2^m>2^m-1.
\]
For \(m=2,3\), the same sufficient determinant size
\(2^m-1\) is at most \(m^c\) directly. The known
Grenet construction, in the checked form F-009, therefore
excludes every \(m\leq2c\). Also M-008 gives
\(r=\ell(\lambda)\geq2c+2\), while positive permanent
multiplicity gives \(r\leq m^2\).

We first prove the assertion at the cuts
\[
              k=2c+4-2s,\qquad 1\leq s\leq c+1,
                                                               \tag{M-014.1}
\]
except for \(s=c+1\) when \(c\) is odd. Suppose
\(\lambda_k\leq s\). Choose nonnegative integers
\(u,v\) with \(u+v=s-1\) and \(|u-v|\leq1\),
and put
\[
                      p=c-2u,\qquad q=c-2v.
\]
The stated range of \(s\), with its one exception,
ensures \(p,q\geq0\). Moreover \(p+q=k-2\).
Thus the F-017 mixed graph applies, with
\(R=r-k+1\), and has size
\[
 S=\binom{m+p}{p}\binom{R+u}{u}
      +\binom{m+q}{q}\binom{R+v}{v}-1.                 \tag{M-014.2}
\]
We show that each binomial product is strictly less than
\(m^c/2\).

Write
\[
 A_p(m)=\frac{\binom{m+p}{p}}{m^p},\qquad
 B_u(r)=\frac{\binom{r-1+u}{u}}{r^u}.
\]
The product formula shows that \(B_u(r)\) decreases
with \(r\). For \(r\geq2\),
\[
              B_u(r)\leq\frac{u+1}{2^u}\leq1.
                                                               \tag{M-014.3}
\]
For \(u\geq3\) and \(r\geq4\), the sharper bound is
\[
              B_u(r)\leq\frac{\binom{u+3}{3}}{4^u}
                                      \leq\frac5{16}.
                                                               \tag{M-014.4}
\]
The last sequence starts at \(20/64\) for \(u=3\)
and decreases, because its successive ratio is
\((u+4)/(4(u+1))<1\).

If \(p\geq3\), then \(m\geq2c\geq2p\) and
\[
 A_p(m)=\frac1{p!}\prod_{i=1}^p(1+i/m)
       \leq\frac{\exp(p(p+1)/(2m))}{p!}
       \leq\frac{\exp((p+1)/4)}{p!}
       \leq\frac e6<\frac12.
\]
The sequence in the penultimate expression decreases for
\(p\geq3\): its successive ratio is
\(e^{1/4}/(p+1)<1\). The elementary bound \(e<3\)
follows directly from its convergent factorial series.
Together with (M-014.3), this gives \(A_pB_u<1/2\).

If \(p=2\), the equation \(p+2u=c\geq6\) gives
\(u\geq2\). Since \(m\geq12\),
\[
 A_2(m)B_u(r)\leq\frac{91}{144}\frac34
                                      =\frac{91}{192}<\frac12.
\]
If \(p=1\), then \(u\geq3\), so (M-014.4) gives
\[
 A_1(m)B_u(r)\leq\frac{13}{12}\frac5{16}
                                      =\frac{65}{192}<\frac12.
\]
If \(p=0\), then \(u\geq3\) and the same bound
\(B_u(r)\leq5/16<1/2\) applies. These cases exhaust
the possibilities. Because \(R\leq r-1\), \(r\leq m^2\)
and \(p+2u=c\), they imply
\[
 \binom{m+p}{p}\binom{R+u}{u}
            \leq A_p(m)B_u(r)m^p r^u<\frac12m^c.
\]
The other product has the identical proof with \(q,v\).
Consequently \(S<m^c\leq n\). F-017 then gives
\(q^m_\lambda\leq K_n(\lambda)\), a contradiction.
Thus \(\lambda_k\geq s+1\) for every cut in
(M-014.1) covered so far.

The omitted odd-\(c\) endpoint is \(k=2\), \(s=c+1\).
It is precisely the jet consequence M-013:
an actual obstruction requires \(\lambda_2\geq c+2\).
All cases of (M-014.1) therefore hold. Substituting
\(j=c+2-s\) gives \eqref{eq:M014stairs}.
\end{proof}

\paragraph{Exact scope and Grothendieck dictionary.}
The result is a uniform numerical consequence of the actual
ideal-multiplicity comparisons F-017 and M-013, with the
known whole-permanent containment bound F-009 used only
to remove small \(m\). It does not infer multiplicity
inequalities from occurrence, nor does it impose the stronger-test
condition \(d>n/m\) on actual obstructions.
The torus-weight supports are the comodule gradings of
SGA~3, I, Proposition~4.7.3, pp.~35--36; their entire
highest-weight spaces arise from the Schur and exterior
operations of SGA~6, VI, Theorem~3.3, pp.~392--393.
The integral graph morphisms and the coefficient-jet map
retain the coordinate-algebra pullbacks. Vanishing of each
complete ideal comodule in characteristic zero is therefore
the precise scheme-theoretic content of the displayed
inequalities. No new purity or saturation input is used.

% END proofs/general_staircase.tex

% BEGIN proofs/cyclic_star_jets.tex
\section{Sharing intermediate vertices in an order-five jet construction}

\subsection*{Finding M-015: order-five jets at matrix size \(r^2\)}
Let \(r\geq4\) and \(n\geq r^2\). The map from
\(n\)-by-\(n\) linear pencils in \(r\) variables to their
determinant coefficient jets through degree five is dominant.
It has an explicit integral point with a unit Jacobian minor.
Consequently
\begin{equation}\label{eq:M015K}
 \lambda\vdash nd,\quad\ell(\lambda)\leq r,\quad\lambda_2\leq5
       \quad\Longrightarrow\quad K_n(\lambda)=a_\lambda(d[n]).
\end{equation}
In particular an actual internally padded permanent obstruction
at \(n\geq m^4\), \(m\geq2\), requires
\begin{equation}\label{eq:M015per}
                              \lambda_2\geq6.
\end{equation}
The earlier IP bound assumed \(n>3m^4\). This result
removes that factor for the second-row conclusion alone;
the other IP staircase inequalities retain their own hypotheses.

\paragraph{The pencil and its core.}
Set \(R=r-1\geq3\), write the variables as
\(x,z_1,\ldots,z_R\), and let
\(\sigma=(1\ 2\ \cdots\ R)\). The directed graph has
the following distinct vertices:
\[
 0;\qquad v_i\ (1\leq i\leq R);\qquad
 p_{ij},q_{ij}\ (1\leq i\leq j\leq R).
\]
Its size is \(1+R+2\binom{R+1}{2}=(R+1)^2=r^2\).
Its edges, and no others, are
\[
 \begin{array}{c|c}
 \text{edge}&\text{label}\\\hline
 p_{ij}\longrightarrow v_i&z_j\\
 v_i\longrightarrow0&z_i\\
 0\longrightarrow v_i&z_{\sigma(i)}\\
 v_{\sigma^{-1}(j)}\longrightarrow q_{ij}&z_i.
 \end{array}                                             \tag{M-015.1}
\]
Let \(B\) be its labeled adjacency matrix and put
\(A=xI-B\). Unlike M-013's acyclic graph, this graph has
two-edge cycles through the central vertex. The inward leaves
\(p_{ij}\) have no incoming edges, and the outward leaves
\(q_{ij}\) have no outgoing edges. At \(x=1\), therefore,
\[
        \det A=D:=1-\sum_{i=1}^R z_i z_{\sigma(i)}.
                                                               \tag{M-015.2}
\]
Indeed only the central star contributes nontrivially to
the determinant; its Schur complement is \(D\).

Let \(C\) be the \((R+1)\)-vertex star block at \(x=1\).
The elementary block inverse gives
\[
 \begin{split}
 D(C^{-1})_{00}&=1,\\
 D(C^{-1})_{0k}&=z_{\sigma(k)},\\
 D(C^{-1})_{ik}&=D\delta_{ik}+z_i z_{\sigma(k)}.
 \end{split}                                             \tag{M-015.3}
\]
These rational identities are cofactor polynomial identities
after multiplication by \(D\), so they apply without
restricting to points where \(D\ne0\).
The two leaf edges also give
\[
 D(A^{-1})_{p_{ij},q_{ab}}
     =z_jz_a\left(D\delta_{i,\sigma^{-1}(b)}+z_i z_b\right).
                                                               \tag{M-015.4}
\]
Here \(A^{-1}\) means the inverse after \(x=1\).

\paragraph{Pure monomial cofactors.}
Equation (M-015.3) supplies the constant and every linear
monomial. It also supplies every quadratic monomial from
an off-diagonal entry. For \(z_a z_b\), use row \(a\)
and column \(\sigma^{-1}(b)\) if these differ. Otherwise
reverse \(a,b\). Both choices could fail only if
\(b=\sigma(a)\) and \(a=\sigma(b)\), which would be
a fixed point or two-cycle of \(\sigma\). Neither exists
because \(R\geq3\).

Every monomial of degree four can be written
\(z_i z_j z_a z_b\) with \(i\leq j\), \(a\leq b\)
and \(b\ne\sigma(i)\). Here is a complete check.
If its largest and smallest indices differ by at least two,
put a smallest index into the first pair and a largest into
the second. The first minimum is then \(i\), the second
maximum is \(b\geq i+2\), and \(b\ne\sigma(i)\).
If every index is the same, \(b=i\ne\sigma(i)\).
The remaining case uses just the adjacent indices \(t,t+1\).
For their multiplicities \((3,1)\), take pairs
\((t,t+1)\), \((t,t)\); for \((1,3)\), take
\((t+1,t+1)\), \((t,t+1)\); and for \((2,2)\), take
\((t+1,t+1)\), \((t,t)\). The first two choices use
absence of fixed points. The last uses
\(\sigma(t+1)\ne t\), including the wraparound case
\(t+1=R\), since then \(R-1\ne1\).
These are all degree-four multisets. Equation (M-015.4)
therefore supplies each degree-four monomial as a pure
cofactor with coefficient one.

\begin{proof}[Proof of jet dominance and the multiplicity conclusions]
Varying entry \((u,v)\) of the pencil by a linear form
\(H\) differentiates its determinant to
\(H(1,z)\det A(1,z)(A(1,z)^{-1})_{vu}\).
The pure cofactors above have degrees zero, one, two and
four. Taking \(H=x\) yields those degrees directly;
taking \(H=z_i\) also yields every monomial of degrees
three and five. Thus the Jacobian of the projection onto
all coefficients through degree five is surjective at this
integral pencil. Each chosen elementary coefficient direction
maps to one distinct monomial with coefficient one, so a
Jacobian minor is the identity over \(\mathbb Z\).
For \(n>r^2\), pad by \(xI_{n-r^2}\); after \(x=1\)
the same cofactor directions remain.

Translation by the coefficient vector \(D\) and restriction
to these parameter directions give a polynomial map
\(F(t)=t+O(t^2)\) over \(\mathbb Z\). As in M-013,
the lowest homogeneous term of any nonzero polynomial
survives composition with \(F\). The pullback is injective
after every commutative base change, and the coefficient
map is dominant over \(\mathbb C\).

The support theorem F-016 places every highest-weight
function of second row at most five in this coefficient
algebra. A determinant-ideal vector in that algebra vanishes
after the dominant projection and must be zero. This holds
for the complete highest-weight space, proving
\eqref{eq:M015K} in characteristic zero.

For \eqref{eq:M015per}, a positive permanent multiplicity
has length \(r\leq m^2\). If \(r\geq4\), then
\(n\geq m^4\geq r^2\), and the preceding conclusion
excludes \(\lambda_2\leq5\). If \(r\leq3\), the
full determinant multiplicity results O-006--O-007 already
exclude every actual obstruction under these hypotheses.
This proves the asserted necessary condition.
\end{proof}

\paragraph{Grothendieck dictionary and method.}
The matrix and all cofactor identities are defined over
\(\mathbb Z\). The target is the finite free coefficient
module of \(\mathbb Z[z]/(z)^6\), and the unit Jacobian
minor gives an actual universal injection on polynomial
coordinate algebras. The coefficient weights are the torus
comodules of SGA~3, I, Proposition~4.7.3, pp.~35--36;
the whole Schur-copy support is F-016's use of SGA~6, VI,
Theorem~3.3, pp.~392--393, and Macdonald, I(4.3), p.~63,
Appendix~A(5.3)--(5.6), pp.~155--156.
Sharing intermediate vertices introduces cycles and changes
the base determinant to \(D\). The explicit off-diagonal
cofactor calculation, rather than acyclicity or an exact
truncated-form representation, supplies dominance. Higher
coefficients remain free. The conclusion is an ideal-comodule
vanishing statement over characteristic zero, not a virtual
\(K_0\) equality or a saturation inference.

% END proofs/cyclic_star_jets.tex

% BEGIN agents/formulas/permanent_jet_dominance.tex
\section{Full permanent multiplicities from detected Taylor coefficients}

\subsection*{Finding F-021: the internal-padding permanent jet map}
Let \(r\geq2\), \(R=r-1\), \(\ell\geq0\), and set
\[
 b=2\ell+2,\qquad h=\lceil\sqrt R\rceil,\qquad
 J_b(r)=2h\binom{R+\ell}{\ell}.
\]
Suppose \(n\geq m\geq J_b(r)\). For every outer degree
\(d\) and every partition \(\lambda\vdash nd\) with
\(\ell(\lambda)\leq r\) and \(\lambda_2\leq b\),
\begin{equation}\label{eq:F021}
 q_\lambda(m,n)=K_n(\lambda)=a_\lambda(d[n]).
\end{equation}
Here \(q\) uses the internal padding
\(X_{11}^{n-m}\operatorname{per}_m\), exactly as in the
supplied IP and BIP sources. The statement includes every
highest-weight copy, and imposes no upper bound on \(d\).
The usual setting \(n>m\) is a special case. This is a
variant derived in this run from M-013, without a priority claim.

\begin{proof}
Use the two separate families of rooted monomial trees
and the bridges of M-013. Thus \(B(z)\) is the adjacency
matrix of an acyclic directed graph on \(J=J_b(r)\)
vertices. Its inward-to-outward path from \((\alpha,a)\)
to \((\beta,c)\), when the bridge labeled \(z_{ac}\)
is present, is unique and has product
\[
                         z^\alpha z_{ac}z^\beta.
                                                               \tag{F-021.1}
\]
All nonempty paths have positive labels. Put
\(A_0(x_1,z)=x_1 I+B(z)\), and append an \(x_1 I\)
block if \(m>J\). Choose any one vertex as the distinguished
first vertex, so \((A_0)_{11}=x_1\).
Since the graph is acyclic, the only permutation contributing
to the permanent of \(I+B\) is its diagonal permutation.
Consequently \(\operatorname{per}A_0=x_1^m\).

The permanent cofactor obtained by deleting row \(i\)
and column \(j\) from \(I+B\) has a useful exact
description. If \(i=j\), it equals one. If \(i\ne j\),
it is the sum of the products of all directed paths from
\(j\) to \(i\). Indeed a term in this minor chooses one
outgoing edge at each vertex other than \(i\), and one
incoming edge at each vertex other than \(j\). Its graph
therefore consists of a directed path from \(j\) to \(i\),
loops, and possibly directed cycles. Acyclicity excludes
all nontrivial cycles, and every such path together with
loops on its unused vertices gives exactly one minor term.
All signs are positive. In particular, the cofactors on
the chosen inward-to-outward pairs are the single monomials
in (F-021.1).

Every monomial of degree between one and \(2\ell+1\)
has the form (F-021.1): remove one variable for a bridge
and split the remaining multiset into two parts of degrees
at most \(\ell\). Diagonal cofactors give the constant
monomial. The linear-entry differential
\[
 d(\operatorname{per})_{A_0}(H E_{ij})
       =H(1,z)\operatorname{per}\bigl((I+B)_{\widehat i,
                                                        \widehat j}\bigr)
                                                               \tag{F-021.2}
\]
therefore spans every monomial through degree \(b\), by
taking \(H=x_1\) through degree \(b-1\) and multiplying
a degree-\(b-1\) cofactor by a variable for the final degree.

The distinguished entry must remain exactly \(x_1\),
so none of its coefficients is varied. This does not remove
any required image direction. One has \(J\geq b\geq2\),
and hence there is another diagonal entry to supply the
constant monomial, as well as any linear monomial if needed.
All positive-degree path cofactors used above come from
off-diagonal entries. Distinct monomial image directions
can be chosen from distinct elementary entry-coefficient
coordinates, each with coefficient one. Thus the map
\[
 \{A:\ A_{11}=x_1\}\longrightarrow
 \mathbb Z[z_1,\ldots,z_R]_{\leq b},\qquad
 A\longmapsto\operatorname{per}A(1,z)\pmod{(z)^{b+1}}
                                                               \tag{F-021.3}
\]
has an identity Jacobian minor at \(A_0\), over
\(\mathbb Z\). Restricting to those source coordinates
and translating origins gives \(F(t)=t+O(t^2)\).
The lowest nonzero homogeneous part of a polynomial survives
substitution into \(F\). Its polynomial pullback is therefore
injective, after every commutative base change. In characteristic
zero this in particular proves dominance of (F-021.3).

For the internally padded permanent, substitution of the
pencil \(A\) with \(A_{11}=x_1\) gives
\[
                         x_1^{n-m}\operatorname{per}A(x_1,z).
\]
At \(x_1=1\) its detected coefficients are exactly those
of (F-021.3). All these forms belong to the permanent orbit
closure: a linear substitution into the matrix entries is
an endomorphism substitution, and invertible endomorphisms
are dense. The variable space fits because
\(m\geq2\lceil\sqrt{r-1}\rceil\) implies \(r\leq m^2\).

Now let \(w\) be any highest-weight vector of weight
\(\lambda\) in the permanent ideal. It uses only the
first \(r\) coordinates. F-016 with \(j=2\) says that
every coefficient generator it uses has degree at most
\(\lambda_2\leq b\) outside the first coordinate.
Thus \(w\) is a polynomial in the detected jet coordinates.
It vanishes on the internally padded pencil family, so
the proved pullback injection forces \(w=0\).
The entire ideal multiplicity space vanishes, and hence
\(q_\lambda=a_\lambda\). M-013 applies as well because
\(n\geq m\geq J_b(r)\), giving \(K_n=a_\lambda\).
This proves \eqref{eq:F021}.
\end{proof}

\paragraph{An additional order-five case.}
The cyclic-star construction of M-015 also has a permanent
version. For \(r\geq4\) and \(n\geq m\geq r^2\),
\begin{equation}\label{eq:F021five}
 \ell(\lambda)\leq r,\quad\lambda_2\leq5
       \quad\Longrightarrow\quad q_\lambda=K_n(\lambda)=a_\lambda.
\end{equation}
To verify this variant, keep the same central vertex, core
vertices, leaves, and positive edge labels as M-015, and
use \(A=x_1I+B\). Every nontrivial directed cycle is a
two-edge cycle through the central vertex. Thus at \(x_1=1\)
the base permanent is
\(D_+=1+\sum_i z_i z_{\sigma(i)}\).

The cofactor for the central diagonal is one. A cofactor
whose open path runs from the central vertex to a core
vertex is the corresponding single variable. Between distinct
core vertices \(v_i,v_k\), the only open path is
\(v_i\longrightarrow0\longrightarrow v_k\), with product
\(z_i z_{\sigma(k)}\). The chosen inward-to-outward leaf
cofactors have \(i\ne\sigma^{-1}(b)\), and their only
open path is
\[
 p_{ij}\longrightarrow v_i\longrightarrow0
         \longrightarrow v_{\sigma^{-1}(b)}\longrightarrow q_{ab},
\]
with product \(z_i z_j z_a z_b\). In each positive-degree
case the path uses the central vertex, so no disjoint
nontrivial cycle can contribute. All unused vertices take
their diagonal loops. These are therefore pure permanent
cofactors, with coefficient one. The complete degree-two
and degree-four coverage proved in M-015 applies unchanged.
Multiplication by a linear entry perturbation supplies the
missing degrees three and five.

Choose an inward leaf as the distinguished vertex and hold
its diagonal equal to \(x_1\). None of the selected directions
varies that diagonal: the constant comes from the central
diagonal and all the other selected cofactors are off diagonal.
The same integral unit-Jacobian and internal-padding argument
as above proves \(q_\lambda=a_\lambda\). M-015 gives
\(K_n=a_\lambda\), proving \eqref{eq:F021five}.

\paragraph{A strict improvement over every exact two-tree support cover.}
For the full simplex
\(\Gamma=\{\alpha\in\mathbb Z_{\geq0}^R:|\alpha|\leq4\}\),
\(R\geq4\), the finite cover invariant of F-018 satisfies
\begin{equation}\label{eq:F021cover}
             \kappa(\Gamma)\geq
                          1+\left\lceil\frac{R(R-1)}{12}\right\rceil.
\end{equation}
To prove this, let \(P,Q\) be any admissible downward-closed
sets. For a squarefree degree-four exponent \(\gamma\),
write \(\gamma=\alpha+\beta+\delta\), where
\(\alpha\in P\), \(\beta\in Q\), and
\(\delta\in\{0,e_1,\ldots,e_R\}\). At least one of
\(\alpha,\beta\) has degree at least two, and its
downward-closed set contains a squarefree pair supported
inside \(\gamma\). Let \(E\) be the set of pairs in
\(P\cup Q\). Every four-subset contains a member of \(E\),
so double-counting gives
\[
 |E|\binom{R-2}{2}\geq\binom R4,
          \qquad |E|\geq\left\lceil R(R-1)/12\right\rceil.
\]
Both \(P,Q\) contain zero; consequently
\(|P|+|Q|-1\geq|P\cup Q|\geq1+|E|\), proving the claim.

Take \(r=1024\), \(R=1023\), and \(b=4\). Then
\[
 J_4(1024)=2\cdot32\cdot1024=65536,
                    \qquad \kappa(\Gamma)\geq87127.
\]
With \(m=65536\) and \(n=65537\), F-021 applies to
every type with second row at most four, whereas neither
\(m\) nor \(n\) reaches the exact-cover lower bound for
this simplex. This is a limitation of the F-018 architecture,
not a lower bound for arbitrary permanent or determinant
representations.

This strict threshold comparison contains positive higher
multiplicity examples. Set \(d=2r\) and
\[
                  \lambda=(nd-4(r-1),4^{r-1}).
                                                               \tag{F-021.4}
\]
Its full nested support at either base degree four or any
larger degree is precisely \(\Gamma\). Its ambient
multiplicity is at least two, as follows. Let \(c_0(F)=F(e_1)\)
and let \(W_r(F)\) be the leading \(r\)-by-\(r\)
Hessian determinant at \(e_1\), as in O-011.
The polynomial \(W_r^2\) has degree \(2r\) and weight
\(\lambda\). It is not divisible by \(c_0\): the independent
coefficient variables can prescribe a Hessian with its
\((1,1)\) entry zero, its \((1,2),(2,1)\) entries one,
its \((2,2)\) entry zero, and an identity block in the
remaining coordinates. At this point \(c_0=0\) but
\(W_r=-1\ne0\), after harmless nonzero coefficient scalings.

F-011, with outer degree \(r\) and \(k=4\), supplies
a nonzero highest-weight polynomial \(I\) of degree \(r\)
and weight \((r(n-4)+4,4^{r-1})\). Therefore
\(I c_0^r\) also has degree \(2r\) and weight \(\lambda\).
It is nonzero and divisible by \(c_0\), whereas \(W_r^2\)
is not. These two highest-weight polynomials are linearly
independent. Hence (F-021.4) at the stated parameters has
\[
                         2\leq q_\lambda=K_n(\lambda)=a_\lambda.
\]
No numerical value larger than this lower bound is asserted.

\paragraph{Grothendieck dictionary and remaining scope.}
The permanent minor is a positive sum of path monomials,
and the resulting coefficient maps are integral morphisms
of affine schemes. The truncated polynomial module is finite
locally free. Its coordinate algebra injects into the source
parameter algebra after arbitrary base change by the unit
Jacobian argument. Schur highest weights and the characteristic-zero
quotient decomposition identify the complete multiplicity spaces;
these are the comodule and Schur-functor operations of SGA~3
I4.7.2--4.7.3 and SGA~6 VI3.3. This geometric injection
contains information not supplied by an equality of virtual
\(K_0\)-classes. The conclusion determines additional full
permanent multiplicities; it does not assert a comparison
outside the stated jet-size condition or solve the remaining
general evaluation-rank problem.

% END agents/formulas/permanent_jet_dominance.tex

% BEGIN agents/formulas/shared_core_odd_jets.tex
\section{All odd jet orders from two trees sharing a cyclic core}

\subsection*{Finding F-022: a uniform odd-order jet construction}
Let \(r\geq4\), \(R=r-1\), and \(\ell\geq1\). Put
\[
 b=2\ell+1,\qquad
 L_b(r)=2\binom{r+\ell-1}{\ell}-r.
                                                               \tag{F-022.1}
\]
The following statements hold in characteristic zero.
\begin{enumerate}
\item If \(n\geq L_b(r)\), every partition \(\lambda\vdash nd\)
with \(\ell(\lambda)\leq r\) and \(\lambda_2\leq b\)
satisfies \(K_n(\lambda)=a_\lambda(d[n])\).
\item If \(n\geq m\geq L_b(r)\), every such partition satisfies
\(q_\lambda(m,n)=K_n(\lambda)=a_\lambda(d[n])\), with
the internal padding \(X_{11}^{n-m}\operatorname{per}_m\).
\end{enumerate}
Both conclusions concern the entire multiplicity space, for
every outer degree \(d\). The coefficient maps have unit
Jacobian minors over \(\mathbb Z\), and their polynomial
pullbacks remain injective after arbitrary commutative base
change. This integral assertion is separate from the
characteristic-zero representation decomposition.
The cases \(\ell=1,2\) have sizes
\[
                         L_3(r)=r,\qquad L_5(r)=r^2.
                                                               \tag{F-022.2}
\]
The second recovers M-015 and its permanent variant F-021.
This is a checked generalization derived in this run, without
a priority claim.

\paragraph{The shared vertices and their labels.}
Write the variables as \(x,z_1,\ldots,z_R\), and let
\(\sigma=(1\ 2\ \cdots\ R)\). Take a central vertex
\(0\) and core vertices \(v_1,\ldots,v_R\).
For each nonnegative multi-index \(\alpha\) with
\(1\leq|\alpha|\leq\ell\), take an inward vertex
\(p_\alpha\) and an outward vertex \(q_\alpha\).
Make exactly the identifications
\[
                  p_{e_i}=v_i,\qquad
                  q_{e_j}=v_{\sigma^{-1}(j)}.
                                                               \tag{F-022.3}
\]
All vertices of degree at least two remain distinct, including
between the two families. The number of vertices is
\[
 1+2\left(\binom{R+\ell}{\ell}-1\right)-R=L_b(r).
\]
For \(|\alpha|\geq2\), remove one occurrence of its
largest supported index to obtain its inward parent:
\[
                 p_\alpha\longrightarrow p_{\alpha-e_j}
                    \quad\text{with label }z_j,
                    \qquad j=\max\operatorname{supp}\alpha.
\]
For \(|\beta|\geq2\), the outward edge instead removes
the smallest supported index in the reverse direction:
\[
                 q_{\beta-e_i}\longrightarrow q_\beta
                    \quad\text{with label }z_i,
                    \qquad i=\min\operatorname{supp}\beta.
\]
Finally put edges
\(v_i\longrightarrow0\) labeled \(z_i\) and
\(0\longrightarrow v_i\) labeled \(z_{\sigma(i)}\).
There are no other edges. Every nontrivial directed cycle
is a two-edge cycle through \(0\): inward paths strictly
decrease degree, outward paths strictly increase degree,
and neither family can return to the core after leaving it.

Let \(B\) be the positive-labeled adjacency matrix.
For the determinant use \(A_-=xI-B\), and for the
permanent use \(A_+=xI+B\). At \(x=1\), their values are
\[
 \det A_-=1-\sum_i z_i z_{\sigma(i)},\qquad
 \operatorname{per}A_+=1+\sum_i z_i z_{\sigma(i)}.
                                                               \tag{F-022.4}
\]
There is at most one nontrivial cycle in a cycle cover,
and every other vertex uses its diagonal loop. This proves
both identities directly.

\paragraph{Pure monomial cofactors.}
The cofactor of the central diagonal is one for both
polynomials, since deleting \(0\) leaves an acyclic graph.
For \(1\leq|\gamma|\leq\ell\), the path from \(0\)
to \(q_\gamma\) is unique and has product \(z^\gamma\).
It passes through the core vertex indexed by
\(\sigma^{-1}(\max\operatorname{supp}\gamma)\).

More generally, let \(\alpha,\beta\ne0\), with their
degrees at most \(\ell\), and suppose
\[
        \max\operatorname{supp}\beta
                      \ne\sigma(\min\operatorname{supp}\alpha).
                                                               \tag{F-022.5}
\]
The path from \(p_\alpha\) to \(q_\beta\) goes
inward to \(v_i\), through \(0\), and outward from
\(v_{\sigma^{-1}(j)}\), where
\(i=\min\operatorname{supp}\alpha\) and
\(j=\max\operatorname{supp}\beta\).
Condition (F-022.5) makes these two core vertices distinct.
The path is therefore simple and unique, and its product
is exactly \(z^{\alpha+\beta}\). It uses \(0\),
so no disjoint nontrivial cycle is available.

These paths give pure cofactors of both \(A_-\) and
\(A_+\). For the permanent this is the open-path minor
description in F-021, with no additional cycle terms because
the path occupies \(0\). For the determinant, differentiate
in the reverse entry to close the path into a cycle. If
the path has \(t\) edges, its negative labels contribute
\((-1)^t\), and the permutation sign of the resulting
\((t+1)\)-cycle is also \((-1)^t\). They cancel.
All unused loops are one and no other cycle is disjoint.
The resulting cofactor is the same positive monomial,
with coefficient one.

\begin{proof}[Coverage of every degree through \(2\ell\)]
Only degrees \(k\) with \(\ell<k\leq2\ell\) remain,
since the center-to-outward paths cover the smaller degrees.
Let \(\gamma\) have degree \(k\).

If its largest and smallest indices differ by at least two,
put a smallest index into \(\alpha\) and a largest into
\(\beta\), and distribute the remaining occurrences so
\(|\alpha|=\ell\), \(|\beta|=k-\ell\). Both degrees
are between one and \(\ell\). The smallest index \(i\)
is less than \(R\), and the largest \(j\) is at least
\(i+2\); hence \(j\ne\sigma(i)=i+1\).

If only one index occurs, split the multiplicity into
\(\ell\) and \(k-\ell\). Condition (F-022.5) follows
because \(\sigma\) has no fixed point.

The only remaining case uses adjacent indices \(t,t+1\)
with positive multiplicities \(a,b\). If \(a>\ell\),
put \(\ell\) low indices in \(\beta\), leaving
\(\alpha\) of degree \(k-\ell\). Its minimum and
\(\beta\)'s maximum are both \(t\), so the condition
holds. If \(b>\ell\), put \(\ell\) high indices in
\(\alpha\); its minimum and \(\beta\)'s maximum are
both \(t+1\), again giving the condition. These cases
cannot occur simultaneously because \(k\leq2\ell\).
Otherwise \(a,b\leq\ell\); put all high indices in
\(\alpha\) and all low indices in \(\beta\).
Now \(\sigma(t+1)\ne t\): if \(t+1<R\) its value
is \(t+2\), and at wraparound it is one whereas
\(t=R-1\geq2\). This uses \(R\geq3\), excluding a
two-cycle. Every case yields the required split and pure
cofactor. All degree bounds and the endpoint \(\ell=1\)
are included.
\end{proof}

\begin{proof}[Jet dominance and complete multiplicity spaces]
Vary a matrix entry by a linear form \(H\). The differential
of its determinant or permanent is \(H\) times the
corresponding cofactor. Taking \(H=x\) realizes all the
pure monomials just proved, and taking \(H=z_i\) supplies
every degree-\(2\ell+1\) monomial. Each target jet monomial
thus has an elementary coefficient direction whose derivative
is exactly that monomial. Distinct monomials require distinct
source coordinates. The chosen Jacobian minor is the identity
over \(\mathbb Z\).

The jet degree does not exceed the matrix size: the split
of \(z_i^{2\ell}\) gives a simple path of length
\(2\ell\), hence \(L_b(r)\geq2\ell+1=b\).
For a larger matrix append an \(xI\) block. Translation
by the appropriate coefficient vector in (F-022.4) gives,
on the chosen source subspace, \(F(t)=t+O(t^2)\)
integrally. The lowest nonzero homogeneous term of a
polynomial survives composition, proving injectivity after
arbitrary commutative base change and dominance over
\(\mathbb C\).

For the permanent, hold the diagonal at \(v_1\) fixed
as \(x\), and simultaneously permute its row and column
to the distinguished position \((1,1)\). This removes
none of the selected directions. The constant uses the
central diagonal; every other selected path is between
distinct vertices and uses an off-diagonal parameter.
This remains valid at \(\ell=1\), when no additional
tree vertices exist. Internal padding now gives exactly
\(x^{n-m}\operatorname{per}A_+\), whose dehomogenized
jets are unchanged.

All resulting pencil forms belong to the respective orbit
closures by endomorphism substitution. Every highest-weight
ideal vector with \(\lambda_2\leq b\) is a polynomial
in these jet coefficients by F-016. Pullback injectivity
forces it to be zero. Thus the whole ideal multiplicity
space vanishes, giving the two claimed ambient equalities.
The bounds also ensure that the \(r\)-variable space
fits in each matrix-entry ambient space, since \(L_b(r)\geq r\).
\end{proof}

\paragraph{An actual second-row gate at every even power.}
For every \(m\geq2\) and \(\ell\geq1\), an actual
internally padded permanent obstruction at
\(n\geq m^{2\ell}\) requires
\begin{equation}\label{eq:F022gate}
                              \lambda_2\geq2\ell+2.
\end{equation}
Indeed positive permanent multiplicity gives
\(r=\ell(\lambda)\leq m^2\). For \(r\leq3\),
O-006--O-007 already give full determinant multiplicity.
For \(r\geq4\), (F-022.2) proves
\(L_3(r)\leq m^2\) and \(L_5(r)\leq m^4\).
For \(\ell\geq3\), the proved M-013 estimate gives
\[
 L_{2\ell+1}(r)
    <2\binom{m^2+\ell-1}{\ell}
    \leq \frac58 m^{2\ell}\leq n.
\]
For completeness, its ratio to \(m^{2\ell}\) decreases
with \(m\geq2\), and its value at two is
\(2\binom{\ell+3}{3}/4^\ell\). This equals \(5/8\)
at \(\ell=3\), and decreases with ratio
\((\ell+4)/(4(\ell+1))\). Thus all parameter endpoints
are checked. The ambient determinant equality excludes
\(\lambda_2\leq2\ell+1\), proving the gate.
In particular \(n\geq m^2\) forces \(\lambda_2\geq4\),
and \(n\geq m^4\) forces \(\lambda_2\geq6\).

\paragraph{Grothendieck dictionary and limits.}
The coefficient targets are finite free quotients of the
polynomial algebra by powers of its augmentation ideal.
All matrices, cycle covers, cofactors, and chosen monomial
directions are integral. Their coordinate-algebra injections
are actual maps of affine schemes, universal under base
change. Passing to complete highest-weight spaces uses
F-016's Schur-functor support calculation and the comodule
formalism of SGA~3 I4.7.2--4.7.3 and SGA~6 VI3.3.
The extra information is the explicit integral injection,
not merely positivity in a representation \(K_0\).
Sharing core vertices allows cycles; the path/sign argument
proves precisely why they do not spoil the selected cofactors.
The results determine the stated full multiplicities and
exclude the stated actual gates, without any saturation
claim or comparison of still-undetermined evaluation ranks.

% END agents/formulas/shared_core_odd_jets.tex

% BEGIN proofs/coupled_core_even_jets.tex
\section{Coupled shared cores and the remaining power thresholds}

\subsection*{Finding M-016: even jets from a matrix of root bridges}
Let \(r\geq4\), \(\ell\geq1\), and put
\[
 R=r-1,\qquad b=2\ell+2,\qquad h=\lceil\sqrt r\rceil,
 \qquad
 L=2\binom{r+\ell-1}{\ell}-r,\qquad N=hL.
\]
For every \(n\geq N\), the determinant map from
\(n\)-by-\(n\) linear pencils in \(r\) variables dominates
the whole coefficient jet through degree \(b\).
Its differential admits an explicit integral right inverse.
In particular, in characteristic zero,
\begin{equation}\label{eq:M016K}
 \lambda\vdash nd,\quad\ell(\lambda)\leq r,\quad\lambda_2\leq b
       \quad\Longrightarrow\quad K_n(\lambda)=a_\lambda(d[n]).
\end{equation}
The useful small orders are
\[
       b=4:\ N=\lceil\sqrt r\rceil r,
       \qquad b=6:\ N=\lceil\sqrt r\rceil r^2.
                                                               \tag{M-016.1}
\]
The construction concerns detected coefficients. It asserts
neither exact representation of each truncated form nor
an algebraic section of the full determinant map.

\paragraph{The components and their root matrix.}
Use the shared-core graph of F-022 with variables
\(x,z_1,\ldots,z_R\) and cyclic permutation
\(\sigma=(1\ 2\ \cdots\ R)\). It has one root,
\(R\) shared degree-one vertices, and separate inward and
outward monomial-tree vertices of degrees two through \(\ell\).
Its size is \(L\). For \(\ell=1\), only the root and
shared degree-one vertices occur.
Take \(h\) disjoint copies, labeled \(1,\ldots,h\),
with roots \(o_1,\ldots,o_h\).

Since \(R\leq h^2-1\), inject the \(R\) variable labels
into the \(h^2-1\) ordered pairs \((a,c)\ne(h,h)\).
Let \(C_{ac}\) be the assigned variable, or zero when
that pair is unused. Add a directed bridge
\(o_a\longrightarrow o_c\) labeled \(C_{ac}\) for each
nonzero label. Diagonal bridges are allowed: they add a
linear term to that root's diagonal matrix entry.
In particular \(C_{hh}=0\).

Let \(B(z)\) be the resulting adjacency matrix and
\(A_0=xI-B(z)\). At \(x=1\), eliminate the nonroot
blocks, whose directed paths outside the roots have no cycles.
The shared degree-one star in each component contributes
the same quadratic polynomial
\[
                 s(z)=\sum_{i=1}^R z_i z_{\sigma(i)}.
\]
Thus the Schur complement on the roots is exactly
\[
                         W=I_h-C-sI_h,
                 \qquad \det A_0(1,z)=\det W.           \tag{M-016.2}
\]
All eliminated diagonal entries equal one; the higher tree
vertices cannot contribute a closed path. These observations
also prove (M-016.2) directly by the block determinant formula.

\paragraph{Cofactors detected by the root matrix.}
The pure monomial cofactors in F-022 can be chosen to use
the central root, with no path bypassing it. Consequently,
for every monomial \(z^\gamma\) of degree at most \(2\ell\)
and every component \(a\), an elementary coefficient
direction has cofactor
\[
                         z^\gamma\operatorname{adj}(W)_{aa}.
                                                               \tag{M-016.3}
\]
For degree zero this is the root diagonal cofactor itself.
For degree at most \(\ell\), use the path from the root
to an outward monomial vertex. For larger degrees use
F-022's split into two monomials, each of degree at most
\(\ell\), with distinct shared parent vertices. That
condition removes the direct path bypassing the root.
The local path factors multiply to \(z^\gamma\);
the remaining inverse entry is \((W^{-1})_{aa}\).
Multiplying by \(\det W\) gives (M-016.3).

For distinct components \(a,c\), and monomials
\(z^\alpha,z^\beta\), each of degree exactly \(\ell\),
the analogous cofactor from the inward vertex in component
\(a\) to the outward vertex in component \(c\) is
\[
                     z^{\alpha+\beta}\operatorname{adj}(W)_{ac}.
                                                               \tag{M-016.4}
\]
Different components have no nonroot vertices in common,
so there is no bypass term. With the adjacency convention
that \(B_{uv}\) labels \(u\to v\), the derivative direction
is the reverse matrix entry: its cofactor is
\(\det A_0(A_0^{-1})_{uv}\). This fixes the orientation
in (M-016.4).

The inverse calculations may first be made over
\(\mathbb Q(z)\), where the constant term one makes
\(W\) invertible. After multiplying by the determinant
they are identities of integral polynomials, and do not
restrict the coefficient parameter points.

\paragraph{First terms of the root cofactors.}
Let \(O(z^j)\) mean terms of total degree at least \(j\).
Since \(C\) is linear and \(s\) is quadratic, expansion
of the determinant and inverse at zero gives
\[
 \operatorname{adj}(W)
     =I_h+C-\operatorname{tr}(C)I_h+O(z^2).
\]
It follows that
\[
 \begin{split}
 \operatorname{adj}(W)_{ac}&=C_{ac}+O(z^2)\quad(a\ne c),\\
 \operatorname{adj}(W)_{aa}-\operatorname{adj}(W)_{hh}
                 &=C_{aa}+O(z^2),\\
 \operatorname{adj}(W)_{hh}&=1+O(z).
 \end{split}                                             \tag{M-016.5}
\]
These first terms can also be read directly from the
cofactor permutation expansion; their coefficients are integral.

\begin{proof}[Proof of the surjective integral differential]
Consider a target monomial of degree \(b=2\ell+2\).
Choose one of its variables as a bridge label and another
as the linear factor \(H\) in an entry variation.
The remaining multiset has degree \(2\ell\).

If the chosen bridge variable is in an off-diagonal slot
\((a,c)\), split the remaining multiset into \(\alpha,\beta\)
of degree \(\ell\). By (M-016.4) and (M-016.5),
the entry variation by \(H\) has, after truncation through
degree \(b\), exactly the chosen target monomial.
Every other term has degree at least \(2\ell+3\).

If the bridge variable is in a diagonal slot \((a,a)\),
it has \(a\ne h\). Write the remaining monomial as
\(z^\gamma\), \(|\gamma|=2\ell\). Take the difference
of the two elementary directions in (M-016.3) for components
\(a,h\), both varied by \(H\). The second line of
(M-016.5) again gives exactly the target monomial through
degree \(b\). Thus the differential contains the entire
top homogeneous coefficient space, over \(\mathbb Z\).

For degrees at most \(2\ell\), (M-016.3) in component
\(h\), varied by \(H=x\), has lowest term precisely
\(z^\gamma\) with coefficient one. Degree \(2\ell+1\)
uses a degree-\(2\ell\) cofactor times a variable.
Starting at degree \(b-1\) and proceeding downwards,
subtract integral combinations of already obtained higher
monomials. This yields every homogeneous target monomial
exactly. No division occurs.

We have therefore constructed integral linear combinations
of source coefficient directions which map to the monomial
basis of the full jet module. They give a right inverse to
the differential over \(\mathbb Z\). Restrict the pencil
parameters to the affine space generated by those directions
through \(A_0\), and translate its target by the coefficient
vector of \(\det W\). The restricted polynomial map is
\(F(t)=t+O(t^2)\). The lowest homogeneous term argument
of M-013 proves injectivity of its polynomial pullback after
every commutative base change. In particular it proves
dominance in characteristic zero.

For \(n>N\), append \(xI_{n-N}\); after \(x=1\)
the same directions and differential apply. Finally F-016
puts every highest-weight function with \(\lambda_2\leq b\)
in the detected coefficient algebra. Injective pullback
annihilates no nonzero determinant-ideal vector there.
This proves \eqref{eq:M016K} for every multiplicity copy.
\end{proof}

\paragraph{A complete second-row gate at integral power thresholds.}
For every pair of integers \(m\geq2,c\geq1\), an actual
internally padded permanent obstruction at \(n\geq m^c\)
requires
\begin{equation}\label{eq:M016powers}
                              \lambda_2\geq c+2.
\end{equation}
For even \(c\geq2\) this is F-022. For odd
\(c=2\ell+1\geq3\), a positive permanent multiplicity
has length \(r\leq m^2\). The cases \(r\leq3\)
are already excluded by O-006--O-007. For \(r\geq4\),
the size above is at most
\[
 m\left(2\binom{m^2+\ell-1}{\ell}-m^2\right).
\]
The monotonicity used here follows from
\(L(r+1)-L(r)=2\binom{r+\ell-1}{\ell-1}-1\geq1\).
At \(\ell=1\) and \(\ell=2\) this is exactly
\(m^3\) and \(m^5\), respectively. At \(\ell\geq3\),
M-013.2 bounds it by \((5/8)m^{2\ell+1}\).
Thus \(n\geq m^c\) gives full determinant multiplicity
whenever \(\lambda_2\leq c+1\), proving the assertion.

For \(c=1\), in fact every ambient type with second row
at most two has full determinant multiplicity. F-014 and
the quadratic multiplicity-free decomposition in O-013
show that its only nonzero non-one-row types are
\((nd-2(r-1),2^{r-1})\), each with ambient multiplicity one.
The nonzero determinant Hessian section in O-013 realizes
that copy for every allowed \(r\leq\min(d,n^2)\).
One-row types also have full multiplicity. Hence no actual
obstruction has \(\lambda_2\leq2\), completing
\eqref{eq:M016powers}.

Equivalently, every actual obstruction with second row
\(b=\lambda_2\) satisfies the useful degree-independent
restriction \(n<m^{b-1}\). It is a necessary restriction,
not a sufficient test for an obstruction. At \(c=25\)
it recovers \(\lambda_2\geq27\); the new endpoint
ranges here are \(c=1,2,3,4,5\).

\paragraph{Grothendieck dictionary.}
The root Schur complement and cofactor formulas are
identities over \(\mathbb Z\). The finite jet module is
locally free, and the explicit integral splitting of the
differential produces a universally injective polynomial
coordinate pullback. Its monomial supports are torus
comodules under SGA~3, I, Proposition~4.7.3, pp.~35--36.
The Schur-copy support used in the final characteristic-zero
step is SGA~6, VI, Theorem~3.3, pp.~392--393, as implemented
by F-016's Cauchy flag-minor argument. The resulting
vanishing concerns the entire ideal comodule, in the sense
of SGA~3, I, Proposition~4.7.2, pp.~34--35.
Sharing and coupling vertices retains explicit coordinate
maps; an equality of virtual \(K_0\) classes would not
provide the differential or its integral splitting.

% END proofs/coupled_core_even_jets.tex

% BEGIN agents/jet_review/permanent_coupled_core_jets.tex
\section{Full permanent multiplicities from coupled shared cores}

\subsection*{Finding F-023: the even-order permanent jet map}
Let \(r\geq4\), \(\ell\geq1\), and define
\[
 R=r-1,\qquad b=2\ell+2,\qquad h=\lceil\sqrt r\rceil,
 \qquad L=2\binom{r+\ell-1}{\ell}-r,\qquad N=hL.
\]
If \(n\geq m\geq N\), then every partition
\(\lambda\vdash nd\) with \(\ell(\lambda)\leq r\)
and \(\lambda_2\leq b\) satisfies
\begin{equation}\label{eq:F023full}
           q^m_\lambda(d[n])=K_n(\lambda)=a_\lambda(d[n]).
\end{equation}
Here the permanent is internally padded by
\(X_{11}^{n-m}\). The equality includes every multiplicity
copy and every outer degree \(d\). Its geometric input is an
integrally split differential of the permanent coefficient-jet
map, with the distinguished matrix entry held equal to the
first variable. All representation decompositions below are
in characteristic zero.

\paragraph{A positive graph and its root cycle covers.}
Use exactly the \(h\) coupled shared-core components of
M-016: their local graphs, roots \(o_a\), monomial trees,
cyclic permutation \(\sigma\), and root bridge matrix
\(C=(C_{ac})\), with \(C_{hh}=0\). All nonzero labels
are positive variables. Let \(B(z)\) be the adjacency
matrix and take \(A_+=xI+B(z)\), where
\(z=(z_1,\ldots,z_R)\). Put
\[
 s=\sum_{i=1}^R z_i z_{\sigma(i)},\qquad
                       P=I_h+C+sI_h.
                                                               \tag{F-023.1}
\]
We claim
\[
                      \operatorname{per}A_+(1,z)
                                  =\operatorname{per}P.
                                                               \tag{F-023.2}
\]
Indeed every nontrivial cycle either is a local two-edge
star cycle through one root, or is a cycle of bridges among
roots. Higher inward and outward tree vertices belong to
no cycle. In a cycle cover, a root not used by a bridge
cycle can use its diagonal weight \(1+C_{aa}\), or choose
one local star cycle, whose total weight is \(s\).
All unused nonroot vertices use their diagonal loops.
These choices give precisely the permutation expansion of
\(\operatorname{per}P\). This is a positive cycle-cover
identity over \(\mathbb Z\), without a permanent version
of Schur-complement elimination.

Define the oriented root minors
\[
 Q_{ac}=\operatorname{per}P_{\widehat c,\widehat a},
                                                               \tag{F-023.3}
\]
where the subscript means deleting row \(c\) and column
\(a\). For \(a\ne c\), this is the minor whose open
path runs from root \(a\) to root \(c\). The first terms
of their permutation expansions are
\[
 \begin{split}
 Q_{ac}&=C_{ac}+O(z^2)\qquad(a\ne c),\\
 Q_{aa}&=1+\operatorname{tr}C-C_{aa}+O(z^2),\\
 Q_{hh}-Q_{aa}&=C_{aa}+O(z^2),\qquad
 Q_{hh}=1+O(z).
 \end{split}                                             \tag{F-023.4}
\]
For the second line the linear terms are the choices of one
nonconstant diagonal entry in the remaining roots; an off-diagonal
permutation cycle uses at least two bridges. For the first line,
the sole degree-one open path is the bridge \(a\to c\).
Thus every displayed leading coefficient is exactly one.

\paragraph{The selected full-matrix minors.}
For every monomial \(z^\gamma\) of degree at most
\(2\ell\), F-022 provides, in each component \(a\),
a local simple path of that weight which uses its root.
The degree-zero case means varying the root diagonal.
For positive degree the endpoints are distinct; for a path
with both endpoints nonroot, the shared core vertices at
its inward and outward ends are distinct. This removes
any path bypassing the root.

The corresponding full-matrix permanent cofactor is exactly
\[
                                  z^\gamma Q_{aa}.
                                                               \tag{F-023.5}
\]
To see this directly, a cofactor permutation has one open
path between the chosen endpoints and disjoint cycles.
The local path occupies root \(a\), so that root cannot
be revisited or belong to another cycle. A path leaving and
returning to that root would revisit it and is forbidden.
The remaining choices at all other roots are exactly the
cycle covers of \(P_{\widehat a,\widehat a}\).
All remaining nonroot vertices of component \(a\) are
forced to use diagonal loops. This proves (F-023.5).

For distinct components \(a,c\), take an inward monomial
vertex of weight \(z^\alpha\) in component \(a\) and
an outward one of weight \(z^\beta\) in component \(c\),
with \(|\alpha|=|\beta|=\ell\). The full-matrix
cofactor is
\[
                                  z^{\alpha+\beta}Q_{ac}.
                                                               \tag{F-023.6}
\]
Here the open path starts with the unique local inward path,
uses an open root path \(a\to c\), and ends with the
unique local outward path. It cannot bypass a root because
the endpoints belong to distinct components. On every root
used by that path the local star cycles are unavailable;
on the other roots the same choices as in (F-023.2) remain.
The root open path and its disjoint cycles are exactly the
minor expansion (F-023.3). With the adjacency convention
\(B_{uv}\) for \(u\to v\), the varied matrix entry
is the reverse entry \((v,u)\). Its cofactor deletes row
\(v\) and column \(u\); this explains the deletion order
in (F-023.3) and fixes the orientation of (F-023.6).

\begin{proof}[Integral jet dominance with the padding entry fixed]
For a target monomial of degree \(b=2\ell+2\), remove
one occurrence as a bridge label and another as the linear
entry variation \(H\). The remaining monomial has degree
\(2\ell\). If the bridge occupies an off-diagonal slot
\((a,c)\), split that remainder into two degree-\(\ell\)
monomials. Equations (F-023.6) and (F-023.4), multiplied
by \(H\), give exactly the target monomial modulo terms
of degree greater than \(b\).

If the bridge occupies a diagonal slot \((a,a)\), then
\(a\ne h\). Use the same remainder \(z^\gamma\)
in components \(h,a\), and subtract their two elementary
directions in that order. By (F-023.5), their difference
is \(H z^\gamma(Q_{hh}-Q_{aa})\), whose truncation
is exactly the target monomial. Thus the entire top degree
is in the image of the differential over \(\mathbb Z\).

For every lower degree at most \(2\ell\), use
(F-023.5) in component \(h\), varied by \(H=x\).
The resulting cofactor has the required monomial as its
lowest term, with coefficient one. Degree \(2\ell+1\)
comes from a degree-\(2\ell\) cofactor times a variable.
Descending by degree and subtracting the higher monomials
already obtained gives every target monomial exactly.
Only integral sums and differences are used. This constructs
an integral linear right inverse to the differential.

Hold one nonroot core diagonal, for instance the diagonal
at \(v_1\) in the first component, fixed as \(x\).
None of these selected directions varies it. The constant
cofactor uses a root diagonal, and all positive-degree
local and cross-component paths have distinct endpoints.
Their elementary variations are off diagonal. The diagonal
comparison in the top-degree argument compares two such
off-diagonal cofactor directions, rather than varying a
nonroot diagonal. Simultaneously permute rows and columns
to place the fixed entry at \((1,1)\).

The order of the jet is allowed: F-022 gives
\(L\geq2\ell+1\), and \(h\geq2\), so \(N\geq b\).
For \(m>N\), append an \(xI_{m-N}\) block. Restrict
the source to the integral affine space specified by the
right-inverse directions and translate the target by its
base coefficient vector. The resulting polynomial map is
\(F(t)=t+O(t^2)\) over \(\mathbb Z\). The lowest
nonzero homogeneous term of any polynomial survives this
substitution, so pullback is injective after every commutative
base change. In particular the jet projection is dominant
over \(\mathbb C\).

Substitution into the internally padded permanent is exactly
\(x^{n-m}\operatorname{per}A_+\), since its distinguished
entry is fixed as \(x\). After \(x=1\), this has the
same detected coefficients. Every such substituted form lies
in the permanent orbit closure. A highest-weight ideal vector
with \(\lambda_2\leq b\) is a polynomial in these jet
coordinates by F-016. Injective pullback forces it to vanish
as a polynomial. Hence no ambient highest-weight copy is lost
on the permanent closure, proving \(q^m_\lambda=a_\lambda\).
M-016 gives \(K_n=a_\lambda\) because \(n\geq m\geq N\).
This proves \eqref{eq:F023full}. If the ambient multiplicity
is zero, the same conclusion holds automatically.
\end{proof}

\paragraph{A smaller explicit positive family beyond every exact cover.}
At \(r=1024\), \(b=4\), one has \(h=32\), \(L=r\),
and \(N=32768\). Take
\[
 m=32768,\quad n=32769,\quad d=2048=2r,\qquad
                  \lambda=(nd-4(r-1),4^{r-1}).
                                                               \tag{F-023.7}
\]
The full nested coefficient support is the degree-four
simplex in \(R=1023\) variables. The checked double-counting
bound \eqref{eq:F021cover} in F-021 gives
\[
          \kappa(\Gamma)\geq1+\left\lceil
                         \frac{1023\cdot1022}{12}\right\rceil
                      =87127>n>m.
\]
Thus every exact two-tree support cover of F-018 fails its
size condition, while \eqref{eq:F023full} applies.
The two independent highest-weight polynomials constructed
in F-021---the square of the leading Hessian determinant,
and \(I(F)F(e_1)^r\) from F-011 at \(k=4\)---have
the displayed weight for these values of \(n,d\).
The first is not divisible by \(F(e_1)\), whereas the
second is; their proof of independence is unchanged.
Consequently this example satisfies
\[
                              2\leq q^m_\lambda=K_n=a_\lambda.
\]
The comparison is with every support cover in that explicit
architecture. It is not a lower bound for arbitrary
representations of the forms and does not provide an
actual obstruction.

\paragraph{Grothendieck dictionary.}
The graph, cycle-cover contractions, root minors and affine
coefficient maps are defined over \(\mathbb Z\). The jet
module is finite locally free, and the integral differential
splitting yields an injection of coordinate algebras universal
under base change. Torus weights and ideal comodules are
those of SGA~3, I, Propositions~4.7.2--4.7.3, pp.~34--36.
The passage to all highest-weight copies is F-016's explicit
Schur-functor support calculation, using SGA~6, VI,
Theorem~3.3, pp.~392--393, and Macdonald, I(4.3), p.~63,
Appendix~A(5.3)--(5.6), pp.~155--156.
An equality in representation \(K_0\) alone supplies
neither the coordinate injection nor the fixed internal
padding entry; both are established by the actual integral
morphisms here. The result determines full permanent and
determinant multiplicities in the stated range, without
claiming a general saturation principle or a comparison
outside that range.

% END agents/jet_review/permanent_coupled_core_jets.tex

% BEGIN proofs/prefix_extended_jets.tex
\section{Extending detected jets in selected coordinate directions}

\subsection*{Finding M-017: a combined second-row and later-row jet criterion}
Let \(r\geq2\), \(\ell\geq0\), \(t\geq0\), and
\(2\leq k\leq r+1\) be integers. Define
\[
 b=2\ell+2,\quad R=r-1,\quad p=k-2,\quad
 T=\lceil\sqrt R\rceil\binom{R+\ell}{\ell},\qquad
 N=T\left(1+\binom{p+t}{t}\right).
                                                               \tag{M-017.1}
\]
For every \(n\geq\max\{N,b+t\}\), the determinant
map on \(n\)-by-\(n\) linear pencils dominates the
coefficient projection with nonfirst-variable support
\[
 \Gamma=\left\{\gamma\in\mathbb Z_{\geq0}^{r-1}:
       |\gamma|\leq b+t,
       \ \sum_{i=k}^r\gamma_i\leq b\right\}.
                                                               \tag{M-017.2}
\]
The empty sum at \(k=r+1\) is zero. Consequently,
for every outer degree \(d\),
\begin{equation}\label{eq:M017K}
 \lambda\vdash nd,\quad\ell(\lambda)\leq r,\quad
 \lambda_2\leq b+t,\quad\lambda_k\leq b
       \quad\Longrightarrow\quad K_n(\lambda)=a_\lambda(d[n]).
\end{equation}
Missing parts of a partition are zero. When \(p\geq1\),
the constructed graph itself has a path of length \(b+t-1\),
so \(N\geq b+t\) and the extra maximum in the size
hypothesis is redundant. When \(p=0\), the target is just
the jet of order \(b\), independent of \(t\); in that
case \(n\geq N\) suffices.

\paragraph{An extension at each inward vertex.}
Start with the acyclic graph of M-013 at jet order \(b\).
It has \(T\) inward vertices and \(T\) outward vertices,
and its selected inward-to-outward inverse entries contain
every monomial of degree one through \(b-1\), each as
a single path product. All variables here are
\(x_1,x_2,\ldots,x_r\).

For each original inward vertex \(v\), make a separate
monomial tree on multi-indices \(\alpha\) in the first
\(p\) nonfirst variables \(x_2,\ldots,x_{k-1}\), with
\(|\alpha|\leq t\). Identify this new tree's root
with \(v\), and identify no other vertices. Orient the
new edges toward the root, label them by the variable
removed by the parent rule, and use positive labels.
There are \(\binom{p+t}{t}-1\) new vertices at each
old inward vertex. The resulting total number is exactly
\(N\) in (M-017.1).

The graph remains acyclic. A path beginning in a new tree
travels inward to its old vertex, then follows old inward
edges, crosses at most one old bridge, and travels outward.
No edge leads from the old graph into a new nonroot vertex.
In particular, if an old selected inverse entry is the
monomial \(x^\eta\), the entry from its new vertex
\(\alpha\) to the same outward endpoint is exactly
\[
                              x^\alpha x^\eta.
                                                               \tag{M-017.3}
\]
Uniqueness of the two path segments proves this identity;
the other attached trees introduce no alternate directed path.
For its adjacency matrix \(B\), use \(A_0=x_1I-B\).
Then \(\det A_0=x_1^N\), and after \(x_1=1\)
the selected determinant cofactors are the monomials in
(M-017.3), just as in M-013.

\begin{proof}[Coverage and dominant coefficient projection]
Every \(\gamma\in\Gamma\) with \(|\gamma|\leq b\)
is already a detected direction of M-013. Suppose instead
that \(b<|\gamma|\leq b+t\), and put
\(j=|\gamma|-b\). The suffix bound in (M-017.2)
implies that the total exponent in the early variables is
at least \(j\). Choose a degree-\(j\) submonomial
\(x^\alpha\) there. It has \(j\leq t\).
The remaining monomial has degree exactly \(b\).
Remove one variable from it for the linear entry variation
\(H\), leaving a monomial of degree \(b-1\geq1\).

The old M-013 graph represents that degree-\(b-1\)
monomial as an inward-to-outward inverse entry. Extend its
source by \(\alpha\) in the corresponding new tree.
Equation (M-017.3), multiplied by \(H(1,x)\), gives
exactly \(x^\gamma\) as the derivative of one matrix-entry
coefficient. Thus every target coefficient monomial is an
elementary derivative with coefficient one. Distinct target
monomials use distinct source coordinates, so the selected
Jacobian minor is the identity over \(\mathbb Z\).

The relevant paths also check the degree bound stated above.
If \(p\geq1\), take an early pure power of degree \(t\)
and an old pure-power path of degree \(b-1\).
Their concatenation is a simple path of length \(b+t-1\),
and uses \(b+t\) distinct vertices. Hence \(N\geq b+t\).
For \(p=0\), only the old jet target of degree \(b\)
occurs, and \(N=2T\geq b\).

Append an \(x_1I\) block when \(n>N\). Restricting
to the selected coefficient directions and translating the
target by its value at \(A_0\) yields \(F(t)=t+O(t^2)\)
over \(\mathbb Z\). The lowest homogeneous term proves
injectivity of polynomial pullback after any commutative
base change, and therefore dominance over \(\mathbb C\).

By F-016, a highest-weight polynomial of type \(\lambda\)
under the hypotheses of \eqref{eq:M017K} uses only
coefficient generators with total nonfirst degree at most
\(\lambda_2\leq b+t\) and suffix degree at most
\(\lambda_k\leq b\). It is a polynomial in the target
coordinates (M-017.2). If it lies in the determinant ideal,
it vanishes after the proved dominant projection and must
be zero. This argument applies to every highest-weight
copy, proving the complete ambient multiplicity equality.
\end{proof}

\paragraph{A further family of BIP necessary conditions.}
Define the explicit integer
\[
                  C_*=
       \left\lfloor\frac{190^{24}}{\binom{190^2+11}{12}}\right\rfloor.
                                                               \tag{M-017.4}
\]
Exact integer division gives \(C_*=478126755\).
At \(n\geq m^{25}\), every actual multiplicity obstruction
with \(\lambda_2\leq26+t\) satisfies
\begin{equation}\label{eq:M017BIP}
 \lambda_k\geq27
 \quad\text{whenever}\quad
 2\leq k\leq r+1,\qquad
                      1+\binom{k-2+t}{t}\leq C_*.
\end{equation}
This uses the known permanent-size reduction \(m\geq190\)
for an actual obstruction. It adds a coupled restriction
to the earlier single-row staircase; no source threshold
has been changed silently.

\begin{proof}
Positive permanent multiplicity gives \(r\leq m^2\).
Use (M-017.1) with \(b=26\), \(\ell=12\).
Since \(\lceil\sqrt{r-1}\rceil\leq m\),
\[
            T\leq m\binom{m^2+11}{12}.
\]
The product formula shows that
\(\binom{m^2+11}{12}/m^{24}\) is nonincreasing
for \(m\geq2\): up to the constant \(1/12!\)
it is \(\prod_{i=0}^{11}(1+i/m^2)\).
By definition of \(C_*\), for every \(m\geq190\)
one therefore has
\[
                  C_*\binom{m^2+11}{12}\leq m^{24}.
\]
If the binomial condition in \eqref{eq:M017BIP} holds,
the constructed size is at most \(m^{25}\leq n\).
The path argument checks any needed target degree bound.
If also \(\lambda_k\leq26\), the full multiplicity
equality just proved contradicts an actual obstruction.
This proves \eqref{eq:M017BIP}.
\end{proof}

For example, take \(k=3\). Then the binomial condition
is \(t+2\leq C_*\), so every BIP actual obstruction
must satisfy
\[
                     \lambda_3\geq27
                         \quad\text{or}\quad\lambda_2\geq C_*+25.
                                                               \tag{M-017.5}
\]
At the smallest allowed second row, \(\lambda_2=27\),
put \(t=1\). The condition becomes \(k\leq C_*\).
Thus such a candidate must have at least \(C_*\) rows,
and its rows \(2,\ldots,C_*\) all equal twenty-seven.
If it had fewer rows, taking \(k=r+1\) would contradict
\eqref{eq:M017BIP}. These are restrictions on remaining
candidates, not assertions that any such candidate is an
obstruction or even has positive permanent multiplicity.

\paragraph{Geometric dictionary and distinction from exact containment.}
The selected target is a finite torus-weight coefficient
module defined by two nested degree bounds. The prefix
extensions are integral affine-matrix morphisms, and the
unit Jacobian gives an actual universally injective coordinate
pullback. The relevant gradings are SGA~3, I,
Proposition~4.7.3, pp.~35--36; all Schur copies are
controlled by F-016's use of SGA~6, VI, Theorem~3.3,
pp.~392--393, and Macdonald, I(4.3), p.~63, Appendix~A
(5.3)--(5.6), pp.~155--156.
Higher coefficient coordinates remain free. Dominance of
the projection does not supply exact representations of
every form supported on \(\Gamma\), and does not follow
from a virtual \(K_0\) class. The strict positive-copy
example O-022 separately verifies that this distinction
can exclude BIP shapes missed by every exact support cover
of F-018 and every entire one-cut family.

% END proofs/prefix_extended_jets.tex

% BEGIN proofs/order_ideal_jet_extensions.tex
\section{Jet extensions by arbitrary monomial order ideals}

\subsection*{Finding M-018: a finite support criterion for detected coefficient jets}
Fix \(r\geq2\), \(\ell\geq0\), and put
\[
 R=r-1,\qquad b=2\ell+2,\qquad
 T=\lceil\sqrt R\rceil\binom{R+\ell}{\ell}.
\]
Let \(A\subset\mathbb Z_{\geq0}^R\) be a nonempty finite
monomial order ideal: it contains every componentwise smaller
multi-index than any of its members. Write
\(D_b=\{\eta:|\eta|\leq b\}\), and let
\(\Gamma\subseteq A+D_b\) be any finite target support.
Then the determinant coefficient projection onto \(\Gamma\)
is dominant at every size
\[
                         n\geq N:=T(1+|A|).                 \tag{M-018.1}
\]
The differential has an elementary unit minor over \(\mathbb Z\).
With the internally padded permanent, the same target projection
is dominant whenever \(n\geq m\geq N\).

In particular, let \(\lambda\vdash nd\) have length at most \(r\),
and define its nested coefficient support
\[
 \Gamma_\lambda=\left\{\gamma\in\mathbb Z_{\geq0}^{r-1}:
       \sum_{i=j}^r\gamma_i\leq\lambda_j\quad(2\leq j\leq r)
                   \right\}.
                                                               \tag{M-018.2}
\]
If \(\Gamma_\lambda\subseteq A+D_b\), (M-018.1) implies
\(K_n(\lambda)=a_\lambda(d[n])\). Under the permanent-size
hypothesis it implies \(q^m_\lambda=K_n=a_\lambda\).
The support condition is finite and explicit; it need not have
an efficient optimization algorithm. M-017 is its special case
where \(A\) is a simplex in an initial segment of the variables.

\begin{proof}
Start with M-013's acyclic graph, with \(T\) inward and
\(T\) outward vertices. Its selected inward-to-outward inverse
entries are pure monomials of every degree one through \(b-1\).
At each inward vertex \(v\), attach a separate tree with vertices
indexed by \(A\), identifying its zero vertex with \(v\).
For every nonzero \(\alpha\in A\), choose any supported index
\(i\); its parent \(\alpha-e_i\) lies in \(A\).
Orient the edge toward the parent and label it by \(z_i\).
Different attached trees have no common nonroot vertices.
There are \(T(|A|-1)\) new vertices, proving the size (M-018.1).

The graph is still acyclic, and no old edge enters a new nonroot
vertex. Consequently an old selected monomial inverse entry
\(z^\eta\), extended at its inward source by \(\alpha\),
becomes exactly \(z^{\alpha+\eta}\). There is one directed
path, not a sum of paths. There is also exactly one path from
the new vertex \(\alpha\) to its own old root \(v\), with
monomial \(z^\alpha\). When \(\alpha=0\), the corresponding
diagonal cofactor is one.

For \(\gamma\in\Gamma\), choose a decomposition
\(\gamma=\alpha+\eta\), \(\alpha\in A\), \(|\eta|\leq b\).
If \(|\eta|\geq2\), remove one variable from \(z^\eta\)
and use that variable as the linear entry variation \(H\).
The remaining degree is between one and \(b-1\), hence is
an old selected inverse entry. Its extension by \(\alpha\)
gives derivative exactly \(z^\gamma\).
If \(|\eta|\leq1\), use the path from \(\alpha\) to its
own old root, varying the reverse entry by \(H=z^\eta\)
when \(|\eta|=1\), and by the homogenizing variable \(x\)
when \(\eta=0\). These again give exactly \(z^\gamma\)
after \(x=1\).

For \(xI-B\), an acyclic path gives its positive monomial
determinant cofactor, as proved in M-013. For \(xI+B\),
the same path gives its positive permanent minor, as proved
in F-021. There are no disjoint nontrivial cycles in either
case. Every target monomial therefore is the derivative of
an elementary source coordinate, with coefficient one and
no other monomial. Distinct target monomials cannot use the
same source coordinate. Their Jacobian minor is the identity.

The degree bound is automatic. A longest new path, of length
\(\max_{\alpha\in A}|\alpha|\), can be attached at the
source of an old degree-\(b-1\) path. Their concatenation
is simple, so \(N\geq\max_{\alpha\in A}|\alpha|+b\).
In particular every target monomial has degree at most \(N\).
Append \(xI\) blocks at larger sizes. Translation and restriction
to the selected directions give \(F(t)=t+O(t^2)\) integrally;
the least homogeneous term proves injective coordinate pullback
after every commutative base change.

For the internally padded permanent, hold one old outward
diagonal entry fixed at \(x\), and move it to position \((1,1)\).
Every nonconstant selected path has distinct endpoints; the
constant cofactor can be taken at an inward diagonal. No selected
direction changes the fixed outward diagonal. Thus multiplication
by its \((n-m)\)-th power preserves all dehomogenized target
coefficients and their dominant projection.

Finally, F-016 places every highest-weight polynomial of type
\(\lambda\) in the algebra generated by the coordinates
(M-018.2). The proved pullback is injective on this entire
polynomial algebra, so no nonzero highest-weight ideal vector
of that type vanishes on the corresponding orbit closure.
This proves all multiplicity assertions, copy by copy.
\end{proof}

\paragraph{A concrete further BIP restriction.}
Every actual obstruction at \(n\geq m^{25}\) satisfies
\[
 \lambda_2\geq1000001\quad\text{or}\quad
 \lambda_3\geq29\quad\text{or}\quad\lambda_4\geq27.
                                                               \tag{M-018.3}
\]
Indeed use \(b=26\), and let \(A\) consist of monomials
\(x_2^a x_3^c\) with
\[
                  a+c\leq999974,\qquad c\leq28.
\]
This order ideal has
\[
 |A|=\sum_{c=0}^{28}(999975-c)=28998869,
                \qquad 1+|A|=28998870<478126755=C_*.
\]
If all three inequalities in (M-018.3) fail, every generator
\(\gamma\in\Gamma_\lambda\) has degree at most one million,
suffix degree from coordinate four at most twenty-six, and
individual coordinate-three exponent at most twenty-eight.
When \(|\gamma|>26\), choose an early submonomial of degree
\(|\gamma|-26\); it lies in \(A\). When \(|\gamma|\leq26\),
use the zero member of \(A\). Thus
\(\Gamma_\lambda\subseteq A+D_{26}\).
The uniform size estimate of M-017, with \(m\geq190\),
gives \(N\leq m^{25}\), and the full determinant multiplicity
excludes an obstruction. Partitions with fewer than four rows
were already excluded by O-006--O-007; zero padding gives the
same argument whenever needed.

\paragraph{Geometric dictionary.}
An order ideal indexes a finite free monomial coefficient module.
The graph constructs an integral morphism between affine spaces,
and its elementary unit Jacobian proves universally injective
coordinate pullback. The multigrading is a torus comodule in the
sense of SGA~3, I, Proposition~4.7.3, pp.~35--36. The vanishing
of the complete ideal multiplicity space uses the comodule
formalism of Proposition~4.7.2, pp.~34--35, and the Schur
operations of SGA~6, VI, Theorem~3.3, pp.~392--393, through
F-016. The decisive data are actual maps of coefficient schemes.
The assertion concerns dominance of a projection; it does not
give exact representations of all polynomials with support
\(\Gamma\), nor a saturation theorem for virtual classes.

% END proofs/order_ideal_jet_extensions.tex

% BEGIN proofs/two_sided_jet_extensions.tex
\section{Two-sided extensions of a coefficient-jet core}

\subsection*{Finding M-019: a two-ideal support criterion for full multiplicities}
Let \(r\geq2\), \(\ell\geq0\), and define
\[
 R=r-1,\quad b=2\ell+2,\quad h=\lceil\sqrt r\rceil,
 \quad Q=\binom{R+\ell}{\ell},\quad T=hQ.
\]
Let \(A,B\subset\mathbb Z_{\geq0}^{R}\) be nonempty
finite monomial order ideals. Suppose that a finite coefficient
support satisfies
\[
                         \Gamma\subseteq A+B+D_b,
      \qquad D_b=\{\eta\in\mathbb Z_{\geq0}^{R}:|\eta|\leq b\}.
                                                               \tag{M-019.1}
\]
Then the determinant coefficient projection onto \(\Gamma\)
is dominant whenever
\[
                            n\geq N=T(|A|+|B|).          \tag{M-019.2}
\]
The internally padded permanent has the same dominant projection
whenever \(n\geq m\geq N\). Both differentials have explicit
elementary unit minors over \(\mathbb Z\).

Consequently, if \(\lambda\vdash nd\), \(\ell(\lambda)\leq r\),
and its full nested support \(\Gamma_\lambda\) from M-018
satisfies (M-019.1), then (M-019.2) gives
\(K_n(\lambda)=a_\lambda(d[n])\). Under the permanent-size
hypothesis, \(q^m_\lambda=K_n=a_\lambda\).
These assertions retain every highest-weight copy.

\paragraph{The constant bridge.}
Use \(h\) inward and \(h\) outward copies of the monomial
tree on \(R\) variables \(z_1,\ldots,z_R\), through degree
\(\ell\), just as in M-013. Each copy has \(Q\) vertices.
The inward copies point toward their roots, and the outward
copies point away from their roots. Since \(h^2\geq r=R+1\),
inject the \(R+1\) labels
\[
                             x,z_1,\ldots,z_R
\]
into distinct ordered pairs of an inward root and an outward
root, and add the corresponding directed bridges. Unused
pairs have no edge. The label \(x\) becomes one after
dehomogenization; it is still a linear homogeneous matrix
entry before that step. The graph remains acyclic.

Every monomial of degree between zero and \(b-1=2\ell+1\)
is now a pure inward-to-outward path product after \(x=1\).
For degree zero use the \(x\)-bridge with its two roots.
For positive degree, choose one factor for a variable bridge
and split the remaining factors into two monomials of degrees
at most \(\ell\). Their inward and outward tree paths,
with that bridge, give exactly the required monomial.
For fixed endpoints the path is unique: it must use the
single bridge between their two components. In particular
the added constant bridge introduces no alternate path.

\paragraph{The two attached order ideals.}
At each of the \(T\) old inward vertices, attach a separate
tree indexed by \(A\), identifying its zero vertex with
that old vertex and pointing every new edge toward it.
At each of the \(T\) old outward vertices, attach a separate
tree indexed by \(B\), again identifying the zero vertex,
but pointing every new edge away from the root.
For every nonzero multi-index, choose a supported variable
to remove for its parent; the order-ideal property ensures
that parent exists. Labels are the removed variables.
All nonroot attached vertices are disjoint.

The total size is
\[
        2T+T(|A|-1)+T(|B|-1)=T(|A|+|B|)=N.
\]
A path from an attached inward vertex \(\alpha\) to an
attached outward vertex \(\beta\), using an old selected
path of weight \(z^\eta\), is unique and has weight
\[
                              z^{\alpha+\eta+\beta}
                                                               \tag{M-019.3}
\]
after \(x=1\). No old edge enters a new inward nonroot
vertex, and no new outward nonroot vertex leads back into
the old graph. Thus the path must have precisely these
three segments. There are no directed cycles.

\begin{proof}[Dominance, fixed padding and all copies]
Let \(\mathcal B\) be the homogeneous adjacency matrix.
At the base pencils \(xI-\mathcal B\) and
\(xI+\mathcal B\), acyclicity gives determinant and
permanent both equal to \(x^N\). In the determinant,
the reverse-entry cofactor of a directed path is its
positive monomial product: its edge signs cancel the
cycle permutation sign. In the permanent the open-path
minor has the same product with positive sign. There
are no disjoint nontrivial cycles. These statements
also apply to a path using the constant \(x\)-bridge.

For \(\gamma\in\Gamma\), choose
\(\gamma=\alpha+\beta+\eta\) as in (M-019.1).
If \(|\eta|\leq b-1\), use its selected core path
and multiply the cofactor (M-019.3) by the linear entry
variation \(H=x\). If \(|\eta|=b\), remove one variable
for \(H\) and use the core path of the remaining degree
\(b-1\). Each resulting derivative is exactly
\(z^\gamma\) after \(x=1\), with no extra monomials.
This includes \(\eta=0\), which is precisely the purpose
of the constant bridge. Different target monomials use
different elementary source coordinates, so the selected
Jacobian minor is the identity over \(\mathbb Z\).

The degree bound causes no extra hypothesis. Attach a
maximum-degree member of \(A\) at the source of an old
degree-\(b-1\) path, and a maximum-degree member of \(B\)
at its target. The resulting simple path has
\(\max\deg A+(b-1)+\max\deg B\) edges. Therefore
\[
                       N\geq\max\deg A+\max\deg B+b,
\]
so all target coefficients are valid at degree \(N\).
The degree-zero core path has one edge rather than zero,
which only makes this bound easier in that endpoint case.
For larger matrices append \(xI\) blocks.

Restricting to the selected source coefficient directions
and translating the target by its base value gives
\(F(t)=t+O(t^2)\) over \(\mathbb Z\). The least
homogeneous term of a nonzero polynomial survives pullback.
This proves injection of the polynomial coordinate algebra
after every commutative base change, and in particular
dominance over \(\mathbb C\).

Every selected derivative varies an entry with one inward
and one outward endpoint. These vertex sets are disjoint,
even for the degree-zero core path. Thus all selected
directions are off diagonal. Hold any one matrix diagonal
fixed at \(x\), and simultaneously permute rows and
columns to make it the distinguished entry \(X_{11}\).
The resulting internally padded permanent is exactly
\(x^{n-m}\operatorname{per}(xI+\mathcal B)\) with the
allowed off-diagonal variations. Its dehomogenized detected
coefficients are unchanged. This proves the permanent
claim with the stated internal-padding convention.

All substituted forms belong to the respective orbit
closures by endomorphism substitution. F-016 places every
highest-weight ideal vector of type \(\lambda\) in the
algebra generated by \(\Gamma_\lambda\). The proved
coordinate injection forces such an ideal vector to be
zero. Hence the corresponding entire ideal multiplicity
space vanishes, giving the asserted ambient equalities.
\end{proof}

\paragraph{A useful two-prefix consequence in BIP range.}
For every actual obstruction at \(n\geq m^{25}\),
\[
                   \lambda_4\geq27
                       \quad\text{or}\quad\lambda_2\geq239063403.
                                                               \tag{M-019.4}
\]
Indeed use \(b=26\), and suppose
\(\lambda_2\leq26+t\), \(\lambda_4\leq26\).
Every target monomial can have all but at most twenty-six
factors removed from the first two nonfirst variables.
Its removed part has exponents at most \(t\) in each.
Thus (M-019.1) holds with
\[
             A=\{1,x_2,\ldots,x_2^t\},\qquad
             B=\{1,x_3,\ldots,x_3^t\}.
\]
The known permanent-size reduction gives \(m\geq190\).
Since \(r\leq m^2\), the new value
\(h=\lceil\sqrt r\rceil\) is still at most \(m\).
Hence \(N\leq2(t+1)m\binom{m^2+11}{12}\).
M-017's exact constant \(C_*=478126755\) shows this
fits whenever \(2(t+1)\leq C_*\).
The largest such integer is \(t=239063376\), proving
(M-019.4). The low-length cases were already excluded by
O-006--O-007; the argument uses \(r\geq4\) otherwise.
The general theorem also permits capped or nested ideals
in each direction, rather than only these univariate ones.

\paragraph{Geometric dictionary and scope.}
The two order ideals and the jet core index finite free
torus-weight coefficient modules. Their bridge and tree
matrices are integral, and the unit differential minor
gives actual coordinate-algebra injection under every
base change. The torus grading and ideal comodules are
SGA~3, I, Propositions~4.7.2--4.7.3, pp.~34--36.
F-016 realizes the relevant Schur operations through
SGA~6, VI, Theorem~3.3, pp.~392--393, and Macdonald,
I(4.3), p.~63, Appendix~A(5.3)--(5.6), pp.~155--156.
Projection dominance compares all multiplicity copies;
it does not assert exact representation of all forms with
the target support. The explicit strictness example O-025
separates this two-sided criterion from every one-sided
order-ideal construction of M-018. No general saturation
or actual multiplicity obstruction is asserted.

% END proofs/two_sided_jet_extensions.tex

\part{Unresolved questions and limits of the current results}
\section{No unchecked theorem is used as a proved finding}
The remaining main problem is to compare the positive ranks in F-010
for permanent and determinant at shapes not excluded by M-008, M-009,
the permanent-size gate F-009, and the other explicit families above.
At \(n\geq m^{25}\), occurrence positivity alone does not decide that
rank comparison. This paper does not prove either a general multiplicity
obstruction or a general inequality \(q\leq K\). It does not prove a
general cancellation-free rule for rectangular Kronecker coefficients,
nor ordinary saturation beyond the stated restricted families.

The archive source qualifications are explicit. A-029 verifies the
relevant original cote~133 scans and recovers the definition of \(q_H^n\).
All attached archive pages have now been checked against their original
facsimiles. A classification of all singular local coefficient rings
is still not asserted. Proved
corrections and counterexamples retain their stated interpretations.

The latest timed review was recorded at 2026-11-21T02:00:17.030081+00:00.
The dated ledger records the open rows, current main gates, and the
movement assessment. The prescribed stopping rule remains in force.

\appendix
\section{Archive inventory}
Each row below is an assertion, definition, question, source-status issue,
or explicit scope exclusion from the supplied material. The full machine
readable version is \texttt{agents/archive/ARCHIVE\_INVENTORY.tsv}.
\small
\begin{longtable}{@{}>{\raggedright\arraybackslash}p{0.07\textwidth}>{\raggedright\arraybackslash}p{0.18\textwidth}>{\raggedright\arraybackslash}p{0.36\textwidth}>{\raggedright\arraybackslash}p{0.28\textwidth}@{}}
\toprule ID & Source and locator & Statement & State and resolution\\\midrule\endhead
AR001 & cote133 scan006 lines4-10 & WG is group completion of strictly nonsingular quadratic modules under orthogonal sum & well-defined additive split group; A-005. Not modern GW with metabolic relations; tensor/exterior operations require conventions \\
AR002 & cote133 scan006 lines12-14 & WG is a pre-lambda ring before assuming2 invertible & refuted; A-006. Rank2 hyperbolic plane in characteristic2 has rank1 exterior square \\
AR003 & cote133 scan006 lines12-14 & Rank is augmentation to H0(Z) with binomial operations & proved under corrected hypothesis; A-006. 2 invertible and exterior-form prelambda; ranks locally constant \\
AR004 & cote133 scan006 lines16-19 & Connected base and2 noninvertible force strictly nonsingular ranks even & proved; A-006. Odd polar determinant belongs to ideal(2); constant rank used \\
AR005 & cote133 scan006 lines19-20 & WG has a tensor unit if2 invertible & proved; A-005,A-006. Use equivalence quadratic forms with symmetric bilinear forms, unit bilinear line1 \\
AR006 & cote133 scan006 lines20-32 & Are exterior operations special on split WG when2 invertible? & refuted; A-005. BC3 over F3; Adams3 fails multiplicativity on hyperbolic modules \\
AR007 & cote133 scan006 lines25-32 & Universal product and iteration polynomials have integer coefficients & proved background; SGA6-V2.4. Product formula needs variables through lambda\textasciicircum{}i, not truncated at lambda\textasciicircum{}(i-1); uncertain transcription \\
AR008 & cote133 scan006 lines34-45 & Check product/iteration identities on classifying topoi O(m)xO(n)/O(m) & proved with explicit hypothesis; A-015. Forms must be locally isometric to chosen standard universal form in the topology; true for usual etale/fppf classification with2 invertible; no automatic arbitrary-topos local triviality \\
AR009 & cote133 scan008 lines2-8 & q plus negative q is hyperbolic for modules over a ring & proved; A-008. Finite projective underlying module and nonsingular polar form; no2 invertibility required over affine ring \\
AR010 & cote133 scan008 lines6-8 & Previous assertion fails on general locally ringed topos & proved by counterexample; A-008. Characteristic2 B(Gm semidirect C2); do not authenticate ambiguous name Lazard/Laborde \\
AR011 & cote133 scan008 lines10-20 & Witt quotient equality iff difference becomes hyperbolic after adding hyperbolic module & refuted in stated generality; A-008. 2Q zero in quotient group completion but never stably hyperbolic on characteristic2 classifying topos \\
AR012 & cote133 scan008 lines22-30 & If E=F+G and Witt cancellation is available then hyperbolic stabilization can be cancelled & proved conditional; A-008. Apply the explicitly assumed cancellation theorem; ordinary field clause is standard Witt theory \\
AR013 & cote133 scan008 lines32-33 & Hyperbolic image is ideal and W is ring & proved under corrected hypothesis; A-007. 2 invertible; H(E) tensor M = H(E tensor M) \\
AR014 & cote133 scan008 lines34-36 & Rank mod2 descends but integer rank does not & proved; A-007. Hyperbolic ranks even; hyperbolic line image has rank2 \\
AR015 & cote133 scan008 lines37-38 & Does exterior prelambda descend to Witt quotient? & refuted; A-007. lambda2 of hyperbolic plane is nonzero odd-rank line \\
AR016 & cote133 scan009 lines2-19 & SW classes exist naturally, normalized on lines, multiplicative under orthogonal sum, w0=1 & proved under precise standard hypotheses; A-009,EKV5.1-5.3. Verified etale scheme setting over Z[1/2]; arbitrary strict-local ringed-topos shorthand underspecified \\
AR017 & cote133 scan009 lines22-28 & w(H(E))=1 and SW classes factor through W & refuted; deletion justified; A-009,EKV5.5. P1\_C, E=O(1), w2 nonzero \\
AR018 & cote133 scan009 lines32-35 & Tensor and exterior SW classes given by universal rank-dependent formulas & proved in verified setting; A-009. Injective orthogonal splitting plus symmetric polynomials \\
AR019 & cote133 scan009 lines36-39 & Rank and total SW class form lambda homomorphism & proved with target convention; A-009. Completed total-class target and Chern-root operations required, not componentwise cup-product ring \\
AR020 & cote133 scan019 lines4-6 & lambda\textasciicircum{}i is class of P\_i(E); sigma\textasciicircum{}i is class of Sym\textasciicircum{}i(E) & conditional interpretation; A-001. Originalscan019 visuallyread: notation P\_i visible, no further set definition on that page; subset interpretation supported by displayed sigma2=P2+E and recovered A-029 \\
AR021 & cote133 scan019 lines8-14 & Sigma\_t lambda\_(-t)=1 and symmetric coefficients universal in lambda & refuted for subsets; proved for canonical opposite structure; A-001,A-002. Keep two prelambda structures distinct \\
AR022 & cote133 scan019 lines20-23 & For line L, lambda\_t=1+tL and sigma\_t=(1-tL)\textasciicircum{}-1 & proved under line-element interpretation; A-002. A rank1 representation or lambda line element; a finite set of cardinality1 is the trivial point \\
AR023 & cote133 scan019 lines25-29 & Linearization from G-sets through integral representations to Q/C representations & proved; A-002,A-015. Commutes with symmetric powers and canonical opposite lambda, not subset lambda \\
AR024 & cote133 scan019 lines30-45 & sigma2=E\textasciicircum{}2-lambda2 & proved for canonical lambda only; A-001. Formal coefficient identity, not subset identity \\
AR025 & cote133 scan019 line46 & E x E is equivariant sum of Sym2E and subset P2E & refuted; A-001. C2 free orbit gives unequal fixed-point marks \\
AR026 & cote133 scan019 line47 & Sym2E=P2E+E & proved; A-001. Canonical equivariant decomposition of multisets \\
AR027 & cote133 scan019 line50 & Universal special-lambda product rule on Burnside ring? & refuted for canonical Burnside operations; A-002. V4 explicit marks counterexample \\
AR028 & cote133 scan019 lines53-56 & 2P2E+E=E\textasciicircum{}2; corresponding cardinality identity & first refuted, cardinality proved; A-001. Action information is lost under cardinality \\
AR029 & cote133 scan019 line58 & P\_alpha,beta(E) and indexed alpha\_i & not a determinate theorem; NONE. No definition or identity supplied \\
AR030 & cote133 scan028 lines2-5 & K(X,*) is direct sum of K(BSnX) & well-defined with exact-category conventions; A-010. Induction graded product; finite local ranks, with A-017 removing any global rank bound \\
AR031 & cote133 scan028 lines6-12 & Tensor powers define completed exponential tau on K0 & proved; A-010. Tensor filtration respects short exact sequences \\
AR032 & cote133 scan028 lines14-18 & Integral monomial Cauchy expansion in lambda classes and rho monomials & proved for exact K0 of finite local ranks; A-010,A-017. Universal flag splitting using SGA6VI1.1 and1.14; no denominator n!; product lemma removes rank bound \\
AR033 & cote133 scan028 lines19-24 & Are monomial expansion terms direct summands for effective xi? & refuted; A-004. Standard2-dimensional S3 representation gives psi2=1+E-sign \\
AR034 & cote133 scan028 lines26-34 & rho\_i trivial representation in degreei; products are induction classes & proved; A-010. Internal tensor product and external induction product distinguished \\
AR035 & cote133 scan028 lines35-48 & Sigma\_alpha is integral weighted-homogeneous monomial symmetric polynomial in elementary functions & proved; A-010. Fundamental theorem of symmetric polynomials; weight i on lambda\textasciicircum{}i \\
AR036 & cote43 scan037 lines3-7 & Augmentation1 formal series form affine ring scheme over Z with Witt operations and lambda structure & proved with convention; A-014. Exterior-series convention; special lambda structure; relevant originalfacsimile wording/symbols now visuallyverified \\
AR037 & cote43 scan037 lines8-9 & (1+at)*f(t)=f(at) & proved; A-014. Universal formal-root identity; relevant originalfacsimile wording/symbols now visuallyverified \\
AR038 & cote43 scan037 lines13-16,30 & Coefficient Frobenius is ring endomorphism in characteristicp & proved; A-014. Operations are integral coefficient polynomials; relevant originalfacsimile wording/symbols now visuallyverified \\
AR039 & cote43 scan037 lines18-28 & f(t)\textasciicircum{}p=F(f)(t\textasciicircum{}p); V is additive injection; FV=VF=p & proved; A-014. p on right is Witt addition, not star power; relevant originalfacsimile wording/symbols now visuallyverified \\
AR040 & cote43 scan037 lines32-41 & (1+x t\textasciicircum{}(p\textasciicircum{}i))*f = F\textasciicircum{}i(f)(x t\textasciicircum{}(p\textasciicircum{}i)) & proved; A-014. Adjoining p\textasciicircum{}i root is faithfully flat; no perfectness assumption; relevant originalfacsimile wording/symbols now visuallyverified \\
AR041 & cote43 scan037 lines43-50 & Natural ring map to p-typical Witt ring with line value Teichmuller(a), determined additively & proved with continuity/naturality qualification; A-014. Do not assert abstract additive generation by lines for fixed A; relevant originalfacsimile wording/symbols now visuallyverified \\
AR042 & cote2 scan017 lines2-5 & Pullback preserves tensor product & proved; A-015. Restriction or locally free sheaf pullback; relevant originalfacsimile wording/symbols now visuallyverified \\
AR043 & cote2 scan017 lines6-9 & Projection formula for u\_* & proved with hypotheses; A-015. Explicit induction formula; right adjoint requires dualizable tensor factor; relevant originalfacsimile wording/symbols now visuallyverified \\
AR044 & cote2 scan017 line10 & Exterior and other tensor operations & proved for exterior construction; A-015. No further unspecified tensor-operation theorem inferred; relevant originalfacsimile wording/symbols now visuallyverified \\
AR045 & cote2 scan017 lines12-27 & External tensor, pullback compatibility, diagonal restriction gives internal tensor & proved; A-015. Finite locally free representation category; relevant originalfacsimile wording/symbols now visuallyverified \\
AR046 & cote2 scan017 lines28-38 & Categories with symmetric-group actions; restriction/induction transitivity & proved in standard setting; A-015. Induction preserves finite free objects for subgroup inclusions; arbitrary group homomorphism induction may fail finite local freeness; relevant originalfacsimile wording/symbols now visuallyverified \\
AR047 & cote12 scan041 lines22-24 & Determine lambda operations using finite groups of operators & interpreted, not a theorem; A-015. Antisymmetrizer requires r! invertible; exterior quotient exists integrally; relevant originalfacsimile wording/symbols now visuallyverified \\
AR048 & cote119 scan217 lines25-31 & Lambda operations connected with symmetric-group operations & mathematical dictionary supplied; A-015. Originalcote119 PDFpage217 visuallyread and compared with AI transcription; lambda/symmetric-group wording confirmed. Historical priority not independently asserted. \\
AR049 & cote16 scan038 entire page & Capital Lambda is Lefschetz lowering operator & outside this lambda/Burnside audit; NONE. Not an exterior lambda operation; formulas about primitive cohomology need surrounding pages; relevant originalfacsimile wording/symbols now visuallyverified \\
AR050 & cote4 scan007 line22 and entire page & Capital Lambda is Lefschetz operator in crystalline/etale comparison & outside this lambda/Burnside audit; NONE. No Burnside or representation lambda-ring assertion; relevant originalfacsimile wording/symbols now visuallyverified \\
AR051 & cote79 scan051 lines7-18 & lambda is scalar/group parameter in T\_lambda & outside this lambda/Burnside audit; NONE. Not lambda-ring notation; relevant originalfacsimile wording/symbols now visuallyverified \\
AR052 & cote89 scan018 entire page & Capital Lambda is coefficient ring; wedge is bilinear pairing & outside this lambda/Burnside audit; NONE. Reconstruction assertion not a lambda-ring assertion; relevant originalfacsimile wording/symbols now visuallyverified \\
SA001 & lines491-496 & Subset and inverse-symmetric Burnside formulas agree in degree2 & refuted; A-001. C2 free orbit; originalscan019 now visuallyread and displayed false identities confirmed \\
SA002 & lines500-504 & Injection powers exponential in induction Burnside ring and extend to group completion & proved; A-011. Originalscans020-023 now visuallyread \\
SA003 & SA17line505; originalcote133scan021 lowerhalf & Operations q\_H\textasciicircum{}n & definition recovered and operation proved; A-029. q\_H([E])=[H\textbackslash{}E], q\_H\textasciicircum{}n=q\_H tau\textasciicircum{}n; integer-valued polynomialdegree n, not ordinary scaling-homogeneous for n>=2; original visuallyread \\
SA004 & lines506-510 & Map powers decompose by fiber partitions; type count n!/product alpha\_i!(i!)\textasciicircum{}alpha\_i & proved; A-011. Originalscans020-023 now visuallyread; study proposition verified as stated \\
SA005 & lines511-515 & Induction from S\_m alone is wrong; wreath stabilizers and inflation are needed & proved; A-011. Explicit n2,m1 counterexample and full corrected type formula \\
SA006 & lines520-521 & Representable injection functors connected and components of map functors & proved; A-012. Category of elements has initial identity; components indexed by kernels \\
SA007 & lines522-523 & Finite sums representables equivalently summands of single map functor & refuted; corrected; A-012. Two copies h1 have value2 at singleton while every HJ has value1; finite sums HJ suffice; originalscan025 confirms the single-H\_I claim, even for arbitrarily large I \\
SA008 & lines522-523 & Special functors stable under finite limits and composition & proved; A-012. Finite equality-pattern decomposition and Yoneda \\
SA009 & line524 & Distinguished finite sums hI/H equivalent sequence Xn construction & proved with finite-support condition; A-012. Each Xn finite and only finitely many nonempty degrees \\
SA010 & lines525,529 & Eventual bijections satisfy2-out-of3 and base change & proved; A-012. Maximum threshold and pointwise pullback \\
SA011 & line526 & Functors with injective transitions closed under sum/product/composition & proved; A-012. Pointwise injections and composition \\
SA012 & lines537-541,546-548 & Tensor Cauchy formula is Schur-Weyl identity in characteristic0 & proved; integral exact-K0 refinement; A-010. Original scan028 supplied; unbounded local ranks handled by A-017 \\
SA013 & lines542,549 & Coefficients rho1\textasciicircum{}n and rho\_n are lambda\textasciicircum{}n and Adams\textasciicircum{}n & proved in universal expression; A-010. Not unique extraction in quotient rings where rho monomials satisfy relations \\
SA014 & lines543,550-556 & Is induction representation ring polynomial on rho\_i for every local ring? & refuted; A-013. F2 exact relation; singular characteristic0 counterexample \(A=\mathbb Z_{(2)}[\epsilon]/\allowbreak(2\epsilon,\epsilon^2)\); question visuallyverified on originalscan029 \\
SA015 & lines546-554 & Characteristic0 fields yes; Fp rank shortage in degreep & proved; A-013. Macdonald I7.3,I2.8; James11.5 \\
SA016 & lines555-556 & Do Z\_p lattice exact representation rings give polynomial induction ring? & proved; A-013. Exact K0→generic fiber isomorphism, unitriangular decomposition matrix kills torsion kernel \\
SA017 & line555 & Exact classification of all coefficient rings with polynomial property & not classified; A-013. Specific universal-local claim refuted; Q-local algebras and Zp proved; arbitrary singular locals remain outside asserted classification \\
SA018 & lines574-575 & Hyperbolic SW formula new in this reading & formula verified, priority claim corrected; A-009. EKV1993 Section5.5 Proposition pp183-184 \\
SA019 & lines578-580 & Prelambda fails without2 invertible; product formula must include lambda\textasciicircum{}i & proved; A-006. Originalscan006 visuallyread; literal split WG and missing2 hypothesis confirmed; polynomialendpoint ink is still recorded cautiously \\
SA020 & line585 & Special lambda structure on split WG & refuted; A-005. BC3/F3 explicit exterior/Adams computation \\
BR001 & Section1.1 lines774-807 & Special lambda geometry and Frobenius lifts on flat Z-schemes & used with scope checked; A-016. Canonical Burnside prelambda not special by A-002 \\
BR002 & Section2.5 lines956-977 & Chebychev representation-ring lambda example & used with scope checked; A-016. Adams polynomial x²-2 demonstrates no arbitrary positive-cone preservation \\
BR003 & Section6 theorem labels at1863,1913 & Finite-type lambda-scheme cohomology abelian/potentially Artin-Tate & carried as cited background, no new proof; A-016. No Kronecker realization or saturation theorem follows solely from K0 lambda structure \\
BR004 & lines2680-2694 & Finite-sets section occurs after first end-document & verified; A-016. Inactive draft text, not numbered theorem of compiled paper \\
\bottomrule\end{longtable}\normalsize

\section{Source catalog and reproducibility}
The complete exact-locator catalogs are
\path{agents/archive/SOURCE_LOCATORS.tsv},
\path{agents/formulas/SOURCE_NOTES.md}, and
\path{agents/obstructions/SOURCE_NOTES.md}. They record the supplied
LaTeX labels as well as printed statement numbers. The archive source
hash and run rules are in \texttt{RUN\_METADATA.json}. The attached
source archive was unpacked once.

The computation files use exact integers or rational arithmetic. The
character engine in \texttt{computations/independent\_characters.py}
uses Murnaghan--Nakayama independently of the family formulas, and
checks character orthogonality. The formulas worker's separate
certificates verify the generating identity as an exact integer
polynomial identity, the incidence maps, subset-sum support and the
determinant constructions. These supplement the full proofs and are
not substitutes for them. All files needed to reproduce these checks
are included; external literature is identified by exact locators
and primary URLs rather than copied as complete books.
\end{document}
